\chapter{Fronts and the weather map}\label{p3-ch:fronts}

\textit{This chapter answers Map Q14.}

\section{Air masses and fronts}

An \textbf{air mass} is a very large body of air, hundreds of kilometres across, whose temperature and humidity are much the same throughout. It takes on the character of the surface it has rested on: air that has sat over a cold, dry continent is cold and dry; air from the tropical ocean is warm and humid. A \textbf{front} is the boundary between two unlike air masses. You can recognise one by any of the following:
\begin{itemize}
\item a sharp change of \textbf{temperature};
\item a sharp change of \textbf{moisture} (the dew point);
\item a sharp change of \textbf{wind direction}.
\end{itemize}
Fronts are classified by how they move.

\begin{description}[leftmargin=0pt,font=\sffamily\bfseries]
\item[Warm front] A warm air mass advances over a cooler one. Drawn as a red line with semicircles pointing the way the front moves.
\item[Cold front] A cold air mass advances under a warmer one. Drawn as a blue line with triangles pointing the way the front moves.
\end{description}

\begin{figure}[ht]
\centering
\begin{tikzpicture}[x=1cm,y=1cm,cl/.style={cloud,cloud puffs=10,cloud puff arc=120,aspect=2.2,draw=gray!70,fill=gray!8,inner sep=1pt,font=\scriptsize}]
  % ----- warm front (left) -----
  \begin{scope}
    \node[font=\sffamily\bfseries] at (3.4,3.7) {Warm front};
    \fill[fsky!35] (0,0) -- (7,0) -- (7,0.1) -- (0.5,0.1) -- cycle;
    \fill[fsky!25] (0,0) -- (7,0) -- (7,1.7) .. controls (4.5,0.8) and (2.5,0.35) .. (1.5,0) -- cycle;
    \fill[forange!25] (7,1.7) .. controls (4.5,0.8) and (2.5,0.35) .. (1.5,0) -- (0,0) -- (0,2.4) -- (7,2.4) -- cycle;
    \draw[very thick,fred] (7,1.7) .. controls (4.5,0.8) and (2.5,0.35) .. (1.5,0);
    \draw[thick] (0,0) -- (7,0);
    \node[font=\scriptsize,text=fred!70!black] at (1.1,1.5) {warm air};
    \node[font=\scriptsize,text=fblue!70!black] at (5.6,0.5) {cool air};
    \node[cl] at (4.0,1.5) {nimbostratus}; \node[cl] at (5.9,2.0) {altostratus}; \node[cl,fill=white] at (6.3,2.9) {cirrus};
    \foreach \x in {3.3,3.7,4.1,4.5,4.9,5.3}{\draw[fblue,thin] (\x,1.0) -- (\x-0.08,0.35);}
    \draw[-{Stealth},very thick,fred] (1.2,2.9) -- (2.8,2.9) node[above,font=\scriptsize,pos=0.5]{front moves};
    \node[font=\scriptsize,gray!50!black] at (3.5,-0.3) {slope about 1 in 150: the rain begins hundreds of km ahead};
  \end{scope}
  % ----- cold front (right) -----
  \begin{scope}[xshift=8.6cm]
    \node[font=\sffamily\bfseries] at (3.4,3.7) {Cold front};
    \fill[forange!25] (0,0) -- (7,0) -- (7,2.4) -- (0,2.4) -- cycle;
    \fill[fsky!35] (0,0) -- (4.2,0) .. controls (3.4,0.5) and (2.6,1.4) .. (1.9,2.4) -- (0,2.4) -- cycle;
    \draw[very thick,fblue] (4.2,0) .. controls (3.4,0.5) and (2.6,1.4) .. (1.9,2.4);
    \draw[thick] (0,0) -- (7,0);
    \node[font=\scriptsize,text=fblue!70!black] at (1.0,0.8) {cold air};
    \node[font=\scriptsize,text=fred!70!black] at (5.8,0.6) {warm air};
    \node[cl,aspect=1.6,fill=white] at (3.7,1.7) {cumulonimbus};
    \foreach \x in {3.2,3.5,3.8,4.1}{\draw[fblue,thin] (\x,1.2) -- (\x-0.05,0.3);}
    \draw[-{Stealth},very thick,fblue] (0.6,2.9) -- (2.2,2.9) node[above,font=\scriptsize,pos=0.5]{front moves};
    \node[font=\scriptsize,gray!50!black] at (3.5,-0.3) {slope about 1 in 50: a narrow band of heavy rain};
  \end{scope}
\end{tikzpicture}
\caption{Cross-sections through a warm front (left) and a cold front (right), with the vertical scale greatly exaggerated. In both the front is a sloping surface along which the warm air rises over the cold air.}\label{p3-fig:fronts}
\end{figure}

\begin{table}[ht]
\centering
\caption{What happens as each kind of front passes a station in the Northern Hemisphere.}
\begin{tabular}{@{}p{3.1cm}p{5.8cm}p{5.8cm}@{}}
\toprule
 & \textbf{Warm front} & \textbf{Cold front}\\
\midrule
Before & pressure falls; cloud thickens from cirrus to nimbostratus; wind backs to south-east & pressure falls; warm, humid, south-westerly wind\\
Passage & steady rain or drizzle; wind veers to south-west; temperature rises & heavy showers, perhaps thunder; wind veers sharply to north-west\\
After & warm, humid air, cloud breaks; pressure steady & colder, drier air; pressure rises; clearing sky\\
\bottomrule
\end{tabular}
\end{table}

\textbf{Wind veering.} Notice that in both cases the wind \emph{veers}, that is, turns clockwise, as the front passes. Notice how the wind direction shifts at the frontal boundary.

\section{A low with a warm front and a cold front}\label{p3-sec:fig9}

Figure~\ref{p3-fig:fig9} shows a low-pressure centre (L) with a warm front curving east and south from it and a cold front trailing to the south-west. Four stations sample the air round them, which gives a small case in which to learn how to read a front before tackling the whole map.

\begin{figure}[ht]
\centering
\begin{tikzpicture}[x=0.0235cm,y=-0.0235cm,line cap=round,line join=round]
  \clip (360,950) rectangle (800,1345);
  \draw[fred,line width=1.16pt] (540.0,1082.0) -- (585.0,1081.0) -- (640.0,1087.0) -- (685.0,1105.0) -- (720.0,1127.0) -- (735.0,1165.0);
\fill[fred] (546.8,1081.8) -- (547.0,1079.9) -- (547.6,1078.0) -- (548.6,1076.2) -- (550.0,1074.8) -- (551.7,1073.7) -- (553.6,1073.0) -- (555.5,1072.7) -- (557.5,1072.9) -- (559.4,1073.5) -- (561.2,1074.6) -- (562.6,1075.9) -- (563.7,1077.6) -- (564.4,1079.5) -- (564.7,1081.5) -- cycle;
\fill[fred] (578.4,1080.3) -- (578.8,1078.3) -- (579.7,1076.5) -- (580.9,1075.0) -- (582.5,1073.7) -- (584.2,1072.8) -- (586.2,1072.4) -- (588.2,1072.4) -- (590.1,1072.8) -- (591.9,1073.7) -- (593.5,1074.9) -- (594.8,1076.5) -- (595.6,1078.3) -- (596.1,1080.2) -- (596.1,1082.2) -- cycle;
\fill[fred] (609.7,1083.7) -- (610.1,1081.7) -- (611.0,1079.9) -- (612.2,1078.4) -- (613.8,1077.1) -- (615.6,1076.2) -- (617.5,1075.8) -- (619.5,1075.8) -- (621.5,1076.2) -- (623.3,1077.1) -- (624.8,1078.3) -- (626.1,1079.9) -- (627.0,1081.7) -- (627.4,1083.6) -- (627.4,1085.6) -- cycle;
\fill[fred] (640.9,1087.4) -- (641.9,1085.6) -- (643.2,1084.1) -- (644.8,1082.9) -- (646.6,1082.1) -- (648.6,1081.8) -- (650.6,1081.9) -- (652.5,1082.4) -- (654.3,1083.3) -- (655.8,1084.7) -- (657.0,1086.3) -- (657.7,1088.1) -- (658.1,1090.1) -- (658.0,1092.1) -- (657.5,1094.0) -- cycle;
\fill[fred] (670.2,1099.1) -- (671.1,1097.3) -- (672.4,1095.8) -- (674.0,1094.6) -- (675.9,1093.8) -- (677.8,1093.5) -- (679.8,1093.6) -- (681.8,1094.1) -- (683.5,1095.0) -- (685.0,1096.4) -- (686.2,1098.0) -- (687.0,1099.8) -- (687.4,1101.8) -- (687.3,1103.8) -- (686.7,1105.7) -- cycle;
\fill[fred] (698.1,1113.3) -- (699.4,1111.7) -- (700.9,1110.5) -- (702.8,1109.6) -- (704.7,1109.1) -- (706.7,1109.1) -- (708.6,1109.6) -- (710.4,1110.5) -- (712.0,1111.7) -- (713.3,1113.3) -- (714.1,1115.1) -- (714.6,1117.0) -- (714.6,1119.0) -- (714.1,1121.0) -- (713.3,1122.8) -- cycle;
\fill[fred] (722.1,1132.3) -- (724.0,1131.8) -- (726.0,1131.7) -- (728.0,1132.1) -- (729.8,1132.8) -- (731.4,1134.0) -- (732.7,1135.5) -- (733.7,1137.3) -- (734.2,1139.2) -- (734.3,1141.2) -- (733.9,1143.2) -- (733.1,1145.0) -- (731.9,1146.6) -- (730.4,1148.0) -- (728.6,1148.9) -- cycle;
\draw[fblue,line width=1.16pt] (505.0,1132.0) -- (502.0,1180.0) -- (485.0,1220.0) -- (475.0,1240.0) -- (450.0,1265.0) -- (407.0,1285.0);
\fill[fblue] (504.7,1137.2) -- (517.6,1148.6) -- (503.4,1158.2) -- cycle;
\fill[fblue] (502.7,1168.7) -- (515.7,1180.0) -- (501.4,1189.6) -- cycle;
\fill[fblue] (494.1,1198.6) -- (502.6,1213.6) -- (485.9,1217.9) -- cycle;
\fill[fblue] (481.3,1227.3) -- (488.8,1242.8) -- (471.9,1246.1) -- cycle;
\fill[fblue] (462.7,1252.3) -- (465.0,1269.3) -- (447.9,1267.1) -- cycle;
\fill[fblue] (437.8,1270.7) -- (434.0,1287.5) -- (418.7,1279.5) -- cycle;
\node[lowmark,scale=1.5] at (510,1092) {L};
\draw[black,line width=0.63pt] (631.0,996.0) -- (673.3,1014.0);
\draw[black,line width=0.63pt] (673.3,1014.0) -- (674.5,1032.5);
\draw[black,line width=0.63pt,fill=white] (631.0,996.0) circle (9.8);
\fill[black] (631.0,996.0) circle (9.8);
\node[anchor=east,inner sep=0pt,font=\fontsize{7.0}{7.0}\selectfont,text=black] at (618.4,983.4) {56};
\node[anchor=east,inner sep=0pt,font=\fontsize{7.0}{7.0}\selectfont,text=black] at (618.4,1019.0) {53};
\node[anchor=west,inner sep=0pt,font=\fontsize{7.0}{7.0}\selectfont,text=black] at (643.6,983.4) {019};
\fill[black] (605.0,1001.8) circle (2.5);
\fill[black] (613.3,1001.8) circle (2.5);
\node[anchor=center,inner sep=0pt,font=\fontsize{9.0}{9.0}\selectfont,text=fig9label] at (647.0,1030.0) {\textbf{A}};
\draw[black,line width=0.63pt] (664.0,1152.0) -- (666.4,1197.9);
\draw[black,line width=0.63pt] (666.4,1197.9) -- (650.3,1207.1);
\draw[black,line width=0.63pt,fill=white] (664.0,1152.0) circle (9.8);
\fill[black] (664.0,1152.0) -- (664.0,1142.2) -- (665.3,1142.3) -- (666.5,1142.6) -- (667.7,1143.0) -- (668.9,1143.5) -- (670.0,1144.2) -- (670.9,1145.1) -- (671.8,1146.0) -- (672.5,1147.1) -- (673.0,1148.3) -- (673.4,1149.5) -- (673.7,1150.7) -- (673.8,1152.0) -- cycle;
\node[anchor=east,inner sep=0pt,font=\fontsize{7.0}{7.0}\selectfont,text=black] at (651.4,1139.3) {68};
\node[anchor=east,inner sep=0pt,font=\fontsize{7.0}{7.0}\selectfont,text=black] at (651.4,1175.0) {63};
\node[anchor=west,inner sep=0pt,font=\fontsize{7.0}{7.0}\selectfont,text=black] at (676.6,1139.3) {025};
\node[anchor=center,inner sep=0pt,font=\fontsize{9.0}{9.0}\selectfont,text=fig9label] at (680.0,1186.0) {\textbf{B}};
\draw[black,line width=0.63pt] (582.0,1237.0) -- (571.7,1281.8);
\draw[black,line width=0.63pt] (571.7,1281.8) -- (553.7,1286.2);
\draw[black,line width=0.63pt,fill=white] (582.0,1237.0) circle (9.8);
\fill[black] (582.0,1237.0) -- (582.0,1227.2) -- (583.3,1227.3) -- (584.5,1227.6) -- (585.7,1228.0) -- (586.9,1228.5) -- (588.0,1229.2) -- (588.9,1230.1) -- (589.8,1231.0) -- (590.5,1232.1) -- (591.0,1233.3) -- (591.4,1234.5) -- (591.7,1235.7) -- (591.8,1237.0) -- (591.7,1238.3) -- (591.4,1239.5) -- (591.0,1240.7) -- (590.5,1241.9) -- (589.8,1243.0) -- (588.9,1243.9) -- (588.0,1244.8) -- (586.9,1245.5) -- (585.7,1246.0) -- (584.5,1246.4) -- (583.3,1246.7) -- (582.0,1246.8) -- cycle;
\node[anchor=east,inner sep=0pt,font=\fontsize{7.0}{7.0}\selectfont,text=black] at (569.4,1224.3) {70};
\node[anchor=east,inner sep=0pt,font=\fontsize{7.0}{7.0}\selectfont,text=black] at (569.4,1260.0) {64};
\node[anchor=west,inner sep=0pt,font=\fontsize{7.0}{7.0}\selectfont,text=black] at (594.6,1224.3) {030};
\node[anchor=center,inner sep=0pt,font=\fontsize{9.0}{9.0}\selectfont,text=fig9label] at (598.0,1271.0) {\textbf{C}};
\draw[black,line width=0.63pt] (412.0,1200.0) -- (369.3,1182.8);
\draw[black,line width=0.63pt] (369.3,1182.8) -- (367.9,1164.3);
\draw[black,line width=0.63pt] (375.7,1185.4) -- (374.3,1166.9);
\draw[black,line width=0.63pt,fill=white] (412.0,1200.0) circle (9.8);
\node[anchor=east,inner sep=0pt,font=\fontsize{7.0}{7.0}\selectfont,text=black] at (399.4,1187.3) {50};
\node[anchor=east,inner sep=0pt,font=\fontsize{7.0}{7.0}\selectfont,text=black] at (399.4,1223.0) {38};
\node[anchor=west,inner sep=0pt,font=\fontsize{7.0}{7.0}\selectfont,text=black] at (424.6,1187.3) {041};
\node[anchor=center,inner sep=0pt,font=\fontsize{9.0}{9.0}\selectfont,text=fig9label] at (386.0,1234.0) {\textbf{D}};

  \draw[forange!85!black,very thick,dashed,-{Stealth[length=2.2mm]}] (582,1237) -- (415,1203);
  \draw[forange!85!black,very thick,dashed,-{Stealth[length=2.2mm]}] (631,1000) -- (664,1150);
  \node[font=\scriptsize\bfseries,text=forange!60!black,fill=white,inner sep=1pt] at (498,1160) {cold front};
  \node[font=\scriptsize\bfseries,text=forange!60!black,fill=white,inner sep=1pt] at (702,1075) {warm front};
\end{tikzpicture}
\caption{Station A is north of the warm front, B and C are in the warm sector between the fronts, D is behind the cold front. The dashed orange arrows mark the two crossings studied below: C to D across the cold front, and A to B across the warm front.}\label{p3-fig:fig9}
\end{figure}

\begin{table}[ht]
\centering
\caption{The four stations of Figure~\ref{p3-fig:fig9}.}
\small
\setlength{\tabcolsep}{4pt}
\begin{tabular}{@{}cccccll@{}}
\toprule
\textbf{Stn} & \textbf{Temp.\ (\degF)} & \textbf{Dew pt (\degF)} & \textbf{RH (\%)} & \textbf{p (hPa)} & \textbf{Wind from} & \textbf{Sky}\\
\midrule
A & 56 & 53 & 90 & 1001.9 & ESE 10 kt & overcast, rain\\
B & 68 & 63 & 84 & 1002.5 & S 10 kt & 2--3/10\\
C & 70 & 64 & 81 & 1003.0 & SSW 10 kt & 5/10\\
D & 50 & 38 & 63 & 1004.1 & WNW 20 kt & clear\\
\bottomrule
\end{tabular}
\end{table}

\begin{enumerate}[label=\textbf{Step \arabic*.},leftmargin=*,itemsep=4pt]
\item \textbf{The warm sector (B and C).} The air between the fronts is warm (68 and 70\,\degF) and humid (dew points 63 and 64\,\degF, relative humidity above 80\,\%), and the wind is from the south.
\item \textbf{Behind the cold front (D).} The air is cooler (50\,\degF), much drier (dew point 38\,\degF, RH 63\,\%), the sky is clear and the wind has swung to the west-north-west.
\item \textbf{Crossing the cold front, C to D.} The temperature falls 20\,\degF, the dew point falls 26\,\degF, the wind turns from south-south-west (193\deg) to west-north-west (292\deg): clockwise by about $100\deg$. The pressure rises and the cloud clears.
\item \textbf{Crossing the warm front, A to B.} The temperature rises 12\,\degF, the dew point 10\,\degF, and the wind turns from east-south-east (113\deg) to south (177\deg): clockwise by about $64\deg$. The sky changes from overcast and raining to almost clear.
\end{enumerate}

\begin{answerbox}[What the four stations teach]
A front is marked by \emph{simultaneous} changes of temperature, moisture and wind. At \textbf{both} fronts the wind \textbf{veers} (turns clockwise). The cold front brings the bigger changes in temperature and dew point; the warm front brings the continuous rain.
\end{answerbox}

\section{Question 14: the full analysis}\label{p3-sec:fig10}

\begin{figure}[H]
\centering
\mapbase{0.0196cm}{\draw[black,line width=0.51pt] (476.0,583.0) -- (456.5,551.8);
\draw[black,line width=0.51pt] (456.5,551.8) -- (464.2,539.2);
\draw[black,line width=0.51pt,fill=white] (476.0,583.0) circle (7.8);
\fill[black] (476.0,583.0) -- (476.0,575.2) -- (477.0,575.2) -- (478.0,575.4) -- (479.0,575.8) -- (479.9,576.2) -- (480.8,576.8) -- (481.5,577.5) -- (482.2,578.2) -- (482.8,579.1) -- (483.2,580.0) -- (483.6,581.0) -- (483.8,582.0) -- (483.8,583.0) -- cycle;
\node[anchor=east,inner sep=0pt,font=\fontsize{5.4}{5.4}\selectfont,text=black] at (465.9,572.9) {34};
\node[anchor=east,inner sep=0pt,font=\fontsize{5.4}{5.4}\selectfont,text=black] at (465.9,601.4) {22};
\node[anchor=west,inner sep=0pt,font=\fontsize{5.4}{5.4}\selectfont,text=black] at (486.1,572.9) {110};
\draw[black,line width=0.51pt] (626.0,593.0) -- (662.7,589.8);
\draw[black,line width=0.51pt] (662.7,589.8) -- (670.4,602.4);
\draw[black,line width=0.51pt,fill=white] (626.0,593.0) circle (7.8);
\fill[black] (626.0,593.0) circle (7.8);
\node[anchor=east,inner sep=0pt,font=\fontsize{5.4}{5.4}\selectfont,text=black] at (615.9,582.9) {45};
\node[anchor=east,inner sep=0pt,font=\fontsize{5.4}{5.4}\selectfont,text=black] at (615.9,611.4) {43};
\node[anchor=west,inner sep=0pt,font=\fontsize{5.4}{5.4}\selectfont,text=black] at (636.1,582.9) {096};
\fill[black] (605.2,597.6) circle (2.0);
\fill[black] (611.8,597.6) circle (2.0);
\draw[black,line width=0.51pt] (824.0,581.0) -- (859.9,572.7);
\draw[black,line width=0.51pt] (859.9,572.7) -- (869.3,584.1);
\draw[black,line width=0.51pt] (854.5,574.0) -- (859.2,579.7);
\draw[black,line width=0.51pt,fill=white] (824.0,581.0) circle (7.8);
\fill[black] (824.0,581.0) -- (824.0,573.2) -- (825.0,573.2) -- (826.0,573.4) -- (827.0,573.8) -- (827.9,574.2) -- (828.8,574.8) -- (829.5,575.5) -- (830.2,576.2) -- (830.8,577.1) -- (831.2,578.0) -- (831.6,579.0) -- (831.8,580.0) -- (831.8,581.0) -- (831.8,582.0) -- (831.6,583.0) -- (831.2,584.0) -- (830.8,584.9) -- (830.2,585.8) -- (829.5,586.5) -- (828.8,587.2) -- (827.9,587.8) -- (827.0,588.2) -- (826.0,588.6) -- (825.0,588.8) -- (824.0,588.8) -- (823.0,588.8) -- (822.0,588.6) -- (821.0,588.2) -- (820.1,587.8) -- (819.2,587.2) -- (818.5,586.5) -- (817.8,585.8) -- (817.2,584.9) -- (816.8,584.0) -- (816.4,583.0) -- (816.2,582.0) -- (816.2,581.0) -- cycle;
\node[anchor=east,inner sep=0pt,font=\fontsize{5.4}{5.4}\selectfont,text=black] at (813.9,570.9) {39};
\node[anchor=east,inner sep=0pt,font=\fontsize{5.4}{5.4}\selectfont,text=black] at (813.9,599.4) {39};
\node[anchor=west,inner sep=0pt,font=\fontsize{5.4}{5.4}\selectfont,text=black] at (834.1,570.9) {154};
\draw[black,line width=0.51pt] (465.0,654.0) -- (441.3,625.8);
\draw[black,line width=0.51pt] (441.3,625.8) -- (447.2,612.2);
\draw[black,line width=0.51pt,fill=white] (465.0,654.0) circle (7.8);
\fill[black] (465.0,654.0) -- (465.0,646.2) -- (466.0,646.2) -- (467.0,646.4) -- (468.0,646.8) -- (468.9,647.2) -- (469.8,647.8) -- (470.5,648.5) -- (471.2,649.2) -- (471.8,650.1) -- (472.2,651.0) -- (472.6,652.0) -- (472.8,653.0) -- (472.8,654.0) -- (472.8,655.0) -- (472.6,656.0) -- (472.2,657.0) -- (471.8,657.9) -- (471.2,658.8) -- (470.5,659.5) -- (469.8,660.2) -- (468.9,660.8) -- (468.0,661.2) -- (467.0,661.6) -- (466.0,661.8) -- (465.0,661.8) -- cycle;
\node[anchor=east,inner sep=0pt,font=\fontsize{5.4}{5.4}\selectfont,text=black] at (454.9,643.9) {36};
\node[anchor=east,inner sep=0pt,font=\fontsize{5.4}{5.4}\selectfont,text=black] at (454.9,672.4) {28};
\node[anchor=west,inner sep=0pt,font=\fontsize{5.4}{5.4}\selectfont,text=black] at (475.1,643.9) {124};
\draw[black,line width=0.51pt] (544.0,662.0) -- (519.4,634.7);
\draw[black,line width=0.51pt] (519.4,634.7) -- (524.8,620.9);
\draw[black,line width=0.51pt] (523.1,638.8) -- (528.5,625.0);
\draw[black,line width=0.51pt,fill=white] (544.0,662.0) circle (7.8);
\fill[black] (544.0,662.0) circle (7.8);
\node[anchor=east,inner sep=0pt,font=\fontsize{5.4}{5.4}\selectfont,text=black] at (533.9,651.9) {40};
\node[anchor=east,inner sep=0pt,font=\fontsize{5.4}{5.4}\selectfont,text=black] at (533.9,680.4) {38};
\node[anchor=west,inner sep=0pt,font=\fontsize{5.4}{5.4}\selectfont,text=black] at (554.1,651.9) {088};
\fill[black] (526.5,666.6) circle (2.0);
\draw[black,line width=0.51pt] (738.0,624.0) -- (767.0,646.7);
\draw[black,line width=0.51pt] (767.0,646.7) -- (764.1,661.2);
\draw[black,line width=0.51pt] (762.6,643.3) -- (759.7,657.8);
\draw[black,line width=0.51pt,fill=white] (738.0,624.0) circle (7.8);
\fill[black] (738.0,624.0) circle (7.8);
\node[anchor=east,inner sep=0pt,font=\fontsize{5.4}{5.4}\selectfont,text=black] at (727.9,613.9) {39};
\node[anchor=east,inner sep=0pt,font=\fontsize{5.4}{5.4}\selectfont,text=black] at (727.9,642.4) {39};
\node[anchor=west,inner sep=0pt,font=\fontsize{5.4}{5.4}\selectfont,text=black] at (748.1,613.9) {084};
\fill[black] (717.2,628.6) circle (2.0);
\fill[black] (723.8,628.6) circle (2.0);
\draw[black,line width=0.51pt] (822.0,635.0) -- (857.0,646.4);
\draw[black,line width=0.51pt] (857.0,646.4) -- (859.2,661.0);
\draw[black,line width=0.51pt] (851.7,644.7) -- (854.0,659.3);
\draw[black,line width=0.51pt,fill=white] (822.0,635.0) circle (7.8);
\fill[black] (822.0,635.0) -- (822.0,627.2) -- (823.0,627.2) -- (824.0,627.4) -- (825.0,627.8) -- (825.9,628.2) -- (826.8,628.8) -- (827.5,629.5) -- (828.2,630.2) -- (828.8,631.1) -- (829.2,632.0) -- (829.6,633.0) -- (829.8,634.0) -- (829.8,635.0) -- (829.8,636.0) -- (829.6,637.0) -- (829.2,638.0) -- (828.8,638.9) -- (828.2,639.8) -- (827.5,640.5) -- (826.8,641.2) -- (825.9,641.8) -- (825.0,642.2) -- (824.0,642.6) -- (823.0,642.8) -- (822.0,642.8) -- cycle;
\draw[black,line width=0.51pt] (814.2,635.0) -- (829.8,635.0);
\draw[black,line width=0.51pt] (822.0,627.2) -- (822.0,642.8);
\node[anchor=east,inner sep=0pt,font=\fontsize{5.4}{5.4}\selectfont,text=black] at (811.9,624.9) {43};
\node[anchor=east,inner sep=0pt,font=\fontsize{5.4}{5.4}\selectfont,text=black] at (811.9,653.4) {41};
\node[anchor=west,inner sep=0pt,font=\fontsize{5.4}{5.4}\selectfont,text=black] at (832.1,624.9) {132};
\draw[black,line width=0.51pt] (672.0,681.0) -- (701.0,658.3);
\draw[black,line width=0.51pt] (701.0,658.3) -- (714.4,664.7);
\draw[black,line width=0.51pt,fill=white] (672.0,681.0) circle (7.8);
\fill[black] (672.0,681.0) circle (7.8);
\node[anchor=east,inner sep=0pt,font=\fontsize{5.4}{5.4}\selectfont,text=black] at (661.9,670.9) {48};
\node[anchor=east,inner sep=0pt,font=\fontsize{5.4}{5.4}\selectfont,text=black] at (661.9,699.4) {45};
\node[anchor=west,inner sep=0pt,font=\fontsize{5.4}{5.4}\selectfont,text=black] at (682.1,670.9) {024};
\fill[black] (651.2,685.6) circle (2.0);
\fill[black] (657.8,685.6) circle (2.0);
\draw[black,line width=0.51pt] (600.0,710.0) -- (565.9,696.2);
\draw[black,line width=0.51pt] (565.9,696.2) -- (564.7,681.4);
\draw[black,line width=0.51pt,fill=white] (600.0,710.0) circle (7.8);
\fill[black] (600.0,710.0) circle (7.8);
\node[anchor=east,inner sep=0pt,font=\fontsize{5.4}{5.4}\selectfont,text=black] at (589.9,699.9) {48};
\node[anchor=west,inner sep=0pt,font=\fontsize{5.4}{5.4}\selectfont,text=black] at (610.1,699.9) {026};
\fill[black] (579.2,714.6) circle (2.0);
\fill[black] (585.8,714.6) circle (2.0);
\draw[black,line width=0.51pt] (824.0,683.0) -- (860.4,688.1);
\draw[black,line width=0.51pt] (855.0,687.4) -- (857.3,694.4);
\draw[black,line width=0.51pt,fill=white] (824.0,683.0) circle (7.8);
\fill[black] (824.0,683.0) circle (7.8);
\node[anchor=east,inner sep=0pt,font=\fontsize{5.4}{5.4}\selectfont,text=black] at (813.9,672.9) {46};
\node[anchor=east,inner sep=0pt,font=\fontsize{5.4}{5.4}\selectfont,text=black] at (813.9,701.4) {43};
\node[anchor=west,inner sep=0pt,font=\fontsize{5.4}{5.4}\selectfont,text=black] at (834.1,672.9) {124};
\fill[black] (803.2,687.6) circle (2.0);
\fill[black] (809.8,687.6) circle (2.0);
\draw[black,line width=0.51pt] (908.0,712.0) -- (943.5,721.5);
\draw[black,line width=0.51pt,fill=white] (908.0,712.0) circle (7.8);
\fill[black] (908.0,712.0) circle (7.8);
\node[anchor=east,inner sep=0pt,font=\fontsize{5.4}{5.4}\selectfont,text=black] at (897.9,701.9) {45};
\node[anchor=east,inner sep=0pt,font=\fontsize{5.4}{5.4}\selectfont,text=black] at (897.9,730.4) {43};
\node[anchor=west,inner sep=0pt,font=\fontsize{5.4}{5.4}\selectfont,text=black] at (918.1,701.9) {163};
\fill[black] (890.5,716.6) circle (2.0);
\draw[black,line width=0.51pt] (419.0,722.0) -- (394.4,694.7);
\draw[black,line width=0.51pt] (394.4,694.7) -- (399.8,680.9);
\draw[black,line width=0.51pt,fill=white] (419.0,722.0) circle (7.8);
\node[anchor=east,inner sep=0pt,font=\fontsize{5.4}{5.4}\selectfont,text=black] at (408.9,711.9) {39};
\node[anchor=east,inner sep=0pt,font=\fontsize{5.4}{5.4}\selectfont,text=black] at (408.9,740.4) {31};
\node[anchor=west,inner sep=0pt,font=\fontsize{5.4}{5.4}\selectfont,text=black] at (429.1,711.9) {124};
\draw[black,line width=0.51pt] (714.0,722.0) -- (698.4,755.4);
\draw[black,line width=0.51pt] (698.4,755.4) -- (683.6,755.8);
\draw[black,line width=0.51pt,fill=white] (714.0,722.0) circle (7.8);
\fill[black] (714.0,722.0) -- (714.0,714.2) -- (715.0,714.2) -- (716.0,714.4) -- (717.0,714.8) -- (717.9,715.2) -- (718.8,715.8) -- (719.5,716.5) -- (720.2,717.2) -- (720.8,718.1) -- (721.2,719.0) -- (721.6,720.0) -- (721.8,721.0) -- (721.8,722.0) -- (721.8,723.0) -- (721.6,724.0) -- (721.2,725.0) -- (720.8,725.9) -- (720.2,726.8) -- (719.5,727.5) -- (718.8,728.2) -- (717.9,728.8) -- (717.0,729.2) -- (716.0,729.6) -- (715.0,729.8) -- (714.0,729.8) -- (713.0,729.8) -- (712.0,729.6) -- (711.0,729.2) -- (710.1,728.8) -- (709.2,728.2) -- (708.5,727.5) -- (707.8,726.8) -- (707.2,725.9) -- (706.8,725.0) -- (706.4,724.0) -- (706.2,723.0) -- (706.2,722.0) -- cycle;
\node[anchor=east,inner sep=0pt,font=\fontsize{5.4}{5.4}\selectfont,text=black] at (703.9,711.9) {57};
\node[anchor=east,inner sep=0pt,font=\fontsize{5.4}{5.4}\selectfont,text=black] at (703.9,740.4) {54};
\node[anchor=west,inner sep=0pt,font=\fontsize{5.4}{5.4}\selectfont,text=black] at (724.1,711.9) {029};
\draw[black,line width=0.51pt] (818.0,742.0) -- (854.2,748.4);
\draw[black,line width=0.51pt] (854.2,748.4) -- (858.5,762.6);
\draw[black,line width=0.51pt,fill=white] (818.0,742.0) circle (7.8);
\fill[black] (818.0,742.0) circle (7.8);
\node[anchor=east,inner sep=0pt,font=\fontsize{5.4}{5.4}\selectfont,text=black] at (807.9,731.9) {54};
\node[anchor=east,inner sep=0pt,font=\fontsize{5.4}{5.4}\selectfont,text=black] at (807.9,760.4) {45};
\node[anchor=west,inner sep=0pt,font=\fontsize{5.4}{5.4}\selectfont,text=black] at (828.1,731.9) {097};
\draw[black,line width=0.51pt] (469.0,770.0) -- (438.5,749.4);
\draw[black,line width=0.51pt] (438.5,749.4) -- (440.4,734.7);
\draw[black,line width=0.51pt,fill=white] (469.0,770.0) circle (7.8);
\draw[black,line width=0.51pt] (469.0,762.2) -- (469.0,777.8);
\node[anchor=east,inner sep=0pt,font=\fontsize{5.4}{5.4}\selectfont,text=black] at (458.9,759.9) {41};
\node[anchor=east,inner sep=0pt,font=\fontsize{5.4}{5.4}\selectfont,text=black] at (458.9,788.4) {30};
\node[anchor=west,inner sep=0pt,font=\fontsize{5.4}{5.4}\selectfont,text=black] at (479.1,759.9) {103};
\draw[black,line width=0.51pt] (570.0,768.0) -- (544.9,741.1);
\draw[black,line width=0.51pt] (544.9,741.1) -- (550.1,727.2);
\draw[black,line width=0.51pt] (548.7,745.1) -- (553.8,731.2);
\draw[black,line width=0.51pt] (552.4,749.2) -- (555.0,742.2);
\draw[black,line width=0.51pt,fill=white] (570.0,768.0) circle (7.8);
\fill[black] (570.0,768.0) -- (570.0,760.2) -- (571.0,760.2) -- (572.0,760.4) -- (573.0,760.8) -- (573.9,761.2) -- (574.8,761.8) -- (575.5,762.5) -- (576.2,763.2) -- (576.8,764.1) -- (577.2,765.0) -- (577.6,766.0) -- (577.8,767.0) -- (577.8,768.0) -- (577.8,769.0) -- (577.6,770.0) -- (577.2,771.0) -- (576.8,771.9) -- (576.2,772.8) -- (575.5,773.5) -- (574.8,774.2) -- (573.9,774.8) -- (573.0,775.2) -- (572.0,775.6) -- (571.0,775.8) -- (570.0,775.8) -- cycle;
\node[anchor=east,inner sep=0pt,font=\fontsize{5.4}{5.4}\selectfont,text=black] at (559.9,757.9) {45};
\node[anchor=east,inner sep=0pt,font=\fontsize{5.4}{5.4}\selectfont,text=black] at (559.9,786.4) {40};
\node[anchor=west,inner sep=0pt,font=\fontsize{5.4}{5.4}\selectfont,text=black] at (580.1,757.9) {069};
\draw[black,line width=0.51pt] (637.0,765.0) -- (618.0,796.5);
\draw[black,line width=0.51pt] (618.0,796.5) -- (603.3,795.4);
\draw[black,line width=0.51pt] (620.9,791.8) -- (606.1,790.7);
\draw[black,line width=0.51pt] (623.7,787.1) -- (616.3,786.5);
\draw[black,line width=0.51pt,fill=white] (637.0,765.0) circle (7.8);
\fill[black] (637.0,765.0) circle (7.8);
\node[anchor=east,inner sep=0pt,font=\fontsize{5.4}{5.4}\selectfont,text=black] at (626.9,754.9) {60};
\node[anchor=east,inner sep=0pt,font=\fontsize{5.4}{5.4}\selectfont,text=black] at (626.9,783.4) {55};
\node[anchor=west,inner sep=0pt,font=\fontsize{5.4}{5.4}\selectfont,text=black] at (647.1,754.9) {040};
\draw[black,line width=0.43pt] (616.4,763.3) -- (616.4,775.9);
\draw[black,line width=0.43pt] (616.4,763.3) -- (621.1,763.3) .. controls (625.4,763.3) and (625.4,770.0) .. (621.1,770.0) -- (616.4,770.0);
\draw[black,line width=0.43pt] (620.3,770.0) -- (625.8,775.9);
\fill[black] (626.9,777.0) -- (623.1,776.1) -- (625.6,773.4) -- cycle;
\fill[black] (618.7,759.4) circle (1.9);
\draw[black,line width=0.51pt] (704.0,785.0) -- (685.0,816.5);
\draw[black,line width=0.51pt] (685.0,816.5) -- (670.3,815.4);
\draw[black,line width=0.51pt,fill=white] (704.0,785.0) circle (7.8);
\fill[black] (704.0,785.0) -- (704.0,777.2) -- (705.0,777.2) -- (706.0,777.4) -- (707.0,777.8) -- (707.9,778.2) -- (708.8,778.8) -- (709.5,779.5) -- (710.2,780.2) -- (710.8,781.1) -- (711.2,782.0) -- (711.6,783.0) -- (711.8,784.0) -- (711.8,785.0) -- (711.8,786.0) -- (711.6,787.0) -- (711.2,788.0) -- (710.8,788.9) -- (710.2,789.8) -- (709.5,790.5) -- (708.8,791.2) -- (707.9,791.8) -- (707.0,792.2) -- (706.0,792.6) -- (705.0,792.8) -- (704.0,792.8) -- cycle;
\node[anchor=east,inner sep=0pt,font=\fontsize{5.4}{5.4}\selectfont,text=black] at (693.9,774.9) {61};
\node[anchor=east,inner sep=0pt,font=\fontsize{5.4}{5.4}\selectfont,text=black] at (693.9,803.4) {57};
\node[anchor=west,inner sep=0pt,font=\fontsize{5.4}{5.4}\selectfont,text=black] at (714.1,774.9) {044};
\fill[black] (683.2,789.6) circle (2.0);
\fill[black] (689.8,789.6) circle (2.0);
\draw[black,line width=0.51pt] (787.0,798.0) -- (787.0,834.8);
\draw[black,line width=0.51pt] (787.0,834.8) -- (773.8,841.4);
\draw[black,line width=0.51pt,fill=white] (787.0,798.0) circle (7.8);
\fill[black] (787.0,798.0) -- (787.0,790.2) -- (788.0,790.2) -- (789.0,790.4) -- (790.0,790.8) -- (790.9,791.2) -- (791.8,791.8) -- (792.5,792.5) -- (793.2,793.2) -- (793.8,794.1) -- (794.2,795.0) -- (794.6,796.0) -- (794.8,797.0) -- (794.8,798.0) -- (794.8,799.0) -- (794.6,800.0) -- (794.2,801.0) -- (793.8,801.9) -- (793.2,802.8) -- (792.5,803.5) -- (791.8,804.2) -- (790.9,804.8) -- (790.0,805.2) -- (789.0,805.6) -- (788.0,805.8) -- (787.0,805.8) -- (786.0,805.8) -- (785.0,805.6) -- (784.0,805.2) -- (783.1,804.8) -- (782.2,804.2) -- (781.5,803.5) -- (780.8,802.8) -- (780.2,801.9) -- (779.8,801.0) -- (779.4,800.0) -- (779.2,799.0) -- (779.2,798.0) -- cycle;
\node[anchor=east,inner sep=0pt,font=\fontsize{5.4}{5.4}\selectfont,text=black] at (776.9,787.9) {66};
\node[anchor=east,inner sep=0pt,font=\fontsize{5.4}{5.4}\selectfont,text=black] at (776.9,816.4) {53};
\node[anchor=west,inner sep=0pt,font=\fontsize{5.4}{5.4}\selectfont,text=black] at (797.1,787.9) {099};
\draw[black,line width=0.51pt] (828.0,818.0) -- (828.0,854.8);
\draw[black,line width=0.51pt] (828.0,854.8) -- (814.8,861.4);
\draw[black,line width=0.51pt,fill=white] (828.0,818.0) circle (7.8);
\fill[black] (828.0,818.0) -- (828.0,810.2) -- (829.0,810.2) -- (830.0,810.4) -- (831.0,810.8) -- (831.9,811.2) -- (832.8,811.8) -- (833.5,812.5) -- (834.2,813.2) -- (834.8,814.1) -- (835.2,815.0) -- (835.6,816.0) -- (835.8,817.0) -- (835.8,818.0) -- cycle;
\node[anchor=east,inner sep=0pt,font=\fontsize{5.4}{5.4}\selectfont,text=black] at (817.9,807.9) {66};
\node[anchor=east,inner sep=0pt,font=\fontsize{5.4}{5.4}\selectfont,text=black] at (817.9,836.4) {55};
\node[anchor=west,inner sep=0pt,font=\fontsize{5.4}{5.4}\selectfont,text=black] at (838.1,807.9) {111};
\draw[black,line width=0.51pt] (405.0,838.0) -- (374.9,816.9);
\draw[black,line width=0.51pt] (374.9,816.9) -- (377.0,802.2);
\draw[black,line width=0.51pt,fill=white] (405.0,838.0) circle (7.8);
\node[anchor=east,inner sep=0pt,font=\fontsize{5.4}{5.4}\selectfont,text=black] at (394.9,827.9) {41};
\node[anchor=east,inner sep=0pt,font=\fontsize{5.4}{5.4}\selectfont,text=black] at (394.9,856.4) {33};
\node[anchor=west,inner sep=0pt,font=\fontsize{5.4}{5.4}\selectfont,text=black] at (415.1,827.9) {126};
\draw[black,line width=0.51pt] (496.0,831.0) -- (467.8,807.3);
\draw[black,line width=0.51pt] (467.8,807.3) -- (471.3,792.9);
\draw[black,line width=0.51pt,fill=white] (496.0,831.0) circle (7.8);
\node[anchor=east,inner sep=0pt,font=\fontsize{5.4}{5.4}\selectfont,text=black] at (485.9,820.9) {42};
\node[anchor=east,inner sep=0pt,font=\fontsize{5.4}{5.4}\selectfont,text=black] at (485.9,849.4) {36};
\node[anchor=west,inner sep=0pt,font=\fontsize{5.4}{5.4}\selectfont,text=black] at (506.1,820.9) {109};
\draw[black,line width=0.51pt] (608.0,861.0) -- (588.5,892.2);
\draw[black,line width=0.51pt] (588.5,892.2) -- (573.8,890.8);
\draw[black,line width=0.51pt,fill=white] (608.0,861.0) circle (7.8);
\fill[black] (608.0,861.0) circle (7.8);
\node[anchor=east,inner sep=0pt,font=\fontsize{5.4}{5.4}\selectfont,text=black] at (597.9,850.9) {61};
\node[anchor=east,inner sep=0pt,font=\fontsize{5.4}{5.4}\selectfont,text=black] at (597.9,879.4) {57};
\node[anchor=west,inner sep=0pt,font=\fontsize{5.4}{5.4}\selectfont,text=black] at (618.1,850.9) {075};
\fill[black] (587.2,865.6) circle (2.0);
\fill[black] (593.8,865.6) circle (2.0);
\draw[black,line width=0.51pt] (828.0,885.0) -- (828.0,921.8);
\draw[black,line width=0.51pt] (828.0,921.8) -- (814.8,928.4);
\draw[black,line width=0.51pt,fill=white] (828.0,885.0) circle (7.8);
\fill[black] (828.0,885.0) -- (828.0,877.2) -- (829.0,877.2) -- (830.0,877.4) -- (831.0,877.8) -- (831.9,878.2) -- (832.8,878.8) -- (833.5,879.5) -- (834.2,880.2) -- (834.8,881.1) -- (835.2,882.0) -- (835.6,883.0) -- (835.8,884.0) -- (835.8,885.0) -- (835.8,886.0) -- (835.6,887.0) -- (835.2,888.0) -- (834.8,888.9) -- (834.2,889.8) -- (833.5,890.5) -- (832.8,891.2) -- (831.9,891.8) -- (831.0,892.2) -- (830.0,892.6) -- (829.0,892.8) -- (828.0,892.8) -- cycle;
\node[anchor=east,inner sep=0pt,font=\fontsize{5.4}{5.4}\selectfont,text=black] at (817.9,874.9) {71};
\node[anchor=east,inner sep=0pt,font=\fontsize{5.4}{5.4}\selectfont,text=black] at (817.9,903.4) {60};
\node[anchor=west,inner sep=0pt,font=\fontsize{5.4}{5.4}\selectfont,text=black] at (838.1,874.9) {149};
\draw[black,line width=0.51pt] (498.0,909.0) -- (465.5,891.7);
\draw[black,line width=0.51pt] (465.5,891.7) -- (465.9,876.9);
\draw[black,line width=0.51pt,fill=white] (498.0,909.0) circle (7.8);
\node[anchor=east,inner sep=0pt,font=\fontsize{5.4}{5.4}\selectfont,text=black] at (487.9,898.9) {48};
\node[anchor=east,inner sep=0pt,font=\fontsize{5.4}{5.4}\selectfont,text=black] at (487.9,927.4) {38};
\node[anchor=west,inner sep=0pt,font=\fontsize{5.4}{5.4}\selectfont,text=black] at (508.1,898.9) {126};
\draw[black,line width=0.51pt] (729.0,914.0) -- (729.0,950.8);
\draw[black,line width=0.51pt] (729.0,950.8) -- (715.8,957.4);
\draw[black,line width=0.51pt,fill=white] (729.0,914.0) circle (7.8);
\fill[black] (729.0,914.0) -- (729.0,906.2) -- (730.0,906.2) -- (731.0,906.4) -- (732.0,906.8) -- (732.9,907.2) -- (733.8,907.8) -- (734.5,908.5) -- (735.2,909.2) -- (735.8,910.1) -- (736.2,911.0) -- (736.6,912.0) -- (736.8,913.0) -- (736.8,914.0) -- (736.8,915.0) -- (736.6,916.0) -- (736.2,917.0) -- (735.8,917.9) -- (735.2,918.8) -- (734.5,919.5) -- (733.8,920.2) -- (732.9,920.8) -- (732.0,921.2) -- (731.0,921.6) -- (730.0,921.8) -- (729.0,921.8) -- cycle;
\node[anchor=east,inner sep=0pt,font=\fontsize{5.4}{5.4}\selectfont,text=black] at (718.9,903.9) {70};
\node[anchor=east,inner sep=0pt,font=\fontsize{5.4}{5.4}\selectfont,text=black] at (718.9,932.4) {59};
\node[anchor=west,inner sep=0pt,font=\fontsize{5.4}{5.4}\selectfont,text=black] at (739.1,903.9) {140};
\draw[black,line width=0.51pt] (390.0,946.0) -- (356.6,930.4);
\draw[black,line width=0.51pt] (356.6,930.4) -- (356.2,915.6);
\draw[black,line width=0.51pt,fill=white] (390.0,946.0) circle (7.8);
\node[anchor=east,inner sep=0pt,font=\fontsize{5.4}{5.4}\selectfont,text=black] at (379.9,935.9) {49};
\node[anchor=east,inner sep=0pt,font=\fontsize{5.4}{5.4}\selectfont,text=black] at (379.9,964.4) {38};
\node[anchor=west,inner sep=0pt,font=\fontsize{5.4}{5.4}\selectfont,text=black] at (400.1,935.9) {128};
\draw[black,line width=0.51pt] (603.0,945.0) -- (573.2,966.6);
\draw[black,line width=0.51pt] (573.2,966.6) -- (560.1,959.8);
\draw[black,line width=0.51pt,fill=white] (603.0,945.0) circle (7.8);
\fill[black] (603.0,945.0) -- (603.0,937.2) -- (604.0,937.2) -- (605.0,937.4) -- (606.0,937.8) -- (606.9,938.2) -- (607.8,938.8) -- (608.5,939.5) -- (609.2,940.2) -- (609.8,941.1) -- (610.2,942.0) -- (610.6,943.0) -- (610.8,944.0) -- (610.8,945.0) -- (610.8,946.0) -- (610.6,947.0) -- (610.2,948.0) -- (609.8,948.9) -- (609.2,949.8) -- (608.5,950.5) -- (607.8,951.2) -- (606.9,951.8) -- (606.0,952.2) -- (605.0,952.6) -- (604.0,952.8) -- (603.0,952.8) -- (602.0,952.8) -- (601.0,952.6) -- (600.0,952.2) -- (599.1,951.8) -- (598.2,951.2) -- (597.5,950.5) -- (596.8,949.8) -- (596.2,948.9) -- (595.8,948.0) -- (595.4,947.0) -- (595.2,946.0) -- (595.2,945.0) -- cycle;
\node[anchor=east,inner sep=0pt,font=\fontsize{5.4}{5.4}\selectfont,text=black] at (592.9,934.9) {66};
\node[anchor=east,inner sep=0pt,font=\fontsize{5.4}{5.4}\selectfont,text=black] at (592.9,963.4) {59};
\node[anchor=west,inner sep=0pt,font=\fontsize{5.4}{5.4}\selectfont,text=black] at (613.1,934.9) {114};
\draw[black,line width=0.51pt] (648.0,960.0) -- (631.3,992.8);
\draw[black,line width=0.51pt] (631.3,992.8) -- (616.5,992.7);
\draw[black,line width=0.51pt,fill=white] (648.0,960.0) circle (7.8);
\node[anchor=east,inner sep=0pt,font=\fontsize{5.4}{5.4}\selectfont,text=black] at (637.9,949.9) {66};
\node[anchor=east,inner sep=0pt,font=\fontsize{5.4}{5.4}\selectfont,text=black] at (637.9,978.4) {60};
\node[anchor=west,inner sep=0pt,font=\fontsize{5.4}{5.4}\selectfont,text=black] at (658.1,949.9) {124};
\draw[black,line width=0.51pt] (452.0,1009.0) -- (423.8,985.3);
\draw[black,line width=0.51pt] (423.8,985.3) -- (427.3,970.9);
\draw[black,line width=0.51pt,fill=white] (452.0,1009.0) circle (7.8);
\node[anchor=east,inner sep=0pt,font=\fontsize{5.4}{5.4}\selectfont,text=black] at (441.9,998.9) {50};
\node[anchor=east,inner sep=0pt,font=\fontsize{5.4}{5.4}\selectfont,text=black] at (441.9,1027.4) {38};
\node[anchor=west,inner sep=0pt,font=\fontsize{5.4}{5.4}\selectfont,text=black] at (462.1,998.9) {120};
\draw[black,line width=0.51pt] (507.0,1061.0) -- (471.5,1070.5);
\draw[black,line width=0.51pt] (471.5,1070.5) -- (461.6,1059.4);
\draw[black,line width=0.51pt,fill=white] (507.0,1061.0) circle (7.8);
\fill[black] (507.0,1061.0) circle (7.8);
\node[anchor=east,inner sep=0pt,font=\fontsize{5.4}{5.4}\selectfont,text=black] at (496.9,1050.9) {52};
\node[anchor=east,inner sep=0pt,font=\fontsize{5.4}{5.4}\selectfont,text=black] at (496.9,1079.4) {50};
\node[anchor=west,inner sep=0pt,font=\fontsize{5.4}{5.4}\selectfont,text=black] at (517.1,1050.9) {128};
\draw[black,line width=0.51pt] (613.0,1041.0) -- (606.0,1077.1);
\draw[black,line width=0.51pt] (606.0,1077.1) -- (591.7,1081.1);
\draw[black,line width=0.51pt,fill=white] (613.0,1041.0) circle (7.8);
\fill[black] (613.0,1041.0) circle (7.8);
\node[anchor=east,inner sep=0pt,font=\fontsize{5.4}{5.4}\selectfont,text=black] at (602.9,1030.9) {63};
\node[anchor=east,inner sep=0pt,font=\fontsize{5.4}{5.4}\selectfont,text=black] at (602.9,1059.4) {61};
\node[anchor=west,inner sep=0pt,font=\fontsize{5.4}{5.4}\selectfont,text=black] at (623.1,1030.9) {153};
\draw[black,line width=0.51pt] (725.0,1031.0) -- (710.6,1064.9);
\draw[black,line width=0.51pt] (710.6,1064.9) -- (695.8,1065.8);
\draw[black,line width=0.51pt,fill=white] (725.0,1031.0) circle (7.8);
\fill[black] (725.0,1031.0) -- (725.0,1023.2) -- (726.0,1023.2) -- (727.0,1023.4) -- (728.0,1023.8) -- (728.9,1024.2) -- (729.8,1024.8) -- (730.5,1025.5) -- (731.2,1026.2) -- (731.8,1027.1) -- (732.2,1028.0) -- (732.6,1029.0) -- (732.8,1030.0) -- (732.8,1031.0) -- cycle;
\node[anchor=east,inner sep=0pt,font=\fontsize{5.4}{5.4}\selectfont,text=black] at (714.9,1020.9) {69};
\node[anchor=east,inner sep=0pt,font=\fontsize{5.4}{5.4}\selectfont,text=black] at (714.9,1049.4) {61};
\node[anchor=west,inner sep=0pt,font=\fontsize{5.4}{5.4}\selectfont,text=black] at (735.1,1020.9) {171};
\draw[black,line width=0.51pt] (822.0,1006.0) -- (806.4,1039.4);
\draw[black,line width=0.51pt] (806.4,1039.4) -- (791.6,1039.8);
\draw[black,line width=0.51pt,fill=white] (822.0,1006.0) circle (7.8);
\node[anchor=east,inner sep=0pt,font=\fontsize{5.4}{5.4}\selectfont,text=black] at (811.9,995.9) {72};
\node[anchor=east,inner sep=0pt,font=\fontsize{5.4}{5.4}\selectfont,text=black] at (811.9,1024.4) {64};
\node[anchor=west,inner sep=0pt,font=\fontsize{5.4}{5.4}\selectfont,text=black] at (832.1,995.9) {194};
}
\caption{The 34 stations with all their values: temperature, dew point, pressure, weather, sky cover and wind. They are listed in Appendix~\ref{p3-app:stations}.}\label{p3-fig:fig10}
\end{figure}

\begin{questionbox}
Using the map of Figure~\ref{p3-fig:fig10}: draw isobars at 4\,mb intervals (e.g.\ 1004, 1008, 1012\,mb); label the low-pressure centre with an ``L''; draw the warm and cold fronts; outline the area where cloud cover exceeds 75\,\%; shade the areas receiving precipitation.
\end{questionbox}

Figure~\ref{p3-fig:fig10} shows the same 34 stations as Figure~\ref{p3-fig:fig8}, but with all the numbers: temperature, dew point, pressure and weather, as well as the sky cover and the wind. The recommended order is the order of the bullets; each step helps the next.

\subsection{Step 0: read the pressures}

Decode every pressure with the rule of Chapter~\ref{p3-ch:station} (a three-digit number below 500 gets 10 in front, otherwise 9), and write the full value next to each station.

\begin{figure}[ht]
\centering
\mapbase{0.0196cm}{\fill[black] (476.0,583.0) circle (2.6);
\node[anchor=west,inner sep=0pt,font=\fontsize{6.0}{6.0}\selectfont,text=black] at (481.0,581.0) {1011.0};
\node[anchor=east,inner sep=0pt,font=\fontsize{4.8}{4.8}\selectfont,text=gray] at (471.0,577.0) {1};
\fill[black] (626.0,593.0) circle (2.6);
\node[anchor=west,inner sep=0pt,font=\fontsize{6.0}{6.0}\selectfont,text=black] at (631.0,591.0) {1009.6};
\node[anchor=east,inner sep=0pt,font=\fontsize{4.8}{4.8}\selectfont,text=gray] at (621.0,587.0) {2};
\fill[black] (824.0,581.0) circle (2.6);
\node[anchor=west,inner sep=0pt,font=\fontsize{6.0}{6.0}\selectfont,text=black] at (829.0,579.0) {1015.4};
\node[anchor=east,inner sep=0pt,font=\fontsize{4.8}{4.8}\selectfont,text=gray] at (819.0,575.0) {3};
\fill[black] (465.0,654.0) circle (2.6);
\node[anchor=west,inner sep=0pt,font=\fontsize{6.0}{6.0}\selectfont,text=black] at (470.0,652.0) {1012.4};
\node[anchor=east,inner sep=0pt,font=\fontsize{4.8}{4.8}\selectfont,text=gray] at (460.0,648.0) {4};
\fill[black] (544.0,662.0) circle (2.6);
\node[anchor=west,inner sep=0pt,font=\fontsize{6.0}{6.0}\selectfont,text=black] at (549.0,660.0) {1008.8};
\node[anchor=east,inner sep=0pt,font=\fontsize{4.8}{4.8}\selectfont,text=gray] at (539.0,656.0) {5};
\fill[black] (738.0,624.0) circle (2.6);
\node[anchor=west,inner sep=0pt,font=\fontsize{6.0}{6.0}\selectfont,text=black] at (743.0,622.0) {1008.4};
\node[anchor=east,inner sep=0pt,font=\fontsize{4.8}{4.8}\selectfont,text=gray] at (733.0,618.0) {6};
\fill[black] (822.0,635.0) circle (2.6);
\node[anchor=west,inner sep=0pt,font=\fontsize{6.0}{6.0}\selectfont,text=black] at (827.0,633.0) {1013.2};
\node[anchor=east,inner sep=0pt,font=\fontsize{4.8}{4.8}\selectfont,text=gray] at (817.0,629.0) {7};
\fill[black] (672.0,681.0) circle (2.6);
\node[anchor=west,inner sep=0pt,font=\fontsize{6.0}{6.0}\selectfont,text=black] at (677.0,679.0) {1002.4};
\node[anchor=east,inner sep=0pt,font=\fontsize{4.8}{4.8}\selectfont,text=gray] at (667.0,675.0) {8};
\fill[black] (600.0,710.0) circle (2.6);
\node[anchor=west,inner sep=0pt,font=\fontsize{6.0}{6.0}\selectfont,text=black] at (605.0,708.0) {1002.6};
\node[anchor=east,inner sep=0pt,font=\fontsize{4.8}{4.8}\selectfont,text=gray] at (595.0,704.0) {9};
\fill[black] (824.0,683.0) circle (2.6);
\node[anchor=west,inner sep=0pt,font=\fontsize{6.0}{6.0}\selectfont,text=black] at (829.0,681.0) {1012.4};
\node[anchor=east,inner sep=0pt,font=\fontsize{4.8}{4.8}\selectfont,text=gray] at (819.0,677.0) {10};
\fill[black] (908.0,712.0) circle (2.6);
\node[anchor=west,inner sep=0pt,font=\fontsize{6.0}{6.0}\selectfont,text=black] at (913.0,710.0) {1016.3};
\node[anchor=east,inner sep=0pt,font=\fontsize{4.8}{4.8}\selectfont,text=gray] at (903.0,706.0) {11};
\fill[black] (419.0,722.0) circle (2.6);
\node[anchor=west,inner sep=0pt,font=\fontsize{6.0}{6.0}\selectfont,text=black] at (424.0,720.0) {1012.4};
\node[anchor=east,inner sep=0pt,font=\fontsize{4.8}{4.8}\selectfont,text=gray] at (414.0,716.0) {12};
\fill[black] (714.0,722.0) circle (2.6);
\node[anchor=west,inner sep=0pt,font=\fontsize{6.0}{6.0}\selectfont,text=black] at (719.0,720.0) {1002.9};
\node[anchor=east,inner sep=0pt,font=\fontsize{4.8}{4.8}\selectfont,text=gray] at (709.0,716.0) {13};
\fill[black] (818.0,742.0) circle (2.6);
\node[anchor=west,inner sep=0pt,font=\fontsize{6.0}{6.0}\selectfont,text=black] at (823.0,740.0) {1009.7};
\node[anchor=east,inner sep=0pt,font=\fontsize{4.8}{4.8}\selectfont,text=gray] at (813.0,736.0) {14};
\fill[black] (469.0,770.0) circle (2.6);
\node[anchor=west,inner sep=0pt,font=\fontsize{6.0}{6.0}\selectfont,text=black] at (474.0,768.0) {1010.3};
\node[anchor=east,inner sep=0pt,font=\fontsize{4.8}{4.8}\selectfont,text=gray] at (464.0,764.0) {15};
\fill[black] (570.0,768.0) circle (2.6);
\node[anchor=west,inner sep=0pt,font=\fontsize{6.0}{6.0}\selectfont,text=black] at (575.0,766.0) {1006.9};
\node[anchor=east,inner sep=0pt,font=\fontsize{4.8}{4.8}\selectfont,text=gray] at (565.0,762.0) {16};
\fill[black] (637.0,765.0) circle (2.6);
\node[anchor=west,inner sep=0pt,font=\fontsize{6.0}{6.0}\selectfont,text=black] at (642.0,763.0) {1004.0};
\node[anchor=east,inner sep=0pt,font=\fontsize{4.8}{4.8}\selectfont,text=gray] at (632.0,759.0) {17};
\fill[black] (704.0,785.0) circle (2.6);
\node[anchor=west,inner sep=0pt,font=\fontsize{6.0}{6.0}\selectfont,text=black] at (709.0,783.0) {1004.4};
\node[anchor=east,inner sep=0pt,font=\fontsize{4.8}{4.8}\selectfont,text=gray] at (699.0,779.0) {18};
\fill[black] (787.0,798.0) circle (2.6);
\node[anchor=west,inner sep=0pt,font=\fontsize{6.0}{6.0}\selectfont,text=black] at (792.0,796.0) {1009.9};
\node[anchor=east,inner sep=0pt,font=\fontsize{4.8}{4.8}\selectfont,text=gray] at (782.0,792.0) {19};
\fill[black] (828.0,818.0) circle (2.6);
\node[anchor=west,inner sep=0pt,font=\fontsize{6.0}{6.0}\selectfont,text=black] at (833.0,816.0) {1011.1};
\node[anchor=east,inner sep=0pt,font=\fontsize{4.8}{4.8}\selectfont,text=gray] at (823.0,812.0) {20};
\fill[black] (405.0,838.0) circle (2.6);
\node[anchor=west,inner sep=0pt,font=\fontsize{6.0}{6.0}\selectfont,text=black] at (410.0,836.0) {1012.6};
\node[anchor=east,inner sep=0pt,font=\fontsize{4.8}{4.8}\selectfont,text=gray] at (400.0,832.0) {21};
\fill[black] (496.0,831.0) circle (2.6);
\node[anchor=west,inner sep=0pt,font=\fontsize{6.0}{6.0}\selectfont,text=black] at (501.0,829.0) {1010.9};
\node[anchor=east,inner sep=0pt,font=\fontsize{4.8}{4.8}\selectfont,text=gray] at (491.0,825.0) {22};
\fill[black] (608.0,861.0) circle (2.6);
\node[anchor=west,inner sep=0pt,font=\fontsize{6.0}{6.0}\selectfont,text=black] at (613.0,859.0) {1007.5};
\node[anchor=east,inner sep=0pt,font=\fontsize{4.8}{4.8}\selectfont,text=gray] at (603.0,855.0) {23};
\fill[black] (828.0,885.0) circle (2.6);
\node[anchor=west,inner sep=0pt,font=\fontsize{6.0}{6.0}\selectfont,text=black] at (833.0,883.0) {1014.9};
\node[anchor=east,inner sep=0pt,font=\fontsize{4.8}{4.8}\selectfont,text=gray] at (823.0,879.0) {24};
\fill[black] (498.0,909.0) circle (2.6);
\node[anchor=west,inner sep=0pt,font=\fontsize{6.0}{6.0}\selectfont,text=black] at (503.0,907.0) {1012.6};
\node[anchor=east,inner sep=0pt,font=\fontsize{4.8}{4.8}\selectfont,text=gray] at (493.0,903.0) {25};
\fill[black] (729.0,914.0) circle (2.6);
\node[anchor=west,inner sep=0pt,font=\fontsize{6.0}{6.0}\selectfont,text=black] at (734.0,912.0) {1014.0};
\node[anchor=east,inner sep=0pt,font=\fontsize{4.8}{4.8}\selectfont,text=gray] at (724.0,908.0) {26};
\fill[black] (390.0,946.0) circle (2.6);
\node[anchor=west,inner sep=0pt,font=\fontsize{6.0}{6.0}\selectfont,text=black] at (395.0,944.0) {1012.8};
\node[anchor=east,inner sep=0pt,font=\fontsize{4.8}{4.8}\selectfont,text=gray] at (385.0,940.0) {27};
\fill[black] (603.0,945.0) circle (2.6);
\node[anchor=west,inner sep=0pt,font=\fontsize{6.0}{6.0}\selectfont,text=black] at (608.0,943.0) {1011.4};
\node[anchor=east,inner sep=0pt,font=\fontsize{4.8}{4.8}\selectfont,text=gray] at (598.0,939.0) {28};
\fill[black] (648.0,960.0) circle (2.6);
\node[anchor=west,inner sep=0pt,font=\fontsize{6.0}{6.0}\selectfont,text=black] at (653.0,958.0) {1012.4};
\node[anchor=east,inner sep=0pt,font=\fontsize{4.8}{4.8}\selectfont,text=gray] at (643.0,954.0) {29};
\fill[black] (452.0,1009.0) circle (2.6);
\node[anchor=west,inner sep=0pt,font=\fontsize{6.0}{6.0}\selectfont,text=black] at (457.0,1007.0) {1012.0};
\node[anchor=east,inner sep=0pt,font=\fontsize{4.8}{4.8}\selectfont,text=gray] at (447.0,1003.0) {30};
\fill[black] (507.0,1061.0) circle (2.6);
\node[anchor=west,inner sep=0pt,font=\fontsize{6.0}{6.0}\selectfont,text=black] at (512.0,1059.0) {1012.8};
\node[anchor=east,inner sep=0pt,font=\fontsize{4.8}{4.8}\selectfont,text=gray] at (502.0,1055.0) {31};
\fill[black] (613.0,1041.0) circle (2.6);
\node[anchor=west,inner sep=0pt,font=\fontsize{6.0}{6.0}\selectfont,text=black] at (618.0,1039.0) {1015.3};
\node[anchor=east,inner sep=0pt,font=\fontsize{4.8}{4.8}\selectfont,text=gray] at (608.0,1035.0) {32};
\fill[black] (725.0,1031.0) circle (2.6);
\node[anchor=west,inner sep=0pt,font=\fontsize{6.0}{6.0}\selectfont,text=black] at (730.0,1029.0) {1017.1};
\node[anchor=east,inner sep=0pt,font=\fontsize{4.8}{4.8}\selectfont,text=gray] at (720.0,1025.0) {33};
\fill[black] (822.0,1006.0) circle (2.6);
\node[anchor=west,inner sep=0pt,font=\fontsize{6.0}{6.0}\selectfont,text=black] at (827.0,1004.0) {1019.4};
\node[anchor=east,inner sep=0pt,font=\fontsize{4.8}{4.8}\selectfont,text=gray] at (817.0,1000.0) {34};
}
\caption{Step 0: the decoded sea-level pressures (hPa) at the 34 stations. The lowest values, 1002.4 and 1002.6, are in Iowa and Illinois (stations 8 and 9); the highest, 1019.4, is in Alabama (station 34).}\label{p3-fig:pressures}
\end{figure}

\subsection{Step 1: isobars every 4\,mb}

\textbf{The isobars are 1004, 1008, 1012, 1016 and 1020\,mb.} (The lowest pressure is 1002.4, so there is no 1000\,mb isobar. The highest, 1019.4 in Alabama, falls just short of 1020; the 1020 isobar lies a little beyond it, where the trend continues, and may be left out.)

\textbf{Method: interpolating.}
\begin{enumerate}[label=\textbf{\arabic*.},leftmargin=*,itemsep=3pt]
\item Look at two neighbouring stations and ask whether a multiple of 4 lies between their pressures.
\item If it does, find where it falls on the straight line between them \emph{in proportion}.
\item Mark that point. Do it for every pair of neighbours.
\item Join the points of equal value with a smooth curve. Do the same for the next value.
\end{enumerate}

\begin{workedbox}[Worked example: where does 1008\,mb cross between stations 18 and 19?]
Station 18 reads 1004.4\,hPa and its neighbour station 19 reads 1009.9\,hPa. The pressure changes by $1009.9-1004.4=5.5$\,hPa between them, and 1008 is $1008-1004.4=3.6$\,hPa above the first station:
\[
\text{fraction of the distance}=\frac{3.6}{5.5}=0.65 .
\]
The 1008\,mb isobar crosses the line joining the stations 65\,\% of the way from station 18 towards station 19.
\begin{center}
\begin{tikzpicture}
  \draw[thick] (0,0) -- (10,0);
  \fill (0,0) circle (3pt); \fill (10,0) circle (3pt);
  \node[anchor=north,font=\small,align=center] at (0,-0.2) {station 18\\ 1004.4}; \node[anchor=north,font=\small,align=center] at (10,-0.2) {station 19\\ 1009.9};
  \foreach \x/\v in {3.0/1006,6.5/1008}{\draw (\x*1,-0.12) -- (\x*1,0.12);}
  \fill[forange!90!black] (6.5,0) circle (3.5pt);
  \node[anchor=south,font=\small\bfseries,text=forange!60!black] at (6.5,0.15) {1008 at 65\,\%};
  \draw[<->,gray] (0,0.7) -- (6.5,0.7) node[midway,above,font=\scriptsize]{3.6 hPa}; \draw[<->,gray] (6.5,0.7) -- (10,0.7) node[midway,above,font=\scriptsize]{1.9 hPa};
\end{tikzpicture}
\end{center}
\end{workedbox}

\begin{table}[ht]
\centering
\caption{More interpolations to practise on.}
\begin{tabular}{@{}lllll@{}}
\toprule
\textbf{Stations} & \textbf{Pressures (hPa)} & \textbf{Isobar} & \textbf{Calculation} & \textbf{Fraction}\\
\midrule
9 and 5 & 1002.6, 1008.8 & 1004 & $(1004-1002.6)/6.2$ & 23\,\%\\
9 and 5 & 1002.6, 1008.8 & 1008 & $(1008-1002.6)/6.2$ & 87\,\%\\
16 and 5 & 1006.9, 1008.8 & 1008 & $(1008-1006.9)/1.9$ & 58\,\%\\
20 and 24 & 1011.1, 1014.9 & 1012 & $(1012-1011.1)/3.8$ & 24\,\%\\
\bottomrule
\end{tabular}
\end{table}

\begin{keybox}[Rules for drawing isobars]
\begin{itemize}
\item Isobars are smooth curves. Each connects points of exactly the same pressure.
\item Isobars \textbf{never cross} and never branch.
\item They never end inside the map: they either close into a loop (round a low or a high) or leave the map edge.
\item The pressure rises or falls in steps of 4\,mb from one isobar to the next; if the next isobar is much further away the gradient is gentle there.
\item Label each isobar with its value.
\end{itemize}
\end{keybox}

\begin{figure}[ht]
\centering
\mapbase{0.0196cm}{\draw[isobar] (637.1,658.0) -- (638.6,657.7) -- (641.5,657.1) -- (644.5,656.6) -- (647.5,656.1) -- (650.5,655.7) -- (653.5,655.5) -- (656.5,655.3) -- (659.5,655.2) -- (662.5,655.2) -- (665.5,655.3) -- (668.5,655.5) -- (671.5,655.8) -- (674.5,656.3) -- (677.5,656.8) -- (679.8,657.3) -- (681.4,657.8) -- (682.9,658.2) -- (684.3,658.6) -- (686.5,659.4) -- (689.5,660.6) -- (692.4,661.9) -- (695.2,663.3) -- (696.7,664.1) -- (696.9,664.2) -- (698.5,665.2) -- (701.5,667.2) -- (703.6,668.6) -- (704.8,669.5) -- (706.3,670.8) -- (708.1,672.4) -- (709.7,673.9) -- (711.2,675.3) -- (712.7,676.9) -- (714.2,678.6) -- (715.5,680.1) -- (716.6,681.4) -- (718.1,683.4) -- (720.0,686.1) -- (721.1,687.6) -- (721.2,687.9) -- (722.2,689.5) -- (723.8,692.5) -- (725.2,695.3) -- (726.4,697.8) -- (727.1,699.3) -- (727.3,699.8) -- (727.9,701.5) -- (729.0,704.5) -- (729.8,707.5) -- (730.6,710.5) -- (731.2,713.5) -- (731.6,716.5) -- (731.9,719.5) -- (732.0,722.5) -- (732.0,725.5) -- (731.7,728.5) -- (731.4,731.5) -- (730.8,734.5) -- (730.1,737.5) -- (729.2,740.5) -- (728.3,743.2) -- (727.4,745.5) -- (726.9,747.0) -- (726.6,747.7) -- (725.6,749.5) -- (724.1,752.5) -- (722.7,755.0) -- (721.6,756.9) -- (720.6,758.4) -- (719.9,759.5) -- (718.5,761.5) -- (716.3,764.5) -- (715.1,766.0) -- (715.0,766.1) -- (713.6,767.6) -- (710.9,770.5) -- (709.4,772.1) -- (709.1,772.4) -- (707.5,773.8) -- (704.5,776.2) -- (702.7,777.5) -- (702.2,777.8) -- (700.7,778.6) -- (698.2,779.9) -- (695.5,780.9) -- (692.5,781.8) -- (689.5,782.5) -- (686.5,783.0) -- (683.5,783.2) -- (680.5,783.3) -- (677.5,783.1) -- (674.5,782.8) -- (671.5,782.3) -- (668.5,781.5) -- (665.5,780.6) -- (662.5,779.5) -- (660.5,778.8) -- (659.4,778.3) -- (657.9,777.5) -- (656.0,776.6) -- (653.5,775.2) -- (650.5,773.4) -- (648.8,772.4) -- (648.4,772.1) -- (646.9,771.2) -- (644.3,769.6) -- (641.9,768.1) -- (639.8,766.7) -- (638.3,765.8) -- (637.4,765.3) -- (635.5,764.4) -- (632.5,763.0) -- (629.8,761.7) -- (627.3,760.6) -- (625.8,759.9) -- (625.3,759.7) -- (623.5,758.9) -- (620.5,757.4) -- (617.8,756.0) -- (615.5,754.7) -- (614.0,753.8) -- (613.3,753.5) -- (611.5,752.3) -- (608.5,750.4) -- (606.5,749.1) -- (605.6,748.4) -- (604.1,747.2) -- (602.0,745.6) -- (600.2,744.1) -- (598.7,742.7) -- (597.2,741.3) -- (595.7,739.9) -- (594.3,738.4) -- (592.9,736.8) -- (591.4,735.0) -- (589.8,732.9) -- (588.7,731.4) -- (588.1,730.5) -- (587.0,728.5) -- (585.5,725.5) -- (584.3,722.7) -- (583.4,720.1) -- (582.9,718.6) -- (582.8,718.2) -- (582.6,716.5) -- (582.2,713.5) -- (582.1,710.5) -- (582.2,707.5) -- (582.4,705.1) -- (582.8,703.4) -- (583.1,701.9) -- (583.4,700.6) -- (584.2,698.5) -- (585.4,695.5) -- (586.7,692.6) -- (588.2,689.9) -- (589.1,688.4) -- (589.3,688.1) -- (590.5,686.5) -- (592.9,683.5) -- (594.3,681.7) -- (594.8,681.2) -- (596.3,679.7) -- (598.8,677.2) -- (600.3,675.8) -- (600.8,675.4) -- (602.5,674.1) -- (605.5,671.9) -- (607.3,670.6) -- (607.9,670.2) -- (609.4,669.3) -- (611.8,667.8) -- (614.5,666.4) -- (617.5,664.9) -- (619.1,664.1) -- (619.2,664.0) -- (620.7,663.4) -- (623.6,662.3) -- (626.5,661.2) -- (629.5,660.2) -- (632.5,659.3) -- (635.5,658.4) -- (637.0,658.0) -- (637.1,658.0) -- (637.1,658.0);
\draw[isobar] (644.3,610.0) -- (645.5,609.8) -- (647.8,609.4) -- (650.5,609.0) -- (653.5,608.5) -- (656.5,608.1) -- (659.5,607.8) -- (662.5,607.5) -- (665.5,607.3) -- (668.5,607.2) -- (671.5,607.2) -- (674.5,607.2) -- (677.5,607.3) -- (680.5,607.4) -- (683.5,607.7) -- (686.5,608.0) -- (689.5,608.5) -- (692.5,609.0) -- (695.5,609.6) -- (697.1,609.9) -- (697.3,610.0) -- (698.8,610.4) -- (701.6,611.1) -- (704.5,612.0) -- (707.5,613.0) -- (710.5,614.1) -- (713.5,615.3) -- (715.0,616.0) -- (715.1,616.0) -- (716.6,616.7) -- (719.5,618.2) -- (722.3,619.7) -- (725.0,621.2) -- (726.5,622.1) -- (726.8,622.3) -- (728.5,623.4) -- (731.5,625.5) -- (733.5,626.9) -- (734.5,627.6) -- (736.0,628.8) -- (738.0,630.5) -- (739.7,632.0) -- (741.1,633.3) -- (742.6,634.8) -- (744.2,636.4) -- (745.6,637.9) -- (746.9,639.3) -- (748.4,641.0) -- (750.1,643.0) -- (751.4,644.5) -- (752.2,645.5) -- (753.6,647.5) -- (755.8,650.5) -- (756.9,652.1) -- (757.0,652.2) -- (757.9,653.7) -- (759.7,656.6) -- (761.2,659.1) -- (762.4,661.2) -- (763.2,662.7) -- (763.7,663.6) -- (764.7,665.5) -- (766.2,668.5) -- (767.5,671.1) -- (768.5,673.3) -- (769.2,674.8) -- (769.5,675.6) -- (770.3,677.5) -- (771.6,680.5) -- (772.8,683.5) -- (774.0,686.5) -- (774.7,688.3) -- (774.9,688.9) -- (775.5,690.4) -- (776.4,692.8) -- (777.3,695.5) -- (778.4,698.5) -- (779.5,701.5) -- (780.5,704.5) -- (781.0,706.0) -- (781.0,706.0) -- (781.5,707.5) -- (782.4,710.5) -- (783.4,713.5) -- (784.2,716.5) -- (785.0,719.5) -- (785.8,722.5) -- (786.4,725.2) -- (786.8,727.6) -- (787.1,729.1) -- (787.2,729.7) -- (787.4,731.5) -- (787.7,734.5) -- (787.8,737.5) -- (787.7,740.5) -- (787.5,743.0) -- (787.2,745.0) -- (786.9,746.5) -- (786.7,747.5) -- (786.2,749.5) -- (785.2,752.5) -- (784.0,755.5) -- (782.5,758.5) -- (781.6,760.3) -- (781.2,761.0) -- (780.3,762.5) -- (778.8,764.8) -- (777.3,767.1) -- (775.8,769.2) -- (774.7,770.7) -- (774.1,771.6) -- (772.6,773.5) -- (770.3,776.5) -- (769.1,778.1) -- (769.0,778.2) -- (767.9,779.7) -- (765.7,782.6) -- (764.2,784.5) -- (763.4,785.5) -- (762.2,787.0) -- (760.7,789.0) -- (759.2,790.9) -- (757.7,792.8) -- (756.6,794.3) -- (755.7,795.4) -- (754.2,797.3) -- (752.1,799.9) -- (750.8,801.4) -- (750.5,801.8) -- (749.0,803.5) -- (746.3,806.4) -- (745.0,807.9) -- (745.0,808.0) -- (743.5,809.5) -- (740.5,812.5) -- (737.5,815.4) -- (734.5,818.1) -- (732.8,819.6) -- (732.5,819.9) -- (731.0,821.1) -- (728.3,823.2) -- (726.4,824.7) -- (725.3,825.6) -- (723.8,826.7) -- (721.9,828.0) -- (719.8,829.5) -- (717.3,831.2) -- (715.8,832.2) -- (715.3,832.5) -- (713.5,833.7) -- (710.5,835.6) -- (708.4,836.9) -- (707.1,837.6) -- (705.6,838.5) -- (703.9,839.5) -- (701.5,840.9) -- (698.5,842.6) -- (696.7,843.6) -- (696.2,843.9) -- (694.7,844.7) -- (692.2,846.1) -- (689.5,847.5) -- (686.5,849.2) -- (685.0,850.0) -- (684.9,850.0) -- (683.4,850.8) -- (680.5,852.3) -- (677.6,853.8) -- (674.8,855.3) -- (673.3,856.1) -- (673.1,856.2) -- (671.5,857.0) -- (668.5,858.5) -- (665.6,859.9) -- (662.8,861.3) -- (661.3,862.0) -- (661.1,862.1) -- (659.5,862.9) -- (656.5,864.3) -- (653.5,865.7) -- (650.5,867.0) -- (648.8,867.7) -- (648.3,867.9) -- (646.8,868.5) -- (644.3,869.5) -- (641.5,870.5) -- (638.5,871.4) -- (635.5,872.3) -- (632.5,873.1) -- (630.2,873.6) -- (628.5,873.9) -- (627.0,874.1) -- (625.7,874.3) -- (623.5,874.5) -- (620.5,874.6) -- (617.5,874.5) -- (614.5,874.3) -- (612.8,874.1) -- (612.5,874.0) -- (611.0,873.7) -- (608.3,873.0) -- (605.5,871.9) -- (602.5,870.5) -- (600.3,869.4) -- (598.9,868.5) -- (597.4,867.4) -- (595.8,866.1) -- (594.1,864.6) -- (592.3,862.9) -- (590.8,861.3) -- (589.6,859.9) -- (588.4,858.4) -- (587.2,856.8) -- (585.7,854.7) -- (583.9,852.0) -- (582.9,850.5) -- (582.7,850.2) -- (581.8,848.5) -- (580.1,845.5) -- (578.7,842.8) -- (577.6,840.5) -- (576.8,839.0) -- (576.5,838.3) -- (575.7,836.5) -- (574.4,833.5) -- (573.1,830.5) -- (571.9,827.5) -- (571.2,825.8) -- (571.1,825.5) -- (570.5,824.0) -- (569.4,821.3) -- (568.3,818.5) -- (567.1,815.5) -- (566.2,813.0) -- (565.4,811.0) -- (564.8,809.5) -- (564.4,808.5) -- (563.6,806.5) -- (562.4,803.5) -- (561.2,800.5) -- (560.0,797.5) -- (559.3,795.8) -- (559.1,795.4) -- (558.4,793.9) -- (557.2,791.3) -- (555.9,788.5) -- (554.5,785.5) -- (553.5,783.7) -- (553.2,783.0) -- (552.4,781.5) -- (551.1,779.2) -- (549.6,776.5) -- (547.9,773.6) -- (547.0,772.1) -- (546.9,772.0) -- (546.0,770.5) -- (544.2,767.5) -- (542.8,764.9) -- (541.6,762.8) -- (540.8,761.3) -- (540.3,760.4) -- (539.4,758.5) -- (538.0,755.5) -- (536.8,752.5) -- (535.7,749.5) -- (535.1,747.8) -- (535.0,747.5) -- (534.7,746.0) -- (534.0,743.3) -- (533.4,740.5) -- (532.9,737.5) -- (532.5,734.5) -- (532.3,731.5) -- (532.2,728.5) -- (532.2,725.5) -- (532.3,722.5) -- (532.6,719.5) -- (533.0,716.5) -- (533.4,713.5) -- (534.0,710.6) -- (534.7,707.7) -- (535.0,706.2) -- (535.0,706.1) -- (535.5,704.5) -- (536.3,701.5) -- (537.3,698.5) -- (538.3,695.5) -- (539.4,692.7) -- (540.5,690.1) -- (541.1,688.6) -- (541.3,688.2) -- (542.0,686.5) -- (543.4,683.5) -- (544.9,680.7) -- (546.3,678.0) -- (547.1,676.5) -- (547.2,676.2) -- (548.2,674.5) -- (549.8,671.5) -- (551.2,669.0) -- (552.4,667.0) -- (553.3,665.5) -- (553.8,664.5) -- (555.0,662.5) -- (556.8,659.5) -- (558.0,657.5) -- (558.7,656.4) -- (559.6,654.9) -- (560.9,653.0) -- (562.4,650.8) -- (564.1,648.4) -- (565.2,646.9) -- (565.7,646.3) -- (567.2,644.5) -- (569.7,641.6) -- (571.0,640.1) -- (571.1,640.0) -- (572.6,638.6) -- (575.5,635.9) -- (577.2,634.4) -- (577.5,634.1) -- (579.0,633.0) -- (581.7,631.0) -- (583.8,629.5) -- (585.3,628.5) -- (586.8,627.6) -- (588.3,626.7) -- (590.5,625.5) -- (593.5,624.0) -- (595.7,622.9) -- (597.2,622.3) -- (598.7,621.7) -- (600.2,621.1) -- (602.5,620.2) -- (605.5,619.2) -- (608.5,618.2) -- (611.5,617.3) -- (613.9,616.7) -- (615.6,616.2) -- (617.1,615.8) -- (618.4,615.5) -- (620.5,615.0) -- (623.5,614.3) -- (626.5,613.6) -- (629.5,613.0) -- (632.5,612.3) -- (635.5,611.7) -- (638.5,611.1) -- (641.5,610.5) -- (643.3,610.2) -- (644.0,610.1) -- (644.3,610.0);
\draw[isobar] (415.0,520.0) -- (414.0,521.5) -- (412.2,524.5) -- (410.7,527.3) -- (409.6,530.0) -- (408.9,531.5) -- (408.8,531.8) -- (408.4,533.5) -- (407.7,536.5) -- (407.2,539.5) -- (406.9,542.5) -- (406.9,545.5) -- (407.0,548.5) -- (407.3,551.5) -- (407.8,554.5) -- (408.3,556.9) -- (408.8,558.6) -- (409.2,560.1) -- (409.5,561.4) -- (410.2,563.5) -- (411.4,566.5) -- (412.6,569.5) -- (414.1,572.5) -- (414.9,574.1) -- (415.0,574.2) -- (415.8,575.7) -- (417.5,578.6) -- (419.0,580.9) -- (420.3,582.8) -- (421.4,584.3) -- (422.1,585.4) -- (423.6,587.4) -- (425.9,590.1) -- (427.1,591.6) -- (427.3,591.9) -- (428.8,593.5) -- (431.6,596.5) -- (433.0,598.0) -- (433.0,598.0) -- (434.5,599.5) -- (437.5,602.3) -- (439.1,603.8) -- (439.2,603.9) -- (440.7,605.2) -- (443.6,607.7) -- (445.3,609.2) -- (446.0,609.7) -- (447.5,610.9) -- (449.8,612.8) -- (451.7,614.3) -- (453.1,615.4) -- (454.6,616.6) -- (456.2,617.9) -- (458.0,619.4) -- (460.0,621.1) -- (461.5,622.4) -- (462.5,623.3) -- (464.1,624.8) -- (466.3,626.9) -- (467.8,628.5) -- (468.6,629.4) -- (469.8,630.9) -- (471.3,633.0) -- (472.8,635.3) -- (474.3,637.8) -- (475.1,639.3) -- (475.3,639.8) -- (475.9,641.5) -- (476.8,644.5) -- (477.4,647.5) -- (477.7,650.5) -- (477.7,653.5) -- (477.4,656.5) -- (476.9,659.5) -- (476.2,662.5) -- (475.6,664.5) -- (475.2,665.6) -- (474.6,667.1) -- (473.8,669.0) -- (472.5,671.5) -- (471.0,674.5) -- (469.9,676.5) -- (469.3,677.5) -- (468.3,679.0) -- (467.0,681.0) -- (465.5,683.2) -- (463.8,685.6) -- (462.8,687.1) -- (462.3,687.7) -- (461.0,689.5) -- (458.7,692.5) -- (457.4,694.2) -- (457.1,694.6) -- (455.9,696.1) -- (453.8,698.7) -- (452.3,700.6) -- (451.4,701.7) -- (450.2,703.2) -- (448.7,705.1) -- (447.2,707.0) -- (445.7,708.9) -- (444.6,710.4) -- (443.7,711.5) -- (442.2,713.4) -- (440.1,716.3) -- (438.9,717.8) -- (438.8,717.9) -- (437.7,719.5) -- (435.5,722.5) -- (434.1,724.5) -- (433.4,725.6) -- (432.3,727.1) -- (431.0,729.0) -- (429.5,731.4) -- (427.8,734.1) -- (426.9,735.6) -- (426.8,735.9) -- (425.8,737.5) -- (424.1,740.5) -- (422.7,743.2) -- (421.6,745.7) -- (420.9,747.2) -- (420.6,747.7) -- (419.8,749.5) -- (418.6,752.5) -- (417.6,755.5) -- (416.7,758.5) -- (415.9,761.4) -- (415.3,764.2) -- (415.0,765.7) -- (414.9,765.9) -- (414.7,767.5) -- (414.3,770.5) -- (414.1,773.5) -- (414.0,776.5) -- (414.0,779.5) -- (414.2,782.5) -- (414.5,785.1) -- (414.8,787.3) -- (415.1,788.8) -- (415.2,789.6) -- (415.6,791.5) -- (416.4,794.5) -- (417.3,797.5) -- (418.4,800.5) -- (419.5,803.0) -- (420.5,805.1) -- (421.2,806.6) -- (421.7,807.5) -- (422.8,809.5) -- (424.6,812.5) -- (425.9,814.5) -- (426.6,815.5) -- (427.7,817.0) -- (429.2,819.0) -- (430.7,820.9) -- (432.2,822.6) -- (433.6,824.1) -- (434.7,825.4) -- (436.2,826.9) -- (438.1,828.7) -- (439.6,830.2) -- (440.9,831.4) -- (442.4,832.7) -- (444.1,834.2) -- (445.9,835.7) -- (447.8,837.2) -- (449.3,838.5) -- (450.4,839.4) -- (452.4,840.9) -- (455.2,843.0) -- (456.7,844.1) -- (456.9,844.2) -- (458.5,845.4) -- (461.5,847.6) -- (463.4,849.0) -- (464.3,849.7) -- (465.8,850.8) -- (467.9,852.4) -- (469.9,853.9) -- (471.8,855.3) -- (473.3,856.4) -- (474.4,857.3) -- (476.3,858.8) -- (479.0,860.9) -- (480.5,862.1) -- (480.8,862.4) -- (482.5,863.9) -- (485.5,866.5) -- (487.1,867.8) -- (487.2,867.9) -- (488.7,869.3) -- (491.6,872.0) -- (493.1,873.5) -- (493.4,873.8) -- (494.9,875.4) -- (497.6,878.1) -- (499.1,879.6) -- (499.3,879.9) -- (500.8,881.5) -- (503.5,884.5) -- (504.9,886.0) -- (505.0,886.1) -- (506.2,887.6) -- (508.5,890.5) -- (510.0,892.4) -- (510.7,893.3) -- (511.8,894.8) -- (513.3,896.9) -- (514.8,899.1) -- (516.3,901.4) -- (517.2,902.9) -- (517.7,903.6) -- (518.9,905.5) -- (520.7,908.5) -- (521.9,910.6) -- (522.6,911.9) -- (523.5,913.4) -- (524.6,915.1) -- (526.0,917.5) -- (527.7,920.5) -- (528.7,922.2) -- (528.9,922.5) -- (529.9,924.0) -- (531.6,926.7) -- (533.1,928.9) -- (534.4,930.7) -- (535.5,932.2) -- (536.4,933.4) -- (537.9,935.2) -- (540.0,937.5) -- (541.3,939.0) -- (542.0,939.7) -- (543.5,941.0) -- (545.8,943.1) -- (547.6,944.6) -- (548.9,945.5) -- (550.4,946.6) -- (552.1,947.7) -- (554.5,949.0) -- (557.5,950.6) -- (559.4,951.5) -- (560.1,951.8) -- (561.6,952.4) -- (563.9,953.3) -- (566.5,954.2) -- (569.5,955.1) -- (572.5,955.9) -- (575.5,956.6) -- (578.5,957.2) -- (581.5,957.7) -- (583.0,958.0) -- (583.1,958.0) -- (584.6,958.2) -- (587.5,958.5) -- (590.5,958.8) -- (593.5,959.1) -- (596.5,959.2) -- (599.5,959.3) -- (602.5,959.4) -- (605.5,959.4) -- (608.5,959.4) -- (611.5,959.3) -- (614.5,959.1) -- (617.5,958.9) -- (620.5,958.6) -- (623.5,958.2) -- (625.0,958.0) -- (625.1,958.0) -- (626.6,957.7) -- (629.5,957.2) -- (632.5,956.5) -- (635.5,955.7) -- (638.5,954.8) -- (641.5,953.6) -- (643.5,952.8) -- (644.6,952.3) -- (646.1,951.5) -- (648.0,950.5) -- (650.5,949.0) -- (653.5,947.0) -- (655.0,946.0) -- (655.0,946.0) -- (656.5,944.8) -- (659.5,942.5) -- (661.4,941.0) -- (662.1,940.3) -- (663.6,939.0) -- (665.9,937.0) -- (667.5,935.5) -- (668.6,934.5) -- (670.1,933.1) -- (672.0,931.2) -- (673.6,929.7) -- (674.7,928.6) -- (676.2,927.1) -- (678.1,925.2) -- (679.6,923.7) -- (680.7,922.6) -- (682.2,921.1) -- (684.1,919.2) -- (685.5,917.7) -- (686.6,916.6) -- (688.1,915.0) -- (690.0,913.1) -- (691.5,911.6) -- (692.6,910.5) -- (694.1,909.1) -- (696.0,907.3) -- (697.6,905.8) -- (698.8,904.6) -- (700.3,903.2) -- (702.1,901.6) -- (703.8,900.1) -- (705.4,898.7) -- (706.9,897.4) -- (708.3,896.3) -- (710.2,894.8) -- (712.5,892.9) -- (714.0,891.8) -- (714.7,891.3) -- (716.5,890.0) -- (719.5,887.8) -- (721.3,886.6) -- (721.8,886.2) -- (723.3,885.2) -- (725.8,883.6) -- (728.1,882.1) -- (730.3,880.7) -- (731.8,879.8) -- (732.6,879.3) -- (734.5,878.1) -- (737.5,876.3) -- (739.6,875.1) -- (740.9,874.4) -- (742.4,873.5) -- (744.1,872.5) -- (746.5,871.1) -- (749.5,869.4) -- (751.2,868.4) -- (751.7,868.1) -- (753.2,867.3) -- (755.7,865.9) -- (758.5,864.4) -- (761.4,862.8) -- (762.9,862.0) -- (763.0,861.9) -- (764.5,861.1) -- (767.5,859.5) -- (770.3,858.0) -- (773.0,856.7) -- (774.5,855.9) -- (774.8,855.7) -- (776.5,854.9) -- (779.5,853.4) -- (782.5,852.0) -- (785.5,850.7) -- (787.0,850.0) -- (787.1,850.0) -- (788.6,849.4) -- (791.5,848.2) -- (794.5,847.1) -- (797.5,846.1) -- (800.3,845.2) -- (802.9,844.4) -- (804.4,843.9) -- (804.8,843.8) -- (806.5,843.3) -- (809.5,842.5) -- (812.5,841.7) -- (815.5,840.8) -- (818.5,840.0) -- (821.5,839.0) -- (823.4,838.4) -- (824.2,838.1) -- (825.7,837.6) -- (827.9,836.7) -- (830.5,835.5) -- (833.5,834.0) -- (835.5,832.9) -- (836.4,832.3) -- (837.9,831.2) -- (840.0,829.7) -- (841.8,828.2) -- (843.4,826.7) -- (844.9,825.2) -- (846.3,823.7) -- (847.6,822.2) -- (848.7,820.7) -- (850.2,818.5) -- (852.0,815.5) -- (852.9,813.9) -- (853.0,813.8) -- (853.6,812.3) -- (854.9,809.4) -- (856.0,806.5) -- (857.0,803.5) -- (857.8,800.5) -- (858.5,797.5) -- (858.9,795.7) -- (859.0,795.1) -- (859.2,793.6) -- (859.5,791.2) -- (859.8,788.5) -- (860.0,785.5) -- (860.1,782.5) -- (860.1,779.5) -- (860.0,776.5) -- (859.8,773.5) -- (859.5,770.5) -- (859.2,767.5) -- (859.0,766.0) -- (859.0,765.9) -- (858.7,764.4) -- (858.2,761.5) -- (857.7,758.5) -- (857.0,755.5) -- (856.2,752.5) -- (855.4,749.5) -- (854.5,746.6) -- (853.5,743.8) -- (853.0,742.3) -- (852.9,742.1) -- (852.3,740.5) -- (851.1,737.5) -- (849.7,734.5) -- (848.3,731.5) -- (847.5,729.7) -- (847.2,729.1) -- (846.4,727.6) -- (845.1,725.2) -- (843.6,722.6) -- (841.9,719.7) -- (841.0,718.2) -- (840.9,718.1) -- (839.8,716.5) -- (837.9,713.5) -- (836.4,711.3) -- (835.5,709.9) -- (834.5,708.4) -- (833.4,706.8) -- (831.9,704.6) -- (830.0,701.9) -- (828.9,700.4) -- (828.7,700.1) -- (827.7,698.5) -- (825.7,695.5) -- (824.3,693.3) -- (823.4,691.9) -- (822.6,690.4) -- (821.7,688.8) -- (820.6,686.5) -- (819.2,683.5) -- (818.1,680.9) -- (817.4,678.7) -- (816.9,677.2) -- (816.6,676.4) -- (816.1,674.5) -- (815.3,671.5) -- (814.5,668.5) -- (813.7,665.5) -- (813.0,662.5) -- (812.2,659.5) -- (811.6,657.2) -- (811.2,655.7) -- (810.8,654.2) -- (810.4,652.7) -- (809.7,650.5) -- (808.8,647.5) -- (807.8,644.5) -- (806.6,641.5) -- (805.8,639.4) -- (805.3,638.2) -- (804.6,636.7) -- (803.7,634.9) -- (802.5,632.5) -- (800.9,629.5) -- (799.8,627.5) -- (799.3,626.6) -- (798.4,625.1) -- (797.1,623.0) -- (795.6,620.7) -- (793.9,618.2) -- (792.8,616.7) -- (792.5,616.2) -- (791.2,614.5) -- (788.9,611.5) -- (787.6,609.8) -- (787.2,609.3) -- (786.0,607.8) -- (783.9,605.3) -- (782.4,603.5) -- (781.5,602.4) -- (780.1,600.9) -- (778.3,599.0) -- (776.8,597.4) -- (775.6,596.1) -- (774.1,594.6) -- (772.4,592.9) -- (770.9,591.4) -- (769.6,590.2) -- (768.0,588.7) -- (766.1,586.9) -- (764.6,585.5) -- (763.5,584.6) -- (761.7,583.1) -- (759.1,581.0) -- (757.6,579.8) -- (757.2,579.5) -- (755.5,578.2) -- (752.5,576.0) -- (750.7,574.7) -- (750.0,574.2) -- (748.5,573.2) -- (746.2,571.6) -- (743.8,570.1) -- (741.5,568.7) -- (740.0,567.8) -- (739.3,567.4) -- (737.5,566.4) -- (734.5,564.7) -- (732.0,563.4) -- (730.1,562.5) -- (728.6,561.7) -- (727.5,561.2) -- (725.5,560.3) -- (722.5,559.0) -- (719.5,557.7) -- (716.6,556.6) -- (715.1,556.0) -- (715.0,556.0) -- (713.5,555.4) -- (710.5,554.4) -- (707.5,553.4) -- (704.5,552.5) -- (701.5,551.6) -- (698.5,550.9) -- (696.5,550.4) -- (695.4,550.1) -- (693.9,549.8) -- (692.0,549.3) -- (689.5,548.8) -- (686.5,548.2) -- (683.5,547.7) -- (680.5,547.2) -- (677.5,546.8) -- (674.5,546.4) -- (671.5,546.1) -- (668.5,545.8) -- (665.5,545.5) -- (662.5,545.3) -- (659.5,545.1) -- (656.5,544.9) -- (653.5,544.8) -- (650.5,544.6) -- (647.5,544.5) -- (644.5,544.4) -- (641.5,544.2) -- (638.5,544.1) -- (636.8,544.0) -- (636.5,544.0) -- (635.0,543.9) -- (632.3,543.8) -- (629.5,543.6) -- (626.5,543.4) -- (623.5,543.2) -- (620.5,542.9) -- (617.5,542.6) -- (614.5,542.3) -- (611.5,541.9) -- (608.5,541.5) -- (605.5,541.0) -- (602.5,540.5) -- (599.5,540.0) -- (596.5,539.4) -- (593.7,538.9) -- (591.0,538.3) -- (589.5,538.0) -- (589.2,537.9) -- (587.5,537.5) -- (584.5,536.7) -- (581.5,535.9) -- (578.5,535.1) -- (575.5,534.2) -- (572.5,533.3) -- (570.4,532.6) -- (569.2,532.2) -- (567.7,531.7) -- (565.9,531.1) -- (563.5,530.2) -- (560.5,529.1) -- (557.5,528.0) -- (554.5,526.8) -- (552.8,526.2) -- (552.5,526.1) -- (551.0,525.4) -- (548.3,524.3) -- (545.5,523.2) -- (542.5,521.9) -- (540.3,520.9) -- (538.8,520.3) -- (538.1,520.0);
\draw[isobar] (307.0,1012.5) -- (308.5,1012.9) -- (311.5,1013.7) -- (314.5,1014.5) -- (317.5,1015.2) -- (320.5,1015.8) -- (323.5,1016.4) -- (326.5,1016.9) -- (329.5,1017.4) -- (331.8,1017.7) -- (333.5,1017.9) -- (335.0,1018.1) -- (336.3,1018.2) -- (338.5,1018.5) -- (341.5,1018.7) -- (344.5,1018.9) -- (347.5,1019.1) -- (350.5,1019.2) -- (353.5,1019.2) -- (356.5,1019.2) -- (359.5,1019.1) -- (362.5,1019.0) -- (365.5,1018.8) -- (368.5,1018.5) -- (371.5,1018.2) -- (373.2,1018.1) -- (373.6,1018.0) -- (375.1,1017.8) -- (377.7,1017.5) -- (380.5,1017.1) -- (383.5,1016.7) -- (386.5,1016.2) -- (389.5,1015.7) -- (392.5,1015.1) -- (395.5,1014.5) -- (398.5,1013.8) -- (401.5,1013.2) -- (403.9,1012.6) -- (405.6,1012.2) -- (407.1,1011.9) -- (408.4,1011.6) -- (410.5,1011.1) -- (413.5,1010.5) -- (416.5,1009.9) -- (419.5,1009.2) -- (422.5,1008.7) -- (425.5,1008.1) -- (428.5,1007.6) -- (431.5,1007.2) -- (434.5,1006.9) -- (437.5,1006.7) -- (440.5,1006.7) -- (443.5,1007.0) -- (446.5,1007.5) -- (449.5,1008.4) -- (452.1,1009.7) -- (454.2,1011.2) -- (455.7,1012.5) -- (456.6,1013.6) -- (457.5,1015.1) -- (458.6,1017.0) -- (459.6,1019.5) -- (460.5,1022.5) -- (461.1,1025.5) -- (461.4,1028.5) -- (461.4,1031.5) -- (461.1,1034.5) -- (460.6,1037.5) -- (460.0,1040.5) -- (459.1,1043.5) -- (458.1,1046.5) -- (457.4,1048.3) -- (457.1,1049.0) -- (456.4,1050.5) -- (455.3,1052.8) -- (453.8,1055.5) -- (452.1,1058.5) -- (451.2,1060.1) -- (451.1,1060.3) -- (450.0,1061.8) -- (448.1,1064.6) -- (446.6,1066.7) -- (445.5,1068.2) -- (444.4,1069.7) -- (443.2,1071.2) -- (441.7,1073.1) -- (439.9,1075.3) -- (438.7,1076.8) -- (438.0,1077.6) -- (436.5,1079.4) -- (434.2,1082.1) -- (432.9,1083.6) -- (432.6,1083.9) -- (431.2,1085.5) -- (428.5,1088.5) -- (427.1,1090.0) -- (427.0,1090.1) -- (425.7,1091.6) -- (423.0,1094.5) -- (421.5,1096.2) -- (421.2,1096.5) -- (419.8,1098.0) -- (417.3,1100.7) -- (415.8,1102.3) -- (415.3,1102.9) -- (413.9,1104.4) -- (411.7,1106.8) -- (410.2,1108.5) -- (409.4,1109.4) -- (409.0,1109.8);
\draw[isobar] (784.3,520.0) -- (785.0,520.6) -- (786.3,521.9) -- (787.9,523.4) -- (789.6,525.1) -- (791.1,526.6) -- (792.4,527.9) -- (793.8,529.4) -- (795.5,531.1) -- (797.0,532.7) -- (798.3,534.1) -- (799.7,535.6) -- (801.1,537.2) -- (802.6,538.9) -- (804.2,540.6) -- (805.5,542.1) -- (806.6,543.4) -- (808.1,545.2) -- (810.0,547.5) -- (811.3,549.0) -- (811.8,549.7) -- (813.3,551.5) -- (815.7,554.5) -- (816.9,556.0) -- (817.0,556.1) -- (818.1,557.6) -- (820.4,560.5) -- (821.9,562.5) -- (822.6,563.5) -- (823.7,565.0) -- (825.2,567.0) -- (826.7,569.0) -- (828.2,571.1) -- (829.3,572.6) -- (830.0,573.5) -- (831.5,575.5) -- (833.6,578.5) -- (834.8,580.1) -- (834.9,580.3) -- (836.0,581.8) -- (838.0,584.6) -- (839.5,586.7) -- (840.5,588.1) -- (841.5,589.6) -- (842.6,591.2) -- (844.1,593.5) -- (846.0,596.4) -- (847.0,597.9) -- (847.1,598.0) -- (848.0,599.5) -- (849.8,602.5) -- (851.3,605.0) -- (852.4,607.0) -- (853.3,608.5) -- (853.9,609.5) -- (854.9,611.5) -- (856.5,614.5) -- (857.7,616.9) -- (858.6,618.7) -- (859.3,620.2) -- (859.9,621.4) -- (860.9,623.5) -- (862.2,626.5) -- (863.4,629.2) -- (864.5,631.7) -- (865.1,633.2) -- (865.4,633.7) -- (866.1,635.5) -- (867.4,638.5) -- (868.6,641.5) -- (869.8,644.5) -- (870.6,646.4) -- (870.9,647.1) -- (871.5,648.6) -- (872.5,650.9) -- (873.6,653.5) -- (874.9,656.5) -- (875.9,658.8) -- (876.6,660.5) -- (877.3,662.0) -- (877.9,663.3) -- (879.0,665.5) -- (880.4,668.5) -- (881.6,671.0) -- (882.5,672.9) -- (883.3,674.4) -- (883.8,675.5) -- (884.9,677.5) -- (886.4,680.5) -- (887.6,682.9) -- (888.5,684.6) -- (889.3,686.1) -- (890.0,687.4) -- (891.1,689.5) -- (892.7,692.5) -- (893.9,694.7) -- (894.6,696.1) -- (895.4,697.6) -- (896.2,699.2) -- (897.4,701.5) -- (898.9,704.5) -- (900.0,706.7) -- (900.7,708.0) -- (901.4,709.5) -- (902.3,711.2) -- (903.4,713.5) -- (904.9,716.5) -- (906.0,718.7) -- (906.7,720.2) -- (907.4,721.7) -- (908.1,723.2) -- (909.2,725.5) -- (910.5,728.5) -- (911.6,731.1) -- (912.5,733.3) -- (913.2,734.8) -- (913.5,735.6) -- (914.2,737.5) -- (915.4,740.5) -- (916.5,743.5) -- (917.5,746.5) -- (918.3,748.8) -- (918.8,750.5) -- (919.2,752.0) -- (919.6,753.3) -- (920.2,755.5) -- (921.0,758.5) -- (921.7,761.5) -- (922.3,764.5) -- (922.9,767.5) -- (923.4,770.5) -- (923.9,773.5) -- (924.2,776.5) -- (924.5,779.5) -- (924.8,782.5) -- (924.9,784.9) -- (925.0,786.6) -- (925.0,788.1) -- (925.1,789.4) -- (925.1,791.2) -- (925.0,793.6) -- (925.0,795.1) -- (925.0,795.7) -- (924.9,797.5) -- (924.7,800.5) -- (924.4,803.5) -- (924.0,806.5) -- (923.6,809.5) -- (923.1,812.5) -- (922.5,815.5) -- (921.8,818.5) -- (921.0,821.5) -- (920.1,824.5) -- (919.5,826.5) -- (919.2,827.4) -- (918.6,828.9) -- (917.9,831.0) -- (916.9,833.5) -- (915.6,836.5) -- (914.5,839.0) -- (913.5,841.0) -- (912.7,842.5) -- (912.2,843.5) -- (911.1,845.5) -- (909.3,848.5) -- (908.1,850.5) -- (907.4,851.6) -- (906.3,853.1) -- (905.0,855.0) -- (903.5,857.1) -- (901.8,859.2) -- (900.6,860.7) -- (899.9,861.6) -- (898.4,863.3) -- (896.1,865.8) -- (894.8,867.3) -- (894.3,867.8) -- (892.8,869.2) -- (890.3,871.6) -- (888.7,873.1) -- (888.0,873.7) -- (886.5,875.0) -- (884.2,876.9) -- (882.3,878.4) -- (880.8,879.5) -- (879.3,880.6) -- (877.8,881.7) -- (875.5,883.2) -- (872.6,885.1) -- (871.1,886.0) -- (871.0,886.1) -- (869.5,887.0) -- (866.5,888.7) -- (863.8,890.2) -- (861.5,891.4) -- (860.0,892.2) -- (859.3,892.5) -- (857.5,893.4) -- (854.5,894.8) -- (851.6,896.2) -- (848.8,897.4) -- (847.3,898.0) -- (847.1,898.1) -- (845.5,898.8) -- (842.5,900.1) -- (839.5,901.3) -- (836.5,902.5) -- (834.5,903.4) -- (833.4,903.8) -- (831.9,904.4) -- (830.0,905.2) -- (827.5,906.2) -- (824.5,907.4) -- (821.8,908.5) -- (819.4,909.5) -- (817.9,910.1) -- (817.3,910.4) -- (815.5,911.2) -- (812.5,912.4) -- (809.5,913.7) -- (806.5,915.0) -- (804.8,915.7) -- (804.3,915.9) -- (802.8,916.6) -- (800.3,917.7) -- (797.5,919.0) -- (794.5,920.3) -- (792.5,921.3) -- (791.4,921.8) -- (789.9,922.5) -- (788.0,923.4) -- (785.5,924.7) -- (782.5,926.2) -- (780.5,927.2) -- (779.5,927.7) -- (778.0,928.6) -- (776.0,929.7) -- (773.5,931.2) -- (770.5,933.0) -- (769.0,933.9) -- (768.9,934.0) -- (767.4,935.0) -- (764.5,937.1) -- (762.4,938.6) -- (761.2,939.5) -- (759.7,940.7) -- (757.9,942.2) -- (756.2,943.7) -- (754.6,945.2) -- (753.1,946.7) -- (751.7,948.1) -- (750.3,949.6) -- (748.8,951.2) -- (747.3,952.9) -- (745.8,954.7) -- (744.5,956.2) -- (743.6,957.4) -- (742.1,959.4) -- (740.0,962.1) -- (738.9,963.6) -- (738.7,963.9) -- (737.6,965.5) -- (735.5,968.5) -- (734.1,970.5) -- (733.4,971.6) -- (732.3,973.1) -- (731.0,975.0) -- (729.5,977.2) -- (727.8,979.7) -- (726.8,981.2) -- (726.5,981.7) -- (725.3,983.5) -- (723.2,986.5) -- (721.9,988.4) -- (721.3,989.2) -- (720.2,990.7) -- (718.6,992.9) -- (717.1,994.9) -- (715.7,996.8) -- (714.5,998.3) -- (713.6,999.4) -- (712.1,1001.2) -- (710.0,1003.7) -- (708.8,1005.2) -- (708.3,1005.7) -- (706.8,1007.3) -- (704.3,1010.0) -- (702.8,1011.5) -- (702.5,1011.8) -- (701.0,1013.2) -- (698.3,1015.7) -- (696.7,1017.2) -- (696.1,1017.7) -- (694.6,1019.0) -- (692.2,1021.0) -- (690.3,1022.5) -- (688.9,1023.5) -- (687.4,1024.6) -- (685.8,1025.8) -- (683.6,1027.3) -- (680.9,1029.1) -- (679.4,1030.1) -- (679.1,1030.3) -- (677.5,1031.3) -- (674.5,1033.1) -- (672.1,1034.5) -- (670.3,1035.5) -- (668.8,1036.4) -- (667.6,1037.1) -- (665.5,1038.2) -- (662.5,1039.8) -- (660.3,1041.0) -- (658.9,1041.7) -- (657.4,1042.4) -- (655.8,1043.3) -- (653.5,1044.5) -- (650.5,1046.0) -- (648.4,1047.1) -- (647.2,1047.7) -- (645.7,1048.5) -- (643.9,1049.5) -- (641.5,1050.8) -- (638.5,1052.6) -- (636.8,1053.6) -- (636.3,1053.9) -- (634.8,1054.9) -- (632.3,1056.7) -- (630.2,1058.2) -- (628.7,1059.4) -- (627.2,1060.7) -- (625.7,1062.0) -- (624.2,1063.5) -- (622.7,1065.2) -- (621.2,1066.9) -- (619.7,1068.7) -- (618.6,1070.2) -- (617.7,1071.4) -- (616.2,1073.5) -- (614.3,1076.5) -- (613.2,1078.1) -- (613.1,1078.4) -- (612.2,1079.9) -- (610.7,1082.6) -- (609.2,1085.5) -- (607.7,1088.4) -- (607.0,1089.9) -- (606.9,1090.0) -- (606.2,1091.5) -- (604.9,1094.5) -- (603.5,1097.5) -- (602.2,1100.5) -- (601.4,1102.4) -- (601.1,1103.1) -- (600.5,1104.6) -- (599.6,1106.9) -- (598.6,1109.5) -- (597.4,1112.5) -- (596.3,1115.2) -- (595.4,1117.5) -- (594.9,1119.0) -- (594.6,1119.7) -- (593.9,1121.5) -- (592.8,1124.5) -- (591.7,1127.5) -- (590.6,1130.5) -- (589.7,1132.7) -- (589.2,1134.0) -- (588.7,1135.5) -- (588.1,1137.2) -- (587.2,1139.5) -- (586.1,1142.5) -- (585.6,1144.0);
\draw[isobar] (882.8,520.0) -- (882.9,520.1) -- (883.0,520.2) -- (883.9,521.7) -- (885.8,524.6) -- (887.3,526.9) -- (888.4,528.7) -- (889.4,530.2) -- (890.2,531.4) -- (891.5,533.5) -- (893.4,536.5) -- (894.5,538.3) -- (894.8,538.8) -- (895.8,540.3) -- (897.3,542.8) -- (898.8,545.2) -- (900.3,547.7) -- (901.2,549.2) -- (901.5,549.7) -- (902.6,551.5) -- (904.3,554.5) -- (905.7,556.8) -- (906.6,558.3) -- (907.4,559.8) -- (908.3,561.3) -- (909.6,563.5) -- (911.3,566.5) -- (912.4,568.4) -- (912.8,569.1) -- (913.6,570.6) -- (914.9,572.9) -- (916.4,575.5) -- (918.0,578.5) -- (918.9,580.1) -- (919.0,580.2) -- (919.8,581.7) -- (921.3,584.6) -- (922.8,587.4) -- (924.3,590.1) -- (925.1,591.6) -- (925.2,591.9) -- (926.1,593.5) -- (927.6,596.5) -- (929.1,599.3) -- (930.4,601.8) -- (931.0,603.0);
\draw[isobar] (931.5,604.0) -- (932.3,605.5) -- (933.8,608.5) -- (935.2,611.2) -- (936.4,613.7) -- (937.1,615.2) -- (937.4,615.7) -- (938.3,617.5) -- (939.8,620.5) -- (941.1,623.3) -- (942.4,625.8) -- (943.1,627.3) -- (943.3,627.8) -- (944.2,629.5) -- (945.6,632.5) -- (947.0,635.4) -- (948.3,638.1) -- (949.1,639.6) -- (949.2,639.9) -- (950.0,641.5) -- (951.4,644.5) -- (952.8,647.5) -- (954.2,650.5) -- (955.0,652.0) -- (955.0,652.1) -- (955.7,653.6) -- (957.1,656.5) -- (958.5,659.5) -- (959.9,662.5) -- (960.7,664.2) -- (960.9,664.7) -- (961.6,666.2) -- (962.7,668.7) -- (964.0,671.5) -- (965.3,674.5) -- (966.2,676.6) -- (966.7,677.7) -- (967.4,679.2) -- (968.2,681.1) -- (969.3,683.5) -- (970.6,686.5) -- (971.6,689.1) -- (972.5,691.2) -- (973.2,692.7) -- (973.5,693.6) -- (974.3,695.5) -- (975.5,698.5) -- (976.7,701.5) -- (977.9,704.5) -- (978.6,706.4) -- (978.9,707.1) -- (979.4,708.6) -- (980.3,710.9) -- (981.2,713.5) -- (982.2,716.5) -- (983.3,719.5) -- (984.3,722.5) -- (984.8,724.2) -- (984.9,724.6) -- (985.4,726.1) -- (986.2,728.7) -- (987.0,731.5) -- (987.9,734.5) -- (988.7,737.5) -- (989.5,740.5) -- (990.2,743.1) -- (990.7,745.4) -- (991.1,746.9) -- (991.3,747.6) -- (991.7,749.5) -- (992.3,752.5) -- (992.9,755.5) -- (993.5,758.5) -- (994.0,761.5) -- (994.5,764.5) -- (994.9,767.5) -- (995.3,770.5) -- (995.7,773.5) -- (996.0,776.5) -- (996.1,778.0);
\draw[isobar] (758.2,1132.0) -- (759.2,1130.5) -- (761.0,1127.5) -- (762.2,1125.6) -- (762.7,1124.8) -- (763.7,1123.3) -- (765.0,1121.1) -- (766.5,1118.7) -- (768.2,1116.1) -- (769.1,1114.6) -- (769.4,1114.2) -- (770.5,1112.5) -- (772.4,1109.5) -- (773.8,1107.4) -- (774.6,1106.1) -- (775.6,1104.6) -- (776.7,1102.9) -- (778.2,1100.6) -- (780.1,1097.7) -- (781.0,1096.2) -- (781.1,1096.1) -- (782.1,1094.5) -- (784.0,1091.5) -- (785.5,1089.2) -- (786.5,1087.7) -- (787.5,1086.2) -- (788.4,1084.7) -- (789.8,1082.5) -- (791.8,1079.5) -- (792.8,1077.9) -- (792.9,1077.7) -- (793.9,1076.2) -- (795.7,1073.4) -- (797.2,1071.1) -- (798.4,1069.2) -- (799.4,1067.7) -- (800.1,1066.6) -- (801.5,1064.5) -- (803.4,1061.5) -- (804.6,1059.8) -- (804.9,1059.3) -- (805.8,1057.8) -- (807.5,1055.3) -- (809.0,1053.0) -- (810.3,1051.1) -- (811.4,1049.6) -- (812.1,1048.5) -- (813.5,1046.5) -- (815.6,1043.5) -- (816.8,1041.9) -- (816.9,1041.7) -- (818.0,1040.2) -- (820.1,1037.4) -- (821.6,1035.4) -- (822.5,1034.1) -- (823.7,1032.6) -- (825.1,1030.9) -- (826.6,1029.0) -- (828.2,1027.1) -- (829.4,1025.6) -- (830.3,1024.5) -- (831.8,1022.8) -- (833.9,1020.5) -- (835.3,1019.0) -- (836.0,1018.3) -- (837.5,1016.8) -- (839.8,1014.5) -- (841.3,1013.0) -- (842.0,1012.3) -- (843.5,1011.0) -- (845.8,1008.9) -- (847.5,1007.4) -- (848.5,1006.5) -- (850.0,1005.2) -- (852.0,1003.5) -- (853.8,1002.0) -- (855.5,1000.7) -- (857.0,999.5) -- (858.3,998.4) -- (860.3,996.9) -- (862.8,995.0) -- (864.3,993.8) -- (864.8,993.5) -- (866.5,992.2) -- (869.5,990.0) -- (871.3,988.7) -- (871.9,988.2) -- (873.4,987.1) -- (875.8,985.4) -- (877.9,983.9) -- (879.6,982.6) -- (881.1,981.6) -- (882.4,980.7) -- (884.5,979.2) -- (887.4,977.1) -- (888.9,976.0) -- (889.0,975.9) -- (890.5,974.8) -- (893.5,972.7) -- (895.5,971.2) -- (896.6,970.4) -- (898.1,969.3) -- (900.0,967.9) -- (902.0,966.4) -- (904.1,964.8) -- (905.6,963.7) -- (906.5,963.0) -- (908.5,961.5) -- (911.4,959.2) -- (912.9,958.0) -- (913.0,957.9) -- (914.5,956.7) -- (917.5,954.2) -- (919.0,953.0);
\node[isolabel,rotate=-10] at (673.0,782.6) {1004};
\node[isolabel,rotate=25] at (661.0,862.2) {1008};
\node[isolabel,rotate=45] at (685.0,918.2) {1012};
\node[isolabel,rotate=35] at (877.0,882.2) {1016};
\node[isolabel,rotate=51] at (825.8,1030.0) {1020};
\fill[black] (476.0,583.0) circle (2.6);
\node[anchor=west,inner sep=0pt,font=\fontsize{6.0}{6.0}\selectfont,text=black] at (481.0,581.0) {1011.0};
\node[anchor=east,inner sep=0pt,font=\fontsize{4.8}{4.8}\selectfont,text=gray] at (471.0,577.0) {1};
\fill[black] (626.0,593.0) circle (2.6);
\node[anchor=west,inner sep=0pt,font=\fontsize{6.0}{6.0}\selectfont,text=black] at (631.0,591.0) {1009.6};
\node[anchor=east,inner sep=0pt,font=\fontsize{4.8}{4.8}\selectfont,text=gray] at (621.0,587.0) {2};
\fill[black] (824.0,581.0) circle (2.6);
\node[anchor=west,inner sep=0pt,font=\fontsize{6.0}{6.0}\selectfont,text=black] at (829.0,579.0) {1015.4};
\node[anchor=east,inner sep=0pt,font=\fontsize{4.8}{4.8}\selectfont,text=gray] at (819.0,575.0) {3};
\fill[black] (465.0,654.0) circle (2.6);
\node[anchor=west,inner sep=0pt,font=\fontsize{6.0}{6.0}\selectfont,text=black] at (470.0,652.0) {1012.4};
\node[anchor=east,inner sep=0pt,font=\fontsize{4.8}{4.8}\selectfont,text=gray] at (460.0,648.0) {4};
\fill[black] (544.0,662.0) circle (2.6);
\node[anchor=west,inner sep=0pt,font=\fontsize{6.0}{6.0}\selectfont,text=black] at (549.0,660.0) {1008.8};
\node[anchor=east,inner sep=0pt,font=\fontsize{4.8}{4.8}\selectfont,text=gray] at (539.0,656.0) {5};
\fill[black] (738.0,624.0) circle (2.6);
\node[anchor=west,inner sep=0pt,font=\fontsize{6.0}{6.0}\selectfont,text=black] at (743.0,622.0) {1008.4};
\node[anchor=east,inner sep=0pt,font=\fontsize{4.8}{4.8}\selectfont,text=gray] at (733.0,618.0) {6};
\fill[black] (822.0,635.0) circle (2.6);
\node[anchor=west,inner sep=0pt,font=\fontsize{6.0}{6.0}\selectfont,text=black] at (827.0,633.0) {1013.2};
\node[anchor=east,inner sep=0pt,font=\fontsize{4.8}{4.8}\selectfont,text=gray] at (817.0,629.0) {7};
\fill[black] (672.0,681.0) circle (2.6);
\node[anchor=west,inner sep=0pt,font=\fontsize{6.0}{6.0}\selectfont,text=black] at (677.0,679.0) {1002.4};
\node[anchor=east,inner sep=0pt,font=\fontsize{4.8}{4.8}\selectfont,text=gray] at (667.0,675.0) {8};
\fill[black] (600.0,710.0) circle (2.6);
\node[anchor=west,inner sep=0pt,font=\fontsize{6.0}{6.0}\selectfont,text=black] at (605.0,708.0) {1002.6};
\node[anchor=east,inner sep=0pt,font=\fontsize{4.8}{4.8}\selectfont,text=gray] at (595.0,704.0) {9};
\fill[black] (824.0,683.0) circle (2.6);
\node[anchor=west,inner sep=0pt,font=\fontsize{6.0}{6.0}\selectfont,text=black] at (829.0,681.0) {1012.4};
\node[anchor=east,inner sep=0pt,font=\fontsize{4.8}{4.8}\selectfont,text=gray] at (819.0,677.0) {10};
\fill[black] (908.0,712.0) circle (2.6);
\node[anchor=west,inner sep=0pt,font=\fontsize{6.0}{6.0}\selectfont,text=black] at (913.0,710.0) {1016.3};
\node[anchor=east,inner sep=0pt,font=\fontsize{4.8}{4.8}\selectfont,text=gray] at (903.0,706.0) {11};
\fill[black] (419.0,722.0) circle (2.6);
\node[anchor=west,inner sep=0pt,font=\fontsize{6.0}{6.0}\selectfont,text=black] at (424.0,720.0) {1012.4};
\node[anchor=east,inner sep=0pt,font=\fontsize{4.8}{4.8}\selectfont,text=gray] at (414.0,716.0) {12};
\fill[black] (714.0,722.0) circle (2.6);
\node[anchor=west,inner sep=0pt,font=\fontsize{6.0}{6.0}\selectfont,text=black] at (719.0,720.0) {1002.9};
\node[anchor=east,inner sep=0pt,font=\fontsize{4.8}{4.8}\selectfont,text=gray] at (709.0,716.0) {13};
\fill[black] (818.0,742.0) circle (2.6);
\node[anchor=west,inner sep=0pt,font=\fontsize{6.0}{6.0}\selectfont,text=black] at (823.0,740.0) {1009.7};
\node[anchor=east,inner sep=0pt,font=\fontsize{4.8}{4.8}\selectfont,text=gray] at (813.0,736.0) {14};
\fill[black] (469.0,770.0) circle (2.6);
\node[anchor=west,inner sep=0pt,font=\fontsize{6.0}{6.0}\selectfont,text=black] at (474.0,768.0) {1010.3};
\node[anchor=east,inner sep=0pt,font=\fontsize{4.8}{4.8}\selectfont,text=gray] at (464.0,764.0) {15};
\fill[black] (570.0,768.0) circle (2.6);
\node[anchor=west,inner sep=0pt,font=\fontsize{6.0}{6.0}\selectfont,text=black] at (575.0,766.0) {1006.9};
\node[anchor=east,inner sep=0pt,font=\fontsize{4.8}{4.8}\selectfont,text=gray] at (565.0,762.0) {16};
\fill[black] (637.0,765.0) circle (2.6);
\node[anchor=west,inner sep=0pt,font=\fontsize{6.0}{6.0}\selectfont,text=black] at (642.0,763.0) {1004.0};
\node[anchor=east,inner sep=0pt,font=\fontsize{4.8}{4.8}\selectfont,text=gray] at (632.0,759.0) {17};
\fill[black] (704.0,785.0) circle (2.6);
\node[anchor=west,inner sep=0pt,font=\fontsize{6.0}{6.0}\selectfont,text=black] at (709.0,783.0) {1004.4};
\node[anchor=east,inner sep=0pt,font=\fontsize{4.8}{4.8}\selectfont,text=gray] at (699.0,779.0) {18};
\fill[black] (787.0,798.0) circle (2.6);
\node[anchor=west,inner sep=0pt,font=\fontsize{6.0}{6.0}\selectfont,text=black] at (792.0,796.0) {1009.9};
\node[anchor=east,inner sep=0pt,font=\fontsize{4.8}{4.8}\selectfont,text=gray] at (782.0,792.0) {19};
\fill[black] (828.0,818.0) circle (2.6);
\node[anchor=west,inner sep=0pt,font=\fontsize{6.0}{6.0}\selectfont,text=black] at (833.0,816.0) {1011.1};
\node[anchor=east,inner sep=0pt,font=\fontsize{4.8}{4.8}\selectfont,text=gray] at (823.0,812.0) {20};
\fill[black] (405.0,838.0) circle (2.6);
\node[anchor=west,inner sep=0pt,font=\fontsize{6.0}{6.0}\selectfont,text=black] at (410.0,836.0) {1012.6};
\node[anchor=east,inner sep=0pt,font=\fontsize{4.8}{4.8}\selectfont,text=gray] at (400.0,832.0) {21};
\fill[black] (496.0,831.0) circle (2.6);
\node[anchor=west,inner sep=0pt,font=\fontsize{6.0}{6.0}\selectfont,text=black] at (501.0,829.0) {1010.9};
\node[anchor=east,inner sep=0pt,font=\fontsize{4.8}{4.8}\selectfont,text=gray] at (491.0,825.0) {22};
\fill[black] (608.0,861.0) circle (2.6);
\node[anchor=west,inner sep=0pt,font=\fontsize{6.0}{6.0}\selectfont,text=black] at (613.0,859.0) {1007.5};
\node[anchor=east,inner sep=0pt,font=\fontsize{4.8}{4.8}\selectfont,text=gray] at (603.0,855.0) {23};
\fill[black] (828.0,885.0) circle (2.6);
\node[anchor=west,inner sep=0pt,font=\fontsize{6.0}{6.0}\selectfont,text=black] at (833.0,883.0) {1014.9};
\node[anchor=east,inner sep=0pt,font=\fontsize{4.8}{4.8}\selectfont,text=gray] at (823.0,879.0) {24};
\fill[black] (498.0,909.0) circle (2.6);
\node[anchor=west,inner sep=0pt,font=\fontsize{6.0}{6.0}\selectfont,text=black] at (503.0,907.0) {1012.6};
\node[anchor=east,inner sep=0pt,font=\fontsize{4.8}{4.8}\selectfont,text=gray] at (493.0,903.0) {25};
\fill[black] (729.0,914.0) circle (2.6);
\node[anchor=west,inner sep=0pt,font=\fontsize{6.0}{6.0}\selectfont,text=black] at (734.0,912.0) {1014.0};
\node[anchor=east,inner sep=0pt,font=\fontsize{4.8}{4.8}\selectfont,text=gray] at (724.0,908.0) {26};
\fill[black] (390.0,946.0) circle (2.6);
\node[anchor=west,inner sep=0pt,font=\fontsize{6.0}{6.0}\selectfont,text=black] at (395.0,944.0) {1012.8};
\node[anchor=east,inner sep=0pt,font=\fontsize{4.8}{4.8}\selectfont,text=gray] at (385.0,940.0) {27};
\fill[black] (603.0,945.0) circle (2.6);
\node[anchor=west,inner sep=0pt,font=\fontsize{6.0}{6.0}\selectfont,text=black] at (608.0,943.0) {1011.4};
\node[anchor=east,inner sep=0pt,font=\fontsize{4.8}{4.8}\selectfont,text=gray] at (598.0,939.0) {28};
\fill[black] (648.0,960.0) circle (2.6);
\node[anchor=west,inner sep=0pt,font=\fontsize{6.0}{6.0}\selectfont,text=black] at (653.0,958.0) {1012.4};
\node[anchor=east,inner sep=0pt,font=\fontsize{4.8}{4.8}\selectfont,text=gray] at (643.0,954.0) {29};
\fill[black] (452.0,1009.0) circle (2.6);
\node[anchor=west,inner sep=0pt,font=\fontsize{6.0}{6.0}\selectfont,text=black] at (457.0,1007.0) {1012.0};
\node[anchor=east,inner sep=0pt,font=\fontsize{4.8}{4.8}\selectfont,text=gray] at (447.0,1003.0) {30};
\fill[black] (507.0,1061.0) circle (2.6);
\node[anchor=west,inner sep=0pt,font=\fontsize{6.0}{6.0}\selectfont,text=black] at (512.0,1059.0) {1012.8};
\node[anchor=east,inner sep=0pt,font=\fontsize{4.8}{4.8}\selectfont,text=gray] at (502.0,1055.0) {31};
\fill[black] (613.0,1041.0) circle (2.6);
\node[anchor=west,inner sep=0pt,font=\fontsize{6.0}{6.0}\selectfont,text=black] at (618.0,1039.0) {1015.3};
\node[anchor=east,inner sep=0pt,font=\fontsize{4.8}{4.8}\selectfont,text=gray] at (608.0,1035.0) {32};
\fill[black] (725.0,1031.0) circle (2.6);
\node[anchor=west,inner sep=0pt,font=\fontsize{6.0}{6.0}\selectfont,text=black] at (730.0,1029.0) {1017.1};
\node[anchor=east,inner sep=0pt,font=\fontsize{4.8}{4.8}\selectfont,text=gray] at (720.0,1025.0) {33};
\fill[black] (822.0,1006.0) circle (2.6);
\node[anchor=west,inner sep=0pt,font=\fontsize{6.0}{6.0}\selectfont,text=black] at (827.0,1004.0) {1019.4};
\node[anchor=east,inner sep=0pt,font=\fontsize{4.8}{4.8}\selectfont,text=gray] at (817.0,1000.0) {34};
}
\caption{Step 1: isobars at 4\,mb intervals through the station pressures. The 1004\,mb isobar closes round the lowest pressures; the 1008 and 1012\,mb isobars follow it outwards; 1016 and 1020\,mb lie to the east and south-east. (The curves here were drawn as a smooth surface through all 34 pressures.)}\label{p3-fig:isobars}
\end{figure}

\subsection{Step 2: the L}

The low-pressure centre is where the isobars close round the lowest pressure. The two lowest pressures on the map are 1002.4\,hPa (station 8, northern Illinois) and 1002.6\,hPa (station 9, Iowa). The 1004\,mb isobar encloses them, together with stations 13 and 17 (1002.9 and 1004.0). Write the \textbf{L} between stations 8 and 9, near the centre of the closed isobar, at about 1002\,hPa.

\begin{answerbox}[Answer: the L]
The L goes in northern Illinois and eastern Iowa, between stations 8 and 9, in the middle of the innermost isobar. It is also the centre of the swirl of streamlines you drew in Question~13.
\end{answerbox}

\subsection{Step 3: the fronts}

A front lies where the temperature, the dew point and the wind direction change rapidly. Compare pairs of neighbouring stations: if two neighbours differ strongly in all three, a front runs between them.

\begin{figure}[ht]
\centering
\mapbase{0.0196cm}{\fill[tempcol!5!tempcold] (476.0,583.0) circle (6.5);
\node[anchor=west,inner sep=0pt,font=\fontsize{6.0}{6.0}\selectfont,text=black] at (484.0,582.0) {34};
\fill[tempcol!32!tempcold] (626.0,593.0) circle (6.5);
\node[anchor=west,inner sep=0pt,font=\fontsize{6.0}{6.0}\selectfont,text=black] at (634.0,592.0) {45};
\fill[tempcol!17!tempcold] (824.0,581.0) circle (6.5);
\node[anchor=west,inner sep=0pt,font=\fontsize{6.0}{6.0}\selectfont,text=black] at (832.0,580.0) {39};
\fill[tempcol!10!tempcold] (465.0,654.0) circle (6.5);
\node[anchor=west,inner sep=0pt,font=\fontsize{6.0}{6.0}\selectfont,text=black] at (473.0,653.0) {36};
\fill[tempcol!20!tempcold] (544.0,662.0) circle (6.5);
\node[anchor=west,inner sep=0pt,font=\fontsize{6.0}{6.0}\selectfont,text=black] at (552.0,661.0) {40};
\fill[tempcol!17!tempcold] (738.0,624.0) circle (6.5);
\node[anchor=west,inner sep=0pt,font=\fontsize{6.0}{6.0}\selectfont,text=black] at (746.0,623.0) {39};
\fill[tempcol!27!tempcold] (822.0,635.0) circle (6.5);
\node[anchor=west,inner sep=0pt,font=\fontsize{6.0}{6.0}\selectfont,text=black] at (830.0,634.0) {43};
\fill[tempcol!40!tempcold] (672.0,681.0) circle (6.5);
\node[anchor=west,inner sep=0pt,font=\fontsize{6.0}{6.0}\selectfont,text=black] at (680.0,680.0) {48};
\fill[tempcol!40!tempcold] (600.0,710.0) circle (6.5);
\node[anchor=west,inner sep=0pt,font=\fontsize{6.0}{6.0}\selectfont,text=black] at (608.0,709.0) {48};
\fill[tempcol!35!tempcold] (824.0,683.0) circle (6.5);
\node[anchor=west,inner sep=0pt,font=\fontsize{6.0}{6.0}\selectfont,text=black] at (832.0,682.0) {46};
\fill[tempcol!32!tempcold] (908.0,712.0) circle (6.5);
\node[anchor=west,inner sep=0pt,font=\fontsize{6.0}{6.0}\selectfont,text=black] at (916.0,711.0) {45};
\fill[tempcol!17!tempcold] (419.0,722.0) circle (6.5);
\node[anchor=west,inner sep=0pt,font=\fontsize{6.0}{6.0}\selectfont,text=black] at (427.0,721.0) {39};
\fill[tempcol!62!tempcold] (714.0,722.0) circle (6.5);
\node[anchor=west,inner sep=0pt,font=\fontsize{6.0}{6.0}\selectfont,text=black] at (722.0,721.0) {57};
\fill[tempcol!55!tempcold] (818.0,742.0) circle (6.5);
\node[anchor=west,inner sep=0pt,font=\fontsize{6.0}{6.0}\selectfont,text=black] at (826.0,741.0) {54};
\fill[tempcol!22!tempcold] (469.0,770.0) circle (6.5);
\node[anchor=west,inner sep=0pt,font=\fontsize{6.0}{6.0}\selectfont,text=black] at (477.0,769.0) {41};
\fill[tempcol!32!tempcold] (570.0,768.0) circle (6.5);
\node[anchor=west,inner sep=0pt,font=\fontsize{6.0}{6.0}\selectfont,text=black] at (578.0,767.0) {45};
\fill[tempcol!70!tempcold] (637.0,765.0) circle (6.5);
\node[anchor=west,inner sep=0pt,font=\fontsize{6.0}{6.0}\selectfont,text=black] at (645.0,764.0) {60};
\fill[tempcol!72!tempcold] (704.0,785.0) circle (6.5);
\node[anchor=west,inner sep=0pt,font=\fontsize{6.0}{6.0}\selectfont,text=black] at (712.0,784.0) {61};
\fill[tempcol!85!tempcold] (787.0,798.0) circle (6.5);
\node[anchor=west,inner sep=0pt,font=\fontsize{6.0}{6.0}\selectfont,text=black] at (795.0,797.0) {66};
\fill[tempcol!85!tempcold] (828.0,818.0) circle (6.5);
\node[anchor=west,inner sep=0pt,font=\fontsize{6.0}{6.0}\selectfont,text=black] at (836.0,817.0) {66};
\fill[tempcol!22!tempcold] (405.0,838.0) circle (6.5);
\node[anchor=west,inner sep=0pt,font=\fontsize{6.0}{6.0}\selectfont,text=black] at (413.0,837.0) {41};
\fill[tempcol!25!tempcold] (496.0,831.0) circle (6.5);
\node[anchor=west,inner sep=0pt,font=\fontsize{6.0}{6.0}\selectfont,text=black] at (504.0,830.0) {42};
\fill[tempcol!72!tempcold] (608.0,861.0) circle (6.5);
\node[anchor=west,inner sep=0pt,font=\fontsize{6.0}{6.0}\selectfont,text=black] at (616.0,860.0) {61};
\fill[tempcol!97!tempcold] (828.0,885.0) circle (6.5);
\node[anchor=west,inner sep=0pt,font=\fontsize{6.0}{6.0}\selectfont,text=black] at (836.0,884.0) {71};
\fill[tempcol!40!tempcold] (498.0,909.0) circle (6.5);
\node[anchor=west,inner sep=0pt,font=\fontsize{6.0}{6.0}\selectfont,text=black] at (506.0,908.0) {48};
\fill[tempcol!95!tempcold] (729.0,914.0) circle (6.5);
\node[anchor=west,inner sep=0pt,font=\fontsize{6.0}{6.0}\selectfont,text=black] at (737.0,913.0) {70};
\fill[tempcol!42!tempcold] (390.0,946.0) circle (6.5);
\node[anchor=west,inner sep=0pt,font=\fontsize{6.0}{6.0}\selectfont,text=black] at (398.0,945.0) {49};
\fill[tempcol!85!tempcold] (603.0,945.0) circle (6.5);
\node[anchor=west,inner sep=0pt,font=\fontsize{6.0}{6.0}\selectfont,text=black] at (611.0,944.0) {66};
\fill[tempcol!85!tempcold] (648.0,960.0) circle (6.5);
\node[anchor=west,inner sep=0pt,font=\fontsize{6.0}{6.0}\selectfont,text=black] at (656.0,959.0) {66};
\fill[tempcol!45!tempcold] (452.0,1009.0) circle (6.5);
\node[anchor=west,inner sep=0pt,font=\fontsize{6.0}{6.0}\selectfont,text=black] at (460.0,1008.0) {50};
\fill[tempcol!50!tempcold] (507.0,1061.0) circle (6.5);
\node[anchor=west,inner sep=0pt,font=\fontsize{6.0}{6.0}\selectfont,text=black] at (515.0,1060.0) {52};
\fill[tempcol!77!tempcold] (613.0,1041.0) circle (6.5);
\node[anchor=west,inner sep=0pt,font=\fontsize{6.0}{6.0}\selectfont,text=black] at (621.0,1040.0) {63};
\fill[tempcol!92!tempcold] (725.0,1031.0) circle (6.5);
\node[anchor=west,inner sep=0pt,font=\fontsize{6.0}{6.0}\selectfont,text=black] at (733.0,1030.0) {69};
\fill[tempcol!100!tempcold] (822.0,1006.0) circle (6.5);
\node[anchor=west,inner sep=0pt,font=\fontsize{6.0}{6.0}\selectfont,text=black] at (830.0,1005.0) {72};
\draw[fblue,line width=1.10pt] (655.0,702.0) -- (620.0,740.0) -- (600.0,770.0) -- (572.0,815.0) -- (553.0,865.0) -- (550.0,915.0) -- (540.0,970.0) -- (535.0,1020.0) -- (525.0,1072.0);
\fill[fblue] (651.6,705.7) -- (654.4,721.8) -- (638.1,720.4) -- cycle;
\fill[fblue] (631.3,727.7) -- (634.1,743.9) -- (617.7,742.5) -- cycle;
\fill[fblue] (612.6,751.1) -- (617.9,766.6) -- (601.5,767.7) -- cycle;
\fill[fblue] (596.2,776.2) -- (601.9,791.5) -- (585.6,793.2) -- cycle;
\fill[fblue] (580.3,801.7) -- (586.1,817.0) -- (569.7,818.6) -- cycle;
\fill[fblue] (566.9,828.4) -- (575.5,842.3) -- (559.8,847.0) -- cycle;
\fill[fblue] (553.6,855.8) -- (565.9,866.6) -- (552.4,875.8) -- cycle;
\fill[fblue] (551.8,885.8) -- (564.1,896.5) -- (550.6,905.7) -- cycle;
\fill[fblue] (549.9,915.7) -- (560.9,927.9) -- (546.3,935.4) -- cycle;
\fill[fblue] (544.5,945.2) -- (555.5,957.4) -- (540.9,964.9) -- cycle;
\fill[fblue] (539.5,974.8) -- (551.5,986.0) -- (537.5,994.7) -- cycle;
\fill[fblue] (536.5,1004.6) -- (548.5,1015.9) -- (534.5,1024.5) -- cycle;
\fill[fblue] (532.3,1034.3) -- (543.1,1046.6) -- (528.5,1053.9) -- cycle;
\draw[fred,line width=1.10pt] (655.0,702.0) -- (690.0,704.0) -- (722.0,710.0) -- (760.0,735.0) -- (818.0,768.0) -- (870.0,772.0) -- (940.0,768.0);
\fill[fred] (661.5,702.4) -- (661.8,700.5) -- (662.5,698.7) -- (663.6,697.2) -- (665.1,695.9) -- (666.7,695.0) -- (668.6,694.5) -- (670.5,694.4) -- (672.3,694.7) -- (674.1,695.4) -- (675.6,696.5) -- (676.9,697.9) -- (677.8,699.6) -- (678.4,701.4) -- (678.5,703.3) -- cycle;
\fill[fred] (691.4,704.3) -- (692.0,702.4) -- (692.9,700.8) -- (694.2,699.4) -- (695.8,698.3) -- (697.6,697.6) -- (699.4,697.3) -- (701.3,697.5) -- (703.2,698.0) -- (704.8,699.0) -- (706.2,700.3) -- (707.3,701.8) -- (708.0,703.6) -- (708.3,705.5) -- (708.1,707.4) -- cycle;
\fill[fred] (721.1,709.4) -- (722.3,707.9) -- (723.8,706.8) -- (725.5,706.0) -- (727.4,705.6) -- (729.3,705.6) -- (731.1,706.1) -- (732.8,707.0) -- (734.3,708.2) -- (735.5,709.7) -- (736.2,711.4) -- (736.6,713.3) -- (736.6,715.2) -- (736.1,717.0) -- (735.3,718.7) -- cycle;
\fill[fred] (746.1,725.9) -- (747.3,724.4) -- (748.9,723.3) -- (750.6,722.5) -- (752.5,722.1) -- (754.4,722.1) -- (756.2,722.6) -- (757.9,723.4) -- (759.4,724.7) -- (760.5,726.2) -- (761.3,727.9) -- (761.7,729.8) -- (761.7,731.7) -- (761.2,733.5) -- (760.3,735.2) -- cycle;
\fill[fred] (771.6,741.6) -- (772.8,740.1) -- (774.2,738.8) -- (775.9,737.9) -- (777.7,737.4) -- (779.6,737.3) -- (781.5,737.7) -- (783.2,738.4) -- (784.8,739.6) -- (786.0,741.0) -- (786.9,742.7) -- (787.4,744.5) -- (787.5,746.4) -- (787.2,748.3) -- (786.4,750.0) -- cycle;
\fill[fred] (797.7,756.5) -- (798.8,754.9) -- (800.3,753.7) -- (802.0,752.8) -- (803.8,752.3) -- (805.7,752.2) -- (807.6,752.5) -- (809.3,753.3) -- (810.9,754.4) -- (812.1,755.8) -- (813.0,757.5) -- (813.5,759.3) -- (813.6,761.2) -- (813.2,763.1) -- (812.5,764.9) -- cycle;
\fill[fred] (824.6,768.5) -- (825.0,766.6) -- (825.8,764.9) -- (826.9,763.4) -- (828.3,762.1) -- (830.0,761.2) -- (831.9,760.8) -- (833.8,760.7) -- (835.6,761.0) -- (837.4,761.8) -- (838.9,762.9) -- (840.2,764.4) -- (841.0,766.1) -- (841.5,767.9) -- (841.6,769.8) -- cycle;
\fill[fred] (854.6,770.8) -- (854.9,768.9) -- (855.7,767.2) -- (856.8,765.7) -- (858.3,764.4) -- (859.9,763.5) -- (861.8,763.1) -- (863.7,763.0) -- (865.6,763.3) -- (867.3,764.1) -- (868.8,765.2) -- (870.1,766.7) -- (871.0,768.4) -- (871.4,770.2) -- (871.5,772.1) -- cycle;
\fill[fred] (884.5,771.2) -- (884.6,769.3) -- (885.1,767.4) -- (886.0,765.8) -- (887.3,764.4) -- (888.9,763.3) -- (890.6,762.5) -- (892.5,762.2) -- (894.4,762.3) -- (896.2,762.8) -- (897.9,763.7) -- (899.3,765.0) -- (900.4,766.6) -- (901.1,768.3) -- (901.5,770.2) -- cycle;
\fill[fred] (914.4,769.5) -- (914.5,767.6) -- (915.1,765.7) -- (916.0,764.1) -- (917.3,762.6) -- (918.8,761.5) -- (920.6,760.8) -- (922.4,760.5) -- (924.3,760.6) -- (926.2,761.1) -- (927.8,762.0) -- (929.3,763.3) -- (930.4,764.9) -- (931.1,766.6) -- (931.4,768.5) -- cycle;
\node[lowmark] at (655.0,706.0) {L};
}
\caption{Step 3: air temperature (\degF) at each station, coloured from blue (cold, 32\,\degF) to red (warm, 72\,\degF). The warm air fills the south-east and a tongue reaches towards the low: the warm sector. Cold air covers the north-west. The cold front (blue, triangles) separates them in the west; the warm front (red, semicircles) in the east.}\label{p3-fig:airmass}
\end{figure}

\begin{table}[ht]
\centering
\caption{Evidence for the fronts: pairs of stations on either side of each front (cold side $\to$ warm side).}
\begin{tabular}{@{}llccc@{}}
\toprule
\textbf{Front} & \textbf{Stations} & \textbf{Temperature (\degF)} & \textbf{Dew point (\degF)} & \textbf{Wind from}\\
\midrule
cold & 16 $\to$ 17 & 45 $\to$ 60 & 40 $\to$ 55 & NW $\to$ SSW \\
cold & 22 $\to$ 23 & 42 $\to$ 61 & 36 $\to$ 57 & NW $\to$ SSW \\
cold & 25 $\to$ 28 & 48 $\to$ 66 & 38 $\to$ 59 & WNW $\to$ SW \\
cold & 31 $\to$ 32 & 52 $\to$ 63 & 50 $\to$ 61 & WSW $\to$ S \\
warm & 8 $\to$ 13 & 48 $\to$ 57 & 45 $\to$ 54 & NE $\to$ SSW \\
warm & 14 $\to$ 19 & 54 $\to$ 66 & 45 $\to$ 53 & E $\to$ S \\
warm & 10 $\to$ 20 & 46 $\to$ 66 & 43 $\to$ 55 & E $\to$ S \\
\bottomrule
\end{tabular}

\end{table}

\textbf{The cold front.} West of the low, the stations behind the front are cool (42 to 52\,\degF), drier (dew points 36 to 50\,\degF) and have winds from the west to north-west; those ahead are warm (60 to 66\,\degF), humid (dew points 55 to 61\,\degF) and have winds from the south-south-west to south-west. The pairs 16--17, 22--23, 25--28 and 31--32 straddle it, with temperature jumps of 11 to 19\,\degF. The front runs from the low towards the south-south-west, down through Missouri, Arkansas and east Texas.

\textbf{The warm front.} East of the low, the stations north of the front are cool (46 to 54\,\degF) with winds from the north-east or east; those south of it are warm (57 to 66\,\degF) with winds from the south. The pairs 8--13, 14--19 and 10--20 straddle it, with temperature jumps of 9 to 20\,\degF and the wind turning from east or north-east to south. The front runs east from the low through Indiana and Ohio.

\begin{answerbox}[Answer: the fronts]
A \textbf{cold front} (blue, with triangles pointing east, the way it moves) from the L towards the south-south-west to Texas; a \textbf{warm front} (red, with semicircles pointing north-east) from the L eastwards through Indiana and Ohio. The warm sector, the wedge of warm air between them, lies to the south-east.
\end{answerbox}

\subsection{Step 4: cloud cover above 75\,\%}

Cloud cover above 75\,\% means the stations drawn with a three-quarters, nine-tenths or completely filled circle. Looking at the sky circles: stations 2, 5, 6, 8, 9, 10, 11, 14, 17, 23, 31 and 32 are completely overcast; stations 3, 13, 19 and 28 have 7 to 8 tenths. Draw a line round the group, including the area between them. The result is a broad area wrapped round the low and extending along the cold front to the south-west.

\begin{figure}[ht]
\centering
\mapbase{0.0196cm}{\path[cloudfill,even odd rule] (567.4,526.0) -- (567.7,526.0) -- (568.1,526.0) -- (568.8,525.9) -- (569.7,525.9) -- (570.7,525.8) -- (571.9,525.8) -- (573.2,525.7) -- (574.8,525.7) -- (576.2,525.6) -- (577.8,525.6) -- (579.2,525.5) -- (580.8,525.4) -- (582.2,525.4) -- (583.8,525.3) -- (585.2,525.3) -- (586.8,525.3) -- (588.2,525.2) -- (589.8,525.2) -- (591.2,525.1) -- (592.8,525.1) -- (594.2,525.1) -- (595.8,525.0) -- (597.2,525.0) -- (598.8,525.0) -- (600.2,525.0) -- (601.8,524.9) -- (603.2,524.9) -- (604.8,524.9) -- (606.2,524.9) -- (607.8,524.8) -- (609.2,524.8) -- (610.8,524.8) -- (612.2,524.8) -- (613.8,524.8) -- (615.2,524.8) -- (616.8,524.7) -- (618.2,524.7) -- (619.8,524.7) -- (621.2,524.7) -- (622.8,524.7) -- (624.2,524.7) -- (625.8,524.7) -- (627.2,524.7) -- (628.8,524.6) -- (630.2,524.6) -- (631.8,524.6) -- (633.2,524.6) -- (634.8,524.6) -- (636.2,524.6) -- (637.8,524.6) -- (639.2,524.6) -- (640.8,524.6) -- (642.2,524.6) -- (643.8,524.6) -- (645.2,524.6) -- (646.8,524.6) -- (648.2,524.6) -- (649.8,524.6) -- (651.2,524.6) -- (652.8,524.6) -- (654.2,524.6) -- (655.8,524.6) -- (657.2,524.6) -- (658.8,524.6) -- (660.2,524.6) -- (661.8,524.6) -- (663.2,524.6) -- (664.8,524.6) -- (666.2,524.6) -- (667.8,524.6) -- (669.2,524.6) -- (670.8,524.6) -- (672.2,524.6) -- (673.8,524.6) -- (675.2,524.6) -- (676.8,524.6) -- (678.2,524.6) -- (679.8,524.6) -- (681.2,524.6) -- (682.8,524.6) -- (684.2,524.6) -- (685.8,524.6) -- (687.2,524.7) -- (688.8,524.7) -- (690.0,524.8) -- (690.9,524.9) -- (691.5,525.2) -- (691.8,525.5) -- (692.4,526.0) -- (693.1,526.8) -- (694.1,527.8) -- (695.2,528.9) -- (696.2,529.9) -- (696.9,530.7) -- (697.5,531.2) -- (697.8,531.5) -- (698.3,532.1) -- (699.1,532.8) -- (700.1,533.8) -- (701.2,535.0) -- (702.5,535.9) -- (703.8,536.5) -- (705.2,536.8) -- (706.8,536.8) -- (708.2,536.5) -- (709.5,535.9) -- (710.8,535.0) -- (712.0,533.8) -- (713.0,532.8) -- (713.8,532.1) -- (714.3,531.6) -- (714.6,531.3) -- (715.2,531.1) -- (716.1,531.0) -- (717.2,530.9) -- (718.8,530.9) -- (720.2,530.9) -- (721.8,530.9) -- (723.2,531.0) -- (724.8,531.0) -- (726.2,531.0) -- (727.8,531.0) -- (729.2,531.0) -- (730.8,531.0) -- (732.2,531.1) -- (733.8,531.1) -- (735.2,531.1) -- (736.8,531.1) -- (738.2,531.1) -- (739.8,531.2) -- (741.2,531.2) -- (742.8,531.2) -- (744.2,530.9) -- (745.6,530.2) -- (747.0,529.3) -- (748.3,528.0) -- (749.4,526.9) -- (750.1,526.2) -- (750.6,525.7) -- (750.7,525.6) -- (751.2,525.5) -- (752.1,525.4) -- (753.2,525.3) -- (754.8,525.4) -- (756.2,525.4) -- (757.8,525.4) -- (759.2,525.4) -- (760.8,525.4) -- (762.2,525.5) -- (763.8,525.5) -- (765.2,525.5) -- (766.8,525.5) -- (768.2,525.5) -- (769.8,525.5) -- (771.2,525.6) -- (772.8,525.6) -- (774.2,525.6) -- (775.8,525.6) -- (777.2,525.6) -- (778.8,525.6) -- (780.2,525.7) -- (781.8,525.7) -- (783.2,525.7) -- (784.8,525.7) -- (786.2,525.7) -- (787.8,525.7) -- (789.2,525.7) -- (790.8,525.7) -- (792.2,525.8) -- (793.8,525.8) -- (795.2,525.8) -- (796.8,525.8) -- (798.2,525.8) -- (799.8,525.8) -- (801.2,525.8) -- (802.8,525.8) -- (804.2,525.8) -- (805.8,525.9) -- (807.2,525.9) -- (808.8,525.9) -- (810.2,525.9) -- (811.8,525.9) -- (813.2,525.9) -- (814.8,525.9) -- (816.2,525.9) -- (817.8,525.9) -- (819.2,525.9) -- (820.8,525.9) -- (822.2,526.0) -- (823.8,526.0) -- (825.2,526.0) -- (826.8,526.0) -- (828.2,526.0) -- (829.6,526.0) -- (830.9,526.0) -- (832.2,526.0) -- (833.2,526.3) -- (833.9,526.9) -- (834.4,527.7) -- (834.6,528.9) -- (834.8,529.8) -- (835.0,530.5) -- (835.1,531.1) -- (835.2,531.5) -- (835.4,532.1) -- (835.6,533.0) -- (836.0,534.2) -- (836.4,535.8) -- (836.9,537.2) -- (837.4,538.8) -- (837.9,540.2) -- (838.5,541.8) -- (839.0,543.1) -- (839.5,544.3) -- (840.0,545.4) -- (840.4,546.3) -- (840.8,547.1) -- (841.1,547.8) -- (841.5,548.5) -- (841.8,549.1) -- (842.3,549.9) -- (843.0,551.0) -- (843.8,552.2) -- (844.9,553.8) -- (845.7,554.9) -- (846.3,555.8) -- (846.6,556.3) -- (846.8,556.4) -- (847.2,556.9) -- (848.0,557.7) -- (849.2,558.7) -- (850.6,560.0) -- (851.7,561.0) -- (852.4,561.7) -- (852.8,562.1) -- (852.9,562.1) -- (853.3,562.3) -- (854.1,562.8) -- (855.2,563.4) -- (856.8,564.2) -- (858.2,564.9) -- (859.8,565.6) -- (861.2,566.1) -- (862.8,566.6) -- (864.1,567.0) -- (865.4,567.4) -- (866.6,567.6) -- (867.6,567.8) -- (868.5,567.9) -- (869.3,568.0) -- (869.9,568.1) -- (870.3,568.1) -- (871.0,568.2) -- (872.0,568.2) -- (873.2,568.2) -- (874.8,568.2) -- (876.0,568.2) -- (876.9,568.2) -- (877.6,568.1) -- (878.1,568.1) -- (878.7,568.0) -- (879.4,567.9) -- (880.3,567.8) -- (881.4,567.7) -- (882.6,567.5) -- (883.9,567.2) -- (885.2,566.9) -- (886.8,566.6) -- (888.2,566.2) -- (889.8,565.8) -- (891.2,565.4) -- (892.8,565.0) -- (894.2,564.5) -- (895.8,564.0) -- (897.2,563.5) -- (898.8,563.0) -- (899.9,562.5) -- (900.7,562.3) -- (901.1,562.1) -- (901.2,562.1) -- (901.7,561.9) -- (902.4,561.5) -- (903.5,561.0) -- (904.9,560.4) -- (906.3,559.7) -- (907.8,559.0) -- (909.2,558.3) -- (910.8,557.6) -- (911.9,557.4) -- (912.6,557.7) -- (913.0,558.6) -- (913.0,559.9) -- (913.0,561.3) -- (913.0,562.8) -- (913.0,564.2) -- (913.1,565.8) -- (913.1,567.2) -- (913.1,568.8) -- (913.1,570.2) -- (913.1,571.8) -- (913.1,573.2) -- (913.1,574.8) -- (913.1,576.2) -- (913.1,577.8) -- (913.2,579.2) -- (913.2,580.8) -- (913.2,582.2) -- (913.2,583.8) -- (913.2,585.2) -- (913.2,586.8) -- (913.2,588.2) -- (913.2,589.8) -- (913.3,591.2) -- (913.3,592.8) -- (913.3,594.2) -- (913.3,595.8) -- (913.3,597.2) -- (913.3,598.8) -- (913.4,600.2) -- (913.4,601.8) -- (913.4,603.2) -- (913.4,604.8) -- (913.4,606.2) -- (913.4,607.8) -- (913.5,609.2) -- (913.5,610.8) -- (913.5,612.2) -- (913.5,613.8) -- (913.9,615.2) -- (914.6,616.6) -- (915.6,618.0) -- (917.0,619.3) -- (918.4,620.3) -- (919.8,621.0) -- (921.2,621.3) -- (922.8,621.3) -- (924.2,621.3) -- (925.8,621.3) -- (927.2,621.3) -- (928.8,621.3) -- (930.2,621.3) -- (931.8,621.3) -- (933.2,621.2) -- (934.8,621.2) -- (935.9,621.3) -- (936.8,621.4) -- (937.3,621.5) -- (937.5,621.7) -- (938.0,622.2) -- (938.7,622.9) -- (939.7,623.9) -- (941.0,625.2) -- (942.4,626.1) -- (943.8,626.8) -- (945.2,627.1) -- (946.8,627.1) -- (947.9,627.1) -- (948.8,627.2) -- (949.4,627.4) -- (949.6,627.6) -- (950.1,628.1) -- (950.8,628.9) -- (951.8,629.9) -- (953.1,631.1) -- (954.4,632.0) -- (955.8,632.6) -- (957.2,632.9) -- (958.8,632.9) -- (959.9,633.0) -- (960.8,633.1) -- (961.4,633.3) -- (961.7,633.6) -- (962.2,634.1) -- (962.9,634.8) -- (963.9,635.8) -- (965.2,637.0) -- (966.2,638.0) -- (966.9,638.8) -- (967.4,639.3) -- (967.7,639.6) -- (968.2,640.1) -- (969.0,640.8) -- (970.0,641.8) -- (971.2,643.0) -- (972.2,644.0) -- (972.9,644.7) -- (973.5,645.2) -- (973.8,645.5) -- (974.3,646.1) -- (975.0,646.8) -- (976.0,647.8) -- (977.2,648.9) -- (978.2,649.9) -- (978.9,650.7) -- (979.5,651.2) -- (979.8,651.5) -- (980.3,652.1) -- (981.1,652.8) -- (982.1,653.8) -- (983.2,654.9) -- (984.2,655.9) -- (985.0,656.6) -- (985.5,657.2) -- (985.8,657.5) -- (986.1,658.1) -- (986.3,659.0) -- (986.3,660.2) -- (986.3,661.8) -- (986.6,663.2) -- (987.2,664.5) -- (988.1,665.7) -- (989.3,666.9) -- (990.2,667.8) -- (991.0,668.6) -- (991.5,669.1) -- (991.9,669.5) -- (992.1,670.1) -- (992.3,671.0) -- (992.4,672.2) -- (992.4,673.8) -- (992.4,675.2) -- (992.4,676.8) -- (992.4,678.2) -- (992.4,679.8) -- (992.4,681.2) -- (992.4,682.8) -- (992.4,684.2) -- (992.4,685.8) -- (992.4,687.2) -- (992.4,688.8) -- (992.4,690.2) -- (992.4,691.8) -- (992.4,693.2) -- (992.4,694.8) -- (992.4,696.2) -- (992.4,697.8) -- (992.4,699.2) -- (992.4,700.8) -- (992.4,702.2) -- (992.4,703.8) -- (992.4,705.2) -- (992.4,706.8) -- (992.4,708.2) -- (992.4,709.8) -- (992.4,711.2) -- (992.4,712.8) -- (992.4,714.2) -- (992.4,715.8) -- (992.4,717.2) -- (992.4,718.8) -- (992.4,720.2) -- (992.4,721.8) -- (992.4,723.2) -- (992.4,724.8) -- (992.4,726.2) -- (992.4,727.8) -- (992.4,729.2) -- (992.4,730.8) -- (992.4,732.2) -- (992.4,733.8) -- (992.4,735.2) -- (992.4,736.8) -- (992.4,738.2) -- (992.4,739.8) -- (992.4,741.2) -- (992.4,742.8) -- (992.4,744.2) -- (992.4,745.8) -- (992.3,747.2) -- (992.3,748.8) -- (992.3,750.2) -- (992.3,751.8) -- (992.2,753.0) -- (992.1,753.9) -- (991.8,754.5) -- (991.5,754.8) -- (991.0,755.4) -- (990.2,756.1) -- (989.2,757.1) -- (988.0,758.2) -- (987.1,759.5) -- (986.5,760.8) -- (986.2,762.2) -- (986.2,763.8) -- (986.1,765.0) -- (986.0,765.9) -- (985.8,766.5) -- (985.5,766.8) -- (984.9,767.3) -- (984.2,768.0) -- (983.2,769.0) -- (982.0,770.2) -- (981.0,771.2) -- (980.2,771.9) -- (979.7,772.4) -- (979.4,772.7) -- (978.9,773.2) -- (978.1,774.0) -- (977.1,775.0) -- (975.9,776.2) -- (974.9,777.1) -- (974.1,777.9) -- (973.6,778.4) -- (973.4,778.6) -- (972.9,779.1) -- (972.1,779.9) -- (971.1,780.9) -- (969.8,782.1) -- (968.8,783.1) -- (968.0,783.8) -- (967.5,784.3) -- (967.3,784.5) -- (966.8,785.0) -- (966.0,785.8) -- (965.0,786.8) -- (963.7,788.1) -- (962.6,789.1) -- (961.9,789.8) -- (961.4,790.2) -- (961.2,790.4) -- (960.7,790.5) -- (959.9,790.6) -- (958.8,790.6) -- (957.2,790.6) -- (955.8,790.9) -- (954.3,791.6) -- (952.9,792.6) -- (951.4,794.0) -- (950.3,795.0) -- (949.6,795.7) -- (949.2,796.1) -- (949.1,796.2) -- (948.7,796.2) -- (947.9,796.3) -- (946.8,796.2) -- (945.2,796.2) -- (943.8,796.3) -- (942.2,796.5) -- (940.8,796.9) -- (939.2,797.4) -- (937.9,797.6) -- (936.6,797.6) -- (935.4,797.3) -- (934.4,796.8) -- (933.5,796.3) -- (932.7,796.0) -- (932.1,795.7) -- (931.7,795.4) -- (931.0,795.1) -- (930.0,794.6) -- (928.8,794.0) -- (927.2,793.3) -- (925.8,792.6) -- (924.5,792.0) -- (923.2,791.4) -- (922.0,790.8) -- (921.0,790.4) -- (920.2,790.0) -- (919.7,789.8) -- (919.4,789.7) -- (918.8,789.4) -- (917.9,789.0) -- (916.8,788.5) -- (915.2,787.8) -- (913.8,787.1) -- (912.2,786.5) -- (910.8,785.9) -- (909.2,785.2) -- (908.1,784.7) -- (907.2,784.4) -- (906.7,784.2) -- (906.5,784.1) -- (906.0,783.9) -- (905.3,783.6) -- (904.3,783.2) -- (903.0,782.7) -- (901.6,782.2) -- (900.2,781.6) -- (898.8,781.1) -- (897.2,780.5) -- (895.8,780.0) -- (894.5,779.5) -- (893.2,779.0) -- (892.0,778.6) -- (891.0,778.3) -- (890.2,778.0) -- (889.7,777.8) -- (889.4,777.7) -- (888.8,777.5) -- (887.9,777.3) -- (886.8,776.9) -- (885.2,776.4) -- (883.8,775.9) -- (882.2,775.4) -- (880.8,775.0) -- (879.2,774.5) -- (877.8,774.1) -- (876.2,773.7) -- (874.8,773.2) -- (873.2,772.8) -- (872.1,772.5) -- (871.2,772.3) -- (870.7,772.1) -- (870.4,772.1) -- (869.9,772.0) -- (869.2,771.8) -- (868.2,771.5) -- (866.9,771.2) -- (865.6,770.9) -- (864.2,770.5) -- (862.8,770.2) -- (861.2,769.8) -- (859.8,769.5) -- (858.2,769.2) -- (856.8,768.9) -- (855.2,768.5) -- (853.8,768.2) -- (852.2,767.9) -- (850.8,767.7) -- (849.2,767.4) -- (847.8,767.1) -- (846.3,766.9) -- (844.9,766.6) -- (843.4,766.4) -- (842.3,766.2) -- (841.6,766.1) -- (841.2,766.0) -- (841.1,766.0) -- (840.7,765.9) -- (839.9,765.8) -- (838.8,765.6) -- (837.2,765.3) -- (835.8,765.1) -- (834.2,764.9) -- (832.8,764.7) -- (831.2,764.4) -- (829.8,764.2) -- (828.2,764.0) -- (826.8,763.8) -- (825.2,763.6) -- (823.8,763.4) -- (822.2,763.2) -- (820.8,763.0) -- (819.2,762.8) -- (817.8,762.6) -- (816.2,762.4) -- (814.8,762.2) -- (813.2,762.0) -- (811.8,761.8) -- (810.2,761.6) -- (808.8,761.5) -- (807.2,761.3) -- (805.8,761.1) -- (804.2,760.9) -- (802.8,760.7) -- (801.2,760.5) -- (800.0,760.4) -- (799.0,760.2) -- (798.3,760.1) -- (797.8,760.1) -- (797.2,760.0) -- (796.5,759.9) -- (795.6,759.8) -- (794.5,759.7) -- (793.4,759.5) -- (792.1,759.3) -- (790.8,759.1) -- (789.2,758.9) -- (787.8,758.7) -- (786.2,758.5) -- (784.8,758.3) -- (783.2,758.1) -- (781.8,757.9) -- (780.2,757.7) -- (778.8,757.5) -- (777.2,757.3) -- (775.8,757.2) -- (774.2,757.0) -- (772.8,756.8) -- (771.2,756.6) -- (769.8,756.4) -- (768.2,756.3) -- (766.8,756.1) -- (765.2,755.9) -- (763.8,755.8) -- (762.2,755.6) -- (760.8,755.5) -- (759.2,755.4) -- (757.8,755.3) -- (756.2,755.2) -- (754.8,755.1) -- (753.2,755.1) -- (751.8,755.0) -- (750.2,755.0) -- (748.8,755.0) -- (747.2,755.0) -- (745.8,755.1) -- (744.2,755.1) -- (742.8,755.2) -- (741.2,755.3) -- (739.8,755.4) -- (738.2,755.6) -- (736.8,755.8) -- (735.2,756.0) -- (733.8,756.3) -- (732.2,756.6) -- (730.8,756.9) -- (729.2,757.3) -- (727.8,757.6) -- (726.2,758.1) -- (724.8,758.5) -- (723.2,759.0) -- (722.1,759.4) -- (721.3,759.7) -- (720.8,759.9) -- (720.7,759.9) -- (720.2,760.1) -- (719.5,760.4) -- (718.4,760.9) -- (717.0,761.5) -- (715.6,762.2) -- (714.2,762.9) -- (712.8,763.6) -- (711.2,764.5) -- (710.1,765.1) -- (709.3,765.5) -- (708.8,765.8) -- (708.7,765.9) -- (708.3,766.2) -- (707.5,766.7) -- (706.5,767.4) -- (705.1,768.3) -- (703.9,769.2) -- (702.8,769.9) -- (701.9,770.6) -- (701.2,771.2) -- (700.4,771.7) -- (699.7,772.3) -- (698.9,772.9) -- (698.2,773.6) -- (697.3,774.3) -- (696.5,775.0) -- (695.6,775.8) -- (694.6,776.7) -- (693.8,777.5) -- (693.0,778.2) -- (692.4,778.8) -- (691.8,779.3) -- (691.2,780.0) -- (690.4,780.7) -- (689.6,781.6) -- (688.7,782.5) -- (687.8,783.4) -- (687.1,784.2) -- (686.4,785.0) -- (685.9,785.6) -- (685.2,786.3) -- (684.6,787.1) -- (683.9,787.9) -- (683.2,788.7) -- (682.4,789.6) -- (681.7,790.5) -- (680.9,791.4) -- (680.1,792.4) -- (679.5,793.2) -- (678.9,794.0) -- (678.3,794.6) -- (677.9,795.2) -- (677.3,795.9) -- (676.6,796.9) -- (675.7,798.1) -- (674.6,799.6) -- (673.8,800.7) -- (673.2,801.4) -- (672.9,801.8) -- (672.9,801.9) -- (672.5,802.3) -- (672.0,803.1) -- (671.1,804.2) -- (670.0,805.8) -- (669.1,807.0) -- (668.3,808.0) -- (667.8,808.7) -- (667.5,809.1) -- (667.0,809.7) -- (666.5,810.5) -- (665.8,811.4) -- (664.9,812.4) -- (664.2,813.4) -- (663.4,814.4) -- (662.7,815.3) -- (662.0,816.1) -- (661.4,816.9) -- (660.7,817.7) -- (660.1,818.4) -- (659.6,819.0) -- (658.9,819.8) -- (658.1,820.6) -- (657.3,821.5) -- (656.4,822.4) -- (655.5,823.3) -- (654.8,824.1) -- (654.1,824.7) -- (653.5,825.2) -- (652.8,825.8) -- (652.1,826.4) -- (651.2,827.1) -- (650.3,827.8) -- (649.3,828.5) -- (648.2,829.3) -- (646.9,830.0) -- (645.6,830.8) -- (644.5,831.4) -- (643.8,831.9) -- (643.3,832.1) -- (643.2,832.2) -- (642.7,832.3) -- (641.9,832.5) -- (640.8,832.8) -- (639.2,833.2) -- (637.8,833.4) -- (636.2,833.6) -- (634.8,833.6) -- (633.2,833.5) -- (632.0,833.3) -- (630.9,833.1) -- (630.0,832.8) -- (629.3,832.5) -- (628.6,832.1) -- (627.9,831.7) -- (627.1,831.3) -- (626.2,830.8) -- (625.4,830.2) -- (624.5,829.4) -- (623.7,828.5) -- (622.8,827.5) -- (621.9,826.6) -- (621.2,825.7) -- (620.5,824.8) -- (619.9,824.0) -- (619.4,823.2) -- (618.9,822.5) -- (618.5,821.8) -- (618.1,821.1) -- (617.6,820.2) -- (617.1,819.1) -- (616.5,817.8) -- (615.8,816.2) -- (615.1,814.9) -- (614.6,813.7) -- (614.1,812.7) -- (613.6,811.8) -- (613.3,811.0) -- (612.9,810.2) -- (612.7,809.5) -- (612.4,808.9) -- (612.1,808.1) -- (611.7,807.0) -- (611.2,805.8) -- (610.6,804.2) -- (610.0,802.8) -- (609.4,801.2) -- (608.8,799.8) -- (608.1,798.2) -- (607.6,797.1) -- (607.3,796.3) -- (607.1,795.9) -- (607.1,795.8) -- (606.9,795.4) -- (606.6,794.6) -- (606.1,793.5) -- (605.6,792.1) -- (605.0,790.7) -- (604.3,789.2) -- (603.7,787.8) -- (603.0,786.2) -- (602.4,785.0) -- (601.9,784.1) -- (601.6,783.4) -- (601.3,782.9) -- (601.0,782.3) -- (600.7,781.6) -- (600.3,780.7) -- (599.8,779.6) -- (599.2,778.4) -- (598.5,777.1) -- (597.8,775.8) -- (597.0,774.2) -- (596.4,773.0) -- (595.9,772.1) -- (595.5,771.5) -- (595.3,771.2) -- (595.0,770.6) -- (594.6,769.9) -- (594.1,768.9) -- (593.4,767.7) -- (592.7,766.5) -- (592.0,765.2) -- (591.2,763.9) -- (590.3,762.4) -- (589.6,761.4) -- (589.2,760.6) -- (588.9,760.2) -- (588.9,760.1) -- (588.6,759.7) -- (588.1,758.9) -- (587.4,757.8) -- (586.5,756.2) -- (585.7,754.9) -- (585.0,753.8) -- (584.3,752.8) -- (583.8,752.0) -- (583.3,751.3) -- (582.8,750.5) -- (582.3,749.8) -- (581.9,749.1) -- (581.3,748.2) -- (580.5,747.1) -- (579.6,745.8) -- (578.6,744.2) -- (577.8,743.1) -- (577.3,742.4) -- (577.0,742.0) -- (577.0,742.0) -- (576.8,741.6) -- (576.2,740.8) -- (575.5,739.7) -- (574.5,738.2) -- (573.6,736.9) -- (572.8,735.9) -- (572.2,735.0) -- (571.7,734.3) -- (571.2,733.6) -- (570.7,732.8) -- (570.2,732.0) -- (569.6,731.2) -- (569.0,730.3) -- (568.2,729.2) -- (567.4,728.0) -- (566.4,726.7) -- (565.7,725.6) -- (565.2,724.9) -- (564.8,724.4) -- (564.7,724.2) -- (564.4,723.7) -- (563.8,722.9) -- (563.0,721.8) -- (562.0,720.2) -- (561.2,719.0) -- (560.5,718.0) -- (560.0,717.2) -- (559.6,716.6) -- (559.1,715.9) -- (558.7,715.2) -- (558.1,714.3) -- (557.5,713.4) -- (556.8,712.3) -- (556.1,711.1) -- (555.3,709.8) -- (554.4,708.4) -- (553.7,707.3) -- (553.2,706.5) -- (553.0,706.1) -- (552.9,706.1) -- (552.7,705.7) -- (552.3,704.9) -- (551.6,703.8) -- (550.8,702.2) -- (550.0,700.8) -- (549.3,699.4) -- (548.6,698.1) -- (548.0,696.9) -- (547.5,695.9) -- (547.1,695.1) -- (546.8,694.6) -- (546.7,694.4) -- (546.4,693.8) -- (546.0,692.9) -- (545.5,691.8) -- (544.9,690.2) -- (544.3,688.8) -- (543.6,687.2) -- (543.1,685.8) -- (542.5,684.2) -- (542.0,683.0) -- (541.7,682.0) -- (541.4,681.2) -- (541.2,680.6) -- (541.1,680.0) -- (540.8,679.2) -- (540.6,678.4) -- (540.3,677.4) -- (540.0,676.3) -- (539.7,675.1) -- (539.4,673.8) -- (539.1,672.2) -- (538.8,670.8) -- (538.6,669.2) -- (538.3,667.8) -- (538.1,666.2) -- (537.9,664.8) -- (537.8,663.2) -- (537.7,661.8) -- (537.6,660.2) -- (537.5,658.8) -- (537.4,657.2) -- (537.4,655.8) -- (537.4,654.2) -- (537.5,652.8) -- (537.6,651.2) -- (537.6,649.8) -- (537.8,648.2) -- (537.9,646.8) -- (538.1,645.2) -- (538.2,643.8) -- (538.4,642.2) -- (538.7,640.8) -- (538.9,639.2) -- (539.2,637.8) -- (539.5,636.2) -- (539.8,634.8) -- (540.0,633.5) -- (540.3,632.2) -- (540.6,631.1) -- (540.8,630.1) -- (541.0,629.3) -- (541.1,628.8) -- (541.2,628.5) -- (541.3,627.9) -- (541.6,627.0) -- (541.9,625.8) -- (542.3,624.2) -- (542.7,622.8) -- (543.2,621.2) -- (543.6,619.8) -- (544.1,618.2) -- (544.5,616.8) -- (545.0,615.2) -- (545.4,613.8) -- (545.9,612.2) -- (546.3,611.0) -- (546.5,610.1) -- (546.7,609.5) -- (546.8,609.2) -- (547.0,608.6) -- (547.3,607.9) -- (547.6,606.9) -- (547.9,605.8) -- (548.3,604.5) -- (548.8,603.2) -- (549.2,601.8) -- (549.7,600.2) -- (550.2,598.8) -- (550.6,597.2) -- (551.1,595.8) -- (551.6,594.2) -- (552.0,593.0) -- (552.3,591.9) -- (552.6,591.1) -- (552.7,590.5) -- (552.9,589.9) -- (553.2,589.1) -- (553.4,588.3) -- (553.7,587.4) -- (554.0,586.3) -- (554.4,585.1) -- (554.8,583.8) -- (555.2,582.2) -- (555.6,580.8) -- (556.1,579.2) -- (556.5,577.8) -- (556.9,576.2) -- (557.3,574.8) -- (557.7,573.3) -- (558.1,571.8) -- (558.4,570.4) -- (558.7,569.3) -- (558.9,568.5) -- (559.0,568.1) -- (559.0,568.1) -- (559.1,567.7) -- (559.3,566.9) -- (559.6,565.8) -- (560.0,564.2) -- (560.3,562.8) -- (560.7,561.2) -- (561.1,559.8) -- (561.4,558.2) -- (561.7,556.8) -- (562.1,555.2) -- (562.4,553.8) -- (562.7,552.2) -- (563.0,550.8) -- (563.3,549.2) -- (563.6,547.8) -- (563.9,546.2) -- (564.2,544.9) -- (564.4,543.7) -- (564.6,542.6) -- (564.8,541.7) -- (564.9,540.9) -- (565.0,540.2) -- (565.2,539.5) -- (565.3,538.9) -- (565.4,538.1) -- (565.6,537.0) -- (565.8,535.8) -- (566.1,534.2) -- (566.3,532.8) -- (566.6,531.2) -- (566.8,529.8) -- (567.1,528.2) -- (567.3,527.1) -- (567.4,526.4) -- (567.4,526.0) -- cycle  (554.8,1036.0) -- (555.1,1035.9) -- (555.6,1035.7) -- (556.4,1035.4) -- (557.4,1035.0) -- (558.6,1034.6) -- (559.9,1034.2) -- (561.2,1033.8) -- (562.8,1033.5) -- (564.2,1033.2) -- (565.8,1032.9) -- (567.2,1032.7) -- (568.8,1032.6) -- (570.2,1032.5) -- (571.8,1032.5) -- (573.2,1032.5) -- (574.8,1032.6) -- (576.2,1032.8) -- (577.8,1033.0) -- (579.2,1033.3) -- (580.8,1033.6) -- (582.2,1034.0) -- (583.6,1034.4) -- (585.0,1034.8) -- (586.4,1035.3) -- (587.4,1035.7) -- (588.2,1035.9) -- (588.6,1036.1) -- (588.8,1036.1) -- (589.3,1036.4) -- (590.1,1036.8) -- (591.2,1037.4) -- (592.8,1038.1) -- (594.2,1038.8) -- (595.5,1039.6) -- (596.7,1040.3) -- (597.8,1041.0) -- (598.7,1041.5) -- (599.5,1042.0) -- (600.0,1042.4) -- (600.4,1042.6) -- (601.1,1043.1) -- (602.0,1043.8) -- (603.2,1044.7) -- (604.8,1045.9) -- (605.9,1046.7) -- (606.7,1047.4) -- (607.2,1047.7) -- (607.3,1047.8) -- (607.7,1048.2) -- (608.5,1048.9) -- (609.5,1049.9) -- (610.9,1051.2) -- (612.0,1052.2) -- (612.8,1052.9) -- (613.3,1053.4) -- (613.6,1053.6) -- (614.1,1054.2) -- (614.8,1055.0) -- (615.8,1056.1) -- (617.1,1057.4) -- (618.1,1058.5) -- (618.8,1059.3) -- (619.2,1059.7) -- (619.3,1059.8) -- (619.6,1060.3) -- (620.3,1061.1) -- (621.2,1062.2) -- (622.4,1063.8) -- (623.3,1064.9) -- (624.0,1065.8) -- (624.5,1066.4) -- (624.7,1066.7) -- (625.0,1067.3) -- (625.5,1068.0) -- (626.2,1069.0) -- (627.0,1070.2) -- (627.8,1071.4) -- (628.5,1072.5) -- (629.3,1073.7) -- (630.0,1074.8) -- (630.5,1075.7) -- (631.0,1076.5) -- (631.3,1077.0) -- (631.5,1077.4) -- (631.9,1078.1) -- (632.4,1079.0) -- (633.1,1080.2) -- (633.9,1081.8) -- (634.6,1083.1) -- (635.2,1084.4) -- (635.8,1085.5) -- (636.3,1086.5) -- (636.7,1087.4) -- (637.0,1088.1) -- (637.3,1088.7) -- (637.6,1089.2) -- (637.9,1090.0) -- (638.3,1091.0) -- (638.9,1092.2) -- (639.5,1093.8) -- (640.1,1095.2) -- (640.7,1096.8) -- (641.3,1098.2) -- (641.9,1099.8) -- (642.3,1100.9) -- (642.7,1101.8) -- (642.8,1102.3) -- (642.9,1102.4) -- (643.1,1102.9) -- (643.3,1103.6) -- (643.7,1104.7) -- (644.1,1106.0) -- (644.6,1107.4) -- (645.0,1108.8) -- (645.5,1110.2) -- (645.9,1111.8) -- (646.4,1113.2) -- (646.8,1114.8) -- (647.3,1116.2) -- (647.7,1117.8) -- (648.1,1119.0) -- (648.4,1120.1) -- (648.6,1120.9) -- (648.8,1121.5) -- (648.9,1122.2) -- (649.1,1122.9) -- (649.3,1123.7) -- (649.6,1124.6) -- (649.7,1125.3) -- (649.7,1125.8) -- (649.6,1126.0) -- (649.3,1126.0) -- (648.8,1126.0) -- (647.9,1126.1) -- (646.8,1126.1) -- (645.2,1126.2) -- (643.8,1126.6) -- (642.3,1127.3) -- (640.9,1128.4) -- (639.6,1129.8) -- (638.5,1131.0) -- (637.7,1131.7) -- (637.3,1132.2) -- (637.2,1132.3) -- (636.7,1132.4) -- (635.9,1132.5) -- (634.8,1132.5) -- (633.2,1132.6) -- (631.8,1132.6) -- (630.2,1132.7) -- (628.8,1132.7) -- (627.2,1132.7) -- (625.8,1132.8) -- (624.2,1132.8) -- (622.8,1132.8) -- (621.2,1132.8) -- (619.8,1132.9) -- (618.2,1132.9) -- (616.8,1132.9) -- (615.2,1132.9) -- (613.8,1133.0) -- (612.2,1133.0) -- (610.8,1133.0) -- (609.2,1133.0) -- (607.8,1133.0) -- (606.2,1133.0) -- (604.8,1133.1) -- (603.2,1133.1) -- (601.8,1133.1) -- (600.2,1133.1) -- (598.8,1133.1) -- (597.2,1133.1) -- (595.8,1133.1) -- (594.2,1133.1) -- (592.8,1133.2) -- (591.2,1133.2) -- (590.1,1133.1) -- (589.2,1133.0) -- (588.6,1132.7) -- (588.3,1132.4) -- (587.7,1131.9) -- (587.0,1131.2) -- (586.0,1130.2) -- (584.8,1129.0) -- (583.5,1128.1) -- (582.2,1127.5) -- (580.8,1127.2) -- (579.2,1127.2) -- (577.8,1127.2) -- (576.2,1127.2) -- (574.8,1127.2) -- (573.2,1127.2) -- (571.8,1127.5) -- (570.5,1128.1) -- (569.2,1129.0) -- (568.0,1130.2) -- (567.0,1131.2) -- (566.3,1131.9) -- (565.8,1132.5) -- (565.5,1132.8) -- (564.9,1133.3) -- (564.2,1134.0) -- (563.2,1135.0) -- (562.0,1136.2) -- (561.0,1137.2) -- (560.3,1137.9) -- (559.8,1138.5) -- (559.5,1138.8) -- (558.9,1139.3) -- (558.2,1140.0) -- (557.2,1141.0) -- (556.0,1142.2) -- (555.1,1143.2) -- (554.3,1143.9) -- (553.8,1144.5) -- (553.5,1144.8) -- (552.9,1145.0) -- (552.0,1145.2) -- (550.8,1145.2) -- (549.2,1145.2) -- (547.8,1145.2) -- (546.2,1145.2) -- (544.8,1145.2) -- (543.2,1145.2) -- (541.8,1145.5) -- (540.5,1146.1) -- (539.2,1147.0) -- (538.0,1148.2) -- (537.0,1149.2) -- (536.3,1149.9) -- (535.7,1150.4) -- (535.4,1150.7) -- (534.8,1151.0) -- (533.9,1151.1) -- (532.8,1151.2) -- (531.2,1151.2) -- (529.8,1151.2) -- (528.2,1151.1) -- (526.8,1151.1) -- (525.2,1151.1) -- (523.8,1151.1) -- (522.2,1151.1) -- (520.8,1151.1) -- (519.2,1151.1) -- (517.8,1151.1) -- (516.2,1151.1) -- (514.8,1151.0) -- (513.2,1151.0) -- (511.8,1151.0) -- (510.2,1151.0) -- (508.8,1151.0) -- (507.2,1151.0) -- (505.8,1151.0) -- (504.2,1151.0) -- (502.8,1150.9) -- (501.2,1150.9) -- (499.8,1150.9) -- (498.2,1150.9) -- (496.8,1150.9) -- (495.2,1150.9) -- (493.8,1150.8) -- (492.2,1150.8) -- (490.8,1150.8) -- (489.2,1150.8) -- (487.8,1150.8) -- (486.2,1150.8) -- (484.8,1150.7) -- (483.2,1150.7) -- (481.8,1150.7) -- (480.2,1150.7) -- (478.8,1150.7) -- (477.2,1150.6) -- (476.1,1150.6) -- (475.3,1150.5) -- (474.8,1150.4) -- (474.6,1150.2) -- (474.2,1149.8) -- (473.4,1149.0) -- (472.4,1147.9) -- (471.0,1146.5) -- (469.6,1145.5) -- (468.2,1144.7) -- (466.8,1144.4) -- (465.2,1144.3) -- (464.1,1144.3) -- (463.3,1144.3) -- (462.9,1144.2) -- (462.8,1144.1) -- (462.4,1143.7) -- (461.6,1142.9) -- (460.6,1141.8) -- (459.1,1140.3) -- (458.1,1139.2) -- (457.3,1138.4) -- (457.0,1138.0) -- (457.0,1138.0) -- (457.0,1137.9) -- (457.0,1137.6) -- (457.0,1137.1) -- (457.0,1136.5) -- (457.2,1135.8) -- (457.6,1135.1) -- (458.2,1134.2) -- (459.0,1133.3) -- (459.8,1132.5) -- (460.6,1131.7) -- (461.3,1130.9) -- (462.0,1130.1) -- (462.7,1129.3) -- (463.3,1128.6) -- (464.0,1127.8) -- (464.7,1127.1) -- (465.4,1126.3) -- (466.1,1125.5) -- (466.9,1124.6) -- (467.7,1123.7) -- (468.5,1122.8) -- (469.2,1122.1) -- (469.7,1121.4) -- (470.2,1120.8) -- (470.9,1120.1) -- (471.6,1119.3) -- (472.5,1118.3) -- (473.5,1117.2) -- (474.3,1116.2) -- (475.0,1115.5) -- (475.5,1114.9) -- (475.8,1114.5) -- (476.3,1113.9) -- (477.1,1113.1) -- (478.1,1112.0) -- (479.2,1110.6) -- (480.2,1109.6) -- (480.8,1108.8) -- (481.2,1108.4) -- (481.3,1108.2) -- (481.8,1107.7) -- (482.5,1106.9) -- (483.5,1105.8) -- (484.8,1104.2) -- (485.8,1103.1) -- (486.5,1102.3) -- (486.8,1101.9) -- (486.9,1101.8) -- (487.3,1101.4) -- (487.9,1100.7) -- (488.9,1099.6) -- (490.1,1098.1) -- (491.1,1097.0) -- (491.8,1096.1) -- (492.3,1095.5) -- (492.6,1095.2) -- (493.1,1094.7) -- (493.7,1094.0) -- (494.5,1093.0) -- (495.6,1091.8) -- (496.5,1090.8) -- (497.2,1089.9) -- (497.8,1089.2) -- (498.3,1088.6) -- (498.9,1088.0) -- (499.5,1087.2) -- (500.2,1086.4) -- (501.1,1085.4) -- (501.8,1084.5) -- (502.6,1083.6) -- (503.3,1082.8) -- (504.0,1082.0) -- (504.6,1081.3) -- (505.3,1080.5) -- (505.9,1079.8) -- (506.5,1079.1) -- (507.2,1078.3) -- (508.0,1077.4) -- (508.8,1076.5) -- (509.7,1075.5) -- (510.4,1074.6) -- (511.1,1073.8) -- (511.6,1073.2) -- (512.1,1072.7) -- (512.7,1072.1) -- (513.4,1071.2) -- (514.3,1070.2) -- (515.4,1069.0) -- (516.3,1068.0) -- (516.9,1067.3) -- (517.4,1066.7) -- (517.7,1066.4) -- (518.2,1065.9) -- (518.9,1065.1) -- (519.9,1064.0) -- (521.2,1062.7) -- (522.1,1061.7) -- (522.8,1060.9) -- (523.3,1060.4) -- (523.4,1060.3) -- (523.9,1059.8) -- (524.6,1059.0) -- (525.7,1058.0) -- (527.0,1056.6) -- (528.0,1055.6) -- (528.8,1054.8) -- (529.2,1054.4) -- (529.4,1054.2) -- (529.9,1053.8) -- (530.6,1053.1) -- (531.7,1052.2) -- (533.0,1050.9) -- (534.1,1049.9) -- (534.9,1049.2) -- (535.5,1048.7) -- (535.8,1048.4) -- (536.3,1048.0) -- (537.1,1047.4) -- (538.1,1046.6) -- (539.2,1045.6) -- (540.3,1044.8) -- (541.3,1044.0) -- (542.2,1043.4) -- (543.0,1042.8) -- (543.7,1042.3) -- (544.5,1041.8) -- (545.2,1041.3) -- (545.9,1040.8) -- (546.8,1040.2) -- (547.9,1039.6) -- (549.2,1038.8) -- (550.8,1038.0) -- (552.0,1037.3) -- (553.0,1036.8) -- (553.7,1036.5) -- (554.1,1036.3) -- (554.5,1036.1) -- (554.7,1036.0) -- (554.8,1036.0) -- cycle ;
\draw[cloudline] (567.4,526.0) -- (567.7,526.0) -- (568.1,526.0) -- (568.8,525.9) -- (569.7,525.9) -- (570.7,525.8) -- (571.9,525.8) -- (573.2,525.7) -- (574.8,525.7) -- (576.2,525.6) -- (577.8,525.6) -- (579.2,525.5) -- (580.8,525.4) -- (582.2,525.4) -- (583.8,525.3) -- (585.2,525.3) -- (586.8,525.3) -- (588.2,525.2) -- (589.8,525.2) -- (591.2,525.1) -- (592.8,525.1) -- (594.2,525.1) -- (595.8,525.0) -- (597.2,525.0) -- (598.8,525.0) -- (600.2,525.0) -- (601.8,524.9) -- (603.2,524.9) -- (604.8,524.9) -- (606.2,524.9) -- (607.8,524.8) -- (609.2,524.8) -- (610.8,524.8) -- (612.2,524.8) -- (613.8,524.8) -- (615.2,524.8) -- (616.8,524.7) -- (618.2,524.7) -- (619.8,524.7) -- (621.2,524.7) -- (622.8,524.7) -- (624.2,524.7) -- (625.8,524.7) -- (627.2,524.7) -- (628.8,524.6) -- (630.2,524.6) -- (631.8,524.6) -- (633.2,524.6) -- (634.8,524.6) -- (636.2,524.6) -- (637.8,524.6) -- (639.2,524.6) -- (640.8,524.6) -- (642.2,524.6) -- (643.8,524.6) -- (645.2,524.6) -- (646.8,524.6) -- (648.2,524.6) -- (649.8,524.6) -- (651.2,524.6) -- (652.8,524.6) -- (654.2,524.6) -- (655.8,524.6) -- (657.2,524.6) -- (658.8,524.6) -- (660.2,524.6) -- (661.8,524.6) -- (663.2,524.6) -- (664.8,524.6) -- (666.2,524.6) -- (667.8,524.6) -- (669.2,524.6) -- (670.8,524.6) -- (672.2,524.6) -- (673.8,524.6) -- (675.2,524.6) -- (676.8,524.6) -- (678.2,524.6) -- (679.8,524.6) -- (681.2,524.6) -- (682.8,524.6) -- (684.2,524.6) -- (685.8,524.6) -- (687.2,524.7) -- (688.8,524.7) -- (690.0,524.8) -- (690.9,524.9) -- (691.5,525.2) -- (691.8,525.5) -- (692.4,526.0) -- (693.1,526.8) -- (694.1,527.8) -- (695.2,528.9) -- (696.2,529.9) -- (696.9,530.7) -- (697.5,531.2) -- (697.8,531.5) -- (698.3,532.1) -- (699.1,532.8) -- (700.1,533.8) -- (701.2,535.0) -- (702.5,535.9) -- (703.8,536.5) -- (705.2,536.8) -- (706.8,536.8) -- (708.2,536.5) -- (709.5,535.9) -- (710.8,535.0) -- (712.0,533.8) -- (713.0,532.8) -- (713.8,532.1) -- (714.3,531.6) -- (714.6,531.3) -- (715.2,531.1) -- (716.1,531.0) -- (717.2,530.9) -- (718.8,530.9) -- (720.2,530.9) -- (721.8,530.9) -- (723.2,531.0) -- (724.8,531.0) -- (726.2,531.0) -- (727.8,531.0) -- (729.2,531.0) -- (730.8,531.0) -- (732.2,531.1) -- (733.8,531.1) -- (735.2,531.1) -- (736.8,531.1) -- (738.2,531.1) -- (739.8,531.2) -- (741.2,531.2) -- (742.8,531.2) -- (744.2,530.9) -- (745.6,530.2) -- (747.0,529.3) -- (748.3,528.0) -- (749.4,526.9) -- (750.1,526.2) -- (750.6,525.7) -- (750.7,525.6) -- (751.2,525.5) -- (752.1,525.4) -- (753.2,525.3) -- (754.8,525.4) -- (756.2,525.4) -- (757.8,525.4) -- (759.2,525.4) -- (760.8,525.4) -- (762.2,525.5) -- (763.8,525.5) -- (765.2,525.5) -- (766.8,525.5) -- (768.2,525.5) -- (769.8,525.5) -- (771.2,525.6) -- (772.8,525.6) -- (774.2,525.6) -- (775.8,525.6) -- (777.2,525.6) -- (778.8,525.6) -- (780.2,525.7) -- (781.8,525.7) -- (783.2,525.7) -- (784.8,525.7) -- (786.2,525.7) -- (787.8,525.7) -- (789.2,525.7) -- (790.8,525.7) -- (792.2,525.8) -- (793.8,525.8) -- (795.2,525.8) -- (796.8,525.8) -- (798.2,525.8) -- (799.8,525.8) -- (801.2,525.8) -- (802.8,525.8) -- (804.2,525.8) -- (805.8,525.9) -- (807.2,525.9) -- (808.8,525.9) -- (810.2,525.9) -- (811.8,525.9) -- (813.2,525.9) -- (814.8,525.9) -- (816.2,525.9) -- (817.8,525.9) -- (819.2,525.9) -- (820.8,525.9) -- (822.2,526.0) -- (823.8,526.0) -- (825.2,526.0) -- (826.8,526.0) -- (828.2,526.0) -- (829.6,526.0) -- (830.9,526.0) -- (832.2,526.0) -- (833.2,526.3) -- (833.9,526.9) -- (834.4,527.7) -- (834.6,528.9) -- (834.8,529.8) -- (835.0,530.5) -- (835.1,531.1) -- (835.2,531.5) -- (835.4,532.1) -- (835.6,533.0) -- (836.0,534.2) -- (836.4,535.8) -- (836.9,537.2) -- (837.4,538.8) -- (837.9,540.2) -- (838.5,541.8) -- (839.0,543.1) -- (839.5,544.3) -- (840.0,545.4) -- (840.4,546.3) -- (840.8,547.1) -- (841.1,547.8) -- (841.5,548.5) -- (841.8,549.1) -- (842.3,549.9) -- (843.0,551.0) -- (843.8,552.2) -- (844.9,553.8) -- (845.7,554.9) -- (846.3,555.8) -- (846.6,556.3) -- (846.8,556.4) -- (847.2,556.9) -- (848.0,557.7) -- (849.2,558.7) -- (850.6,560.0) -- (851.7,561.0) -- (852.4,561.7) -- (852.8,562.1) -- (852.9,562.1) -- (853.3,562.3) -- (854.1,562.8) -- (855.2,563.4) -- (856.8,564.2) -- (858.2,564.9) -- (859.8,565.6) -- (861.2,566.1) -- (862.8,566.6) -- (864.1,567.0) -- (865.4,567.4) -- (866.6,567.6) -- (867.6,567.8) -- (868.5,567.9) -- (869.3,568.0) -- (869.9,568.1) -- (870.3,568.1) -- (871.0,568.2) -- (872.0,568.2) -- (873.2,568.2) -- (874.8,568.2) -- (876.0,568.2) -- (876.9,568.2) -- (877.6,568.1) -- (878.1,568.1) -- (878.7,568.0) -- (879.4,567.9) -- (880.3,567.8) -- (881.4,567.7) -- (882.6,567.5) -- (883.9,567.2) -- (885.2,566.9) -- (886.8,566.6) -- (888.2,566.2) -- (889.8,565.8) -- (891.2,565.4) -- (892.8,565.0) -- (894.2,564.5) -- (895.8,564.0) -- (897.2,563.5) -- (898.8,563.0) -- (899.9,562.5) -- (900.7,562.3) -- (901.1,562.1) -- (901.2,562.1) -- (901.7,561.9) -- (902.4,561.5) -- (903.5,561.0) -- (904.9,560.4) -- (906.3,559.7) -- (907.8,559.0) -- (909.2,558.3) -- (910.8,557.6) -- (911.9,557.4) -- (912.6,557.7) -- (913.0,558.6) -- (913.0,559.9) -- (913.0,561.3) -- (913.0,562.8) -- (913.0,564.2) -- (913.1,565.8) -- (913.1,567.2) -- (913.1,568.8) -- (913.1,570.2) -- (913.1,571.8) -- (913.1,573.2) -- (913.1,574.8) -- (913.1,576.2) -- (913.1,577.8) -- (913.2,579.2) -- (913.2,580.8) -- (913.2,582.2) -- (913.2,583.8) -- (913.2,585.2) -- (913.2,586.8) -- (913.2,588.2) -- (913.2,589.8) -- (913.3,591.2) -- (913.3,592.8) -- (913.3,594.2) -- (913.3,595.8) -- (913.3,597.2) -- (913.3,598.8) -- (913.4,600.2) -- (913.4,601.8) -- (913.4,603.2) -- (913.4,604.8) -- (913.4,606.2) -- (913.4,607.8) -- (913.5,609.2) -- (913.5,610.8) -- (913.5,612.2) -- (913.5,613.8) -- (913.9,615.2) -- (914.6,616.6) -- (915.6,618.0) -- (917.0,619.3) -- (918.4,620.3) -- (919.8,621.0) -- (921.2,621.3) -- (922.8,621.3) -- (924.2,621.3) -- (925.8,621.3) -- (927.2,621.3) -- (928.8,621.3) -- (930.2,621.3) -- (931.8,621.3) -- (933.2,621.2) -- (934.8,621.2) -- (935.9,621.3) -- (936.8,621.4) -- (937.3,621.5) -- (937.5,621.7) -- (938.0,622.2) -- (938.7,622.9) -- (939.7,623.9) -- (941.0,625.2) -- (942.4,626.1) -- (943.8,626.8) -- (945.2,627.1) -- (946.8,627.1) -- (947.9,627.1) -- (948.8,627.2) -- (949.4,627.4) -- (949.6,627.6) -- (950.1,628.1) -- (950.8,628.9) -- (951.8,629.9) -- (953.1,631.1) -- (954.4,632.0) -- (955.8,632.6) -- (957.2,632.9) -- (958.8,632.9) -- (959.9,633.0) -- (960.8,633.1) -- (961.4,633.3) -- (961.7,633.6) -- (962.2,634.1) -- (962.9,634.8) -- (963.9,635.8) -- (965.2,637.0) -- (966.2,638.0) -- (966.9,638.8) -- (967.4,639.3) -- (967.7,639.6) -- (968.2,640.1) -- (969.0,640.8) -- (970.0,641.8) -- (971.2,643.0) -- (972.2,644.0) -- (972.9,644.7) -- (973.5,645.2) -- (973.8,645.5) -- (974.3,646.1) -- (975.0,646.8) -- (976.0,647.8) -- (977.2,648.9) -- (978.2,649.9) -- (978.9,650.7) -- (979.5,651.2) -- (979.8,651.5) -- (980.3,652.1) -- (981.1,652.8) -- (982.1,653.8) -- (983.2,654.9) -- (984.2,655.9) -- (985.0,656.6) -- (985.5,657.2) -- (985.8,657.5) -- (986.1,658.1) -- (986.3,659.0) -- (986.3,660.2) -- (986.3,661.8) -- (986.6,663.2) -- (987.2,664.5) -- (988.1,665.7) -- (989.3,666.9) -- (990.2,667.8) -- (991.0,668.6) -- (991.5,669.1) -- (991.9,669.5) -- (992.1,670.1) -- (992.3,671.0) -- (992.4,672.2) -- (992.4,673.8) -- (992.4,675.2) -- (992.4,676.8) -- (992.4,678.2) -- (992.4,679.8) -- (992.4,681.2) -- (992.4,682.8) -- (992.4,684.2) -- (992.4,685.8) -- (992.4,687.2) -- (992.4,688.8) -- (992.4,690.2) -- (992.4,691.8) -- (992.4,693.2) -- (992.4,694.8) -- (992.4,696.2) -- (992.4,697.8) -- (992.4,699.2) -- (992.4,700.8) -- (992.4,702.2) -- (992.4,703.8) -- (992.4,705.2) -- (992.4,706.8) -- (992.4,708.2) -- (992.4,709.8) -- (992.4,711.2) -- (992.4,712.8) -- (992.4,714.2) -- (992.4,715.8) -- (992.4,717.2) -- (992.4,718.8) -- (992.4,720.2) -- (992.4,721.8) -- (992.4,723.2) -- (992.4,724.8) -- (992.4,726.2) -- (992.4,727.8) -- (992.4,729.2) -- (992.4,730.8) -- (992.4,732.2) -- (992.4,733.8) -- (992.4,735.2) -- (992.4,736.8) -- (992.4,738.2) -- (992.4,739.8) -- (992.4,741.2) -- (992.4,742.8) -- (992.4,744.2) -- (992.4,745.8) -- (992.3,747.2) -- (992.3,748.8) -- (992.3,750.2) -- (992.3,751.8) -- (992.2,753.0) -- (992.1,753.9) -- (991.8,754.5) -- (991.5,754.8) -- (991.0,755.4) -- (990.2,756.1) -- (989.2,757.1) -- (988.0,758.2) -- (987.1,759.5) -- (986.5,760.8) -- (986.2,762.2) -- (986.2,763.8) -- (986.1,765.0) -- (986.0,765.9) -- (985.8,766.5) -- (985.5,766.8) -- (984.9,767.3) -- (984.2,768.0) -- (983.2,769.0) -- (982.0,770.2) -- (981.0,771.2) -- (980.2,771.9) -- (979.7,772.4) -- (979.4,772.7) -- (978.9,773.2) -- (978.1,774.0) -- (977.1,775.0) -- (975.9,776.2) -- (974.9,777.1) -- (974.1,777.9) -- (973.6,778.4) -- (973.4,778.6) -- (972.9,779.1) -- (972.1,779.9) -- (971.1,780.9) -- (969.8,782.1) -- (968.8,783.1) -- (968.0,783.8) -- (967.5,784.3) -- (967.3,784.5) -- (966.8,785.0) -- (966.0,785.8) -- (965.0,786.8) -- (963.7,788.1) -- (962.6,789.1) -- (961.9,789.8) -- (961.4,790.2) -- (961.2,790.4) -- (960.7,790.5) -- (959.9,790.6) -- (958.8,790.6) -- (957.2,790.6) -- (955.8,790.9) -- (954.3,791.6) -- (952.9,792.6) -- (951.4,794.0) -- (950.3,795.0) -- (949.6,795.7) -- (949.2,796.1) -- (949.1,796.2) -- (948.7,796.2) -- (947.9,796.3) -- (946.8,796.2) -- (945.2,796.2) -- (943.8,796.3) -- (942.2,796.5) -- (940.8,796.9) -- (939.2,797.4) -- (937.9,797.6) -- (936.6,797.6) -- (935.4,797.3) -- (934.4,796.8) -- (933.5,796.3) -- (932.7,796.0) -- (932.1,795.7) -- (931.7,795.4) -- (931.0,795.1) -- (930.0,794.6) -- (928.8,794.0) -- (927.2,793.3) -- (925.8,792.6) -- (924.5,792.0) -- (923.2,791.4) -- (922.0,790.8) -- (921.0,790.4) -- (920.2,790.0) -- (919.7,789.8) -- (919.4,789.7) -- (918.8,789.4) -- (917.9,789.0) -- (916.8,788.5) -- (915.2,787.8) -- (913.8,787.1) -- (912.2,786.5) -- (910.8,785.9) -- (909.2,785.2) -- (908.1,784.7) -- (907.2,784.4) -- (906.7,784.2) -- (906.5,784.1) -- (906.0,783.9) -- (905.3,783.6) -- (904.3,783.2) -- (903.0,782.7) -- (901.6,782.2) -- (900.2,781.6) -- (898.8,781.1) -- (897.2,780.5) -- (895.8,780.0) -- (894.5,779.5) -- (893.2,779.0) -- (892.0,778.6) -- (891.0,778.3) -- (890.2,778.0) -- (889.7,777.8) -- (889.4,777.7) -- (888.8,777.5) -- (887.9,777.3) -- (886.8,776.9) -- (885.2,776.4) -- (883.8,775.9) -- (882.2,775.4) -- (880.8,775.0) -- (879.2,774.5) -- (877.8,774.1) -- (876.2,773.7) -- (874.8,773.2) -- (873.2,772.8) -- (872.1,772.5) -- (871.2,772.3) -- (870.7,772.1) -- (870.4,772.1) -- (869.9,772.0) -- (869.2,771.8) -- (868.2,771.5) -- (866.9,771.2) -- (865.6,770.9) -- (864.2,770.5) -- (862.8,770.2) -- (861.2,769.8) -- (859.8,769.5) -- (858.2,769.2) -- (856.8,768.9) -- (855.2,768.5) -- (853.8,768.2) -- (852.2,767.9) -- (850.8,767.7) -- (849.2,767.4) -- (847.8,767.1) -- (846.3,766.9) -- (844.9,766.6) -- (843.4,766.4) -- (842.3,766.2) -- (841.6,766.1) -- (841.2,766.0) -- (841.1,766.0) -- (840.7,765.9) -- (839.9,765.8) -- (838.8,765.6) -- (837.2,765.3) -- (835.8,765.1) -- (834.2,764.9) -- (832.8,764.7) -- (831.2,764.4) -- (829.8,764.2) -- (828.2,764.0) -- (826.8,763.8) -- (825.2,763.6) -- (823.8,763.4) -- (822.2,763.2) -- (820.8,763.0) -- (819.2,762.8) -- (817.8,762.6) -- (816.2,762.4) -- (814.8,762.2) -- (813.2,762.0) -- (811.8,761.8) -- (810.2,761.6) -- (808.8,761.5) -- (807.2,761.3) -- (805.8,761.1) -- (804.2,760.9) -- (802.8,760.7) -- (801.2,760.5) -- (800.0,760.4) -- (799.0,760.2) -- (798.3,760.1) -- (797.8,760.1) -- (797.2,760.0) -- (796.5,759.9) -- (795.6,759.8) -- (794.5,759.7) -- (793.4,759.5) -- (792.1,759.3) -- (790.8,759.1) -- (789.2,758.9) -- (787.8,758.7) -- (786.2,758.5) -- (784.8,758.3) -- (783.2,758.1) -- (781.8,757.9) -- (780.2,757.7) -- (778.8,757.5) -- (777.2,757.3) -- (775.8,757.2) -- (774.2,757.0) -- (772.8,756.8) -- (771.2,756.6) -- (769.8,756.4) -- (768.2,756.3) -- (766.8,756.1) -- (765.2,755.9) -- (763.8,755.8) -- (762.2,755.6) -- (760.8,755.5) -- (759.2,755.4) -- (757.8,755.3) -- (756.2,755.2) -- (754.8,755.1) -- (753.2,755.1) -- (751.8,755.0) -- (750.2,755.0) -- (748.8,755.0) -- (747.2,755.0) -- (745.8,755.1) -- (744.2,755.1) -- (742.8,755.2) -- (741.2,755.3) -- (739.8,755.4) -- (738.2,755.6) -- (736.8,755.8) -- (735.2,756.0) -- (733.8,756.3) -- (732.2,756.6) -- (730.8,756.9) -- (729.2,757.3) -- (727.8,757.6) -- (726.2,758.1) -- (724.8,758.5) -- (723.2,759.0) -- (722.1,759.4) -- (721.3,759.7) -- (720.8,759.9) -- (720.7,759.9) -- (720.2,760.1) -- (719.5,760.4) -- (718.4,760.9) -- (717.0,761.5) -- (715.6,762.2) -- (714.2,762.9) -- (712.8,763.6) -- (711.2,764.5) -- (710.1,765.1) -- (709.3,765.5) -- (708.8,765.8) -- (708.7,765.9) -- (708.3,766.2) -- (707.5,766.7) -- (706.5,767.4) -- (705.1,768.3) -- (703.9,769.2) -- (702.8,769.9) -- (701.9,770.6) -- (701.2,771.2) -- (700.4,771.7) -- (699.7,772.3) -- (698.9,772.9) -- (698.2,773.6) -- (697.3,774.3) -- (696.5,775.0) -- (695.6,775.8) -- (694.6,776.7) -- (693.8,777.5) -- (693.0,778.2) -- (692.4,778.8) -- (691.8,779.3) -- (691.2,780.0) -- (690.4,780.7) -- (689.6,781.6) -- (688.7,782.5) -- (687.8,783.4) -- (687.1,784.2) -- (686.4,785.0) -- (685.9,785.6) -- (685.2,786.3) -- (684.6,787.1) -- (683.9,787.9) -- (683.2,788.7) -- (682.4,789.6) -- (681.7,790.5) -- (680.9,791.4) -- (680.1,792.4) -- (679.5,793.2) -- (678.9,794.0) -- (678.3,794.6) -- (677.9,795.2) -- (677.3,795.9) -- (676.6,796.9) -- (675.7,798.1) -- (674.6,799.6) -- (673.8,800.7) -- (673.2,801.4) -- (672.9,801.8) -- (672.9,801.9) -- (672.5,802.3) -- (672.0,803.1) -- (671.1,804.2) -- (670.0,805.8) -- (669.1,807.0) -- (668.3,808.0) -- (667.8,808.7) -- (667.5,809.1) -- (667.0,809.7) -- (666.5,810.5) -- (665.8,811.4) -- (664.9,812.4) -- (664.2,813.4) -- (663.4,814.4) -- (662.7,815.3) -- (662.0,816.1) -- (661.4,816.9) -- (660.7,817.7) -- (660.1,818.4) -- (659.6,819.0) -- (658.9,819.8) -- (658.1,820.6) -- (657.3,821.5) -- (656.4,822.4) -- (655.5,823.3) -- (654.8,824.1) -- (654.1,824.7) -- (653.5,825.2) -- (652.8,825.8) -- (652.1,826.4) -- (651.2,827.1) -- (650.3,827.8) -- (649.3,828.5) -- (648.2,829.3) -- (646.9,830.0) -- (645.6,830.8) -- (644.5,831.4) -- (643.8,831.9) -- (643.3,832.1) -- (643.2,832.2) -- (642.7,832.3) -- (641.9,832.5) -- (640.8,832.8) -- (639.2,833.2) -- (637.8,833.4) -- (636.2,833.6) -- (634.8,833.6) -- (633.2,833.5) -- (632.0,833.3) -- (630.9,833.1) -- (630.0,832.8) -- (629.3,832.5) -- (628.6,832.1) -- (627.9,831.7) -- (627.1,831.3) -- (626.2,830.8) -- (625.4,830.2) -- (624.5,829.4) -- (623.7,828.5) -- (622.8,827.5) -- (621.9,826.6) -- (621.2,825.7) -- (620.5,824.8) -- (619.9,824.0) -- (619.4,823.2) -- (618.9,822.5) -- (618.5,821.8) -- (618.1,821.1) -- (617.6,820.2) -- (617.1,819.1) -- (616.5,817.8) -- (615.8,816.2) -- (615.1,814.9) -- (614.6,813.7) -- (614.1,812.7) -- (613.6,811.8) -- (613.3,811.0) -- (612.9,810.2) -- (612.7,809.5) -- (612.4,808.9) -- (612.1,808.1) -- (611.7,807.0) -- (611.2,805.8) -- (610.6,804.2) -- (610.0,802.8) -- (609.4,801.2) -- (608.8,799.8) -- (608.1,798.2) -- (607.6,797.1) -- (607.3,796.3) -- (607.1,795.9) -- (607.1,795.8) -- (606.9,795.4) -- (606.6,794.6) -- (606.1,793.5) -- (605.6,792.1) -- (605.0,790.7) -- (604.3,789.2) -- (603.7,787.8) -- (603.0,786.2) -- (602.4,785.0) -- (601.9,784.1) -- (601.6,783.4) -- (601.3,782.9) -- (601.0,782.3) -- (600.7,781.6) -- (600.3,780.7) -- (599.8,779.6) -- (599.2,778.4) -- (598.5,777.1) -- (597.8,775.8) -- (597.0,774.2) -- (596.4,773.0) -- (595.9,772.1) -- (595.5,771.5) -- (595.3,771.2) -- (595.0,770.6) -- (594.6,769.9) -- (594.1,768.9) -- (593.4,767.7) -- (592.7,766.5) -- (592.0,765.2) -- (591.2,763.9) -- (590.3,762.4) -- (589.6,761.4) -- (589.2,760.6) -- (588.9,760.2) -- (588.9,760.1) -- (588.6,759.7) -- (588.1,758.9) -- (587.4,757.8) -- (586.5,756.2) -- (585.7,754.9) -- (585.0,753.8) -- (584.3,752.8) -- (583.8,752.0) -- (583.3,751.3) -- (582.8,750.5) -- (582.3,749.8) -- (581.9,749.1) -- (581.3,748.2) -- (580.5,747.1) -- (579.6,745.8) -- (578.6,744.2) -- (577.8,743.1) -- (577.3,742.4) -- (577.0,742.0) -- (577.0,742.0) -- (576.8,741.6) -- (576.2,740.8) -- (575.5,739.7) -- (574.5,738.2) -- (573.6,736.9) -- (572.8,735.9) -- (572.2,735.0) -- (571.7,734.3) -- (571.2,733.6) -- (570.7,732.8) -- (570.2,732.0) -- (569.6,731.2) -- (569.0,730.3) -- (568.2,729.2) -- (567.4,728.0) -- (566.4,726.7) -- (565.7,725.6) -- (565.2,724.9) -- (564.8,724.4) -- (564.7,724.2) -- (564.4,723.7) -- (563.8,722.9) -- (563.0,721.8) -- (562.0,720.2) -- (561.2,719.0) -- (560.5,718.0) -- (560.0,717.2) -- (559.6,716.6) -- (559.1,715.9) -- (558.7,715.2) -- (558.1,714.3) -- (557.5,713.4) -- (556.8,712.3) -- (556.1,711.1) -- (555.3,709.8) -- (554.4,708.4) -- (553.7,707.3) -- (553.2,706.5) -- (553.0,706.1) -- (552.9,706.1) -- (552.7,705.7) -- (552.3,704.9) -- (551.6,703.8) -- (550.8,702.2) -- (550.0,700.8) -- (549.3,699.4) -- (548.6,698.1) -- (548.0,696.9) -- (547.5,695.9) -- (547.1,695.1) -- (546.8,694.6) -- (546.7,694.4) -- (546.4,693.8) -- (546.0,692.9) -- (545.5,691.8) -- (544.9,690.2) -- (544.3,688.8) -- (543.6,687.2) -- (543.1,685.8) -- (542.5,684.2) -- (542.0,683.0) -- (541.7,682.0) -- (541.4,681.2) -- (541.2,680.6) -- (541.1,680.0) -- (540.8,679.2) -- (540.6,678.4) -- (540.3,677.4) -- (540.0,676.3) -- (539.7,675.1) -- (539.4,673.8) -- (539.1,672.2) -- (538.8,670.8) -- (538.6,669.2) -- (538.3,667.8) -- (538.1,666.2) -- (537.9,664.8) -- (537.8,663.2) -- (537.7,661.8) -- (537.6,660.2) -- (537.5,658.8) -- (537.4,657.2) -- (537.4,655.8) -- (537.4,654.2) -- (537.5,652.8) -- (537.6,651.2) -- (537.6,649.8) -- (537.8,648.2) -- (537.9,646.8) -- (538.1,645.2) -- (538.2,643.8) -- (538.4,642.2) -- (538.7,640.8) -- (538.9,639.2) -- (539.2,637.8) -- (539.5,636.2) -- (539.8,634.8) -- (540.0,633.5) -- (540.3,632.2) -- (540.6,631.1) -- (540.8,630.1) -- (541.0,629.3) -- (541.1,628.8) -- (541.2,628.5) -- (541.3,627.9) -- (541.6,627.0) -- (541.9,625.8) -- (542.3,624.2) -- (542.7,622.8) -- (543.2,621.2) -- (543.6,619.8) -- (544.1,618.2) -- (544.5,616.8) -- (545.0,615.2) -- (545.4,613.8) -- (545.9,612.2) -- (546.3,611.0) -- (546.5,610.1) -- (546.7,609.5) -- (546.8,609.2) -- (547.0,608.6) -- (547.3,607.9) -- (547.6,606.9) -- (547.9,605.8) -- (548.3,604.5) -- (548.8,603.2) -- (549.2,601.8) -- (549.7,600.2) -- (550.2,598.8) -- (550.6,597.2) -- (551.1,595.8) -- (551.6,594.2) -- (552.0,593.0) -- (552.3,591.9) -- (552.6,591.1) -- (552.7,590.5) -- (552.9,589.9) -- (553.2,589.1) -- (553.4,588.3) -- (553.7,587.4) -- (554.0,586.3) -- (554.4,585.1) -- (554.8,583.8) -- (555.2,582.2) -- (555.6,580.8) -- (556.1,579.2) -- (556.5,577.8) -- (556.9,576.2) -- (557.3,574.8) -- (557.7,573.3) -- (558.1,571.8) -- (558.4,570.4) -- (558.7,569.3) -- (558.9,568.5) -- (559.0,568.1) -- (559.0,568.1) -- (559.1,567.7) -- (559.3,566.9) -- (559.6,565.8) -- (560.0,564.2) -- (560.3,562.8) -- (560.7,561.2) -- (561.1,559.8) -- (561.4,558.2) -- (561.7,556.8) -- (562.1,555.2) -- (562.4,553.8) -- (562.7,552.2) -- (563.0,550.8) -- (563.3,549.2) -- (563.6,547.8) -- (563.9,546.2) -- (564.2,544.9) -- (564.4,543.7) -- (564.6,542.6) -- (564.8,541.7) -- (564.9,540.9) -- (565.0,540.2) -- (565.2,539.5) -- (565.3,538.9) -- (565.4,538.1) -- (565.6,537.0) -- (565.8,535.8) -- (566.1,534.2) -- (566.3,532.8) -- (566.6,531.2) -- (566.8,529.8) -- (567.1,528.2) -- (567.3,527.1) -- (567.4,526.4) -- (567.4,526.0);
\draw[cloudline] (554.8,1036.0) -- (555.1,1035.9) -- (555.6,1035.7) -- (556.4,1035.4) -- (557.4,1035.0) -- (558.6,1034.6) -- (559.9,1034.2) -- (561.2,1033.8) -- (562.8,1033.5) -- (564.2,1033.2) -- (565.8,1032.9) -- (567.2,1032.7) -- (568.8,1032.6) -- (570.2,1032.5) -- (571.8,1032.5) -- (573.2,1032.5) -- (574.8,1032.6) -- (576.2,1032.8) -- (577.8,1033.0) -- (579.2,1033.3) -- (580.8,1033.6) -- (582.2,1034.0) -- (583.6,1034.4) -- (585.0,1034.8) -- (586.4,1035.3) -- (587.4,1035.7) -- (588.2,1035.9) -- (588.6,1036.1) -- (588.8,1036.1) -- (589.3,1036.4) -- (590.1,1036.8) -- (591.2,1037.4) -- (592.8,1038.1) -- (594.2,1038.8) -- (595.5,1039.6) -- (596.7,1040.3) -- (597.8,1041.0) -- (598.7,1041.5) -- (599.5,1042.0) -- (600.0,1042.4) -- (600.4,1042.6) -- (601.1,1043.1) -- (602.0,1043.8) -- (603.2,1044.7) -- (604.8,1045.9) -- (605.9,1046.7) -- (606.7,1047.4) -- (607.2,1047.7) -- (607.3,1047.8) -- (607.7,1048.2) -- (608.5,1048.9) -- (609.5,1049.9) -- (610.9,1051.2) -- (612.0,1052.2) -- (612.8,1052.9) -- (613.3,1053.4) -- (613.6,1053.6) -- (614.1,1054.2) -- (614.8,1055.0) -- (615.8,1056.1) -- (617.1,1057.4) -- (618.1,1058.5) -- (618.8,1059.3) -- (619.2,1059.7) -- (619.3,1059.8) -- (619.6,1060.3) -- (620.3,1061.1) -- (621.2,1062.2) -- (622.4,1063.8) -- (623.3,1064.9) -- (624.0,1065.8) -- (624.5,1066.4) -- (624.7,1066.7) -- (625.0,1067.3) -- (625.5,1068.0) -- (626.2,1069.0) -- (627.0,1070.2) -- (627.8,1071.4) -- (628.5,1072.5) -- (629.3,1073.7) -- (630.0,1074.8) -- (630.5,1075.7) -- (631.0,1076.5) -- (631.3,1077.0) -- (631.5,1077.4) -- (631.9,1078.1) -- (632.4,1079.0) -- (633.1,1080.2) -- (633.9,1081.8) -- (634.6,1083.1) -- (635.2,1084.4) -- (635.8,1085.5) -- (636.3,1086.5) -- (636.7,1087.4) -- (637.0,1088.1) -- (637.3,1088.7) -- (637.6,1089.2) -- (637.9,1090.0) -- (638.3,1091.0) -- (638.9,1092.2) -- (639.5,1093.8) -- (640.1,1095.2) -- (640.7,1096.8) -- (641.3,1098.2) -- (641.9,1099.8) -- (642.3,1100.9) -- (642.7,1101.8) -- (642.8,1102.3) -- (642.9,1102.4) -- (643.1,1102.9) -- (643.3,1103.6) -- (643.7,1104.7) -- (644.1,1106.0) -- (644.6,1107.4) -- (645.0,1108.8) -- (645.5,1110.2) -- (645.9,1111.8) -- (646.4,1113.2) -- (646.8,1114.8) -- (647.3,1116.2) -- (647.7,1117.8) -- (648.1,1119.0) -- (648.4,1120.1) -- (648.6,1120.9) -- (648.8,1121.5) -- (648.9,1122.2) -- (649.1,1122.9) -- (649.3,1123.7) -- (649.6,1124.6) -- (649.7,1125.3) -- (649.7,1125.8) -- (649.6,1126.0) -- (649.3,1126.0) -- (648.8,1126.0) -- (647.9,1126.1) -- (646.8,1126.1) -- (645.2,1126.2) -- (643.8,1126.6) -- (642.3,1127.3) -- (640.9,1128.4) -- (639.6,1129.8) -- (638.5,1131.0) -- (637.7,1131.7) -- (637.3,1132.2) -- (637.2,1132.3) -- (636.7,1132.4) -- (635.9,1132.5) -- (634.8,1132.5) -- (633.2,1132.6) -- (631.8,1132.6) -- (630.2,1132.7) -- (628.8,1132.7) -- (627.2,1132.7) -- (625.8,1132.8) -- (624.2,1132.8) -- (622.8,1132.8) -- (621.2,1132.8) -- (619.8,1132.9) -- (618.2,1132.9) -- (616.8,1132.9) -- (615.2,1132.9) -- (613.8,1133.0) -- (612.2,1133.0) -- (610.8,1133.0) -- (609.2,1133.0) -- (607.8,1133.0) -- (606.2,1133.0) -- (604.8,1133.1) -- (603.2,1133.1) -- (601.8,1133.1) -- (600.2,1133.1) -- (598.8,1133.1) -- (597.2,1133.1) -- (595.8,1133.1) -- (594.2,1133.1) -- (592.8,1133.2) -- (591.2,1133.2) -- (590.1,1133.1) -- (589.2,1133.0) -- (588.6,1132.7) -- (588.3,1132.4) -- (587.7,1131.9) -- (587.0,1131.2) -- (586.0,1130.2) -- (584.8,1129.0) -- (583.5,1128.1) -- (582.2,1127.5) -- (580.8,1127.2) -- (579.2,1127.2) -- (577.8,1127.2) -- (576.2,1127.2) -- (574.8,1127.2) -- (573.2,1127.2) -- (571.8,1127.5) -- (570.5,1128.1) -- (569.2,1129.0) -- (568.0,1130.2) -- (567.0,1131.2) -- (566.3,1131.9) -- (565.8,1132.5) -- (565.5,1132.8) -- (564.9,1133.3) -- (564.2,1134.0) -- (563.2,1135.0) -- (562.0,1136.2) -- (561.0,1137.2) -- (560.3,1137.9) -- (559.8,1138.5) -- (559.5,1138.8) -- (558.9,1139.3) -- (558.2,1140.0) -- (557.2,1141.0) -- (556.0,1142.2) -- (555.1,1143.2) -- (554.3,1143.9) -- (553.8,1144.5) -- (553.5,1144.8) -- (552.9,1145.0) -- (552.0,1145.2) -- (550.8,1145.2) -- (549.2,1145.2) -- (547.8,1145.2) -- (546.2,1145.2) -- (544.8,1145.2) -- (543.2,1145.2) -- (541.8,1145.5) -- (540.5,1146.1) -- (539.2,1147.0) -- (538.0,1148.2) -- (537.0,1149.2) -- (536.3,1149.9) -- (535.7,1150.4) -- (535.4,1150.7) -- (534.8,1151.0) -- (533.9,1151.1) -- (532.8,1151.2) -- (531.2,1151.2) -- (529.8,1151.2) -- (528.2,1151.1) -- (526.8,1151.1) -- (525.2,1151.1) -- (523.8,1151.1) -- (522.2,1151.1) -- (520.8,1151.1) -- (519.2,1151.1) -- (517.8,1151.1) -- (516.2,1151.1) -- (514.8,1151.0) -- (513.2,1151.0) -- (511.8,1151.0) -- (510.2,1151.0) -- (508.8,1151.0) -- (507.2,1151.0) -- (505.8,1151.0) -- (504.2,1151.0) -- (502.8,1150.9) -- (501.2,1150.9) -- (499.8,1150.9) -- (498.2,1150.9) -- (496.8,1150.9) -- (495.2,1150.9) -- (493.8,1150.8) -- (492.2,1150.8) -- (490.8,1150.8) -- (489.2,1150.8) -- (487.8,1150.8) -- (486.2,1150.8) -- (484.8,1150.7) -- (483.2,1150.7) -- (481.8,1150.7) -- (480.2,1150.7) -- (478.8,1150.7) -- (477.2,1150.6) -- (476.1,1150.6) -- (475.3,1150.5) -- (474.8,1150.4) -- (474.6,1150.2) -- (474.2,1149.8) -- (473.4,1149.0) -- (472.4,1147.9) -- (471.0,1146.5) -- (469.6,1145.5) -- (468.2,1144.7) -- (466.8,1144.4) -- (465.2,1144.3) -- (464.1,1144.3) -- (463.3,1144.3) -- (462.9,1144.2) -- (462.8,1144.1) -- (462.4,1143.7) -- (461.6,1142.9) -- (460.6,1141.8) -- (459.1,1140.3) -- (458.1,1139.2) -- (457.3,1138.4) -- (457.0,1138.0) -- (457.0,1138.0) -- (457.0,1137.9) -- (457.0,1137.6) -- (457.0,1137.1) -- (457.0,1136.5) -- (457.2,1135.8) -- (457.6,1135.1) -- (458.2,1134.2) -- (459.0,1133.3) -- (459.8,1132.5) -- (460.6,1131.7) -- (461.3,1130.9) -- (462.0,1130.1) -- (462.7,1129.3) -- (463.3,1128.6) -- (464.0,1127.8) -- (464.7,1127.1) -- (465.4,1126.3) -- (466.1,1125.5) -- (466.9,1124.6) -- (467.7,1123.7) -- (468.5,1122.8) -- (469.2,1122.1) -- (469.7,1121.4) -- (470.2,1120.8) -- (470.9,1120.1) -- (471.6,1119.3) -- (472.5,1118.3) -- (473.5,1117.2) -- (474.3,1116.2) -- (475.0,1115.5) -- (475.5,1114.9) -- (475.8,1114.5) -- (476.3,1113.9) -- (477.1,1113.1) -- (478.1,1112.0) -- (479.2,1110.6) -- (480.2,1109.6) -- (480.8,1108.8) -- (481.2,1108.4) -- (481.3,1108.2) -- (481.8,1107.7) -- (482.5,1106.9) -- (483.5,1105.8) -- (484.8,1104.2) -- (485.8,1103.1) -- (486.5,1102.3) -- (486.8,1101.9) -- (486.9,1101.8) -- (487.3,1101.4) -- (487.9,1100.7) -- (488.9,1099.6) -- (490.1,1098.1) -- (491.1,1097.0) -- (491.8,1096.1) -- (492.3,1095.5) -- (492.6,1095.2) -- (493.1,1094.7) -- (493.7,1094.0) -- (494.5,1093.0) -- (495.6,1091.8) -- (496.5,1090.8) -- (497.2,1089.9) -- (497.8,1089.2) -- (498.3,1088.6) -- (498.9,1088.0) -- (499.5,1087.2) -- (500.2,1086.4) -- (501.1,1085.4) -- (501.8,1084.5) -- (502.6,1083.6) -- (503.3,1082.8) -- (504.0,1082.0) -- (504.6,1081.3) -- (505.3,1080.5) -- (505.9,1079.8) -- (506.5,1079.1) -- (507.2,1078.3) -- (508.0,1077.4) -- (508.8,1076.5) -- (509.7,1075.5) -- (510.4,1074.6) -- (511.1,1073.8) -- (511.6,1073.2) -- (512.1,1072.7) -- (512.7,1072.1) -- (513.4,1071.2) -- (514.3,1070.2) -- (515.4,1069.0) -- (516.3,1068.0) -- (516.9,1067.3) -- (517.4,1066.7) -- (517.7,1066.4) -- (518.2,1065.9) -- (518.9,1065.1) -- (519.9,1064.0) -- (521.2,1062.7) -- (522.1,1061.7) -- (522.8,1060.9) -- (523.3,1060.4) -- (523.4,1060.3) -- (523.9,1059.8) -- (524.6,1059.0) -- (525.7,1058.0) -- (527.0,1056.6) -- (528.0,1055.6) -- (528.8,1054.8) -- (529.2,1054.4) -- (529.4,1054.2) -- (529.9,1053.8) -- (530.6,1053.1) -- (531.7,1052.2) -- (533.0,1050.9) -- (534.1,1049.9) -- (534.9,1049.2) -- (535.5,1048.7) -- (535.8,1048.4) -- (536.3,1048.0) -- (537.1,1047.4) -- (538.1,1046.6) -- (539.2,1045.6) -- (540.3,1044.8) -- (541.3,1044.0) -- (542.2,1043.4) -- (543.0,1042.8) -- (543.7,1042.3) -- (544.5,1041.8) -- (545.2,1041.3) -- (545.9,1040.8) -- (546.8,1040.2) -- (547.9,1039.6) -- (549.2,1038.8) -- (550.8,1038.0) -- (552.0,1037.3) -- (553.0,1036.8) -- (553.7,1036.5) -- (554.1,1036.3) -- (554.5,1036.1) -- (554.7,1036.0) -- (554.8,1036.0);
\draw[black,line width=0.51pt] (476.0,583.0) -- (456.5,551.8);
\draw[black,line width=0.51pt] (456.5,551.8) -- (464.2,539.2);
\draw[black,line width=0.51pt,fill=white] (476.0,583.0) circle (7.8);
\fill[black] (476.0,583.0) -- (476.0,575.2) -- (477.0,575.2) -- (478.0,575.4) -- (479.0,575.8) -- (479.9,576.2) -- (480.8,576.8) -- (481.5,577.5) -- (482.2,578.2) -- (482.8,579.1) -- (483.2,580.0) -- (483.6,581.0) -- (483.8,582.0) -- (483.8,583.0) -- cycle;
\node[anchor=east,inner sep=0pt,font=\fontsize{5.4}{5.4}\selectfont,text=black] at (465.9,572.9) {34};
\node[anchor=east,inner sep=0pt,font=\fontsize{5.4}{5.4}\selectfont,text=black] at (465.9,601.4) {22};
\node[anchor=west,inner sep=0pt,font=\fontsize{5.4}{5.4}\selectfont,text=black] at (486.1,572.9) {110};
\draw[black,line width=0.51pt] (626.0,593.0) -- (662.7,589.8);
\draw[black,line width=0.51pt] (662.7,589.8) -- (670.4,602.4);
\draw[black,line width=0.51pt,fill=white] (626.0,593.0) circle (7.8);
\fill[black] (626.0,593.0) circle (7.8);
\node[anchor=east,inner sep=0pt,font=\fontsize{5.4}{5.4}\selectfont,text=black] at (615.9,582.9) {45};
\node[anchor=east,inner sep=0pt,font=\fontsize{5.4}{5.4}\selectfont,text=black] at (615.9,611.4) {43};
\node[anchor=west,inner sep=0pt,font=\fontsize{5.4}{5.4}\selectfont,text=black] at (636.1,582.9) {096};
\fill[black] (605.2,597.6) circle (2.0);
\fill[black] (611.8,597.6) circle (2.0);
\draw[black,line width=0.51pt] (824.0,581.0) -- (859.9,572.7);
\draw[black,line width=0.51pt] (859.9,572.7) -- (869.3,584.1);
\draw[black,line width=0.51pt] (854.5,574.0) -- (859.2,579.7);
\draw[black,line width=0.51pt,fill=white] (824.0,581.0) circle (7.8);
\fill[black] (824.0,581.0) -- (824.0,573.2) -- (825.0,573.2) -- (826.0,573.4) -- (827.0,573.8) -- (827.9,574.2) -- (828.8,574.8) -- (829.5,575.5) -- (830.2,576.2) -- (830.8,577.1) -- (831.2,578.0) -- (831.6,579.0) -- (831.8,580.0) -- (831.8,581.0) -- (831.8,582.0) -- (831.6,583.0) -- (831.2,584.0) -- (830.8,584.9) -- (830.2,585.8) -- (829.5,586.5) -- (828.8,587.2) -- (827.9,587.8) -- (827.0,588.2) -- (826.0,588.6) -- (825.0,588.8) -- (824.0,588.8) -- (823.0,588.8) -- (822.0,588.6) -- (821.0,588.2) -- (820.1,587.8) -- (819.2,587.2) -- (818.5,586.5) -- (817.8,585.8) -- (817.2,584.9) -- (816.8,584.0) -- (816.4,583.0) -- (816.2,582.0) -- (816.2,581.0) -- cycle;
\node[anchor=east,inner sep=0pt,font=\fontsize{5.4}{5.4}\selectfont,text=black] at (813.9,570.9) {39};
\node[anchor=east,inner sep=0pt,font=\fontsize{5.4}{5.4}\selectfont,text=black] at (813.9,599.4) {39};
\node[anchor=west,inner sep=0pt,font=\fontsize{5.4}{5.4}\selectfont,text=black] at (834.1,570.9) {154};
\draw[black,line width=0.51pt] (465.0,654.0) -- (441.3,625.8);
\draw[black,line width=0.51pt] (441.3,625.8) -- (447.2,612.2);
\draw[black,line width=0.51pt,fill=white] (465.0,654.0) circle (7.8);
\fill[black] (465.0,654.0) -- (465.0,646.2) -- (466.0,646.2) -- (467.0,646.4) -- (468.0,646.8) -- (468.9,647.2) -- (469.8,647.8) -- (470.5,648.5) -- (471.2,649.2) -- (471.8,650.1) -- (472.2,651.0) -- (472.6,652.0) -- (472.8,653.0) -- (472.8,654.0) -- (472.8,655.0) -- (472.6,656.0) -- (472.2,657.0) -- (471.8,657.9) -- (471.2,658.8) -- (470.5,659.5) -- (469.8,660.2) -- (468.9,660.8) -- (468.0,661.2) -- (467.0,661.6) -- (466.0,661.8) -- (465.0,661.8) -- cycle;
\node[anchor=east,inner sep=0pt,font=\fontsize{5.4}{5.4}\selectfont,text=black] at (454.9,643.9) {36};
\node[anchor=east,inner sep=0pt,font=\fontsize{5.4}{5.4}\selectfont,text=black] at (454.9,672.4) {28};
\node[anchor=west,inner sep=0pt,font=\fontsize{5.4}{5.4}\selectfont,text=black] at (475.1,643.9) {124};
\draw[black,line width=0.51pt] (544.0,662.0) -- (519.4,634.7);
\draw[black,line width=0.51pt] (519.4,634.7) -- (524.8,620.9);
\draw[black,line width=0.51pt] (523.1,638.8) -- (528.5,625.0);
\draw[black,line width=0.51pt,fill=white] (544.0,662.0) circle (7.8);
\fill[black] (544.0,662.0) circle (7.8);
\node[anchor=east,inner sep=0pt,font=\fontsize{5.4}{5.4}\selectfont,text=black] at (533.9,651.9) {40};
\node[anchor=east,inner sep=0pt,font=\fontsize{5.4}{5.4}\selectfont,text=black] at (533.9,680.4) {38};
\node[anchor=west,inner sep=0pt,font=\fontsize{5.4}{5.4}\selectfont,text=black] at (554.1,651.9) {088};
\fill[black] (526.5,666.6) circle (2.0);
\draw[black,line width=0.51pt] (738.0,624.0) -- (767.0,646.7);
\draw[black,line width=0.51pt] (767.0,646.7) -- (764.1,661.2);
\draw[black,line width=0.51pt] (762.6,643.3) -- (759.7,657.8);
\draw[black,line width=0.51pt,fill=white] (738.0,624.0) circle (7.8);
\fill[black] (738.0,624.0) circle (7.8);
\node[anchor=east,inner sep=0pt,font=\fontsize{5.4}{5.4}\selectfont,text=black] at (727.9,613.9) {39};
\node[anchor=east,inner sep=0pt,font=\fontsize{5.4}{5.4}\selectfont,text=black] at (727.9,642.4) {39};
\node[anchor=west,inner sep=0pt,font=\fontsize{5.4}{5.4}\selectfont,text=black] at (748.1,613.9) {084};
\fill[black] (717.2,628.6) circle (2.0);
\fill[black] (723.8,628.6) circle (2.0);
\draw[black,line width=0.51pt] (822.0,635.0) -- (857.0,646.4);
\draw[black,line width=0.51pt] (857.0,646.4) -- (859.2,661.0);
\draw[black,line width=0.51pt] (851.7,644.7) -- (854.0,659.3);
\draw[black,line width=0.51pt,fill=white] (822.0,635.0) circle (7.8);
\fill[black] (822.0,635.0) -- (822.0,627.2) -- (823.0,627.2) -- (824.0,627.4) -- (825.0,627.8) -- (825.9,628.2) -- (826.8,628.8) -- (827.5,629.5) -- (828.2,630.2) -- (828.8,631.1) -- (829.2,632.0) -- (829.6,633.0) -- (829.8,634.0) -- (829.8,635.0) -- (829.8,636.0) -- (829.6,637.0) -- (829.2,638.0) -- (828.8,638.9) -- (828.2,639.8) -- (827.5,640.5) -- (826.8,641.2) -- (825.9,641.8) -- (825.0,642.2) -- (824.0,642.6) -- (823.0,642.8) -- (822.0,642.8) -- cycle;
\draw[black,line width=0.51pt] (814.2,635.0) -- (829.8,635.0);
\draw[black,line width=0.51pt] (822.0,627.2) -- (822.0,642.8);
\node[anchor=east,inner sep=0pt,font=\fontsize{5.4}{5.4}\selectfont,text=black] at (811.9,624.9) {43};
\node[anchor=east,inner sep=0pt,font=\fontsize{5.4}{5.4}\selectfont,text=black] at (811.9,653.4) {41};
\node[anchor=west,inner sep=0pt,font=\fontsize{5.4}{5.4}\selectfont,text=black] at (832.1,624.9) {132};
\draw[black,line width=0.51pt] (672.0,681.0) -- (701.0,658.3);
\draw[black,line width=0.51pt] (701.0,658.3) -- (714.4,664.7);
\draw[black,line width=0.51pt,fill=white] (672.0,681.0) circle (7.8);
\fill[black] (672.0,681.0) circle (7.8);
\node[anchor=east,inner sep=0pt,font=\fontsize{5.4}{5.4}\selectfont,text=black] at (661.9,670.9) {48};
\node[anchor=east,inner sep=0pt,font=\fontsize{5.4}{5.4}\selectfont,text=black] at (661.9,699.4) {45};
\node[anchor=west,inner sep=0pt,font=\fontsize{5.4}{5.4}\selectfont,text=black] at (682.1,670.9) {024};
\fill[black] (651.2,685.6) circle (2.0);
\fill[black] (657.8,685.6) circle (2.0);
\draw[black,line width=0.51pt] (600.0,710.0) -- (565.9,696.2);
\draw[black,line width=0.51pt] (565.9,696.2) -- (564.7,681.4);
\draw[black,line width=0.51pt,fill=white] (600.0,710.0) circle (7.8);
\fill[black] (600.0,710.0) circle (7.8);
\node[anchor=east,inner sep=0pt,font=\fontsize{5.4}{5.4}\selectfont,text=black] at (589.9,699.9) {48};
\node[anchor=west,inner sep=0pt,font=\fontsize{5.4}{5.4}\selectfont,text=black] at (610.1,699.9) {026};
\fill[black] (579.2,714.6) circle (2.0);
\fill[black] (585.8,714.6) circle (2.0);
\draw[black,line width=0.51pt] (824.0,683.0) -- (860.4,688.1);
\draw[black,line width=0.51pt] (855.0,687.4) -- (857.3,694.4);
\draw[black,line width=0.51pt,fill=white] (824.0,683.0) circle (7.8);
\fill[black] (824.0,683.0) circle (7.8);
\node[anchor=east,inner sep=0pt,font=\fontsize{5.4}{5.4}\selectfont,text=black] at (813.9,672.9) {46};
\node[anchor=east,inner sep=0pt,font=\fontsize{5.4}{5.4}\selectfont,text=black] at (813.9,701.4) {43};
\node[anchor=west,inner sep=0pt,font=\fontsize{5.4}{5.4}\selectfont,text=black] at (834.1,672.9) {124};
\fill[black] (803.2,687.6) circle (2.0);
\fill[black] (809.8,687.6) circle (2.0);
\draw[black,line width=0.51pt] (908.0,712.0) -- (943.5,721.5);
\draw[black,line width=0.51pt,fill=white] (908.0,712.0) circle (7.8);
\fill[black] (908.0,712.0) circle (7.8);
\node[anchor=east,inner sep=0pt,font=\fontsize{5.4}{5.4}\selectfont,text=black] at (897.9,701.9) {45};
\node[anchor=east,inner sep=0pt,font=\fontsize{5.4}{5.4}\selectfont,text=black] at (897.9,730.4) {43};
\node[anchor=west,inner sep=0pt,font=\fontsize{5.4}{5.4}\selectfont,text=black] at (918.1,701.9) {163};
\fill[black] (890.5,716.6) circle (2.0);
\draw[black,line width=0.51pt] (419.0,722.0) -- (394.4,694.7);
\draw[black,line width=0.51pt] (394.4,694.7) -- (399.8,680.9);
\draw[black,line width=0.51pt,fill=white] (419.0,722.0) circle (7.8);
\node[anchor=east,inner sep=0pt,font=\fontsize{5.4}{5.4}\selectfont,text=black] at (408.9,711.9) {39};
\node[anchor=east,inner sep=0pt,font=\fontsize{5.4}{5.4}\selectfont,text=black] at (408.9,740.4) {31};
\node[anchor=west,inner sep=0pt,font=\fontsize{5.4}{5.4}\selectfont,text=black] at (429.1,711.9) {124};
\draw[black,line width=0.51pt] (714.0,722.0) -- (698.4,755.4);
\draw[black,line width=0.51pt] (698.4,755.4) -- (683.6,755.8);
\draw[black,line width=0.51pt,fill=white] (714.0,722.0) circle (7.8);
\fill[black] (714.0,722.0) -- (714.0,714.2) -- (715.0,714.2) -- (716.0,714.4) -- (717.0,714.8) -- (717.9,715.2) -- (718.8,715.8) -- (719.5,716.5) -- (720.2,717.2) -- (720.8,718.1) -- (721.2,719.0) -- (721.6,720.0) -- (721.8,721.0) -- (721.8,722.0) -- (721.8,723.0) -- (721.6,724.0) -- (721.2,725.0) -- (720.8,725.9) -- (720.2,726.8) -- (719.5,727.5) -- (718.8,728.2) -- (717.9,728.8) -- (717.0,729.2) -- (716.0,729.6) -- (715.0,729.8) -- (714.0,729.8) -- (713.0,729.8) -- (712.0,729.6) -- (711.0,729.2) -- (710.1,728.8) -- (709.2,728.2) -- (708.5,727.5) -- (707.8,726.8) -- (707.2,725.9) -- (706.8,725.0) -- (706.4,724.0) -- (706.2,723.0) -- (706.2,722.0) -- cycle;
\node[anchor=east,inner sep=0pt,font=\fontsize{5.4}{5.4}\selectfont,text=black] at (703.9,711.9) {57};
\node[anchor=east,inner sep=0pt,font=\fontsize{5.4}{5.4}\selectfont,text=black] at (703.9,740.4) {54};
\node[anchor=west,inner sep=0pt,font=\fontsize{5.4}{5.4}\selectfont,text=black] at (724.1,711.9) {029};
\draw[black,line width=0.51pt] (818.0,742.0) -- (854.2,748.4);
\draw[black,line width=0.51pt] (854.2,748.4) -- (858.5,762.6);
\draw[black,line width=0.51pt,fill=white] (818.0,742.0) circle (7.8);
\fill[black] (818.0,742.0) circle (7.8);
\node[anchor=east,inner sep=0pt,font=\fontsize{5.4}{5.4}\selectfont,text=black] at (807.9,731.9) {54};
\node[anchor=east,inner sep=0pt,font=\fontsize{5.4}{5.4}\selectfont,text=black] at (807.9,760.4) {45};
\node[anchor=west,inner sep=0pt,font=\fontsize{5.4}{5.4}\selectfont,text=black] at (828.1,731.9) {097};
\draw[black,line width=0.51pt] (469.0,770.0) -- (438.5,749.4);
\draw[black,line width=0.51pt] (438.5,749.4) -- (440.4,734.7);
\draw[black,line width=0.51pt,fill=white] (469.0,770.0) circle (7.8);
\draw[black,line width=0.51pt] (469.0,762.2) -- (469.0,777.8);
\node[anchor=east,inner sep=0pt,font=\fontsize{5.4}{5.4}\selectfont,text=black] at (458.9,759.9) {41};
\node[anchor=east,inner sep=0pt,font=\fontsize{5.4}{5.4}\selectfont,text=black] at (458.9,788.4) {30};
\node[anchor=west,inner sep=0pt,font=\fontsize{5.4}{5.4}\selectfont,text=black] at (479.1,759.9) {103};
\draw[black,line width=0.51pt] (570.0,768.0) -- (544.9,741.1);
\draw[black,line width=0.51pt] (544.9,741.1) -- (550.1,727.2);
\draw[black,line width=0.51pt] (548.7,745.1) -- (553.8,731.2);
\draw[black,line width=0.51pt] (552.4,749.2) -- (555.0,742.2);
\draw[black,line width=0.51pt,fill=white] (570.0,768.0) circle (7.8);
\fill[black] (570.0,768.0) -- (570.0,760.2) -- (571.0,760.2) -- (572.0,760.4) -- (573.0,760.8) -- (573.9,761.2) -- (574.8,761.8) -- (575.5,762.5) -- (576.2,763.2) -- (576.8,764.1) -- (577.2,765.0) -- (577.6,766.0) -- (577.8,767.0) -- (577.8,768.0) -- (577.8,769.0) -- (577.6,770.0) -- (577.2,771.0) -- (576.8,771.9) -- (576.2,772.8) -- (575.5,773.5) -- (574.8,774.2) -- (573.9,774.8) -- (573.0,775.2) -- (572.0,775.6) -- (571.0,775.8) -- (570.0,775.8) -- cycle;
\node[anchor=east,inner sep=0pt,font=\fontsize{5.4}{5.4}\selectfont,text=black] at (559.9,757.9) {45};
\node[anchor=east,inner sep=0pt,font=\fontsize{5.4}{5.4}\selectfont,text=black] at (559.9,786.4) {40};
\node[anchor=west,inner sep=0pt,font=\fontsize{5.4}{5.4}\selectfont,text=black] at (580.1,757.9) {069};
\draw[black,line width=0.51pt] (637.0,765.0) -- (618.0,796.5);
\draw[black,line width=0.51pt] (618.0,796.5) -- (603.3,795.4);
\draw[black,line width=0.51pt] (620.9,791.8) -- (606.1,790.7);
\draw[black,line width=0.51pt] (623.7,787.1) -- (616.3,786.5);
\draw[black,line width=0.51pt,fill=white] (637.0,765.0) circle (7.8);
\fill[black] (637.0,765.0) circle (7.8);
\node[anchor=east,inner sep=0pt,font=\fontsize{5.4}{5.4}\selectfont,text=black] at (626.9,754.9) {60};
\node[anchor=east,inner sep=0pt,font=\fontsize{5.4}{5.4}\selectfont,text=black] at (626.9,783.4) {55};
\node[anchor=west,inner sep=0pt,font=\fontsize{5.4}{5.4}\selectfont,text=black] at (647.1,754.9) {040};
\draw[black,line width=0.43pt] (616.4,763.3) -- (616.4,775.9);
\draw[black,line width=0.43pt] (616.4,763.3) -- (621.1,763.3) .. controls (625.4,763.3) and (625.4,770.0) .. (621.1,770.0) -- (616.4,770.0);
\draw[black,line width=0.43pt] (620.3,770.0) -- (625.8,775.9);
\fill[black] (626.9,777.0) -- (623.1,776.1) -- (625.6,773.4) -- cycle;
\fill[black] (618.7,759.4) circle (1.9);
\draw[black,line width=0.51pt] (704.0,785.0) -- (685.0,816.5);
\draw[black,line width=0.51pt] (685.0,816.5) -- (670.3,815.4);
\draw[black,line width=0.51pt,fill=white] (704.0,785.0) circle (7.8);
\fill[black] (704.0,785.0) -- (704.0,777.2) -- (705.0,777.2) -- (706.0,777.4) -- (707.0,777.8) -- (707.9,778.2) -- (708.8,778.8) -- (709.5,779.5) -- (710.2,780.2) -- (710.8,781.1) -- (711.2,782.0) -- (711.6,783.0) -- (711.8,784.0) -- (711.8,785.0) -- (711.8,786.0) -- (711.6,787.0) -- (711.2,788.0) -- (710.8,788.9) -- (710.2,789.8) -- (709.5,790.5) -- (708.8,791.2) -- (707.9,791.8) -- (707.0,792.2) -- (706.0,792.6) -- (705.0,792.8) -- (704.0,792.8) -- cycle;
\node[anchor=east,inner sep=0pt,font=\fontsize{5.4}{5.4}\selectfont,text=black] at (693.9,774.9) {61};
\node[anchor=east,inner sep=0pt,font=\fontsize{5.4}{5.4}\selectfont,text=black] at (693.9,803.4) {57};
\node[anchor=west,inner sep=0pt,font=\fontsize{5.4}{5.4}\selectfont,text=black] at (714.1,774.9) {044};
\fill[black] (683.2,789.6) circle (2.0);
\fill[black] (689.8,789.6) circle (2.0);
\draw[black,line width=0.51pt] (787.0,798.0) -- (787.0,834.8);
\draw[black,line width=0.51pt] (787.0,834.8) -- (773.8,841.4);
\draw[black,line width=0.51pt,fill=white] (787.0,798.0) circle (7.8);
\fill[black] (787.0,798.0) -- (787.0,790.2) -- (788.0,790.2) -- (789.0,790.4) -- (790.0,790.8) -- (790.9,791.2) -- (791.8,791.8) -- (792.5,792.5) -- (793.2,793.2) -- (793.8,794.1) -- (794.2,795.0) -- (794.6,796.0) -- (794.8,797.0) -- (794.8,798.0) -- (794.8,799.0) -- (794.6,800.0) -- (794.2,801.0) -- (793.8,801.9) -- (793.2,802.8) -- (792.5,803.5) -- (791.8,804.2) -- (790.9,804.8) -- (790.0,805.2) -- (789.0,805.6) -- (788.0,805.8) -- (787.0,805.8) -- (786.0,805.8) -- (785.0,805.6) -- (784.0,805.2) -- (783.1,804.8) -- (782.2,804.2) -- (781.5,803.5) -- (780.8,802.8) -- (780.2,801.9) -- (779.8,801.0) -- (779.4,800.0) -- (779.2,799.0) -- (779.2,798.0) -- cycle;
\node[anchor=east,inner sep=0pt,font=\fontsize{5.4}{5.4}\selectfont,text=black] at (776.9,787.9) {66};
\node[anchor=east,inner sep=0pt,font=\fontsize{5.4}{5.4}\selectfont,text=black] at (776.9,816.4) {53};
\node[anchor=west,inner sep=0pt,font=\fontsize{5.4}{5.4}\selectfont,text=black] at (797.1,787.9) {099};
\draw[black,line width=0.51pt] (828.0,818.0) -- (828.0,854.8);
\draw[black,line width=0.51pt] (828.0,854.8) -- (814.8,861.4);
\draw[black,line width=0.51pt,fill=white] (828.0,818.0) circle (7.8);
\fill[black] (828.0,818.0) -- (828.0,810.2) -- (829.0,810.2) -- (830.0,810.4) -- (831.0,810.8) -- (831.9,811.2) -- (832.8,811.8) -- (833.5,812.5) -- (834.2,813.2) -- (834.8,814.1) -- (835.2,815.0) -- (835.6,816.0) -- (835.8,817.0) -- (835.8,818.0) -- cycle;
\node[anchor=east,inner sep=0pt,font=\fontsize{5.4}{5.4}\selectfont,text=black] at (817.9,807.9) {66};
\node[anchor=east,inner sep=0pt,font=\fontsize{5.4}{5.4}\selectfont,text=black] at (817.9,836.4) {55};
\node[anchor=west,inner sep=0pt,font=\fontsize{5.4}{5.4}\selectfont,text=black] at (838.1,807.9) {111};
\draw[black,line width=0.51pt] (405.0,838.0) -- (374.9,816.9);
\draw[black,line width=0.51pt] (374.9,816.9) -- (377.0,802.2);
\draw[black,line width=0.51pt,fill=white] (405.0,838.0) circle (7.8);
\node[anchor=east,inner sep=0pt,font=\fontsize{5.4}{5.4}\selectfont,text=black] at (394.9,827.9) {41};
\node[anchor=east,inner sep=0pt,font=\fontsize{5.4}{5.4}\selectfont,text=black] at (394.9,856.4) {33};
\node[anchor=west,inner sep=0pt,font=\fontsize{5.4}{5.4}\selectfont,text=black] at (415.1,827.9) {126};
\draw[black,line width=0.51pt] (496.0,831.0) -- (467.8,807.3);
\draw[black,line width=0.51pt] (467.8,807.3) -- (471.3,792.9);
\draw[black,line width=0.51pt,fill=white] (496.0,831.0) circle (7.8);
\node[anchor=east,inner sep=0pt,font=\fontsize{5.4}{5.4}\selectfont,text=black] at (485.9,820.9) {42};
\node[anchor=east,inner sep=0pt,font=\fontsize{5.4}{5.4}\selectfont,text=black] at (485.9,849.4) {36};
\node[anchor=west,inner sep=0pt,font=\fontsize{5.4}{5.4}\selectfont,text=black] at (506.1,820.9) {109};
\draw[black,line width=0.51pt] (608.0,861.0) -- (588.5,892.2);
\draw[black,line width=0.51pt] (588.5,892.2) -- (573.8,890.8);
\draw[black,line width=0.51pt,fill=white] (608.0,861.0) circle (7.8);
\fill[black] (608.0,861.0) circle (7.8);
\node[anchor=east,inner sep=0pt,font=\fontsize{5.4}{5.4}\selectfont,text=black] at (597.9,850.9) {61};
\node[anchor=east,inner sep=0pt,font=\fontsize{5.4}{5.4}\selectfont,text=black] at (597.9,879.4) {57};
\node[anchor=west,inner sep=0pt,font=\fontsize{5.4}{5.4}\selectfont,text=black] at (618.1,850.9) {075};
\fill[black] (587.2,865.6) circle (2.0);
\fill[black] (593.8,865.6) circle (2.0);
\draw[black,line width=0.51pt] (828.0,885.0) -- (828.0,921.8);
\draw[black,line width=0.51pt] (828.0,921.8) -- (814.8,928.4);
\draw[black,line width=0.51pt,fill=white] (828.0,885.0) circle (7.8);
\fill[black] (828.0,885.0) -- (828.0,877.2) -- (829.0,877.2) -- (830.0,877.4) -- (831.0,877.8) -- (831.9,878.2) -- (832.8,878.8) -- (833.5,879.5) -- (834.2,880.2) -- (834.8,881.1) -- (835.2,882.0) -- (835.6,883.0) -- (835.8,884.0) -- (835.8,885.0) -- (835.8,886.0) -- (835.6,887.0) -- (835.2,888.0) -- (834.8,888.9) -- (834.2,889.8) -- (833.5,890.5) -- (832.8,891.2) -- (831.9,891.8) -- (831.0,892.2) -- (830.0,892.6) -- (829.0,892.8) -- (828.0,892.8) -- cycle;
\node[anchor=east,inner sep=0pt,font=\fontsize{5.4}{5.4}\selectfont,text=black] at (817.9,874.9) {71};
\node[anchor=east,inner sep=0pt,font=\fontsize{5.4}{5.4}\selectfont,text=black] at (817.9,903.4) {60};
\node[anchor=west,inner sep=0pt,font=\fontsize{5.4}{5.4}\selectfont,text=black] at (838.1,874.9) {149};
\draw[black,line width=0.51pt] (498.0,909.0) -- (465.5,891.7);
\draw[black,line width=0.51pt] (465.5,891.7) -- (465.9,876.9);
\draw[black,line width=0.51pt,fill=white] (498.0,909.0) circle (7.8);
\node[anchor=east,inner sep=0pt,font=\fontsize{5.4}{5.4}\selectfont,text=black] at (487.9,898.9) {48};
\node[anchor=east,inner sep=0pt,font=\fontsize{5.4}{5.4}\selectfont,text=black] at (487.9,927.4) {38};
\node[anchor=west,inner sep=0pt,font=\fontsize{5.4}{5.4}\selectfont,text=black] at (508.1,898.9) {126};
\draw[black,line width=0.51pt] (729.0,914.0) -- (729.0,950.8);
\draw[black,line width=0.51pt] (729.0,950.8) -- (715.8,957.4);
\draw[black,line width=0.51pt,fill=white] (729.0,914.0) circle (7.8);
\fill[black] (729.0,914.0) -- (729.0,906.2) -- (730.0,906.2) -- (731.0,906.4) -- (732.0,906.8) -- (732.9,907.2) -- (733.8,907.8) -- (734.5,908.5) -- (735.2,909.2) -- (735.8,910.1) -- (736.2,911.0) -- (736.6,912.0) -- (736.8,913.0) -- (736.8,914.0) -- (736.8,915.0) -- (736.6,916.0) -- (736.2,917.0) -- (735.8,917.9) -- (735.2,918.8) -- (734.5,919.5) -- (733.8,920.2) -- (732.9,920.8) -- (732.0,921.2) -- (731.0,921.6) -- (730.0,921.8) -- (729.0,921.8) -- cycle;
\node[anchor=east,inner sep=0pt,font=\fontsize{5.4}{5.4}\selectfont,text=black] at (718.9,903.9) {70};
\node[anchor=east,inner sep=0pt,font=\fontsize{5.4}{5.4}\selectfont,text=black] at (718.9,932.4) {59};
\node[anchor=west,inner sep=0pt,font=\fontsize{5.4}{5.4}\selectfont,text=black] at (739.1,903.9) {140};
\draw[black,line width=0.51pt] (390.0,946.0) -- (356.6,930.4);
\draw[black,line width=0.51pt] (356.6,930.4) -- (356.2,915.6);
\draw[black,line width=0.51pt,fill=white] (390.0,946.0) circle (7.8);
\node[anchor=east,inner sep=0pt,font=\fontsize{5.4}{5.4}\selectfont,text=black] at (379.9,935.9) {49};
\node[anchor=east,inner sep=0pt,font=\fontsize{5.4}{5.4}\selectfont,text=black] at (379.9,964.4) {38};
\node[anchor=west,inner sep=0pt,font=\fontsize{5.4}{5.4}\selectfont,text=black] at (400.1,935.9) {128};
\draw[black,line width=0.51pt] (603.0,945.0) -- (573.2,966.6);
\draw[black,line width=0.51pt] (573.2,966.6) -- (560.1,959.8);
\draw[black,line width=0.51pt,fill=white] (603.0,945.0) circle (7.8);
\fill[black] (603.0,945.0) -- (603.0,937.2) -- (604.0,937.2) -- (605.0,937.4) -- (606.0,937.8) -- (606.9,938.2) -- (607.8,938.8) -- (608.5,939.5) -- (609.2,940.2) -- (609.8,941.1) -- (610.2,942.0) -- (610.6,943.0) -- (610.8,944.0) -- (610.8,945.0) -- (610.8,946.0) -- (610.6,947.0) -- (610.2,948.0) -- (609.8,948.9) -- (609.2,949.8) -- (608.5,950.5) -- (607.8,951.2) -- (606.9,951.8) -- (606.0,952.2) -- (605.0,952.6) -- (604.0,952.8) -- (603.0,952.8) -- (602.0,952.8) -- (601.0,952.6) -- (600.0,952.2) -- (599.1,951.8) -- (598.2,951.2) -- (597.5,950.5) -- (596.8,949.8) -- (596.2,948.9) -- (595.8,948.0) -- (595.4,947.0) -- (595.2,946.0) -- (595.2,945.0) -- cycle;
\node[anchor=east,inner sep=0pt,font=\fontsize{5.4}{5.4}\selectfont,text=black] at (592.9,934.9) {66};
\node[anchor=east,inner sep=0pt,font=\fontsize{5.4}{5.4}\selectfont,text=black] at (592.9,963.4) {59};
\node[anchor=west,inner sep=0pt,font=\fontsize{5.4}{5.4}\selectfont,text=black] at (613.1,934.9) {114};
\draw[black,line width=0.51pt] (648.0,960.0) -- (631.3,992.8);
\draw[black,line width=0.51pt] (631.3,992.8) -- (616.5,992.7);
\draw[black,line width=0.51pt,fill=white] (648.0,960.0) circle (7.8);
\node[anchor=east,inner sep=0pt,font=\fontsize{5.4}{5.4}\selectfont,text=black] at (637.9,949.9) {66};
\node[anchor=east,inner sep=0pt,font=\fontsize{5.4}{5.4}\selectfont,text=black] at (637.9,978.4) {60};
\node[anchor=west,inner sep=0pt,font=\fontsize{5.4}{5.4}\selectfont,text=black] at (658.1,949.9) {124};
\draw[black,line width=0.51pt] (452.0,1009.0) -- (423.8,985.3);
\draw[black,line width=0.51pt] (423.8,985.3) -- (427.3,970.9);
\draw[black,line width=0.51pt,fill=white] (452.0,1009.0) circle (7.8);
\node[anchor=east,inner sep=0pt,font=\fontsize{5.4}{5.4}\selectfont,text=black] at (441.9,998.9) {50};
\node[anchor=east,inner sep=0pt,font=\fontsize{5.4}{5.4}\selectfont,text=black] at (441.9,1027.4) {38};
\node[anchor=west,inner sep=0pt,font=\fontsize{5.4}{5.4}\selectfont,text=black] at (462.1,998.9) {120};
\draw[black,line width=0.51pt] (507.0,1061.0) -- (471.5,1070.5);
\draw[black,line width=0.51pt] (471.5,1070.5) -- (461.6,1059.4);
\draw[black,line width=0.51pt,fill=white] (507.0,1061.0) circle (7.8);
\fill[black] (507.0,1061.0) circle (7.8);
\node[anchor=east,inner sep=0pt,font=\fontsize{5.4}{5.4}\selectfont,text=black] at (496.9,1050.9) {52};
\node[anchor=east,inner sep=0pt,font=\fontsize{5.4}{5.4}\selectfont,text=black] at (496.9,1079.4) {50};
\node[anchor=west,inner sep=0pt,font=\fontsize{5.4}{5.4}\selectfont,text=black] at (517.1,1050.9) {128};
\draw[black,line width=0.51pt] (613.0,1041.0) -- (606.0,1077.1);
\draw[black,line width=0.51pt] (606.0,1077.1) -- (591.7,1081.1);
\draw[black,line width=0.51pt,fill=white] (613.0,1041.0) circle (7.8);
\fill[black] (613.0,1041.0) circle (7.8);
\node[anchor=east,inner sep=0pt,font=\fontsize{5.4}{5.4}\selectfont,text=black] at (602.9,1030.9) {63};
\node[anchor=east,inner sep=0pt,font=\fontsize{5.4}{5.4}\selectfont,text=black] at (602.9,1059.4) {61};
\node[anchor=west,inner sep=0pt,font=\fontsize{5.4}{5.4}\selectfont,text=black] at (623.1,1030.9) {153};
\draw[black,line width=0.51pt] (725.0,1031.0) -- (710.6,1064.9);
\draw[black,line width=0.51pt] (710.6,1064.9) -- (695.8,1065.8);
\draw[black,line width=0.51pt,fill=white] (725.0,1031.0) circle (7.8);
\fill[black] (725.0,1031.0) -- (725.0,1023.2) -- (726.0,1023.2) -- (727.0,1023.4) -- (728.0,1023.8) -- (728.9,1024.2) -- (729.8,1024.8) -- (730.5,1025.5) -- (731.2,1026.2) -- (731.8,1027.1) -- (732.2,1028.0) -- (732.6,1029.0) -- (732.8,1030.0) -- (732.8,1031.0) -- cycle;
\node[anchor=east,inner sep=0pt,font=\fontsize{5.4}{5.4}\selectfont,text=black] at (714.9,1020.9) {69};
\node[anchor=east,inner sep=0pt,font=\fontsize{5.4}{5.4}\selectfont,text=black] at (714.9,1049.4) {61};
\node[anchor=west,inner sep=0pt,font=\fontsize{5.4}{5.4}\selectfont,text=black] at (735.1,1020.9) {171};
\draw[black,line width=0.51pt] (822.0,1006.0) -- (806.4,1039.4);
\draw[black,line width=0.51pt] (806.4,1039.4) -- (791.6,1039.8);
\draw[black,line width=0.51pt,fill=white] (822.0,1006.0) circle (7.8);
\node[anchor=east,inner sep=0pt,font=\fontsize{5.4}{5.4}\selectfont,text=black] at (811.9,995.9) {72};
\node[anchor=east,inner sep=0pt,font=\fontsize{5.4}{5.4}\selectfont,text=black] at (811.9,1024.4) {64};
\node[anchor=west,inner sep=0pt,font=\fontsize{5.4}{5.4}\selectfont,text=black] at (832.1,995.9) {194};
}
\caption{Step 4: the area where the cloud cover exceeds 75\,\% (dashed outline and grey tint). It wraps round the low on its north and east sides and tails away to the south-west along the cold front.}
\end{figure}

\subsection{Step 5: the precipitation}

Shade round every station that reports precipitation: the weather symbol is a dot or two dots (rain), or the thunderstorm symbol with a dot. These are stations 2, 5, 6, 8, 9, 10, 11, 17, 18 and 23: a band to the north and east of the low, along the warm front, and a narrow band along the cold front in Missouri.

\begin{figure}[ht]
\centering
\mapbase{0.0196cm}{\path[precipfill,even odd rule] (541.0,526.0) -- (541.0,526.0) -- (541.0,526.0) -- (541.0,526.0) -- (541.0,526.0) -- (541.4,525.9) -- (542.1,525.7) -- (543.2,525.5) -- (544.8,525.1) -- (546.0,525.0) -- (546.9,525.0) -- (547.5,525.1) -- (547.9,525.5) -- (548.4,526.0) -- (549.1,526.6) -- (550.1,527.4) -- (551.3,528.3) -- (552.3,528.7) -- (553.3,528.7) -- (554.2,528.3) -- (554.9,527.4) -- (555.7,526.5) -- (556.4,525.7) -- (557.2,524.9) -- (557.9,524.2) -- (558.8,523.6) -- (559.9,523.2) -- (561.2,522.9) -- (562.8,522.8) -- (564.2,522.7) -- (565.8,522.6) -- (567.2,522.5) -- (568.8,522.5) -- (570.2,522.4) -- (571.8,522.3) -- (573.2,522.3) -- (574.8,522.2) -- (576.2,522.2) -- (577.8,522.2) -- (579.2,522.1) -- (580.8,522.1) -- (582.2,522.1) -- (583.8,522.0) -- (585.2,522.0) -- (586.8,522.0) -- (588.2,522.0) -- (589.8,522.0) -- (591.2,521.9) -- (592.8,521.9) -- (594.2,521.9) -- (595.8,521.9) -- (597.2,521.9) -- (598.8,521.9) -- (600.2,521.9) -- (601.8,521.9) -- (603.2,521.9) -- (604.8,521.9) -- (606.2,521.8) -- (607.8,521.8) -- (609.2,521.8) -- (610.8,521.8) -- (612.2,521.8) -- (613.8,521.8) -- (615.2,521.8) -- (616.8,521.8) -- (618.2,521.8) -- (619.8,521.8) -- (621.2,521.8) -- (622.8,521.8) -- (624.2,521.8) -- (625.8,521.8) -- (627.2,521.8) -- (628.8,521.8) -- (630.2,521.8) -- (631.8,521.8) -- (633.2,521.8) -- (634.8,521.8) -- (636.2,521.8) -- (637.8,521.8) -- (639.2,521.8) -- (640.8,521.8) -- (642.2,521.8) -- (643.8,521.8) -- (645.2,521.8) -- (646.8,521.8) -- (648.2,521.8) -- (649.8,521.8) -- (651.2,521.8) -- (652.8,521.8) -- (654.2,521.8) -- (655.8,521.8) -- (657.2,521.8) -- (658.8,521.8) -- (660.2,521.8) -- (661.8,521.8) -- (663.2,521.8) -- (664.8,521.8) -- (666.2,521.8) -- (667.8,521.8) -- (669.2,521.8) -- (670.8,521.8) -- (672.2,521.8) -- (673.8,521.8) -- (675.2,521.8) -- (676.8,521.8) -- (678.2,521.8) -- (679.8,521.8) -- (681.2,521.8) -- (682.8,521.8) -- (684.2,521.8) -- (685.8,521.9) -- (687.2,521.9) -- (688.8,521.9) -- (690.1,522.1) -- (691.4,522.7) -- (692.5,523.4) -- (693.6,524.5) -- (694.5,525.3) -- (695.2,526.1) -- (695.8,526.7) -- (696.3,527.2) -- (696.9,527.8) -- (697.6,528.6) -- (698.5,529.4) -- (699.6,530.5) -- (700.4,531.4) -- (701.2,532.1) -- (701.8,532.7) -- (702.3,533.2) -- (703.0,533.6) -- (704.0,533.9) -- (705.2,534.0) -- (706.8,534.0) -- (708.0,533.9) -- (709.0,533.7) -- (709.8,533.3) -- (710.4,532.8) -- (711.0,532.1) -- (711.8,531.4) -- (712.6,530.6) -- (713.6,529.6) -- (714.7,528.9) -- (715.9,528.5) -- (717.2,528.2) -- (718.8,528.3) -- (720.2,528.3) -- (721.8,528.4) -- (723.2,528.4) -- (724.8,528.5) -- (726.2,528.5) -- (727.8,528.6) -- (729.2,528.6) -- (730.8,528.7) -- (732.2,528.8) -- (733.8,528.9) -- (735.2,528.9) -- (736.8,529.0) -- (738.2,529.1) -- (739.8,529.2) -- (741.2,529.3) -- (742.8,529.4) -- (744.1,529.3) -- (745.4,528.9) -- (746.6,528.2) -- (747.7,527.3) -- (748.7,526.6) -- (749.4,525.9) -- (750.0,525.4) -- (750.4,525.0) -- (751.1,524.7) -- (752.0,524.6) -- (753.2,524.7) -- (754.8,524.9) -- (756.2,525.1) -- (757.6,525.3) -- (759.0,525.5) -- (760.3,525.7) -- (761.3,526.0) -- (762.0,526.3) -- (762.5,526.8) -- (762.7,527.3) -- (763.0,527.9) -- (763.3,528.7) -- (763.6,529.5) -- (764.0,530.5) -- (764.5,531.6) -- (765.0,532.9) -- (765.6,534.2) -- (766.2,535.8) -- (766.8,537.2) -- (767.3,538.5) -- (767.8,539.8) -- (768.3,541.0) -- (768.7,542.0) -- (769.0,542.8) -- (769.2,543.3) -- (769.3,543.6) -- (769.6,544.2) -- (769.9,545.1) -- (770.4,546.2) -- (771.0,547.8) -- (771.7,549.2) -- (772.3,550.8) -- (772.9,552.2) -- (773.5,553.8) -- (773.9,555.0) -- (774.3,555.9) -- (774.6,556.6) -- (774.8,557.0) -- (775.0,557.6) -- (775.3,558.4) -- (775.7,559.3) -- (776.1,560.4) -- (776.6,561.6) -- (777.2,562.9) -- (777.7,564.2) -- (778.3,565.8) -- (778.9,567.1) -- (779.4,568.4) -- (779.9,569.6) -- (780.3,570.7) -- (780.7,571.7) -- (781.0,572.4) -- (781.3,573.0) -- (781.4,573.4) -- (781.7,574.1) -- (782.1,575.0) -- (782.7,576.2) -- (783.3,577.8) -- (784.0,579.2) -- (784.6,580.8) -- (785.2,582.2) -- (785.9,583.8) -- (786.4,584.9) -- (786.7,585.7) -- (786.9,586.2) -- (786.9,586.3) -- (787.1,586.7) -- (787.5,587.5) -- (788.0,588.5) -- (788.7,589.9) -- (789.4,591.3) -- (790.1,592.8) -- (790.8,594.2) -- (791.5,595.8) -- (792.1,596.9) -- (792.5,597.8) -- (792.7,598.3) -- (792.8,598.5) -- (793.1,599.0) -- (793.5,599.7) -- (794.1,600.7) -- (794.9,602.0) -- (795.7,603.3) -- (796.4,604.6) -- (797.2,605.8) -- (797.9,606.9) -- (798.5,607.9) -- (799.0,608.7) -- (799.4,609.2) -- (799.6,609.5) -- (800.1,610.1) -- (800.8,611.0) -- (801.8,612.2) -- (803.0,613.8) -- (804.0,614.9) -- (804.6,615.7) -- (804.9,616.1) -- (805.0,616.1) -- (805.4,616.5) -- (806.1,617.1) -- (807.2,618.1) -- (808.8,619.5) -- (809.9,620.5) -- (810.8,621.2) -- (811.3,621.7) -- (811.5,621.8) -- (812.0,622.1) -- (812.7,622.6) -- (813.8,623.2) -- (815.1,624.0) -- (816.4,624.7) -- (817.8,625.4) -- (819.2,626.1) -- (820.8,626.7) -- (821.9,627.2) -- (822.8,627.6) -- (823.4,627.8) -- (823.7,627.9) -- (824.3,628.0) -- (825.0,628.2) -- (826.0,628.5) -- (827.2,628.8) -- (828.5,629.1) -- (829.8,629.3) -- (831.2,629.5) -- (832.8,629.8) -- (834.2,629.9) -- (835.8,630.1) -- (837.2,630.2) -- (838.8,630.2) -- (840.2,630.3) -- (841.8,630.3) -- (843.2,630.2) -- (844.8,630.2) -- (846.2,630.1) -- (847.8,630.0) -- (849.2,629.8) -- (850.8,629.7) -- (852.2,629.5) -- (853.8,629.3) -- (855.2,629.0) -- (856.8,628.8) -- (858.0,628.5) -- (858.9,628.4) -- (859.6,628.2) -- (860.0,628.1) -- (860.5,628.0) -- (861.3,627.8) -- (862.2,627.6) -- (863.3,627.3) -- (864.5,627.0) -- (865.8,626.7) -- (867.2,626.3) -- (868.8,625.8) -- (870.2,625.4) -- (871.8,624.9) -- (873.2,624.4) -- (874.8,623.9) -- (876.1,623.4) -- (877.2,623.0) -- (878.1,622.7) -- (878.8,622.4) -- (879.5,622.1) -- (880.3,621.8) -- (881.1,621.5) -- (881.8,621.2) -- (882.8,620.8) -- (883.9,620.4) -- (885.2,619.8) -- (886.8,619.2) -- (888.2,618.5) -- (889.5,617.9) -- (890.7,617.4) -- (891.9,616.8) -- (892.9,616.4) -- (893.6,616.0) -- (894.2,615.8) -- (894.5,615.6) -- (895.1,615.3) -- (896.0,614.8) -- (897.2,614.2) -- (898.8,613.4) -- (900.1,612.7) -- (901.4,612.1) -- (902.6,611.4) -- (903.7,610.9) -- (904.6,610.4) -- (905.4,610.0) -- (905.9,609.7) -- (906.4,609.4) -- (907.1,609.0) -- (908.0,608.5) -- (909.2,607.8) -- (910.8,607.0) -- (911.9,606.6) -- (912.8,606.7) -- (913.3,607.3) -- (913.5,608.4) -- (913.6,609.6) -- (913.8,610.9) -- (914.0,612.2) -- (914.2,613.8) -- (914.7,615.1) -- (915.3,616.3) -- (916.2,617.3) -- (917.3,618.2) -- (918.5,618.9) -- (919.8,619.3) -- (921.2,619.4) -- (922.8,619.4) -- (924.2,619.3) -- (925.8,619.2) -- (927.2,619.2) -- (928.8,619.1) -- (930.2,619.0) -- (931.8,619.0) -- (933.2,618.9) -- (934.8,618.9) -- (936.1,619.0) -- (937.2,619.4) -- (938.2,620.0) -- (939.0,620.8) -- (939.8,621.5) -- (940.5,622.2) -- (941.3,622.9) -- (942.0,623.5) -- (942.9,623.9) -- (944.0,624.2) -- (945.2,624.4) -- (946.8,624.3) -- (948.1,624.5) -- (949.3,625.0) -- (950.4,625.7) -- (951.3,626.6) -- (952.2,627.4) -- (952.9,628.2) -- (953.6,628.8) -- (954.1,629.3) -- (954.9,629.7) -- (956.0,630.0) -- (957.2,630.1) -- (958.8,630.1) -- (960.1,630.3) -- (961.4,630.8) -- (962.5,631.5) -- (963.5,632.5) -- (964.3,633.4) -- (965.1,634.1) -- (965.7,634.7) -- (966.2,635.2) -- (966.9,635.8) -- (967.6,636.6) -- (968.5,637.5) -- (969.5,638.5) -- (970.4,639.4) -- (971.2,640.1) -- (971.8,640.7) -- (972.3,641.2) -- (972.9,641.8) -- (973.6,642.5) -- (974.5,643.4) -- (975.6,644.5) -- (976.5,645.3) -- (977.2,646.1) -- (977.8,646.7) -- (978.3,647.2) -- (978.9,647.8) -- (979.7,648.5) -- (980.6,649.4) -- (981.6,650.4) -- (982.5,651.3) -- (983.2,652.1) -- (983.8,652.7) -- (984.3,653.1) -- (984.9,653.8) -- (985.7,654.5) -- (986.6,655.4) -- (987.6,656.4) -- (988.4,657.6) -- (988.9,658.9) -- (989.2,660.2) -- (989.2,661.8) -- (989.3,663.0) -- (989.5,664.0) -- (989.9,664.7) -- (990.3,665.1) -- (990.9,665.7) -- (991.7,666.5) -- (992.6,667.4) -- (993.6,668.4) -- (994.4,669.6) -- (994.9,670.9) -- (995.2,672.2) -- (995.2,673.8) -- (995.2,675.2) -- (995.2,676.8) -- (995.2,678.2) -- (995.2,679.8) -- (995.2,681.2) -- (995.2,682.8) -- (995.2,684.2) -- (995.2,685.8) -- (995.2,687.2) -- (995.2,688.8) -- (995.2,690.2) -- (995.2,691.8) -- (995.2,693.2) -- (995.2,694.8) -- (995.2,696.2) -- (995.2,697.8) -- (995.2,699.2) -- (995.2,700.8) -- (995.2,702.2) -- (995.2,703.8) -- (995.2,705.2) -- (995.2,706.8) -- (995.2,708.2) -- (995.2,709.8) -- (995.2,711.2) -- (995.2,712.8) -- (995.2,714.2) -- (995.2,715.8) -- (995.2,717.2) -- (995.2,718.8) -- (995.2,720.2) -- (995.2,721.8) -- (995.2,723.2) -- (995.2,724.8) -- (995.2,726.2) -- (995.2,727.8) -- (995.2,729.2) -- (995.2,730.8) -- (995.2,732.2) -- (995.2,733.8) -- (995.2,735.2) -- (995.2,736.8) -- (995.2,738.2) -- (995.2,739.8) -- (995.2,741.2) -- (995.2,742.8) -- (995.2,744.2) -- (995.2,745.8) -- (995.2,747.2) -- (995.2,748.8) -- (995.2,750.2) -- (995.2,751.8) -- (994.9,753.1) -- (994.4,754.4) -- (993.6,755.6) -- (992.6,756.6) -- (991.7,757.5) -- (990.9,758.3) -- (990.3,758.9) -- (989.9,759.3) -- (989.5,760.0) -- (989.3,761.0) -- (989.2,762.2) -- (989.2,763.8) -- (988.9,765.1) -- (988.4,766.4) -- (987.6,767.6) -- (986.6,768.6) -- (985.7,769.5) -- (984.9,770.2) -- (984.3,770.8) -- (983.8,771.3) -- (983.2,771.9) -- (982.5,772.7) -- (981.6,773.5) -- (980.5,774.6) -- (979.7,775.5) -- (978.9,776.2) -- (978.3,776.8) -- (977.8,777.3) -- (977.2,777.9) -- (976.4,778.7) -- (975.6,779.5) -- (974.5,780.6) -- (973.6,781.4) -- (972.9,782.2) -- (972.3,782.8) -- (971.8,783.3) -- (971.1,783.9) -- (970.4,784.6) -- (969.5,785.5) -- (968.5,786.5) -- (967.6,787.4) -- (966.9,788.1) -- (966.2,788.8) -- (965.7,789.3) -- (965.1,789.9) -- (964.3,790.6) -- (963.4,791.5) -- (962.5,792.4) -- (961.4,793.2) -- (960.1,793.6) -- (958.8,793.9) -- (957.2,793.9) -- (956.0,794.0) -- (954.9,794.2) -- (954.1,794.6) -- (953.5,795.2) -- (952.8,795.8) -- (952.1,796.5) -- (951.2,797.3) -- (950.3,798.2) -- (949.3,798.9) -- (948.1,799.3) -- (946.8,799.5) -- (945.2,799.5) -- (943.9,799.6) -- (942.8,799.9) -- (941.8,800.4) -- (941.1,801.0) -- (940.3,801.7) -- (939.5,802.4) -- (938.8,803.1) -- (938.1,803.8) -- (937.2,804.3) -- (936.1,804.7) -- (934.8,804.8) -- (933.2,804.7) -- (931.8,804.6) -- (930.2,804.5) -- (928.8,804.4) -- (927.2,804.3) -- (925.8,804.4) -- (924.4,804.9) -- (923.0,805.6) -- (921.7,806.5) -- (920.6,807.5) -- (919.9,808.4) -- (919.4,809.3) -- (919.3,810.2) -- (918.8,810.7) -- (918.2,810.7) -- (917.3,810.2) -- (916.1,809.3) -- (915.2,808.6) -- (914.4,808.0) -- (913.9,807.6) -- (913.5,807.3) -- (912.9,806.8) -- (912.0,806.1) -- (910.8,805.1) -- (909.2,803.9) -- (908.1,803.0) -- (907.4,802.3) -- (907.0,802.0) -- (907.0,802.0) -- (906.6,801.7) -- (905.8,801.1) -- (904.7,800.2) -- (903.2,798.9) -- (902.0,797.9) -- (901.1,797.2) -- (900.5,796.7) -- (900.2,796.4) -- (899.7,795.9) -- (898.9,795.3) -- (898.0,794.5) -- (896.8,793.4) -- (895.7,792.6) -- (894.9,791.8) -- (894.2,791.2) -- (893.6,790.7) -- (893.0,790.1) -- (892.2,789.5) -- (891.4,788.7) -- (890.4,787.9) -- (889.5,787.1) -- (888.7,786.3) -- (887.9,785.6) -- (887.2,785.0) -- (886.5,784.3) -- (885.7,783.6) -- (884.9,782.9) -- (884.2,782.2) -- (883.4,781.5) -- (882.6,780.7) -- (881.7,779.9) -- (880.9,779.2) -- (880.1,778.4) -- (879.3,777.7) -- (878.6,777.1) -- (878.0,776.5) -- (877.2,775.8) -- (876.4,775.0) -- (875.5,774.2) -- (874.6,773.3) -- (873.7,772.5) -- (872.9,771.9) -- (872.3,771.3) -- (871.8,770.8) -- (871.1,770.1) -- (870.3,769.4) -- (869.3,768.5) -- (868.2,767.5) -- (867.3,766.7) -- (866.5,766.0) -- (866.0,765.5) -- (865.6,765.1) -- (865.0,764.5) -- (864.1,763.8) -- (863.0,762.9) -- (861.7,761.7) -- (860.7,760.8) -- (859.9,760.2) -- (859.5,759.7) -- (859.3,759.6) -- (858.8,759.1) -- (857.9,758.4) -- (856.8,757.5) -- (855.2,756.2) -- (854.1,755.3) -- (853.3,754.6) -- (852.8,754.2) -- (852.7,754.1) -- (852.2,753.8) -- (851.5,753.2) -- (850.4,752.4) -- (849.1,751.4) -- (847.9,750.5) -- (846.8,749.8) -- (846.0,749.1) -- (845.3,748.7) -- (844.6,748.2) -- (843.8,747.7) -- (843.1,747.1) -- (842.2,746.6) -- (841.2,745.9) -- (840.1,745.2) -- (838.8,744.4) -- (837.2,743.4) -- (836.1,742.8) -- (835.4,742.3) -- (835.0,742.0) -- (834.9,742.0) -- (834.5,741.8) -- (833.8,741.4) -- (832.7,740.8) -- (831.2,740.0) -- (829.7,739.3) -- (828.2,738.5) -- (826.8,737.8) -- (825.2,737.1) -- (824.1,736.6) -- (823.3,736.2) -- (822.9,736.0) -- (822.9,736.0) -- (822.5,735.9) -- (821.8,735.6) -- (820.6,735.1) -- (819.2,734.5) -- (817.7,734.0) -- (816.2,733.4) -- (814.8,732.9) -- (813.2,732.3) -- (811.8,731.9) -- (810.4,731.4) -- (809.0,731.0) -- (807.7,730.6) -- (806.6,730.3) -- (805.9,730.1) -- (805.4,729.9) -- (805.3,729.9) -- (804.8,729.7) -- (803.9,729.5) -- (802.8,729.3) -- (801.2,728.9) -- (799.8,728.5) -- (798.2,728.2) -- (796.8,727.9) -- (795.2,727.6) -- (793.8,727.4) -- (792.2,727.2) -- (790.8,726.9) -- (789.2,726.8) -- (787.8,726.6) -- (786.2,726.5) -- (784.8,726.4) -- (783.2,726.3) -- (781.8,726.3) -- (780.2,726.3) -- (778.8,726.4) -- (777.2,726.5) -- (775.8,726.7) -- (774.2,726.9) -- (772.8,727.2) -- (771.2,727.6) -- (769.9,727.9) -- (768.6,728.4) -- (767.5,728.8) -- (766.5,729.3) -- (765.7,729.7) -- (764.9,730.1) -- (764.3,730.5) -- (763.8,730.9) -- (763.1,731.4) -- (762.4,732.2) -- (761.5,733.1) -- (760.4,734.3) -- (759.6,735.4) -- (758.8,736.4) -- (758.2,737.4) -- (757.7,738.4) -- (757.3,739.2) -- (757.0,740.0) -- (756.7,740.6) -- (756.5,741.2) -- (756.3,742.0) -- (756.0,743.0) -- (755.7,744.2) -- (755.4,745.8) -- (755.1,747.2) -- (754.9,748.8) -- (754.7,750.2) -- (754.6,751.8) -- (754.5,753.2) -- (754.4,754.8) -- (754.3,756.2) -- (754.2,757.8) -- (754.2,759.2) -- (754.1,760.8) -- (754.1,762.2) -- (754.1,763.8) -- (754.1,765.2) -- (754.1,766.8) -- (754.0,768.2) -- (754.0,769.8) -- (754.0,771.2) -- (754.0,772.8) -- (754.0,774.2) -- (753.9,775.8) -- (753.9,777.2) -- (753.9,778.8) -- (753.8,780.2) -- (753.8,781.8) -- (753.7,783.2) -- (753.6,784.8) -- (753.5,786.2) -- (753.5,787.8) -- (753.4,789.2) -- (753.2,790.8) -- (753.1,792.2) -- (753.0,793.8) -- (752.8,795.2) -- (752.7,796.8) -- (752.5,798.2) -- (752.3,799.8) -- (752.1,801.2) -- (751.9,802.7) -- (751.6,804.2) -- (751.4,805.6) -- (751.2,806.8) -- (751.1,807.5) -- (751.0,807.9) -- (751.0,807.9) -- (750.9,808.3) -- (750.7,809.1) -- (750.5,810.2) -- (750.2,811.8) -- (749.9,813.2) -- (749.5,814.8) -- (749.1,816.2) -- (748.7,817.8) -- (748.3,819.2) -- (747.9,820.8) -- (747.4,822.2) -- (746.9,823.8) -- (746.4,825.0) -- (746.0,826.1) -- (745.7,827.0) -- (745.4,827.6) -- (745.1,828.3) -- (744.8,829.1) -- (744.5,829.9) -- (744.1,830.7) -- (743.6,831.7) -- (743.1,832.9) -- (742.4,834.2) -- (741.6,835.8) -- (740.9,837.0) -- (740.3,838.0) -- (739.9,838.8) -- (739.5,839.3) -- (739.1,839.9) -- (738.6,840.7) -- (738.0,841.5) -- (737.3,842.5) -- (736.6,843.5) -- (735.8,844.5) -- (735.0,845.4) -- (734.2,846.4) -- (733.5,847.2) -- (732.8,848.0) -- (732.3,848.6) -- (731.8,849.2) -- (731.1,849.8) -- (730.4,850.6) -- (729.5,851.4) -- (728.5,852.3) -- (727.6,853.2) -- (726.7,853.9) -- (726.0,854.6) -- (725.3,855.1) -- (724.6,855.7) -- (723.9,856.3) -- (723.1,857.0) -- (722.2,857.6) -- (721.3,858.3) -- (720.3,859.0) -- (719.2,859.8) -- (718.0,860.7) -- (717.0,861.4) -- (716.2,861.9) -- (715.7,862.3) -- (715.4,862.5) -- (714.8,862.9) -- (713.9,863.5) -- (712.8,864.3) -- (711.2,865.2) -- (710.0,866.1) -- (709.0,866.7) -- (708.2,867.2) -- (707.7,867.5) -- (707.1,867.9) -- (706.4,868.4) -- (705.5,869.0) -- (704.5,869.6) -- (703.4,870.3) -- (702.2,871.1) -- (701.0,871.9) -- (699.7,872.7) -- (698.6,873.4) -- (697.9,873.9) -- (697.4,874.2) -- (697.2,874.3) -- (696.7,874.6) -- (695.9,875.2) -- (694.8,876.0) -- (693.2,877.0) -- (692.0,877.8) -- (691.0,878.5) -- (690.2,879.1) -- (689.6,879.4) -- (688.9,879.9) -- (688.2,880.4) -- (687.3,881.0) -- (686.4,881.6) -- (685.4,882.4) -- (684.3,883.1) -- (683.1,883.9) -- (681.9,884.7) -- (680.9,885.4) -- (680.1,885.9) -- (679.6,886.3) -- (679.4,886.4) -- (678.8,886.8) -- (677.9,887.4) -- (676.8,888.2) -- (675.2,889.2) -- (674.0,890.0) -- (673.0,890.7) -- (672.2,891.2) -- (671.6,891.5) -- (671.0,891.9) -- (670.2,892.4) -- (669.4,892.9) -- (668.4,893.5) -- (667.3,894.1) -- (666.1,894.8) -- (664.8,895.6) -- (663.2,896.4) -- (662.1,897.0) -- (661.2,897.5) -- (660.7,897.8) -- (660.5,897.9) -- (660.0,898.1) -- (659.3,898.5) -- (658.3,899.0) -- (657.0,899.6) -- (655.6,900.2) -- (654.2,900.9) -- (652.8,901.5) -- (651.2,902.1) -- (649.9,902.6) -- (648.8,903.1) -- (647.9,903.4) -- (647.2,903.6) -- (646.5,903.9) -- (645.8,904.2) -- (645.0,904.4) -- (644.2,904.7) -- (643.2,905.1) -- (642.1,905.4) -- (640.8,905.8) -- (639.2,906.3) -- (637.8,906.7) -- (636.2,907.1) -- (634.8,907.5) -- (633.2,907.9) -- (631.8,908.3) -- (630.2,908.6) -- (628.8,908.9) -- (627.2,909.2) -- (626.0,909.4) -- (624.9,909.6) -- (624.1,909.8) -- (623.6,909.9) -- (622.9,910.0) -- (622.2,910.1) -- (621.3,910.2) -- (620.4,910.4) -- (619.3,910.6) -- (618.1,910.7) -- (616.8,910.9) -- (615.2,911.1) -- (613.8,911.3) -- (612.2,911.4) -- (610.8,911.6) -- (609.2,911.7) -- (607.8,911.8) -- (606.2,911.9) -- (604.8,911.9) -- (603.2,912.0) -- (601.8,912.0) -- (600.2,912.0) -- (598.8,912.0) -- (597.2,912.0) -- (595.8,911.9) -- (594.2,911.8) -- (592.8,911.7) -- (591.2,911.6) -- (589.8,911.4) -- (588.2,911.2) -- (586.8,911.0) -- (585.2,910.7) -- (584.0,910.5) -- (583.1,910.3) -- (582.5,910.2) -- (582.2,910.1) -- (581.7,910.0) -- (580.9,909.8) -- (580.0,909.5) -- (578.8,909.2) -- (577.5,908.7) -- (576.2,908.3) -- (574.8,907.7) -- (573.2,907.0) -- (571.9,906.4) -- (570.8,905.8) -- (569.8,905.3) -- (568.9,904.8) -- (568.1,904.3) -- (567.4,903.8) -- (566.7,903.3) -- (566.0,902.8) -- (565.2,902.2) -- (564.3,901.5) -- (563.3,900.6) -- (562.2,899.6) -- (561.3,898.7) -- (560.6,897.9) -- (560.0,897.3) -- (559.6,896.9) -- (559.1,896.3) -- (558.5,895.5) -- (557.8,894.6) -- (557.1,893.6) -- (556.3,892.5) -- (555.6,891.3) -- (554.8,890.0) -- (554.1,888.7) -- (553.5,887.7) -- (553.1,887.0) -- (552.9,886.5) -- (552.8,886.3) -- (552.5,885.8) -- (552.2,884.9) -- (551.7,883.8) -- (551.1,882.2) -- (550.6,880.8) -- (550.1,879.2) -- (549.7,877.8) -- (549.2,876.2) -- (548.9,874.8) -- (548.6,873.2) -- (548.3,871.8) -- (548.0,870.2) -- (547.8,868.8) -- (547.6,867.2) -- (547.5,865.8) -- (547.4,864.2) -- (547.3,862.8) -- (547.2,861.2) -- (547.2,859.8) -- (547.2,858.2) -- (547.2,856.8) -- (547.3,855.2) -- (547.4,853.8) -- (547.5,852.2) -- (547.6,850.8) -- (547.8,849.2) -- (548.0,847.8) -- (548.2,846.2) -- (548.4,844.8) -- (548.6,843.2) -- (548.9,841.8) -- (549.2,840.2) -- (549.5,838.8) -- (549.8,837.2) -- (550.2,835.8) -- (550.6,834.2) -- (550.9,832.8) -- (551.3,831.2) -- (551.7,829.8) -- (552.2,828.2) -- (552.5,827.1) -- (552.7,826.2) -- (552.9,825.7) -- (552.9,825.5) -- (553.1,825.1) -- (553.3,824.3) -- (553.6,823.3) -- (553.9,822.0) -- (554.3,820.6) -- (554.7,819.2) -- (555.1,817.8) -- (555.5,816.2) -- (555.9,814.8) -- (556.3,813.2) -- (556.6,811.8) -- (557.0,810.2) -- (557.3,808.8) -- (557.6,807.2) -- (557.9,805.8) -- (558.2,804.2) -- (558.4,802.8) -- (558.6,801.4) -- (558.7,800.0) -- (558.8,798.7) -- (558.9,797.6) -- (559.0,796.9) -- (559.0,796.4) -- (559.0,796.2) -- (559.0,796.0) -- (559.0,795.7) -- (559.0,795.3) -- (559.0,794.9) -- (559.0,794.3) -- (559.0,793.5) -- (559.0,792.6) -- (558.9,791.6) -- (558.8,790.4) -- (558.7,789.1) -- (558.5,787.8) -- (558.3,786.2) -- (558.0,784.8) -- (557.6,783.2) -- (557.3,781.8) -- (556.8,780.2) -- (556.3,778.8) -- (555.8,777.2) -- (555.2,775.8) -- (554.5,774.2) -- (554.0,773.1) -- (553.6,772.2) -- (553.4,771.6) -- (553.2,771.4) -- (552.9,770.9) -- (552.5,770.1) -- (551.9,769.1) -- (551.1,767.9) -- (550.4,766.7) -- (549.6,765.5) -- (548.9,764.4) -- (548.1,763.4) -- (547.5,762.5) -- (546.9,761.8) -- (546.5,761.2) -- (546.1,760.7) -- (545.6,760.0) -- (544.8,759.1) -- (543.9,758.0) -- (542.7,756.7) -- (541.8,755.7) -- (541.2,755.0) -- (540.7,754.5) -- (540.6,754.3) -- (540.1,753.8) -- (539.3,753.1) -- (538.3,752.0) -- (537.0,750.7) -- (536.0,749.7) -- (535.2,748.9) -- (534.7,748.4) -- (534.6,748.3) -- (534.1,747.8) -- (533.4,747.1) -- (532.3,746.1) -- (531.0,744.9) -- (529.9,743.9) -- (529.1,743.2) -- (528.6,742.7) -- (528.3,742.4) -- (527.8,741.9) -- (527.1,741.2) -- (526.1,740.3) -- (524.8,739.2) -- (523.8,738.2) -- (523.0,737.5) -- (522.4,736.9) -- (522.1,736.5) -- (521.5,736.0) -- (520.8,735.3) -- (519.8,734.3) -- (518.7,733.2) -- (517.8,732.3) -- (517.0,731.6) -- (516.4,731.0) -- (516.0,730.6) -- (515.5,730.0) -- (514.7,729.2) -- (513.8,728.2) -- (512.7,727.0) -- (511.8,726.0) -- (511.1,725.2) -- (510.6,724.7) -- (510.4,724.4) -- (509.9,723.8) -- (509.2,722.9) -- (508.2,721.8) -- (507.1,720.2) -- (506.2,719.1) -- (505.5,718.3) -- (505.2,717.8) -- (505.1,717.7) -- (504.8,717.3) -- (504.3,716.6) -- (503.6,715.5) -- (502.6,714.1) -- (501.8,712.8) -- (501.0,711.6) -- (500.4,710.5) -- (499.8,709.5) -- (499.3,708.6) -- (498.9,707.9) -- (498.6,707.3) -- (498.3,706.8) -- (497.9,706.0) -- (497.4,705.0) -- (496.9,703.8) -- (496.2,702.2) -- (495.5,700.8) -- (494.9,699.2) -- (494.4,697.8) -- (493.9,696.2) -- (493.5,695.1) -- (493.2,694.3) -- (493.1,693.9) -- (493.0,693.8) -- (492.9,693.4) -- (492.7,692.6) -- (492.4,691.5) -- (492.0,690.1) -- (491.7,688.7) -- (491.4,687.2) -- (491.1,685.8) -- (490.8,684.2) -- (490.6,682.8) -- (490.3,681.2) -- (490.2,679.8) -- (490.0,678.2) -- (489.9,676.8) -- (489.8,675.2) -- (489.8,673.8) -- (489.7,672.2) -- (489.7,670.8) -- (489.7,669.2) -- (489.8,667.8) -- (489.9,666.2) -- (490.0,664.8) -- (490.1,663.2) -- (490.2,661.8) -- (490.4,660.2) -- (490.6,658.8) -- (490.8,657.2) -- (491.1,655.8) -- (491.4,654.2) -- (491.7,652.8) -- (492.0,651.5) -- (492.3,650.3) -- (492.6,649.1) -- (492.8,648.2) -- (493.0,647.4) -- (493.1,646.9) -- (493.2,646.5) -- (493.4,645.9) -- (493.7,645.0) -- (494.0,643.8) -- (494.5,642.2) -- (495.0,640.8) -- (495.5,639.2) -- (496.0,637.8) -- (496.6,636.2) -- (497.1,634.9) -- (497.6,633.7) -- (498.0,632.7) -- (498.4,631.8) -- (498.8,631.0) -- (499.1,630.2) -- (499.4,629.6) -- (499.7,628.9) -- (500.0,628.1) -- (500.5,627.0) -- (501.1,625.8) -- (501.8,624.2) -- (502.5,622.9) -- (503.1,621.7) -- (503.7,620.6) -- (504.2,619.6) -- (504.7,618.8) -- (505.1,618.0) -- (505.4,617.4) -- (505.7,616.8) -- (506.1,616.0) -- (506.7,615.0) -- (507.4,613.8) -- (508.2,612.2) -- (508.9,611.0) -- (509.5,609.9) -- (510.0,609.0) -- (510.4,608.4) -- (510.8,607.7) -- (511.3,607.0) -- (511.7,606.1) -- (512.2,605.3) -- (512.8,604.3) -- (513.5,603.1) -- (514.2,601.8) -- (515.1,600.2) -- (515.8,599.1) -- (516.3,598.2) -- (516.7,597.7) -- (516.8,597.4) -- (517.1,596.9) -- (517.5,596.2) -- (518.0,595.2) -- (518.7,593.9) -- (519.5,592.6) -- (520.2,591.2) -- (521.0,589.8) -- (521.8,588.4) -- (522.4,587.3) -- (522.8,586.5) -- (523.0,586.1) -- (523.1,586.1) -- (523.3,585.7) -- (523.6,584.9) -- (524.2,583.8) -- (524.9,582.2) -- (525.6,580.8) -- (526.3,579.2) -- (527.0,577.8) -- (527.7,576.2) -- (528.2,575.1) -- (528.6,574.2) -- (528.8,573.7) -- (528.9,573.6) -- (529.1,573.1) -- (529.4,572.4) -- (529.7,571.3) -- (530.2,570.0) -- (530.7,568.6) -- (531.2,567.2) -- (531.8,565.8) -- (532.3,564.2) -- (532.8,562.8) -- (533.3,561.2) -- (533.8,559.8) -- (534.2,558.2) -- (534.6,557.1) -- (534.8,556.3) -- (535.0,555.9) -- (535.0,555.9) -- (535.1,555.5) -- (535.3,554.7) -- (535.5,553.6) -- (535.9,552.2) -- (536.2,550.7) -- (536.6,549.2) -- (536.9,547.8) -- (537.2,546.2) -- (537.6,544.8) -- (537.9,543.2) -- (538.2,541.8) -- (538.5,540.2) -- (538.8,538.8) -- (539.0,537.2) -- (539.3,535.8) -- (539.6,534.2) -- (539.9,532.8) -- (540.1,531.2) -- (540.4,529.8) -- (540.6,528.2) -- (540.8,527.1) -- (540.9,526.4) -- (541.0,526.0) -- cycle ;
\draw[black,line width=0.51pt] (476.0,583.0) -- (456.5,551.8);
\draw[black,line width=0.51pt] (456.5,551.8) -- (464.2,539.2);
\draw[black,line width=0.51pt,fill=white] (476.0,583.0) circle (7.8);
\fill[black] (476.0,583.0) -- (476.0,575.2) -- (477.0,575.2) -- (478.0,575.4) -- (479.0,575.8) -- (479.9,576.2) -- (480.8,576.8) -- (481.5,577.5) -- (482.2,578.2) -- (482.8,579.1) -- (483.2,580.0) -- (483.6,581.0) -- (483.8,582.0) -- (483.8,583.0) -- cycle;
\node[anchor=east,inner sep=0pt,font=\fontsize{5.4}{5.4}\selectfont,text=black] at (465.9,572.9) {34};
\node[anchor=east,inner sep=0pt,font=\fontsize{5.4}{5.4}\selectfont,text=black] at (465.9,601.4) {22};
\node[anchor=west,inner sep=0pt,font=\fontsize{5.4}{5.4}\selectfont,text=black] at (486.1,572.9) {110};
\draw[black,line width=0.51pt] (626.0,593.0) -- (662.7,589.8);
\draw[black,line width=0.51pt] (662.7,589.8) -- (670.4,602.4);
\draw[black,line width=0.51pt,fill=white] (626.0,593.0) circle (7.8);
\fill[black] (626.0,593.0) circle (7.8);
\node[anchor=east,inner sep=0pt,font=\fontsize{5.4}{5.4}\selectfont,text=black] at (615.9,582.9) {45};
\node[anchor=east,inner sep=0pt,font=\fontsize{5.4}{5.4}\selectfont,text=black] at (615.9,611.4) {43};
\node[anchor=west,inner sep=0pt,font=\fontsize{5.4}{5.4}\selectfont,text=black] at (636.1,582.9) {096};
\fill[black] (605.2,597.6) circle (2.0);
\fill[black] (611.8,597.6) circle (2.0);
\draw[black,line width=0.51pt] (824.0,581.0) -- (859.9,572.7);
\draw[black,line width=0.51pt] (859.9,572.7) -- (869.3,584.1);
\draw[black,line width=0.51pt] (854.5,574.0) -- (859.2,579.7);
\draw[black,line width=0.51pt,fill=white] (824.0,581.0) circle (7.8);
\fill[black] (824.0,581.0) -- (824.0,573.2) -- (825.0,573.2) -- (826.0,573.4) -- (827.0,573.8) -- (827.9,574.2) -- (828.8,574.8) -- (829.5,575.5) -- (830.2,576.2) -- (830.8,577.1) -- (831.2,578.0) -- (831.6,579.0) -- (831.8,580.0) -- (831.8,581.0) -- (831.8,582.0) -- (831.6,583.0) -- (831.2,584.0) -- (830.8,584.9) -- (830.2,585.8) -- (829.5,586.5) -- (828.8,587.2) -- (827.9,587.8) -- (827.0,588.2) -- (826.0,588.6) -- (825.0,588.8) -- (824.0,588.8) -- (823.0,588.8) -- (822.0,588.6) -- (821.0,588.2) -- (820.1,587.8) -- (819.2,587.2) -- (818.5,586.5) -- (817.8,585.8) -- (817.2,584.9) -- (816.8,584.0) -- (816.4,583.0) -- (816.2,582.0) -- (816.2,581.0) -- cycle;
\node[anchor=east,inner sep=0pt,font=\fontsize{5.4}{5.4}\selectfont,text=black] at (813.9,570.9) {39};
\node[anchor=east,inner sep=0pt,font=\fontsize{5.4}{5.4}\selectfont,text=black] at (813.9,599.4) {39};
\node[anchor=west,inner sep=0pt,font=\fontsize{5.4}{5.4}\selectfont,text=black] at (834.1,570.9) {154};
\draw[black,line width=0.51pt] (465.0,654.0) -- (441.3,625.8);
\draw[black,line width=0.51pt] (441.3,625.8) -- (447.2,612.2);
\draw[black,line width=0.51pt,fill=white] (465.0,654.0) circle (7.8);
\fill[black] (465.0,654.0) -- (465.0,646.2) -- (466.0,646.2) -- (467.0,646.4) -- (468.0,646.8) -- (468.9,647.2) -- (469.8,647.8) -- (470.5,648.5) -- (471.2,649.2) -- (471.8,650.1) -- (472.2,651.0) -- (472.6,652.0) -- (472.8,653.0) -- (472.8,654.0) -- (472.8,655.0) -- (472.6,656.0) -- (472.2,657.0) -- (471.8,657.9) -- (471.2,658.8) -- (470.5,659.5) -- (469.8,660.2) -- (468.9,660.8) -- (468.0,661.2) -- (467.0,661.6) -- (466.0,661.8) -- (465.0,661.8) -- cycle;
\node[anchor=east,inner sep=0pt,font=\fontsize{5.4}{5.4}\selectfont,text=black] at (454.9,643.9) {36};
\node[anchor=east,inner sep=0pt,font=\fontsize{5.4}{5.4}\selectfont,text=black] at (454.9,672.4) {28};
\node[anchor=west,inner sep=0pt,font=\fontsize{5.4}{5.4}\selectfont,text=black] at (475.1,643.9) {124};
\draw[black,line width=0.51pt] (544.0,662.0) -- (519.4,634.7);
\draw[black,line width=0.51pt] (519.4,634.7) -- (524.8,620.9);
\draw[black,line width=0.51pt] (523.1,638.8) -- (528.5,625.0);
\draw[black,line width=0.51pt,fill=white] (544.0,662.0) circle (7.8);
\fill[black] (544.0,662.0) circle (7.8);
\node[anchor=east,inner sep=0pt,font=\fontsize{5.4}{5.4}\selectfont,text=black] at (533.9,651.9) {40};
\node[anchor=east,inner sep=0pt,font=\fontsize{5.4}{5.4}\selectfont,text=black] at (533.9,680.4) {38};
\node[anchor=west,inner sep=0pt,font=\fontsize{5.4}{5.4}\selectfont,text=black] at (554.1,651.9) {088};
\fill[black] (526.5,666.6) circle (2.0);
\draw[black,line width=0.51pt] (738.0,624.0) -- (767.0,646.7);
\draw[black,line width=0.51pt] (767.0,646.7) -- (764.1,661.2);
\draw[black,line width=0.51pt] (762.6,643.3) -- (759.7,657.8);
\draw[black,line width=0.51pt,fill=white] (738.0,624.0) circle (7.8);
\fill[black] (738.0,624.0) circle (7.8);
\node[anchor=east,inner sep=0pt,font=\fontsize{5.4}{5.4}\selectfont,text=black] at (727.9,613.9) {39};
\node[anchor=east,inner sep=0pt,font=\fontsize{5.4}{5.4}\selectfont,text=black] at (727.9,642.4) {39};
\node[anchor=west,inner sep=0pt,font=\fontsize{5.4}{5.4}\selectfont,text=black] at (748.1,613.9) {084};
\fill[black] (717.2,628.6) circle (2.0);
\fill[black] (723.8,628.6) circle (2.0);
\draw[black,line width=0.51pt] (822.0,635.0) -- (857.0,646.4);
\draw[black,line width=0.51pt] (857.0,646.4) -- (859.2,661.0);
\draw[black,line width=0.51pt] (851.7,644.7) -- (854.0,659.3);
\draw[black,line width=0.51pt,fill=white] (822.0,635.0) circle (7.8);
\fill[black] (822.0,635.0) -- (822.0,627.2) -- (823.0,627.2) -- (824.0,627.4) -- (825.0,627.8) -- (825.9,628.2) -- (826.8,628.8) -- (827.5,629.5) -- (828.2,630.2) -- (828.8,631.1) -- (829.2,632.0) -- (829.6,633.0) -- (829.8,634.0) -- (829.8,635.0) -- (829.8,636.0) -- (829.6,637.0) -- (829.2,638.0) -- (828.8,638.9) -- (828.2,639.8) -- (827.5,640.5) -- (826.8,641.2) -- (825.9,641.8) -- (825.0,642.2) -- (824.0,642.6) -- (823.0,642.8) -- (822.0,642.8) -- cycle;
\draw[black,line width=0.51pt] (814.2,635.0) -- (829.8,635.0);
\draw[black,line width=0.51pt] (822.0,627.2) -- (822.0,642.8);
\node[anchor=east,inner sep=0pt,font=\fontsize{5.4}{5.4}\selectfont,text=black] at (811.9,624.9) {43};
\node[anchor=east,inner sep=0pt,font=\fontsize{5.4}{5.4}\selectfont,text=black] at (811.9,653.4) {41};
\node[anchor=west,inner sep=0pt,font=\fontsize{5.4}{5.4}\selectfont,text=black] at (832.1,624.9) {132};
\draw[black,line width=0.51pt] (672.0,681.0) -- (701.0,658.3);
\draw[black,line width=0.51pt] (701.0,658.3) -- (714.4,664.7);
\draw[black,line width=0.51pt,fill=white] (672.0,681.0) circle (7.8);
\fill[black] (672.0,681.0) circle (7.8);
\node[anchor=east,inner sep=0pt,font=\fontsize{5.4}{5.4}\selectfont,text=black] at (661.9,670.9) {48};
\node[anchor=east,inner sep=0pt,font=\fontsize{5.4}{5.4}\selectfont,text=black] at (661.9,699.4) {45};
\node[anchor=west,inner sep=0pt,font=\fontsize{5.4}{5.4}\selectfont,text=black] at (682.1,670.9) {024};
\fill[black] (651.2,685.6) circle (2.0);
\fill[black] (657.8,685.6) circle (2.0);
\draw[black,line width=0.51pt] (600.0,710.0) -- (565.9,696.2);
\draw[black,line width=0.51pt] (565.9,696.2) -- (564.7,681.4);
\draw[black,line width=0.51pt,fill=white] (600.0,710.0) circle (7.8);
\fill[black] (600.0,710.0) circle (7.8);
\node[anchor=east,inner sep=0pt,font=\fontsize{5.4}{5.4}\selectfont,text=black] at (589.9,699.9) {48};
\node[anchor=west,inner sep=0pt,font=\fontsize{5.4}{5.4}\selectfont,text=black] at (610.1,699.9) {026};
\fill[black] (579.2,714.6) circle (2.0);
\fill[black] (585.8,714.6) circle (2.0);
\draw[black,line width=0.51pt] (824.0,683.0) -- (860.4,688.1);
\draw[black,line width=0.51pt] (855.0,687.4) -- (857.3,694.4);
\draw[black,line width=0.51pt,fill=white] (824.0,683.0) circle (7.8);
\fill[black] (824.0,683.0) circle (7.8);
\node[anchor=east,inner sep=0pt,font=\fontsize{5.4}{5.4}\selectfont,text=black] at (813.9,672.9) {46};
\node[anchor=east,inner sep=0pt,font=\fontsize{5.4}{5.4}\selectfont,text=black] at (813.9,701.4) {43};
\node[anchor=west,inner sep=0pt,font=\fontsize{5.4}{5.4}\selectfont,text=black] at (834.1,672.9) {124};
\fill[black] (803.2,687.6) circle (2.0);
\fill[black] (809.8,687.6) circle (2.0);
\draw[black,line width=0.51pt] (908.0,712.0) -- (943.5,721.5);
\draw[black,line width=0.51pt,fill=white] (908.0,712.0) circle (7.8);
\fill[black] (908.0,712.0) circle (7.8);
\node[anchor=east,inner sep=0pt,font=\fontsize{5.4}{5.4}\selectfont,text=black] at (897.9,701.9) {45};
\node[anchor=east,inner sep=0pt,font=\fontsize{5.4}{5.4}\selectfont,text=black] at (897.9,730.4) {43};
\node[anchor=west,inner sep=0pt,font=\fontsize{5.4}{5.4}\selectfont,text=black] at (918.1,701.9) {163};
\fill[black] (890.5,716.6) circle (2.0);
\draw[black,line width=0.51pt] (419.0,722.0) -- (394.4,694.7);
\draw[black,line width=0.51pt] (394.4,694.7) -- (399.8,680.9);
\draw[black,line width=0.51pt,fill=white] (419.0,722.0) circle (7.8);
\node[anchor=east,inner sep=0pt,font=\fontsize{5.4}{5.4}\selectfont,text=black] at (408.9,711.9) {39};
\node[anchor=east,inner sep=0pt,font=\fontsize{5.4}{5.4}\selectfont,text=black] at (408.9,740.4) {31};
\node[anchor=west,inner sep=0pt,font=\fontsize{5.4}{5.4}\selectfont,text=black] at (429.1,711.9) {124};
\draw[black,line width=0.51pt] (714.0,722.0) -- (698.4,755.4);
\draw[black,line width=0.51pt] (698.4,755.4) -- (683.6,755.8);
\draw[black,line width=0.51pt,fill=white] (714.0,722.0) circle (7.8);
\fill[black] (714.0,722.0) -- (714.0,714.2) -- (715.0,714.2) -- (716.0,714.4) -- (717.0,714.8) -- (717.9,715.2) -- (718.8,715.8) -- (719.5,716.5) -- (720.2,717.2) -- (720.8,718.1) -- (721.2,719.0) -- (721.6,720.0) -- (721.8,721.0) -- (721.8,722.0) -- (721.8,723.0) -- (721.6,724.0) -- (721.2,725.0) -- (720.8,725.9) -- (720.2,726.8) -- (719.5,727.5) -- (718.8,728.2) -- (717.9,728.8) -- (717.0,729.2) -- (716.0,729.6) -- (715.0,729.8) -- (714.0,729.8) -- (713.0,729.8) -- (712.0,729.6) -- (711.0,729.2) -- (710.1,728.8) -- (709.2,728.2) -- (708.5,727.5) -- (707.8,726.8) -- (707.2,725.9) -- (706.8,725.0) -- (706.4,724.0) -- (706.2,723.0) -- (706.2,722.0) -- cycle;
\node[anchor=east,inner sep=0pt,font=\fontsize{5.4}{5.4}\selectfont,text=black] at (703.9,711.9) {57};
\node[anchor=east,inner sep=0pt,font=\fontsize{5.4}{5.4}\selectfont,text=black] at (703.9,740.4) {54};
\node[anchor=west,inner sep=0pt,font=\fontsize{5.4}{5.4}\selectfont,text=black] at (724.1,711.9) {029};
\draw[black,line width=0.51pt] (818.0,742.0) -- (854.2,748.4);
\draw[black,line width=0.51pt] (854.2,748.4) -- (858.5,762.6);
\draw[black,line width=0.51pt,fill=white] (818.0,742.0) circle (7.8);
\fill[black] (818.0,742.0) circle (7.8);
\node[anchor=east,inner sep=0pt,font=\fontsize{5.4}{5.4}\selectfont,text=black] at (807.9,731.9) {54};
\node[anchor=east,inner sep=0pt,font=\fontsize{5.4}{5.4}\selectfont,text=black] at (807.9,760.4) {45};
\node[anchor=west,inner sep=0pt,font=\fontsize{5.4}{5.4}\selectfont,text=black] at (828.1,731.9) {097};
\draw[black,line width=0.51pt] (469.0,770.0) -- (438.5,749.4);
\draw[black,line width=0.51pt] (438.5,749.4) -- (440.4,734.7);
\draw[black,line width=0.51pt,fill=white] (469.0,770.0) circle (7.8);
\draw[black,line width=0.51pt] (469.0,762.2) -- (469.0,777.8);
\node[anchor=east,inner sep=0pt,font=\fontsize{5.4}{5.4}\selectfont,text=black] at (458.9,759.9) {41};
\node[anchor=east,inner sep=0pt,font=\fontsize{5.4}{5.4}\selectfont,text=black] at (458.9,788.4) {30};
\node[anchor=west,inner sep=0pt,font=\fontsize{5.4}{5.4}\selectfont,text=black] at (479.1,759.9) {103};
\draw[black,line width=0.51pt] (570.0,768.0) -- (544.9,741.1);
\draw[black,line width=0.51pt] (544.9,741.1) -- (550.1,727.2);
\draw[black,line width=0.51pt] (548.7,745.1) -- (553.8,731.2);
\draw[black,line width=0.51pt] (552.4,749.2) -- (555.0,742.2);
\draw[black,line width=0.51pt,fill=white] (570.0,768.0) circle (7.8);
\fill[black] (570.0,768.0) -- (570.0,760.2) -- (571.0,760.2) -- (572.0,760.4) -- (573.0,760.8) -- (573.9,761.2) -- (574.8,761.8) -- (575.5,762.5) -- (576.2,763.2) -- (576.8,764.1) -- (577.2,765.0) -- (577.6,766.0) -- (577.8,767.0) -- (577.8,768.0) -- (577.8,769.0) -- (577.6,770.0) -- (577.2,771.0) -- (576.8,771.9) -- (576.2,772.8) -- (575.5,773.5) -- (574.8,774.2) -- (573.9,774.8) -- (573.0,775.2) -- (572.0,775.6) -- (571.0,775.8) -- (570.0,775.8) -- cycle;
\node[anchor=east,inner sep=0pt,font=\fontsize{5.4}{5.4}\selectfont,text=black] at (559.9,757.9) {45};
\node[anchor=east,inner sep=0pt,font=\fontsize{5.4}{5.4}\selectfont,text=black] at (559.9,786.4) {40};
\node[anchor=west,inner sep=0pt,font=\fontsize{5.4}{5.4}\selectfont,text=black] at (580.1,757.9) {069};
\draw[black,line width=0.51pt] (637.0,765.0) -- (618.0,796.5);
\draw[black,line width=0.51pt] (618.0,796.5) -- (603.3,795.4);
\draw[black,line width=0.51pt] (620.9,791.8) -- (606.1,790.7);
\draw[black,line width=0.51pt] (623.7,787.1) -- (616.3,786.5);
\draw[black,line width=0.51pt,fill=white] (637.0,765.0) circle (7.8);
\fill[black] (637.0,765.0) circle (7.8);
\node[anchor=east,inner sep=0pt,font=\fontsize{5.4}{5.4}\selectfont,text=black] at (626.9,754.9) {60};
\node[anchor=east,inner sep=0pt,font=\fontsize{5.4}{5.4}\selectfont,text=black] at (626.9,783.4) {55};
\node[anchor=west,inner sep=0pt,font=\fontsize{5.4}{5.4}\selectfont,text=black] at (647.1,754.9) {040};
\draw[black,line width=0.43pt] (616.4,763.3) -- (616.4,775.9);
\draw[black,line width=0.43pt] (616.4,763.3) -- (621.1,763.3) .. controls (625.4,763.3) and (625.4,770.0) .. (621.1,770.0) -- (616.4,770.0);
\draw[black,line width=0.43pt] (620.3,770.0) -- (625.8,775.9);
\fill[black] (626.9,777.0) -- (623.1,776.1) -- (625.6,773.4) -- cycle;
\fill[black] (618.7,759.4) circle (1.9);
\draw[black,line width=0.51pt] (704.0,785.0) -- (685.0,816.5);
\draw[black,line width=0.51pt] (685.0,816.5) -- (670.3,815.4);
\draw[black,line width=0.51pt,fill=white] (704.0,785.0) circle (7.8);
\fill[black] (704.0,785.0) -- (704.0,777.2) -- (705.0,777.2) -- (706.0,777.4) -- (707.0,777.8) -- (707.9,778.2) -- (708.8,778.8) -- (709.5,779.5) -- (710.2,780.2) -- (710.8,781.1) -- (711.2,782.0) -- (711.6,783.0) -- (711.8,784.0) -- (711.8,785.0) -- (711.8,786.0) -- (711.6,787.0) -- (711.2,788.0) -- (710.8,788.9) -- (710.2,789.8) -- (709.5,790.5) -- (708.8,791.2) -- (707.9,791.8) -- (707.0,792.2) -- (706.0,792.6) -- (705.0,792.8) -- (704.0,792.8) -- cycle;
\node[anchor=east,inner sep=0pt,font=\fontsize{5.4}{5.4}\selectfont,text=black] at (693.9,774.9) {61};
\node[anchor=east,inner sep=0pt,font=\fontsize{5.4}{5.4}\selectfont,text=black] at (693.9,803.4) {57};
\node[anchor=west,inner sep=0pt,font=\fontsize{5.4}{5.4}\selectfont,text=black] at (714.1,774.9) {044};
\fill[black] (683.2,789.6) circle (2.0);
\fill[black] (689.8,789.6) circle (2.0);
\draw[black,line width=0.51pt] (787.0,798.0) -- (787.0,834.8);
\draw[black,line width=0.51pt] (787.0,834.8) -- (773.8,841.4);
\draw[black,line width=0.51pt,fill=white] (787.0,798.0) circle (7.8);
\fill[black] (787.0,798.0) -- (787.0,790.2) -- (788.0,790.2) -- (789.0,790.4) -- (790.0,790.8) -- (790.9,791.2) -- (791.8,791.8) -- (792.5,792.5) -- (793.2,793.2) -- (793.8,794.1) -- (794.2,795.0) -- (794.6,796.0) -- (794.8,797.0) -- (794.8,798.0) -- (794.8,799.0) -- (794.6,800.0) -- (794.2,801.0) -- (793.8,801.9) -- (793.2,802.8) -- (792.5,803.5) -- (791.8,804.2) -- (790.9,804.8) -- (790.0,805.2) -- (789.0,805.6) -- (788.0,805.8) -- (787.0,805.8) -- (786.0,805.8) -- (785.0,805.6) -- (784.0,805.2) -- (783.1,804.8) -- (782.2,804.2) -- (781.5,803.5) -- (780.8,802.8) -- (780.2,801.9) -- (779.8,801.0) -- (779.4,800.0) -- (779.2,799.0) -- (779.2,798.0) -- cycle;
\node[anchor=east,inner sep=0pt,font=\fontsize{5.4}{5.4}\selectfont,text=black] at (776.9,787.9) {66};
\node[anchor=east,inner sep=0pt,font=\fontsize{5.4}{5.4}\selectfont,text=black] at (776.9,816.4) {53};
\node[anchor=west,inner sep=0pt,font=\fontsize{5.4}{5.4}\selectfont,text=black] at (797.1,787.9) {099};
\draw[black,line width=0.51pt] (828.0,818.0) -- (828.0,854.8);
\draw[black,line width=0.51pt] (828.0,854.8) -- (814.8,861.4);
\draw[black,line width=0.51pt,fill=white] (828.0,818.0) circle (7.8);
\fill[black] (828.0,818.0) -- (828.0,810.2) -- (829.0,810.2) -- (830.0,810.4) -- (831.0,810.8) -- (831.9,811.2) -- (832.8,811.8) -- (833.5,812.5) -- (834.2,813.2) -- (834.8,814.1) -- (835.2,815.0) -- (835.6,816.0) -- (835.8,817.0) -- (835.8,818.0) -- cycle;
\node[anchor=east,inner sep=0pt,font=\fontsize{5.4}{5.4}\selectfont,text=black] at (817.9,807.9) {66};
\node[anchor=east,inner sep=0pt,font=\fontsize{5.4}{5.4}\selectfont,text=black] at (817.9,836.4) {55};
\node[anchor=west,inner sep=0pt,font=\fontsize{5.4}{5.4}\selectfont,text=black] at (838.1,807.9) {111};
\draw[black,line width=0.51pt] (405.0,838.0) -- (374.9,816.9);
\draw[black,line width=0.51pt] (374.9,816.9) -- (377.0,802.2);
\draw[black,line width=0.51pt,fill=white] (405.0,838.0) circle (7.8);
\node[anchor=east,inner sep=0pt,font=\fontsize{5.4}{5.4}\selectfont,text=black] at (394.9,827.9) {41};
\node[anchor=east,inner sep=0pt,font=\fontsize{5.4}{5.4}\selectfont,text=black] at (394.9,856.4) {33};
\node[anchor=west,inner sep=0pt,font=\fontsize{5.4}{5.4}\selectfont,text=black] at (415.1,827.9) {126};
\draw[black,line width=0.51pt] (496.0,831.0) -- (467.8,807.3);
\draw[black,line width=0.51pt] (467.8,807.3) -- (471.3,792.9);
\draw[black,line width=0.51pt,fill=white] (496.0,831.0) circle (7.8);
\node[anchor=east,inner sep=0pt,font=\fontsize{5.4}{5.4}\selectfont,text=black] at (485.9,820.9) {42};
\node[anchor=east,inner sep=0pt,font=\fontsize{5.4}{5.4}\selectfont,text=black] at (485.9,849.4) {36};
\node[anchor=west,inner sep=0pt,font=\fontsize{5.4}{5.4}\selectfont,text=black] at (506.1,820.9) {109};
\draw[black,line width=0.51pt] (608.0,861.0) -- (588.5,892.2);
\draw[black,line width=0.51pt] (588.5,892.2) -- (573.8,890.8);
\draw[black,line width=0.51pt,fill=white] (608.0,861.0) circle (7.8);
\fill[black] (608.0,861.0) circle (7.8);
\node[anchor=east,inner sep=0pt,font=\fontsize{5.4}{5.4}\selectfont,text=black] at (597.9,850.9) {61};
\node[anchor=east,inner sep=0pt,font=\fontsize{5.4}{5.4}\selectfont,text=black] at (597.9,879.4) {57};
\node[anchor=west,inner sep=0pt,font=\fontsize{5.4}{5.4}\selectfont,text=black] at (618.1,850.9) {075};
\fill[black] (587.2,865.6) circle (2.0);
\fill[black] (593.8,865.6) circle (2.0);
\draw[black,line width=0.51pt] (828.0,885.0) -- (828.0,921.8);
\draw[black,line width=0.51pt] (828.0,921.8) -- (814.8,928.4);
\draw[black,line width=0.51pt,fill=white] (828.0,885.0) circle (7.8);
\fill[black] (828.0,885.0) -- (828.0,877.2) -- (829.0,877.2) -- (830.0,877.4) -- (831.0,877.8) -- (831.9,878.2) -- (832.8,878.8) -- (833.5,879.5) -- (834.2,880.2) -- (834.8,881.1) -- (835.2,882.0) -- (835.6,883.0) -- (835.8,884.0) -- (835.8,885.0) -- (835.8,886.0) -- (835.6,887.0) -- (835.2,888.0) -- (834.8,888.9) -- (834.2,889.8) -- (833.5,890.5) -- (832.8,891.2) -- (831.9,891.8) -- (831.0,892.2) -- (830.0,892.6) -- (829.0,892.8) -- (828.0,892.8) -- cycle;
\node[anchor=east,inner sep=0pt,font=\fontsize{5.4}{5.4}\selectfont,text=black] at (817.9,874.9) {71};
\node[anchor=east,inner sep=0pt,font=\fontsize{5.4}{5.4}\selectfont,text=black] at (817.9,903.4) {60};
\node[anchor=west,inner sep=0pt,font=\fontsize{5.4}{5.4}\selectfont,text=black] at (838.1,874.9) {149};
\draw[black,line width=0.51pt] (498.0,909.0) -- (465.5,891.7);
\draw[black,line width=0.51pt] (465.5,891.7) -- (465.9,876.9);
\draw[black,line width=0.51pt,fill=white] (498.0,909.0) circle (7.8);
\node[anchor=east,inner sep=0pt,font=\fontsize{5.4}{5.4}\selectfont,text=black] at (487.9,898.9) {48};
\node[anchor=east,inner sep=0pt,font=\fontsize{5.4}{5.4}\selectfont,text=black] at (487.9,927.4) {38};
\node[anchor=west,inner sep=0pt,font=\fontsize{5.4}{5.4}\selectfont,text=black] at (508.1,898.9) {126};
\draw[black,line width=0.51pt] (729.0,914.0) -- (729.0,950.8);
\draw[black,line width=0.51pt] (729.0,950.8) -- (715.8,957.4);
\draw[black,line width=0.51pt,fill=white] (729.0,914.0) circle (7.8);
\fill[black] (729.0,914.0) -- (729.0,906.2) -- (730.0,906.2) -- (731.0,906.4) -- (732.0,906.8) -- (732.9,907.2) -- (733.8,907.8) -- (734.5,908.5) -- (735.2,909.2) -- (735.8,910.1) -- (736.2,911.0) -- (736.6,912.0) -- (736.8,913.0) -- (736.8,914.0) -- (736.8,915.0) -- (736.6,916.0) -- (736.2,917.0) -- (735.8,917.9) -- (735.2,918.8) -- (734.5,919.5) -- (733.8,920.2) -- (732.9,920.8) -- (732.0,921.2) -- (731.0,921.6) -- (730.0,921.8) -- (729.0,921.8) -- cycle;
\node[anchor=east,inner sep=0pt,font=\fontsize{5.4}{5.4}\selectfont,text=black] at (718.9,903.9) {70};
\node[anchor=east,inner sep=0pt,font=\fontsize{5.4}{5.4}\selectfont,text=black] at (718.9,932.4) {59};
\node[anchor=west,inner sep=0pt,font=\fontsize{5.4}{5.4}\selectfont,text=black] at (739.1,903.9) {140};
\draw[black,line width=0.51pt] (390.0,946.0) -- (356.6,930.4);
\draw[black,line width=0.51pt] (356.6,930.4) -- (356.2,915.6);
\draw[black,line width=0.51pt,fill=white] (390.0,946.0) circle (7.8);
\node[anchor=east,inner sep=0pt,font=\fontsize{5.4}{5.4}\selectfont,text=black] at (379.9,935.9) {49};
\node[anchor=east,inner sep=0pt,font=\fontsize{5.4}{5.4}\selectfont,text=black] at (379.9,964.4) {38};
\node[anchor=west,inner sep=0pt,font=\fontsize{5.4}{5.4}\selectfont,text=black] at (400.1,935.9) {128};
\draw[black,line width=0.51pt] (603.0,945.0) -- (573.2,966.6);
\draw[black,line width=0.51pt] (573.2,966.6) -- (560.1,959.8);
\draw[black,line width=0.51pt,fill=white] (603.0,945.0) circle (7.8);
\fill[black] (603.0,945.0) -- (603.0,937.2) -- (604.0,937.2) -- (605.0,937.4) -- (606.0,937.8) -- (606.9,938.2) -- (607.8,938.8) -- (608.5,939.5) -- (609.2,940.2) -- (609.8,941.1) -- (610.2,942.0) -- (610.6,943.0) -- (610.8,944.0) -- (610.8,945.0) -- (610.8,946.0) -- (610.6,947.0) -- (610.2,948.0) -- (609.8,948.9) -- (609.2,949.8) -- (608.5,950.5) -- (607.8,951.2) -- (606.9,951.8) -- (606.0,952.2) -- (605.0,952.6) -- (604.0,952.8) -- (603.0,952.8) -- (602.0,952.8) -- (601.0,952.6) -- (600.0,952.2) -- (599.1,951.8) -- (598.2,951.2) -- (597.5,950.5) -- (596.8,949.8) -- (596.2,948.9) -- (595.8,948.0) -- (595.4,947.0) -- (595.2,946.0) -- (595.2,945.0) -- cycle;
\node[anchor=east,inner sep=0pt,font=\fontsize{5.4}{5.4}\selectfont,text=black] at (592.9,934.9) {66};
\node[anchor=east,inner sep=0pt,font=\fontsize{5.4}{5.4}\selectfont,text=black] at (592.9,963.4) {59};
\node[anchor=west,inner sep=0pt,font=\fontsize{5.4}{5.4}\selectfont,text=black] at (613.1,934.9) {114};
\draw[black,line width=0.51pt] (648.0,960.0) -- (631.3,992.8);
\draw[black,line width=0.51pt] (631.3,992.8) -- (616.5,992.7);
\draw[black,line width=0.51pt,fill=white] (648.0,960.0) circle (7.8);
\node[anchor=east,inner sep=0pt,font=\fontsize{5.4}{5.4}\selectfont,text=black] at (637.9,949.9) {66};
\node[anchor=east,inner sep=0pt,font=\fontsize{5.4}{5.4}\selectfont,text=black] at (637.9,978.4) {60};
\node[anchor=west,inner sep=0pt,font=\fontsize{5.4}{5.4}\selectfont,text=black] at (658.1,949.9) {124};
\draw[black,line width=0.51pt] (452.0,1009.0) -- (423.8,985.3);
\draw[black,line width=0.51pt] (423.8,985.3) -- (427.3,970.9);
\draw[black,line width=0.51pt,fill=white] (452.0,1009.0) circle (7.8);
\node[anchor=east,inner sep=0pt,font=\fontsize{5.4}{5.4}\selectfont,text=black] at (441.9,998.9) {50};
\node[anchor=east,inner sep=0pt,font=\fontsize{5.4}{5.4}\selectfont,text=black] at (441.9,1027.4) {38};
\node[anchor=west,inner sep=0pt,font=\fontsize{5.4}{5.4}\selectfont,text=black] at (462.1,998.9) {120};
\draw[black,line width=0.51pt] (507.0,1061.0) -- (471.5,1070.5);
\draw[black,line width=0.51pt] (471.5,1070.5) -- (461.6,1059.4);
\draw[black,line width=0.51pt,fill=white] (507.0,1061.0) circle (7.8);
\fill[black] (507.0,1061.0) circle (7.8);
\node[anchor=east,inner sep=0pt,font=\fontsize{5.4}{5.4}\selectfont,text=black] at (496.9,1050.9) {52};
\node[anchor=east,inner sep=0pt,font=\fontsize{5.4}{5.4}\selectfont,text=black] at (496.9,1079.4) {50};
\node[anchor=west,inner sep=0pt,font=\fontsize{5.4}{5.4}\selectfont,text=black] at (517.1,1050.9) {128};
\draw[black,line width=0.51pt] (613.0,1041.0) -- (606.0,1077.1);
\draw[black,line width=0.51pt] (606.0,1077.1) -- (591.7,1081.1);
\draw[black,line width=0.51pt,fill=white] (613.0,1041.0) circle (7.8);
\fill[black] (613.0,1041.0) circle (7.8);
\node[anchor=east,inner sep=0pt,font=\fontsize{5.4}{5.4}\selectfont,text=black] at (602.9,1030.9) {63};
\node[anchor=east,inner sep=0pt,font=\fontsize{5.4}{5.4}\selectfont,text=black] at (602.9,1059.4) {61};
\node[anchor=west,inner sep=0pt,font=\fontsize{5.4}{5.4}\selectfont,text=black] at (623.1,1030.9) {153};
\draw[black,line width=0.51pt] (725.0,1031.0) -- (710.6,1064.9);
\draw[black,line width=0.51pt] (710.6,1064.9) -- (695.8,1065.8);
\draw[black,line width=0.51pt,fill=white] (725.0,1031.0) circle (7.8);
\fill[black] (725.0,1031.0) -- (725.0,1023.2) -- (726.0,1023.2) -- (727.0,1023.4) -- (728.0,1023.8) -- (728.9,1024.2) -- (729.8,1024.8) -- (730.5,1025.5) -- (731.2,1026.2) -- (731.8,1027.1) -- (732.2,1028.0) -- (732.6,1029.0) -- (732.8,1030.0) -- (732.8,1031.0) -- cycle;
\node[anchor=east,inner sep=0pt,font=\fontsize{5.4}{5.4}\selectfont,text=black] at (714.9,1020.9) {69};
\node[anchor=east,inner sep=0pt,font=\fontsize{5.4}{5.4}\selectfont,text=black] at (714.9,1049.4) {61};
\node[anchor=west,inner sep=0pt,font=\fontsize{5.4}{5.4}\selectfont,text=black] at (735.1,1020.9) {171};
\draw[black,line width=0.51pt] (822.0,1006.0) -- (806.4,1039.4);
\draw[black,line width=0.51pt] (806.4,1039.4) -- (791.6,1039.8);
\draw[black,line width=0.51pt,fill=white] (822.0,1006.0) circle (7.8);
\node[anchor=east,inner sep=0pt,font=\fontsize{5.4}{5.4}\selectfont,text=black] at (811.9,995.9) {72};
\node[anchor=east,inner sep=0pt,font=\fontsize{5.4}{5.4}\selectfont,text=black] at (811.9,1024.4) {64};
\node[anchor=west,inner sep=0pt,font=\fontsize{5.4}{5.4}\selectfont,text=black] at (832.1,995.9) {194};
}
\caption{Step 5: the areas receiving precipitation (blue shading) at the ten stations that report rain or thunderstorms.}
\end{figure}

\begin{landscape}
\begin{figure}[p]
\centering
\mapbase{0.0225cm}{\path[cloudfill,even odd rule] (567.4,526.0) -- (567.7,526.0) -- (568.1,526.0) -- (568.8,525.9) -- (569.7,525.9) -- (570.7,525.8) -- (571.9,525.8) -- (573.2,525.7) -- (574.8,525.7) -- (576.2,525.6) -- (577.8,525.6) -- (579.2,525.5) -- (580.8,525.4) -- (582.2,525.4) -- (583.8,525.3) -- (585.2,525.3) -- (586.8,525.3) -- (588.2,525.2) -- (589.8,525.2) -- (591.2,525.1) -- (592.8,525.1) -- (594.2,525.1) -- (595.8,525.0) -- (597.2,525.0) -- (598.8,525.0) -- (600.2,525.0) -- (601.8,524.9) -- (603.2,524.9) -- (604.8,524.9) -- (606.2,524.9) -- (607.8,524.8) -- (609.2,524.8) -- (610.8,524.8) -- (612.2,524.8) -- (613.8,524.8) -- (615.2,524.8) -- (616.8,524.7) -- (618.2,524.7) -- (619.8,524.7) -- (621.2,524.7) -- (622.8,524.7) -- (624.2,524.7) -- (625.8,524.7) -- (627.2,524.7) -- (628.8,524.6) -- (630.2,524.6) -- (631.8,524.6) -- (633.2,524.6) -- (634.8,524.6) -- (636.2,524.6) -- (637.8,524.6) -- (639.2,524.6) -- (640.8,524.6) -- (642.2,524.6) -- (643.8,524.6) -- (645.2,524.6) -- (646.8,524.6) -- (648.2,524.6) -- (649.8,524.6) -- (651.2,524.6) -- (652.8,524.6) -- (654.2,524.6) -- (655.8,524.6) -- (657.2,524.6) -- (658.8,524.6) -- (660.2,524.6) -- (661.8,524.6) -- (663.2,524.6) -- (664.8,524.6) -- (666.2,524.6) -- (667.8,524.6) -- (669.2,524.6) -- (670.8,524.6) -- (672.2,524.6) -- (673.8,524.6) -- (675.2,524.6) -- (676.8,524.6) -- (678.2,524.6) -- (679.8,524.6) -- (681.2,524.6) -- (682.8,524.6) -- (684.2,524.6) -- (685.8,524.6) -- (687.2,524.7) -- (688.8,524.7) -- (690.0,524.8) -- (690.9,524.9) -- (691.5,525.2) -- (691.8,525.5) -- (692.4,526.0) -- (693.1,526.8) -- (694.1,527.8) -- (695.2,528.9) -- (696.2,529.9) -- (696.9,530.7) -- (697.5,531.2) -- (697.8,531.5) -- (698.3,532.1) -- (699.1,532.8) -- (700.1,533.8) -- (701.2,535.0) -- (702.5,535.9) -- (703.8,536.5) -- (705.2,536.8) -- (706.8,536.8) -- (708.2,536.5) -- (709.5,535.9) -- (710.8,535.0) -- (712.0,533.8) -- (713.0,532.8) -- (713.8,532.1) -- (714.3,531.6) -- (714.6,531.3) -- (715.2,531.1) -- (716.1,531.0) -- (717.2,530.9) -- (718.8,530.9) -- (720.2,530.9) -- (721.8,530.9) -- (723.2,531.0) -- (724.8,531.0) -- (726.2,531.0) -- (727.8,531.0) -- (729.2,531.0) -- (730.8,531.0) -- (732.2,531.1) -- (733.8,531.1) -- (735.2,531.1) -- (736.8,531.1) -- (738.2,531.1) -- (739.8,531.2) -- (741.2,531.2) -- (742.8,531.2) -- (744.2,530.9) -- (745.6,530.2) -- (747.0,529.3) -- (748.3,528.0) -- (749.4,526.9) -- (750.1,526.2) -- (750.6,525.7) -- (750.7,525.6) -- (751.2,525.5) -- (752.1,525.4) -- (753.2,525.3) -- (754.8,525.4) -- (756.2,525.4) -- (757.8,525.4) -- (759.2,525.4) -- (760.8,525.4) -- (762.2,525.5) -- (763.8,525.5) -- (765.2,525.5) -- (766.8,525.5) -- (768.2,525.5) -- (769.8,525.5) -- (771.2,525.6) -- (772.8,525.6) -- (774.2,525.6) -- (775.8,525.6) -- (777.2,525.6) -- (778.8,525.6) -- (780.2,525.7) -- (781.8,525.7) -- (783.2,525.7) -- (784.8,525.7) -- (786.2,525.7) -- (787.8,525.7) -- (789.2,525.7) -- (790.8,525.7) -- (792.2,525.8) -- (793.8,525.8) -- (795.2,525.8) -- (796.8,525.8) -- (798.2,525.8) -- (799.8,525.8) -- (801.2,525.8) -- (802.8,525.8) -- (804.2,525.8) -- (805.8,525.9) -- (807.2,525.9) -- (808.8,525.9) -- (810.2,525.9) -- (811.8,525.9) -- (813.2,525.9) -- (814.8,525.9) -- (816.2,525.9) -- (817.8,525.9) -- (819.2,525.9) -- (820.8,525.9) -- (822.2,526.0) -- (823.8,526.0) -- (825.2,526.0) -- (826.8,526.0) -- (828.2,526.0) -- (829.6,526.0) -- (830.9,526.0) -- (832.2,526.0) -- (833.2,526.3) -- (833.9,526.9) -- (834.4,527.7) -- (834.6,528.9) -- (834.8,529.8) -- (835.0,530.5) -- (835.1,531.1) -- (835.2,531.5) -- (835.4,532.1) -- (835.6,533.0) -- (836.0,534.2) -- (836.4,535.8) -- (836.9,537.2) -- (837.4,538.8) -- (837.9,540.2) -- (838.5,541.8) -- (839.0,543.1) -- (839.5,544.3) -- (840.0,545.4) -- (840.4,546.3) -- (840.8,547.1) -- (841.1,547.8) -- (841.5,548.5) -- (841.8,549.1) -- (842.3,549.9) -- (843.0,551.0) -- (843.8,552.2) -- (844.9,553.8) -- (845.7,554.9) -- (846.3,555.8) -- (846.6,556.3) -- (846.8,556.4) -- (847.2,556.9) -- (848.0,557.7) -- (849.2,558.7) -- (850.6,560.0) -- (851.7,561.0) -- (852.4,561.7) -- (852.8,562.1) -- (852.9,562.1) -- (853.3,562.3) -- (854.1,562.8) -- (855.2,563.4) -- (856.8,564.2) -- (858.2,564.9) -- (859.8,565.6) -- (861.2,566.1) -- (862.8,566.6) -- (864.1,567.0) -- (865.4,567.4) -- (866.6,567.6) -- (867.6,567.8) -- (868.5,567.9) -- (869.3,568.0) -- (869.9,568.1) -- (870.3,568.1) -- (871.0,568.2) -- (872.0,568.2) -- (873.2,568.2) -- (874.8,568.2) -- (876.0,568.2) -- (876.9,568.2) -- (877.6,568.1) -- (878.1,568.1) -- (878.7,568.0) -- (879.4,567.9) -- (880.3,567.8) -- (881.4,567.7) -- (882.6,567.5) -- (883.9,567.2) -- (885.2,566.9) -- (886.8,566.6) -- (888.2,566.2) -- (889.8,565.8) -- (891.2,565.4) -- (892.8,565.0) -- (894.2,564.5) -- (895.8,564.0) -- (897.2,563.5) -- (898.8,563.0) -- (899.9,562.5) -- (900.7,562.3) -- (901.1,562.1) -- (901.2,562.1) -- (901.7,561.9) -- (902.4,561.5) -- (903.5,561.0) -- (904.9,560.4) -- (906.3,559.7) -- (907.8,559.0) -- (909.2,558.3) -- (910.8,557.6) -- (911.9,557.4) -- (912.6,557.7) -- (913.0,558.6) -- (913.0,559.9) -- (913.0,561.3) -- (913.0,562.8) -- (913.0,564.2) -- (913.1,565.8) -- (913.1,567.2) -- (913.1,568.8) -- (913.1,570.2) -- (913.1,571.8) -- (913.1,573.2) -- (913.1,574.8) -- (913.1,576.2) -- (913.1,577.8) -- (913.2,579.2) -- (913.2,580.8) -- (913.2,582.2) -- (913.2,583.8) -- (913.2,585.2) -- (913.2,586.8) -- (913.2,588.2) -- (913.2,589.8) -- (913.3,591.2) -- (913.3,592.8) -- (913.3,594.2) -- (913.3,595.8) -- (913.3,597.2) -- (913.3,598.8) -- (913.4,600.2) -- (913.4,601.8) -- (913.4,603.2) -- (913.4,604.8) -- (913.4,606.2) -- (913.4,607.8) -- (913.5,609.2) -- (913.5,610.8) -- (913.5,612.2) -- (913.5,613.8) -- (913.9,615.2) -- (914.6,616.6) -- (915.6,618.0) -- (917.0,619.3) -- (918.4,620.3) -- (919.8,621.0) -- (921.2,621.3) -- (922.8,621.3) -- (924.2,621.3) -- (925.8,621.3) -- (927.2,621.3) -- (928.8,621.3) -- (930.2,621.3) -- (931.8,621.3) -- (933.2,621.2) -- (934.8,621.2) -- (935.9,621.3) -- (936.8,621.4) -- (937.3,621.5) -- (937.5,621.7) -- (938.0,622.2) -- (938.7,622.9) -- (939.7,623.9) -- (941.0,625.2) -- (942.4,626.1) -- (943.8,626.8) -- (945.2,627.1) -- (946.8,627.1) -- (947.9,627.1) -- (948.8,627.2) -- (949.4,627.4) -- (949.6,627.6) -- (950.1,628.1) -- (950.8,628.9) -- (951.8,629.9) -- (953.1,631.1) -- (954.4,632.0) -- (955.8,632.6) -- (957.2,632.9) -- (958.8,632.9) -- (959.9,633.0) -- (960.8,633.1) -- (961.4,633.3) -- (961.7,633.6) -- (962.2,634.1) -- (962.9,634.8) -- (963.9,635.8) -- (965.2,637.0) -- (966.2,638.0) -- (966.9,638.8) -- (967.4,639.3) -- (967.7,639.6) -- (968.2,640.1) -- (969.0,640.8) -- (970.0,641.8) -- (971.2,643.0) -- (972.2,644.0) -- (972.9,644.7) -- (973.5,645.2) -- (973.8,645.5) -- (974.3,646.1) -- (975.0,646.8) -- (976.0,647.8) -- (977.2,648.9) -- (978.2,649.9) -- (978.9,650.7) -- (979.5,651.2) -- (979.8,651.5) -- (980.3,652.1) -- (981.1,652.8) -- (982.1,653.8) -- (983.2,654.9) -- (984.2,655.9) -- (985.0,656.6) -- (985.5,657.2) -- (985.8,657.5) -- (986.1,658.1) -- (986.3,659.0) -- (986.3,660.2) -- (986.3,661.8) -- (986.6,663.2) -- (987.2,664.5) -- (988.1,665.7) -- (989.3,666.9) -- (990.2,667.8) -- (991.0,668.6) -- (991.5,669.1) -- (991.9,669.5) -- (992.1,670.1) -- (992.3,671.0) -- (992.4,672.2) -- (992.4,673.8) -- (992.4,675.2) -- (992.4,676.8) -- (992.4,678.2) -- (992.4,679.8) -- (992.4,681.2) -- (992.4,682.8) -- (992.4,684.2) -- (992.4,685.8) -- (992.4,687.2) -- (992.4,688.8) -- (992.4,690.2) -- (992.4,691.8) -- (992.4,693.2) -- (992.4,694.8) -- (992.4,696.2) -- (992.4,697.8) -- (992.4,699.2) -- (992.4,700.8) -- (992.4,702.2) -- (992.4,703.8) -- (992.4,705.2) -- (992.4,706.8) -- (992.4,708.2) -- (992.4,709.8) -- (992.4,711.2) -- (992.4,712.8) -- (992.4,714.2) -- (992.4,715.8) -- (992.4,717.2) -- (992.4,718.8) -- (992.4,720.2) -- (992.4,721.8) -- (992.4,723.2) -- (992.4,724.8) -- (992.4,726.2) -- (992.4,727.8) -- (992.4,729.2) -- (992.4,730.8) -- (992.4,732.2) -- (992.4,733.8) -- (992.4,735.2) -- (992.4,736.8) -- (992.4,738.2) -- (992.4,739.8) -- (992.4,741.2) -- (992.4,742.8) -- (992.4,744.2) -- (992.4,745.8) -- (992.3,747.2) -- (992.3,748.8) -- (992.3,750.2) -- (992.3,751.8) -- (992.2,753.0) -- (992.1,753.9) -- (991.8,754.5) -- (991.5,754.8) -- (991.0,755.4) -- (990.2,756.1) -- (989.2,757.1) -- (988.0,758.2) -- (987.1,759.5) -- (986.5,760.8) -- (986.2,762.2) -- (986.2,763.8) -- (986.1,765.0) -- (986.0,765.9) -- (985.8,766.5) -- (985.5,766.8) -- (984.9,767.3) -- (984.2,768.0) -- (983.2,769.0) -- (982.0,770.2) -- (981.0,771.2) -- (980.2,771.9) -- (979.7,772.4) -- (979.4,772.7) -- (978.9,773.2) -- (978.1,774.0) -- (977.1,775.0) -- (975.9,776.2) -- (974.9,777.1) -- (974.1,777.9) -- (973.6,778.4) -- (973.4,778.6) -- (972.9,779.1) -- (972.1,779.9) -- (971.1,780.9) -- (969.8,782.1) -- (968.8,783.1) -- (968.0,783.8) -- (967.5,784.3) -- (967.3,784.5) -- (966.8,785.0) -- (966.0,785.8) -- (965.0,786.8) -- (963.7,788.1) -- (962.6,789.1) -- (961.9,789.8) -- (961.4,790.2) -- (961.2,790.4) -- (960.7,790.5) -- (959.9,790.6) -- (958.8,790.6) -- (957.2,790.6) -- (955.8,790.9) -- (954.3,791.6) -- (952.9,792.6) -- (951.4,794.0) -- (950.3,795.0) -- (949.6,795.7) -- (949.2,796.1) -- (949.1,796.2) -- (948.7,796.2) -- (947.9,796.3) -- (946.8,796.2) -- (945.2,796.2) -- (943.8,796.3) -- (942.2,796.5) -- (940.8,796.9) -- (939.2,797.4) -- (937.9,797.6) -- (936.6,797.6) -- (935.4,797.3) -- (934.4,796.8) -- (933.5,796.3) -- (932.7,796.0) -- (932.1,795.7) -- (931.7,795.4) -- (931.0,795.1) -- (930.0,794.6) -- (928.8,794.0) -- (927.2,793.3) -- (925.8,792.6) -- (924.5,792.0) -- (923.2,791.4) -- (922.0,790.8) -- (921.0,790.4) -- (920.2,790.0) -- (919.7,789.8) -- (919.4,789.7) -- (918.8,789.4) -- (917.9,789.0) -- (916.8,788.5) -- (915.2,787.8) -- (913.8,787.1) -- (912.2,786.5) -- (910.8,785.9) -- (909.2,785.2) -- (908.1,784.7) -- (907.2,784.4) -- (906.7,784.2) -- (906.5,784.1) -- (906.0,783.9) -- (905.3,783.6) -- (904.3,783.2) -- (903.0,782.7) -- (901.6,782.2) -- (900.2,781.6) -- (898.8,781.1) -- (897.2,780.5) -- (895.8,780.0) -- (894.5,779.5) -- (893.2,779.0) -- (892.0,778.6) -- (891.0,778.3) -- (890.2,778.0) -- (889.7,777.8) -- (889.4,777.7) -- (888.8,777.5) -- (887.9,777.3) -- (886.8,776.9) -- (885.2,776.4) -- (883.8,775.9) -- (882.2,775.4) -- (880.8,775.0) -- (879.2,774.5) -- (877.8,774.1) -- (876.2,773.7) -- (874.8,773.2) -- (873.2,772.8) -- (872.1,772.5) -- (871.2,772.3) -- (870.7,772.1) -- (870.4,772.1) -- (869.9,772.0) -- (869.2,771.8) -- (868.2,771.5) -- (866.9,771.2) -- (865.6,770.9) -- (864.2,770.5) -- (862.8,770.2) -- (861.2,769.8) -- (859.8,769.5) -- (858.2,769.2) -- (856.8,768.9) -- (855.2,768.5) -- (853.8,768.2) -- (852.2,767.9) -- (850.8,767.7) -- (849.2,767.4) -- (847.8,767.1) -- (846.3,766.9) -- (844.9,766.6) -- (843.4,766.4) -- (842.3,766.2) -- (841.6,766.1) -- (841.2,766.0) -- (841.1,766.0) -- (840.7,765.9) -- (839.9,765.8) -- (838.8,765.6) -- (837.2,765.3) -- (835.8,765.1) -- (834.2,764.9) -- (832.8,764.7) -- (831.2,764.4) -- (829.8,764.2) -- (828.2,764.0) -- (826.8,763.8) -- (825.2,763.6) -- (823.8,763.4) -- (822.2,763.2) -- (820.8,763.0) -- (819.2,762.8) -- (817.8,762.6) -- (816.2,762.4) -- (814.8,762.2) -- (813.2,762.0) -- (811.8,761.8) -- (810.2,761.6) -- (808.8,761.5) -- (807.2,761.3) -- (805.8,761.1) -- (804.2,760.9) -- (802.8,760.7) -- (801.2,760.5) -- (800.0,760.4) -- (799.0,760.2) -- (798.3,760.1) -- (797.8,760.1) -- (797.2,760.0) -- (796.5,759.9) -- (795.6,759.8) -- (794.5,759.7) -- (793.4,759.5) -- (792.1,759.3) -- (790.8,759.1) -- (789.2,758.9) -- (787.8,758.7) -- (786.2,758.5) -- (784.8,758.3) -- (783.2,758.1) -- (781.8,757.9) -- (780.2,757.7) -- (778.8,757.5) -- (777.2,757.3) -- (775.8,757.2) -- (774.2,757.0) -- (772.8,756.8) -- (771.2,756.6) -- (769.8,756.4) -- (768.2,756.3) -- (766.8,756.1) -- (765.2,755.9) -- (763.8,755.8) -- (762.2,755.6) -- (760.8,755.5) -- (759.2,755.4) -- (757.8,755.3) -- (756.2,755.2) -- (754.8,755.1) -- (753.2,755.1) -- (751.8,755.0) -- (750.2,755.0) -- (748.8,755.0) -- (747.2,755.0) -- (745.8,755.1) -- (744.2,755.1) -- (742.8,755.2) -- (741.2,755.3) -- (739.8,755.4) -- (738.2,755.6) -- (736.8,755.8) -- (735.2,756.0) -- (733.8,756.3) -- (732.2,756.6) -- (730.8,756.9) -- (729.2,757.3) -- (727.8,757.6) -- (726.2,758.1) -- (724.8,758.5) -- (723.2,759.0) -- (722.1,759.4) -- (721.3,759.7) -- (720.8,759.9) -- (720.7,759.9) -- (720.2,760.1) -- (719.5,760.4) -- (718.4,760.9) -- (717.0,761.5) -- (715.6,762.2) -- (714.2,762.9) -- (712.8,763.6) -- (711.2,764.5) -- (710.1,765.1) -- (709.3,765.5) -- (708.8,765.8) -- (708.7,765.9) -- (708.3,766.2) -- (707.5,766.7) -- (706.5,767.4) -- (705.1,768.3) -- (703.9,769.2) -- (702.8,769.9) -- (701.9,770.6) -- (701.2,771.2) -- (700.4,771.7) -- (699.7,772.3) -- (698.9,772.9) -- (698.2,773.6) -- (697.3,774.3) -- (696.5,775.0) -- (695.6,775.8) -- (694.6,776.7) -- (693.8,777.5) -- (693.0,778.2) -- (692.4,778.8) -- (691.8,779.3) -- (691.2,780.0) -- (690.4,780.7) -- (689.6,781.6) -- (688.7,782.5) -- (687.8,783.4) -- (687.1,784.2) -- (686.4,785.0) -- (685.9,785.6) -- (685.2,786.3) -- (684.6,787.1) -- (683.9,787.9) -- (683.2,788.7) -- (682.4,789.6) -- (681.7,790.5) -- (680.9,791.4) -- (680.1,792.4) -- (679.5,793.2) -- (678.9,794.0) -- (678.3,794.6) -- (677.9,795.2) -- (677.3,795.9) -- (676.6,796.9) -- (675.7,798.1) -- (674.6,799.6) -- (673.8,800.7) -- (673.2,801.4) -- (672.9,801.8) -- (672.9,801.9) -- (672.5,802.3) -- (672.0,803.1) -- (671.1,804.2) -- (670.0,805.8) -- (669.1,807.0) -- (668.3,808.0) -- (667.8,808.7) -- (667.5,809.1) -- (667.0,809.7) -- (666.5,810.5) -- (665.8,811.4) -- (664.9,812.4) -- (664.2,813.4) -- (663.4,814.4) -- (662.7,815.3) -- (662.0,816.1) -- (661.4,816.9) -- (660.7,817.7) -- (660.1,818.4) -- (659.6,819.0) -- (658.9,819.8) -- (658.1,820.6) -- (657.3,821.5) -- (656.4,822.4) -- (655.5,823.3) -- (654.8,824.1) -- (654.1,824.7) -- (653.5,825.2) -- (652.8,825.8) -- (652.1,826.4) -- (651.2,827.1) -- (650.3,827.8) -- (649.3,828.5) -- (648.2,829.3) -- (646.9,830.0) -- (645.6,830.8) -- (644.5,831.4) -- (643.8,831.9) -- (643.3,832.1) -- (643.2,832.2) -- (642.7,832.3) -- (641.9,832.5) -- (640.8,832.8) -- (639.2,833.2) -- (637.8,833.4) -- (636.2,833.6) -- (634.8,833.6) -- (633.2,833.5) -- (632.0,833.3) -- (630.9,833.1) -- (630.0,832.8) -- (629.3,832.5) -- (628.6,832.1) -- (627.9,831.7) -- (627.1,831.3) -- (626.2,830.8) -- (625.4,830.2) -- (624.5,829.4) -- (623.7,828.5) -- (622.8,827.5) -- (621.9,826.6) -- (621.2,825.7) -- (620.5,824.8) -- (619.9,824.0) -- (619.4,823.2) -- (618.9,822.5) -- (618.5,821.8) -- (618.1,821.1) -- (617.6,820.2) -- (617.1,819.1) -- (616.5,817.8) -- (615.8,816.2) -- (615.1,814.9) -- (614.6,813.7) -- (614.1,812.7) -- (613.6,811.8) -- (613.3,811.0) -- (612.9,810.2) -- (612.7,809.5) -- (612.4,808.9) -- (612.1,808.1) -- (611.7,807.0) -- (611.2,805.8) -- (610.6,804.2) -- (610.0,802.8) -- (609.4,801.2) -- (608.8,799.8) -- (608.1,798.2) -- (607.6,797.1) -- (607.3,796.3) -- (607.1,795.9) -- (607.1,795.8) -- (606.9,795.4) -- (606.6,794.6) -- (606.1,793.5) -- (605.6,792.1) -- (605.0,790.7) -- (604.3,789.2) -- (603.7,787.8) -- (603.0,786.2) -- (602.4,785.0) -- (601.9,784.1) -- (601.6,783.4) -- (601.3,782.9) -- (601.0,782.3) -- (600.7,781.6) -- (600.3,780.7) -- (599.8,779.6) -- (599.2,778.4) -- (598.5,777.1) -- (597.8,775.8) -- (597.0,774.2) -- (596.4,773.0) -- (595.9,772.1) -- (595.5,771.5) -- (595.3,771.2) -- (595.0,770.6) -- (594.6,769.9) -- (594.1,768.9) -- (593.4,767.7) -- (592.7,766.5) -- (592.0,765.2) -- (591.2,763.9) -- (590.3,762.4) -- (589.6,761.4) -- (589.2,760.6) -- (588.9,760.2) -- (588.9,760.1) -- (588.6,759.7) -- (588.1,758.9) -- (587.4,757.8) -- (586.5,756.2) -- (585.7,754.9) -- (585.0,753.8) -- (584.3,752.8) -- (583.8,752.0) -- (583.3,751.3) -- (582.8,750.5) -- (582.3,749.8) -- (581.9,749.1) -- (581.3,748.2) -- (580.5,747.1) -- (579.6,745.8) -- (578.6,744.2) -- (577.8,743.1) -- (577.3,742.4) -- (577.0,742.0) -- (577.0,742.0) -- (576.8,741.6) -- (576.2,740.8) -- (575.5,739.7) -- (574.5,738.2) -- (573.6,736.9) -- (572.8,735.9) -- (572.2,735.0) -- (571.7,734.3) -- (571.2,733.6) -- (570.7,732.8) -- (570.2,732.0) -- (569.6,731.2) -- (569.0,730.3) -- (568.2,729.2) -- (567.4,728.0) -- (566.4,726.7) -- (565.7,725.6) -- (565.2,724.9) -- (564.8,724.4) -- (564.7,724.2) -- (564.4,723.7) -- (563.8,722.9) -- (563.0,721.8) -- (562.0,720.2) -- (561.2,719.0) -- (560.5,718.0) -- (560.0,717.2) -- (559.6,716.6) -- (559.1,715.9) -- (558.7,715.2) -- (558.1,714.3) -- (557.5,713.4) -- (556.8,712.3) -- (556.1,711.1) -- (555.3,709.8) -- (554.4,708.4) -- (553.7,707.3) -- (553.2,706.5) -- (553.0,706.1) -- (552.9,706.1) -- (552.7,705.7) -- (552.3,704.9) -- (551.6,703.8) -- (550.8,702.2) -- (550.0,700.8) -- (549.3,699.4) -- (548.6,698.1) -- (548.0,696.9) -- (547.5,695.9) -- (547.1,695.1) -- (546.8,694.6) -- (546.7,694.4) -- (546.4,693.8) -- (546.0,692.9) -- (545.5,691.8) -- (544.9,690.2) -- (544.3,688.8) -- (543.6,687.2) -- (543.1,685.8) -- (542.5,684.2) -- (542.0,683.0) -- (541.7,682.0) -- (541.4,681.2) -- (541.2,680.6) -- (541.1,680.0) -- (540.8,679.2) -- (540.6,678.4) -- (540.3,677.4) -- (540.0,676.3) -- (539.7,675.1) -- (539.4,673.8) -- (539.1,672.2) -- (538.8,670.8) -- (538.6,669.2) -- (538.3,667.8) -- (538.1,666.2) -- (537.9,664.8) -- (537.8,663.2) -- (537.7,661.8) -- (537.6,660.2) -- (537.5,658.8) -- (537.4,657.2) -- (537.4,655.8) -- (537.4,654.2) -- (537.5,652.8) -- (537.6,651.2) -- (537.6,649.8) -- (537.8,648.2) -- (537.9,646.8) -- (538.1,645.2) -- (538.2,643.8) -- (538.4,642.2) -- (538.7,640.8) -- (538.9,639.2) -- (539.2,637.8) -- (539.5,636.2) -- (539.8,634.8) -- (540.0,633.5) -- (540.3,632.2) -- (540.6,631.1) -- (540.8,630.1) -- (541.0,629.3) -- (541.1,628.8) -- (541.2,628.5) -- (541.3,627.9) -- (541.6,627.0) -- (541.9,625.8) -- (542.3,624.2) -- (542.7,622.8) -- (543.2,621.2) -- (543.6,619.8) -- (544.1,618.2) -- (544.5,616.8) -- (545.0,615.2) -- (545.4,613.8) -- (545.9,612.2) -- (546.3,611.0) -- (546.5,610.1) -- (546.7,609.5) -- (546.8,609.2) -- (547.0,608.6) -- (547.3,607.9) -- (547.6,606.9) -- (547.9,605.8) -- (548.3,604.5) -- (548.8,603.2) -- (549.2,601.8) -- (549.7,600.2) -- (550.2,598.8) -- (550.6,597.2) -- (551.1,595.8) -- (551.6,594.2) -- (552.0,593.0) -- (552.3,591.9) -- (552.6,591.1) -- (552.7,590.5) -- (552.9,589.9) -- (553.2,589.1) -- (553.4,588.3) -- (553.7,587.4) -- (554.0,586.3) -- (554.4,585.1) -- (554.8,583.8) -- (555.2,582.2) -- (555.6,580.8) -- (556.1,579.2) -- (556.5,577.8) -- (556.9,576.2) -- (557.3,574.8) -- (557.7,573.3) -- (558.1,571.8) -- (558.4,570.4) -- (558.7,569.3) -- (558.9,568.5) -- (559.0,568.1) -- (559.0,568.1) -- (559.1,567.7) -- (559.3,566.9) -- (559.6,565.8) -- (560.0,564.2) -- (560.3,562.8) -- (560.7,561.2) -- (561.1,559.8) -- (561.4,558.2) -- (561.7,556.8) -- (562.1,555.2) -- (562.4,553.8) -- (562.7,552.2) -- (563.0,550.8) -- (563.3,549.2) -- (563.6,547.8) -- (563.9,546.2) -- (564.2,544.9) -- (564.4,543.7) -- (564.6,542.6) -- (564.8,541.7) -- (564.9,540.9) -- (565.0,540.2) -- (565.2,539.5) -- (565.3,538.9) -- (565.4,538.1) -- (565.6,537.0) -- (565.8,535.8) -- (566.1,534.2) -- (566.3,532.8) -- (566.6,531.2) -- (566.8,529.8) -- (567.1,528.2) -- (567.3,527.1) -- (567.4,526.4) -- (567.4,526.0) -- cycle  (554.8,1036.0) -- (555.1,1035.9) -- (555.6,1035.7) -- (556.4,1035.4) -- (557.4,1035.0) -- (558.6,1034.6) -- (559.9,1034.2) -- (561.2,1033.8) -- (562.8,1033.5) -- (564.2,1033.2) -- (565.8,1032.9) -- (567.2,1032.7) -- (568.8,1032.6) -- (570.2,1032.5) -- (571.8,1032.5) -- (573.2,1032.5) -- (574.8,1032.6) -- (576.2,1032.8) -- (577.8,1033.0) -- (579.2,1033.3) -- (580.8,1033.6) -- (582.2,1034.0) -- (583.6,1034.4) -- (585.0,1034.8) -- (586.4,1035.3) -- (587.4,1035.7) -- (588.2,1035.9) -- (588.6,1036.1) -- (588.8,1036.1) -- (589.3,1036.4) -- (590.1,1036.8) -- (591.2,1037.4) -- (592.8,1038.1) -- (594.2,1038.8) -- (595.5,1039.6) -- (596.7,1040.3) -- (597.8,1041.0) -- (598.7,1041.5) -- (599.5,1042.0) -- (600.0,1042.4) -- (600.4,1042.6) -- (601.1,1043.1) -- (602.0,1043.8) -- (603.2,1044.7) -- (604.8,1045.9) -- (605.9,1046.7) -- (606.7,1047.4) -- (607.2,1047.7) -- (607.3,1047.8) -- (607.7,1048.2) -- (608.5,1048.9) -- (609.5,1049.9) -- (610.9,1051.2) -- (612.0,1052.2) -- (612.8,1052.9) -- (613.3,1053.4) -- (613.6,1053.6) -- (614.1,1054.2) -- (614.8,1055.0) -- (615.8,1056.1) -- (617.1,1057.4) -- (618.1,1058.5) -- (618.8,1059.3) -- (619.2,1059.7) -- (619.3,1059.8) -- (619.6,1060.3) -- (620.3,1061.1) -- (621.2,1062.2) -- (622.4,1063.8) -- (623.3,1064.9) -- (624.0,1065.8) -- (624.5,1066.4) -- (624.7,1066.7) -- (625.0,1067.3) -- (625.5,1068.0) -- (626.2,1069.0) -- (627.0,1070.2) -- (627.8,1071.4) -- (628.5,1072.5) -- (629.3,1073.7) -- (630.0,1074.8) -- (630.5,1075.7) -- (631.0,1076.5) -- (631.3,1077.0) -- (631.5,1077.4) -- (631.9,1078.1) -- (632.4,1079.0) -- (633.1,1080.2) -- (633.9,1081.8) -- (634.6,1083.1) -- (635.2,1084.4) -- (635.8,1085.5) -- (636.3,1086.5) -- (636.7,1087.4) -- (637.0,1088.1) -- (637.3,1088.7) -- (637.6,1089.2) -- (637.9,1090.0) -- (638.3,1091.0) -- (638.9,1092.2) -- (639.5,1093.8) -- (640.1,1095.2) -- (640.7,1096.8) -- (641.3,1098.2) -- (641.9,1099.8) -- (642.3,1100.9) -- (642.7,1101.8) -- (642.8,1102.3) -- (642.9,1102.4) -- (643.1,1102.9) -- (643.3,1103.6) -- (643.7,1104.7) -- (644.1,1106.0) -- (644.6,1107.4) -- (645.0,1108.8) -- (645.5,1110.2) -- (645.9,1111.8) -- (646.4,1113.2) -- (646.8,1114.8) -- (647.3,1116.2) -- (647.7,1117.8) -- (648.1,1119.0) -- (648.4,1120.1) -- (648.6,1120.9) -- (648.8,1121.5) -- (648.9,1122.2) -- (649.1,1122.9) -- (649.3,1123.7) -- (649.6,1124.6) -- (649.7,1125.3) -- (649.7,1125.8) -- (649.6,1126.0) -- (649.3,1126.0) -- (648.8,1126.0) -- (647.9,1126.1) -- (646.8,1126.1) -- (645.2,1126.2) -- (643.8,1126.6) -- (642.3,1127.3) -- (640.9,1128.4) -- (639.6,1129.8) -- (638.5,1131.0) -- (637.7,1131.7) -- (637.3,1132.2) -- (637.2,1132.3) -- (636.7,1132.4) -- (635.9,1132.5) -- (634.8,1132.5) -- (633.2,1132.6) -- (631.8,1132.6) -- (630.2,1132.7) -- (628.8,1132.7) -- (627.2,1132.7) -- (625.8,1132.8) -- (624.2,1132.8) -- (622.8,1132.8) -- (621.2,1132.8) -- (619.8,1132.9) -- (618.2,1132.9) -- (616.8,1132.9) -- (615.2,1132.9) -- (613.8,1133.0) -- (612.2,1133.0) -- (610.8,1133.0) -- (609.2,1133.0) -- (607.8,1133.0) -- (606.2,1133.0) -- (604.8,1133.1) -- (603.2,1133.1) -- (601.8,1133.1) -- (600.2,1133.1) -- (598.8,1133.1) -- (597.2,1133.1) -- (595.8,1133.1) -- (594.2,1133.1) -- (592.8,1133.2) -- (591.2,1133.2) -- (590.1,1133.1) -- (589.2,1133.0) -- (588.6,1132.7) -- (588.3,1132.4) -- (587.7,1131.9) -- (587.0,1131.2) -- (586.0,1130.2) -- (584.8,1129.0) -- (583.5,1128.1) -- (582.2,1127.5) -- (580.8,1127.2) -- (579.2,1127.2) -- (577.8,1127.2) -- (576.2,1127.2) -- (574.8,1127.2) -- (573.2,1127.2) -- (571.8,1127.5) -- (570.5,1128.1) -- (569.2,1129.0) -- (568.0,1130.2) -- (567.0,1131.2) -- (566.3,1131.9) -- (565.8,1132.5) -- (565.5,1132.8) -- (564.9,1133.3) -- (564.2,1134.0) -- (563.2,1135.0) -- (562.0,1136.2) -- (561.0,1137.2) -- (560.3,1137.9) -- (559.8,1138.5) -- (559.5,1138.8) -- (558.9,1139.3) -- (558.2,1140.0) -- (557.2,1141.0) -- (556.0,1142.2) -- (555.1,1143.2) -- (554.3,1143.9) -- (553.8,1144.5) -- (553.5,1144.8) -- (552.9,1145.0) -- (552.0,1145.2) -- (550.8,1145.2) -- (549.2,1145.2) -- (547.8,1145.2) -- (546.2,1145.2) -- (544.8,1145.2) -- (543.2,1145.2) -- (541.8,1145.5) -- (540.5,1146.1) -- (539.2,1147.0) -- (538.0,1148.2) -- (537.0,1149.2) -- (536.3,1149.9) -- (535.7,1150.4) -- (535.4,1150.7) -- (534.8,1151.0) -- (533.9,1151.1) -- (532.8,1151.2) -- (531.2,1151.2) -- (529.8,1151.2) -- (528.2,1151.1) -- (526.8,1151.1) -- (525.2,1151.1) -- (523.8,1151.1) -- (522.2,1151.1) -- (520.8,1151.1) -- (519.2,1151.1) -- (517.8,1151.1) -- (516.2,1151.1) -- (514.8,1151.0) -- (513.2,1151.0) -- (511.8,1151.0) -- (510.2,1151.0) -- (508.8,1151.0) -- (507.2,1151.0) -- (505.8,1151.0) -- (504.2,1151.0) -- (502.8,1150.9) -- (501.2,1150.9) -- (499.8,1150.9) -- (498.2,1150.9) -- (496.8,1150.9) -- (495.2,1150.9) -- (493.8,1150.8) -- (492.2,1150.8) -- (490.8,1150.8) -- (489.2,1150.8) -- (487.8,1150.8) -- (486.2,1150.8) -- (484.8,1150.7) -- (483.2,1150.7) -- (481.8,1150.7) -- (480.2,1150.7) -- (478.8,1150.7) -- (477.2,1150.6) -- (476.1,1150.6) -- (475.3,1150.5) -- (474.8,1150.4) -- (474.6,1150.2) -- (474.2,1149.8) -- (473.4,1149.0) -- (472.4,1147.9) -- (471.0,1146.5) -- (469.6,1145.5) -- (468.2,1144.7) -- (466.8,1144.4) -- (465.2,1144.3) -- (464.1,1144.3) -- (463.3,1144.3) -- (462.9,1144.2) -- (462.8,1144.1) -- (462.4,1143.7) -- (461.6,1142.9) -- (460.6,1141.8) -- (459.1,1140.3) -- (458.1,1139.2) -- (457.3,1138.4) -- (457.0,1138.0) -- (457.0,1138.0) -- (457.0,1137.9) -- (457.0,1137.6) -- (457.0,1137.1) -- (457.0,1136.5) -- (457.2,1135.8) -- (457.6,1135.1) -- (458.2,1134.2) -- (459.0,1133.3) -- (459.8,1132.5) -- (460.6,1131.7) -- (461.3,1130.9) -- (462.0,1130.1) -- (462.7,1129.3) -- (463.3,1128.6) -- (464.0,1127.8) -- (464.7,1127.1) -- (465.4,1126.3) -- (466.1,1125.5) -- (466.9,1124.6) -- (467.7,1123.7) -- (468.5,1122.8) -- (469.2,1122.1) -- (469.7,1121.4) -- (470.2,1120.8) -- (470.9,1120.1) -- (471.6,1119.3) -- (472.5,1118.3) -- (473.5,1117.2) -- (474.3,1116.2) -- (475.0,1115.5) -- (475.5,1114.9) -- (475.8,1114.5) -- (476.3,1113.9) -- (477.1,1113.1) -- (478.1,1112.0) -- (479.2,1110.6) -- (480.2,1109.6) -- (480.8,1108.8) -- (481.2,1108.4) -- (481.3,1108.2) -- (481.8,1107.7) -- (482.5,1106.9) -- (483.5,1105.8) -- (484.8,1104.2) -- (485.8,1103.1) -- (486.5,1102.3) -- (486.8,1101.9) -- (486.9,1101.8) -- (487.3,1101.4) -- (487.9,1100.7) -- (488.9,1099.6) -- (490.1,1098.1) -- (491.1,1097.0) -- (491.8,1096.1) -- (492.3,1095.5) -- (492.6,1095.2) -- (493.1,1094.7) -- (493.7,1094.0) -- (494.5,1093.0) -- (495.6,1091.8) -- (496.5,1090.8) -- (497.2,1089.9) -- (497.8,1089.2) -- (498.3,1088.6) -- (498.9,1088.0) -- (499.5,1087.2) -- (500.2,1086.4) -- (501.1,1085.4) -- (501.8,1084.5) -- (502.6,1083.6) -- (503.3,1082.8) -- (504.0,1082.0) -- (504.6,1081.3) -- (505.3,1080.5) -- (505.9,1079.8) -- (506.5,1079.1) -- (507.2,1078.3) -- (508.0,1077.4) -- (508.8,1076.5) -- (509.7,1075.5) -- (510.4,1074.6) -- (511.1,1073.8) -- (511.6,1073.2) -- (512.1,1072.7) -- (512.7,1072.1) -- (513.4,1071.2) -- (514.3,1070.2) -- (515.4,1069.0) -- (516.3,1068.0) -- (516.9,1067.3) -- (517.4,1066.7) -- (517.7,1066.4) -- (518.2,1065.9) -- (518.9,1065.1) -- (519.9,1064.0) -- (521.2,1062.7) -- (522.1,1061.7) -- (522.8,1060.9) -- (523.3,1060.4) -- (523.4,1060.3) -- (523.9,1059.8) -- (524.6,1059.0) -- (525.7,1058.0) -- (527.0,1056.6) -- (528.0,1055.6) -- (528.8,1054.8) -- (529.2,1054.4) -- (529.4,1054.2) -- (529.9,1053.8) -- (530.6,1053.1) -- (531.7,1052.2) -- (533.0,1050.9) -- (534.1,1049.9) -- (534.9,1049.2) -- (535.5,1048.7) -- (535.8,1048.4) -- (536.3,1048.0) -- (537.1,1047.4) -- (538.1,1046.6) -- (539.2,1045.6) -- (540.3,1044.8) -- (541.3,1044.0) -- (542.2,1043.4) -- (543.0,1042.8) -- (543.7,1042.3) -- (544.5,1041.8) -- (545.2,1041.3) -- (545.9,1040.8) -- (546.8,1040.2) -- (547.9,1039.6) -- (549.2,1038.8) -- (550.8,1038.0) -- (552.0,1037.3) -- (553.0,1036.8) -- (553.7,1036.5) -- (554.1,1036.3) -- (554.5,1036.1) -- (554.7,1036.0) -- (554.8,1036.0) -- cycle ;
\draw[cloudline] (567.4,526.0) -- (567.7,526.0) -- (568.1,526.0) -- (568.8,525.9) -- (569.7,525.9) -- (570.7,525.8) -- (571.9,525.8) -- (573.2,525.7) -- (574.8,525.7) -- (576.2,525.6) -- (577.8,525.6) -- (579.2,525.5) -- (580.8,525.4) -- (582.2,525.4) -- (583.8,525.3) -- (585.2,525.3) -- (586.8,525.3) -- (588.2,525.2) -- (589.8,525.2) -- (591.2,525.1) -- (592.8,525.1) -- (594.2,525.1) -- (595.8,525.0) -- (597.2,525.0) -- (598.8,525.0) -- (600.2,525.0) -- (601.8,524.9) -- (603.2,524.9) -- (604.8,524.9) -- (606.2,524.9) -- (607.8,524.8) -- (609.2,524.8) -- (610.8,524.8) -- (612.2,524.8) -- (613.8,524.8) -- (615.2,524.8) -- (616.8,524.7) -- (618.2,524.7) -- (619.8,524.7) -- (621.2,524.7) -- (622.8,524.7) -- (624.2,524.7) -- (625.8,524.7) -- (627.2,524.7) -- (628.8,524.6) -- (630.2,524.6) -- (631.8,524.6) -- (633.2,524.6) -- (634.8,524.6) -- (636.2,524.6) -- (637.8,524.6) -- (639.2,524.6) -- (640.8,524.6) -- (642.2,524.6) -- (643.8,524.6) -- (645.2,524.6) -- (646.8,524.6) -- (648.2,524.6) -- (649.8,524.6) -- (651.2,524.6) -- (652.8,524.6) -- (654.2,524.6) -- (655.8,524.6) -- (657.2,524.6) -- (658.8,524.6) -- (660.2,524.6) -- (661.8,524.6) -- (663.2,524.6) -- (664.8,524.6) -- (666.2,524.6) -- (667.8,524.6) -- (669.2,524.6) -- (670.8,524.6) -- (672.2,524.6) -- (673.8,524.6) -- (675.2,524.6) -- (676.8,524.6) -- (678.2,524.6) -- (679.8,524.6) -- (681.2,524.6) -- (682.8,524.6) -- (684.2,524.6) -- (685.8,524.6) -- (687.2,524.7) -- (688.8,524.7) -- (690.0,524.8) -- (690.9,524.9) -- (691.5,525.2) -- (691.8,525.5) -- (692.4,526.0) -- (693.1,526.8) -- (694.1,527.8) -- (695.2,528.9) -- (696.2,529.9) -- (696.9,530.7) -- (697.5,531.2) -- (697.8,531.5) -- (698.3,532.1) -- (699.1,532.8) -- (700.1,533.8) -- (701.2,535.0) -- (702.5,535.9) -- (703.8,536.5) -- (705.2,536.8) -- (706.8,536.8) -- (708.2,536.5) -- (709.5,535.9) -- (710.8,535.0) -- (712.0,533.8) -- (713.0,532.8) -- (713.8,532.1) -- (714.3,531.6) -- (714.6,531.3) -- (715.2,531.1) -- (716.1,531.0) -- (717.2,530.9) -- (718.8,530.9) -- (720.2,530.9) -- (721.8,530.9) -- (723.2,531.0) -- (724.8,531.0) -- (726.2,531.0) -- (727.8,531.0) -- (729.2,531.0) -- (730.8,531.0) -- (732.2,531.1) -- (733.8,531.1) -- (735.2,531.1) -- (736.8,531.1) -- (738.2,531.1) -- (739.8,531.2) -- (741.2,531.2) -- (742.8,531.2) -- (744.2,530.9) -- (745.6,530.2) -- (747.0,529.3) -- (748.3,528.0) -- (749.4,526.9) -- (750.1,526.2) -- (750.6,525.7) -- (750.7,525.6) -- (751.2,525.5) -- (752.1,525.4) -- (753.2,525.3) -- (754.8,525.4) -- (756.2,525.4) -- (757.8,525.4) -- (759.2,525.4) -- (760.8,525.4) -- (762.2,525.5) -- (763.8,525.5) -- (765.2,525.5) -- (766.8,525.5) -- (768.2,525.5) -- (769.8,525.5) -- (771.2,525.6) -- (772.8,525.6) -- (774.2,525.6) -- (775.8,525.6) -- (777.2,525.6) -- (778.8,525.6) -- (780.2,525.7) -- (781.8,525.7) -- (783.2,525.7) -- (784.8,525.7) -- (786.2,525.7) -- (787.8,525.7) -- (789.2,525.7) -- (790.8,525.7) -- (792.2,525.8) -- (793.8,525.8) -- (795.2,525.8) -- (796.8,525.8) -- (798.2,525.8) -- (799.8,525.8) -- (801.2,525.8) -- (802.8,525.8) -- (804.2,525.8) -- (805.8,525.9) -- (807.2,525.9) -- (808.8,525.9) -- (810.2,525.9) -- (811.8,525.9) -- (813.2,525.9) -- (814.8,525.9) -- (816.2,525.9) -- (817.8,525.9) -- (819.2,525.9) -- (820.8,525.9) -- (822.2,526.0) -- (823.8,526.0) -- (825.2,526.0) -- (826.8,526.0) -- (828.2,526.0) -- (829.6,526.0) -- (830.9,526.0) -- (832.2,526.0) -- (833.2,526.3) -- (833.9,526.9) -- (834.4,527.7) -- (834.6,528.9) -- (834.8,529.8) -- (835.0,530.5) -- (835.1,531.1) -- (835.2,531.5) -- (835.4,532.1) -- (835.6,533.0) -- (836.0,534.2) -- (836.4,535.8) -- (836.9,537.2) -- (837.4,538.8) -- (837.9,540.2) -- (838.5,541.8) -- (839.0,543.1) -- (839.5,544.3) -- (840.0,545.4) -- (840.4,546.3) -- (840.8,547.1) -- (841.1,547.8) -- (841.5,548.5) -- (841.8,549.1) -- (842.3,549.9) -- (843.0,551.0) -- (843.8,552.2) -- (844.9,553.8) -- (845.7,554.9) -- (846.3,555.8) -- (846.6,556.3) -- (846.8,556.4) -- (847.2,556.9) -- (848.0,557.7) -- (849.2,558.7) -- (850.6,560.0) -- (851.7,561.0) -- (852.4,561.7) -- (852.8,562.1) -- (852.9,562.1) -- (853.3,562.3) -- (854.1,562.8) -- (855.2,563.4) -- (856.8,564.2) -- (858.2,564.9) -- (859.8,565.6) -- (861.2,566.1) -- (862.8,566.6) -- (864.1,567.0) -- (865.4,567.4) -- (866.6,567.6) -- (867.6,567.8) -- (868.5,567.9) -- (869.3,568.0) -- (869.9,568.1) -- (870.3,568.1) -- (871.0,568.2) -- (872.0,568.2) -- (873.2,568.2) -- (874.8,568.2) -- (876.0,568.2) -- (876.9,568.2) -- (877.6,568.1) -- (878.1,568.1) -- (878.7,568.0) -- (879.4,567.9) -- (880.3,567.8) -- (881.4,567.7) -- (882.6,567.5) -- (883.9,567.2) -- (885.2,566.9) -- (886.8,566.6) -- (888.2,566.2) -- (889.8,565.8) -- (891.2,565.4) -- (892.8,565.0) -- (894.2,564.5) -- (895.8,564.0) -- (897.2,563.5) -- (898.8,563.0) -- (899.9,562.5) -- (900.7,562.3) -- (901.1,562.1) -- (901.2,562.1) -- (901.7,561.9) -- (902.4,561.5) -- (903.5,561.0) -- (904.9,560.4) -- (906.3,559.7) -- (907.8,559.0) -- (909.2,558.3) -- (910.8,557.6) -- (911.9,557.4) -- (912.6,557.7) -- (913.0,558.6) -- (913.0,559.9) -- (913.0,561.3) -- (913.0,562.8) -- (913.0,564.2) -- (913.1,565.8) -- (913.1,567.2) -- (913.1,568.8) -- (913.1,570.2) -- (913.1,571.8) -- (913.1,573.2) -- (913.1,574.8) -- (913.1,576.2) -- (913.1,577.8) -- (913.2,579.2) -- (913.2,580.8) -- (913.2,582.2) -- (913.2,583.8) -- (913.2,585.2) -- (913.2,586.8) -- (913.2,588.2) -- (913.2,589.8) -- (913.3,591.2) -- (913.3,592.8) -- (913.3,594.2) -- (913.3,595.8) -- (913.3,597.2) -- (913.3,598.8) -- (913.4,600.2) -- (913.4,601.8) -- (913.4,603.2) -- (913.4,604.8) -- (913.4,606.2) -- (913.4,607.8) -- (913.5,609.2) -- (913.5,610.8) -- (913.5,612.2) -- (913.5,613.8) -- (913.9,615.2) -- (914.6,616.6) -- (915.6,618.0) -- (917.0,619.3) -- (918.4,620.3) -- (919.8,621.0) -- (921.2,621.3) -- (922.8,621.3) -- (924.2,621.3) -- (925.8,621.3) -- (927.2,621.3) -- (928.8,621.3) -- (930.2,621.3) -- (931.8,621.3) -- (933.2,621.2) -- (934.8,621.2) -- (935.9,621.3) -- (936.8,621.4) -- (937.3,621.5) -- (937.5,621.7) -- (938.0,622.2) -- (938.7,622.9) -- (939.7,623.9) -- (941.0,625.2) -- (942.4,626.1) -- (943.8,626.8) -- (945.2,627.1) -- (946.8,627.1) -- (947.9,627.1) -- (948.8,627.2) -- (949.4,627.4) -- (949.6,627.6) -- (950.1,628.1) -- (950.8,628.9) -- (951.8,629.9) -- (953.1,631.1) -- (954.4,632.0) -- (955.8,632.6) -- (957.2,632.9) -- (958.8,632.9) -- (959.9,633.0) -- (960.8,633.1) -- (961.4,633.3) -- (961.7,633.6) -- (962.2,634.1) -- (962.9,634.8) -- (963.9,635.8) -- (965.2,637.0) -- (966.2,638.0) -- (966.9,638.8) -- (967.4,639.3) -- (967.7,639.6) -- (968.2,640.1) -- (969.0,640.8) -- (970.0,641.8) -- (971.2,643.0) -- (972.2,644.0) -- (972.9,644.7) -- (973.5,645.2) -- (973.8,645.5) -- (974.3,646.1) -- (975.0,646.8) -- (976.0,647.8) -- (977.2,648.9) -- (978.2,649.9) -- (978.9,650.7) -- (979.5,651.2) -- (979.8,651.5) -- (980.3,652.1) -- (981.1,652.8) -- (982.1,653.8) -- (983.2,654.9) -- (984.2,655.9) -- (985.0,656.6) -- (985.5,657.2) -- (985.8,657.5) -- (986.1,658.1) -- (986.3,659.0) -- (986.3,660.2) -- (986.3,661.8) -- (986.6,663.2) -- (987.2,664.5) -- (988.1,665.7) -- (989.3,666.9) -- (990.2,667.8) -- (991.0,668.6) -- (991.5,669.1) -- (991.9,669.5) -- (992.1,670.1) -- (992.3,671.0) -- (992.4,672.2) -- (992.4,673.8) -- (992.4,675.2) -- (992.4,676.8) -- (992.4,678.2) -- (992.4,679.8) -- (992.4,681.2) -- (992.4,682.8) -- (992.4,684.2) -- (992.4,685.8) -- (992.4,687.2) -- (992.4,688.8) -- (992.4,690.2) -- (992.4,691.8) -- (992.4,693.2) -- (992.4,694.8) -- (992.4,696.2) -- (992.4,697.8) -- (992.4,699.2) -- (992.4,700.8) -- (992.4,702.2) -- (992.4,703.8) -- (992.4,705.2) -- (992.4,706.8) -- (992.4,708.2) -- (992.4,709.8) -- (992.4,711.2) -- (992.4,712.8) -- (992.4,714.2) -- (992.4,715.8) -- (992.4,717.2) -- (992.4,718.8) -- (992.4,720.2) -- (992.4,721.8) -- (992.4,723.2) -- (992.4,724.8) -- (992.4,726.2) -- (992.4,727.8) -- (992.4,729.2) -- (992.4,730.8) -- (992.4,732.2) -- (992.4,733.8) -- (992.4,735.2) -- (992.4,736.8) -- (992.4,738.2) -- (992.4,739.8) -- (992.4,741.2) -- (992.4,742.8) -- (992.4,744.2) -- (992.4,745.8) -- (992.3,747.2) -- (992.3,748.8) -- (992.3,750.2) -- (992.3,751.8) -- (992.2,753.0) -- (992.1,753.9) -- (991.8,754.5) -- (991.5,754.8) -- (991.0,755.4) -- (990.2,756.1) -- (989.2,757.1) -- (988.0,758.2) -- (987.1,759.5) -- (986.5,760.8) -- (986.2,762.2) -- (986.2,763.8) -- (986.1,765.0) -- (986.0,765.9) -- (985.8,766.5) -- (985.5,766.8) -- (984.9,767.3) -- (984.2,768.0) -- (983.2,769.0) -- (982.0,770.2) -- (981.0,771.2) -- (980.2,771.9) -- (979.7,772.4) -- (979.4,772.7) -- (978.9,773.2) -- (978.1,774.0) -- (977.1,775.0) -- (975.9,776.2) -- (974.9,777.1) -- (974.1,777.9) -- (973.6,778.4) -- (973.4,778.6) -- (972.9,779.1) -- (972.1,779.9) -- (971.1,780.9) -- (969.8,782.1) -- (968.8,783.1) -- (968.0,783.8) -- (967.5,784.3) -- (967.3,784.5) -- (966.8,785.0) -- (966.0,785.8) -- (965.0,786.8) -- (963.7,788.1) -- (962.6,789.1) -- (961.9,789.8) -- (961.4,790.2) -- (961.2,790.4) -- (960.7,790.5) -- (959.9,790.6) -- (958.8,790.6) -- (957.2,790.6) -- (955.8,790.9) -- (954.3,791.6) -- (952.9,792.6) -- (951.4,794.0) -- (950.3,795.0) -- (949.6,795.7) -- (949.2,796.1) -- (949.1,796.2) -- (948.7,796.2) -- (947.9,796.3) -- (946.8,796.2) -- (945.2,796.2) -- (943.8,796.3) -- (942.2,796.5) -- (940.8,796.9) -- (939.2,797.4) -- (937.9,797.6) -- (936.6,797.6) -- (935.4,797.3) -- (934.4,796.8) -- (933.5,796.3) -- (932.7,796.0) -- (932.1,795.7) -- (931.7,795.4) -- (931.0,795.1) -- (930.0,794.6) -- (928.8,794.0) -- (927.2,793.3) -- (925.8,792.6) -- (924.5,792.0) -- (923.2,791.4) -- (922.0,790.8) -- (921.0,790.4) -- (920.2,790.0) -- (919.7,789.8) -- (919.4,789.7) -- (918.8,789.4) -- (917.9,789.0) -- (916.8,788.5) -- (915.2,787.8) -- (913.8,787.1) -- (912.2,786.5) -- (910.8,785.9) -- (909.2,785.2) -- (908.1,784.7) -- (907.2,784.4) -- (906.7,784.2) -- (906.5,784.1) -- (906.0,783.9) -- (905.3,783.6) -- (904.3,783.2) -- (903.0,782.7) -- (901.6,782.2) -- (900.2,781.6) -- (898.8,781.1) -- (897.2,780.5) -- (895.8,780.0) -- (894.5,779.5) -- (893.2,779.0) -- (892.0,778.6) -- (891.0,778.3) -- (890.2,778.0) -- (889.7,777.8) -- (889.4,777.7) -- (888.8,777.5) -- (887.9,777.3) -- (886.8,776.9) -- (885.2,776.4) -- (883.8,775.9) -- (882.2,775.4) -- (880.8,775.0) -- (879.2,774.5) -- (877.8,774.1) -- (876.2,773.7) -- (874.8,773.2) -- (873.2,772.8) -- (872.1,772.5) -- (871.2,772.3) -- (870.7,772.1) -- (870.4,772.1) -- (869.9,772.0) -- (869.2,771.8) -- (868.2,771.5) -- (866.9,771.2) -- (865.6,770.9) -- (864.2,770.5) -- (862.8,770.2) -- (861.2,769.8) -- (859.8,769.5) -- (858.2,769.2) -- (856.8,768.9) -- (855.2,768.5) -- (853.8,768.2) -- (852.2,767.9) -- (850.8,767.7) -- (849.2,767.4) -- (847.8,767.1) -- (846.3,766.9) -- (844.9,766.6) -- (843.4,766.4) -- (842.3,766.2) -- (841.6,766.1) -- (841.2,766.0) -- (841.1,766.0) -- (840.7,765.9) -- (839.9,765.8) -- (838.8,765.6) -- (837.2,765.3) -- (835.8,765.1) -- (834.2,764.9) -- (832.8,764.7) -- (831.2,764.4) -- (829.8,764.2) -- (828.2,764.0) -- (826.8,763.8) -- (825.2,763.6) -- (823.8,763.4) -- (822.2,763.2) -- (820.8,763.0) -- (819.2,762.8) -- (817.8,762.6) -- (816.2,762.4) -- (814.8,762.2) -- (813.2,762.0) -- (811.8,761.8) -- (810.2,761.6) -- (808.8,761.5) -- (807.2,761.3) -- (805.8,761.1) -- (804.2,760.9) -- (802.8,760.7) -- (801.2,760.5) -- (800.0,760.4) -- (799.0,760.2) -- (798.3,760.1) -- (797.8,760.1) -- (797.2,760.0) -- (796.5,759.9) -- (795.6,759.8) -- (794.5,759.7) -- (793.4,759.5) -- (792.1,759.3) -- (790.8,759.1) -- (789.2,758.9) -- (787.8,758.7) -- (786.2,758.5) -- (784.8,758.3) -- (783.2,758.1) -- (781.8,757.9) -- (780.2,757.7) -- (778.8,757.5) -- (777.2,757.3) -- (775.8,757.2) -- (774.2,757.0) -- (772.8,756.8) -- (771.2,756.6) -- (769.8,756.4) -- (768.2,756.3) -- (766.8,756.1) -- (765.2,755.9) -- (763.8,755.8) -- (762.2,755.6) -- (760.8,755.5) -- (759.2,755.4) -- (757.8,755.3) -- (756.2,755.2) -- (754.8,755.1) -- (753.2,755.1) -- (751.8,755.0) -- (750.2,755.0) -- (748.8,755.0) -- (747.2,755.0) -- (745.8,755.1) -- (744.2,755.1) -- (742.8,755.2) -- (741.2,755.3) -- (739.8,755.4) -- (738.2,755.6) -- (736.8,755.8) -- (735.2,756.0) -- (733.8,756.3) -- (732.2,756.6) -- (730.8,756.9) -- (729.2,757.3) -- (727.8,757.6) -- (726.2,758.1) -- (724.8,758.5) -- (723.2,759.0) -- (722.1,759.4) -- (721.3,759.7) -- (720.8,759.9) -- (720.7,759.9) -- (720.2,760.1) -- (719.5,760.4) -- (718.4,760.9) -- (717.0,761.5) -- (715.6,762.2) -- (714.2,762.9) -- (712.8,763.6) -- (711.2,764.5) -- (710.1,765.1) -- (709.3,765.5) -- (708.8,765.8) -- (708.7,765.9) -- (708.3,766.2) -- (707.5,766.7) -- (706.5,767.4) -- (705.1,768.3) -- (703.9,769.2) -- (702.8,769.9) -- (701.9,770.6) -- (701.2,771.2) -- (700.4,771.7) -- (699.7,772.3) -- (698.9,772.9) -- (698.2,773.6) -- (697.3,774.3) -- (696.5,775.0) -- (695.6,775.8) -- (694.6,776.7) -- (693.8,777.5) -- (693.0,778.2) -- (692.4,778.8) -- (691.8,779.3) -- (691.2,780.0) -- (690.4,780.7) -- (689.6,781.6) -- (688.7,782.5) -- (687.8,783.4) -- (687.1,784.2) -- (686.4,785.0) -- (685.9,785.6) -- (685.2,786.3) -- (684.6,787.1) -- (683.9,787.9) -- (683.2,788.7) -- (682.4,789.6) -- (681.7,790.5) -- (680.9,791.4) -- (680.1,792.4) -- (679.5,793.2) -- (678.9,794.0) -- (678.3,794.6) -- (677.9,795.2) -- (677.3,795.9) -- (676.6,796.9) -- (675.7,798.1) -- (674.6,799.6) -- (673.8,800.7) -- (673.2,801.4) -- (672.9,801.8) -- (672.9,801.9) -- (672.5,802.3) -- (672.0,803.1) -- (671.1,804.2) -- (670.0,805.8) -- (669.1,807.0) -- (668.3,808.0) -- (667.8,808.7) -- (667.5,809.1) -- (667.0,809.7) -- (666.5,810.5) -- (665.8,811.4) -- (664.9,812.4) -- (664.2,813.4) -- (663.4,814.4) -- (662.7,815.3) -- (662.0,816.1) -- (661.4,816.9) -- (660.7,817.7) -- (660.1,818.4) -- (659.6,819.0) -- (658.9,819.8) -- (658.1,820.6) -- (657.3,821.5) -- (656.4,822.4) -- (655.5,823.3) -- (654.8,824.1) -- (654.1,824.7) -- (653.5,825.2) -- (652.8,825.8) -- (652.1,826.4) -- (651.2,827.1) -- (650.3,827.8) -- (649.3,828.5) -- (648.2,829.3) -- (646.9,830.0) -- (645.6,830.8) -- (644.5,831.4) -- (643.8,831.9) -- (643.3,832.1) -- (643.2,832.2) -- (642.7,832.3) -- (641.9,832.5) -- (640.8,832.8) -- (639.2,833.2) -- (637.8,833.4) -- (636.2,833.6) -- (634.8,833.6) -- (633.2,833.5) -- (632.0,833.3) -- (630.9,833.1) -- (630.0,832.8) -- (629.3,832.5) -- (628.6,832.1) -- (627.9,831.7) -- (627.1,831.3) -- (626.2,830.8) -- (625.4,830.2) -- (624.5,829.4) -- (623.7,828.5) -- (622.8,827.5) -- (621.9,826.6) -- (621.2,825.7) -- (620.5,824.8) -- (619.9,824.0) -- (619.4,823.2) -- (618.9,822.5) -- (618.5,821.8) -- (618.1,821.1) -- (617.6,820.2) -- (617.1,819.1) -- (616.5,817.8) -- (615.8,816.2) -- (615.1,814.9) -- (614.6,813.7) -- (614.1,812.7) -- (613.6,811.8) -- (613.3,811.0) -- (612.9,810.2) -- (612.7,809.5) -- (612.4,808.9) -- (612.1,808.1) -- (611.7,807.0) -- (611.2,805.8) -- (610.6,804.2) -- (610.0,802.8) -- (609.4,801.2) -- (608.8,799.8) -- (608.1,798.2) -- (607.6,797.1) -- (607.3,796.3) -- (607.1,795.9) -- (607.1,795.8) -- (606.9,795.4) -- (606.6,794.6) -- (606.1,793.5) -- (605.6,792.1) -- (605.0,790.7) -- (604.3,789.2) -- (603.7,787.8) -- (603.0,786.2) -- (602.4,785.0) -- (601.9,784.1) -- (601.6,783.4) -- (601.3,782.9) -- (601.0,782.3) -- (600.7,781.6) -- (600.3,780.7) -- (599.8,779.6) -- (599.2,778.4) -- (598.5,777.1) -- (597.8,775.8) -- (597.0,774.2) -- (596.4,773.0) -- (595.9,772.1) -- (595.5,771.5) -- (595.3,771.2) -- (595.0,770.6) -- (594.6,769.9) -- (594.1,768.9) -- (593.4,767.7) -- (592.7,766.5) -- (592.0,765.2) -- (591.2,763.9) -- (590.3,762.4) -- (589.6,761.4) -- (589.2,760.6) -- (588.9,760.2) -- (588.9,760.1) -- (588.6,759.7) -- (588.1,758.9) -- (587.4,757.8) -- (586.5,756.2) -- (585.7,754.9) -- (585.0,753.8) -- (584.3,752.8) -- (583.8,752.0) -- (583.3,751.3) -- (582.8,750.5) -- (582.3,749.8) -- (581.9,749.1) -- (581.3,748.2) -- (580.5,747.1) -- (579.6,745.8) -- (578.6,744.2) -- (577.8,743.1) -- (577.3,742.4) -- (577.0,742.0) -- (577.0,742.0) -- (576.8,741.6) -- (576.2,740.8) -- (575.5,739.7) -- (574.5,738.2) -- (573.6,736.9) -- (572.8,735.9) -- (572.2,735.0) -- (571.7,734.3) -- (571.2,733.6) -- (570.7,732.8) -- (570.2,732.0) -- (569.6,731.2) -- (569.0,730.3) -- (568.2,729.2) -- (567.4,728.0) -- (566.4,726.7) -- (565.7,725.6) -- (565.2,724.9) -- (564.8,724.4) -- (564.7,724.2) -- (564.4,723.7) -- (563.8,722.9) -- (563.0,721.8) -- (562.0,720.2) -- (561.2,719.0) -- (560.5,718.0) -- (560.0,717.2) -- (559.6,716.6) -- (559.1,715.9) -- (558.7,715.2) -- (558.1,714.3) -- (557.5,713.4) -- (556.8,712.3) -- (556.1,711.1) -- (555.3,709.8) -- (554.4,708.4) -- (553.7,707.3) -- (553.2,706.5) -- (553.0,706.1) -- (552.9,706.1) -- (552.7,705.7) -- (552.3,704.9) -- (551.6,703.8) -- (550.8,702.2) -- (550.0,700.8) -- (549.3,699.4) -- (548.6,698.1) -- (548.0,696.9) -- (547.5,695.9) -- (547.1,695.1) -- (546.8,694.6) -- (546.7,694.4) -- (546.4,693.8) -- (546.0,692.9) -- (545.5,691.8) -- (544.9,690.2) -- (544.3,688.8) -- (543.6,687.2) -- (543.1,685.8) -- (542.5,684.2) -- (542.0,683.0) -- (541.7,682.0) -- (541.4,681.2) -- (541.2,680.6) -- (541.1,680.0) -- (540.8,679.2) -- (540.6,678.4) -- (540.3,677.4) -- (540.0,676.3) -- (539.7,675.1) -- (539.4,673.8) -- (539.1,672.2) -- (538.8,670.8) -- (538.6,669.2) -- (538.3,667.8) -- (538.1,666.2) -- (537.9,664.8) -- (537.8,663.2) -- (537.7,661.8) -- (537.6,660.2) -- (537.5,658.8) -- (537.4,657.2) -- (537.4,655.8) -- (537.4,654.2) -- (537.5,652.8) -- (537.6,651.2) -- (537.6,649.8) -- (537.8,648.2) -- (537.9,646.8) -- (538.1,645.2) -- (538.2,643.8) -- (538.4,642.2) -- (538.7,640.8) -- (538.9,639.2) -- (539.2,637.8) -- (539.5,636.2) -- (539.8,634.8) -- (540.0,633.5) -- (540.3,632.2) -- (540.6,631.1) -- (540.8,630.1) -- (541.0,629.3) -- (541.1,628.8) -- (541.2,628.5) -- (541.3,627.9) -- (541.6,627.0) -- (541.9,625.8) -- (542.3,624.2) -- (542.7,622.8) -- (543.2,621.2) -- (543.6,619.8) -- (544.1,618.2) -- (544.5,616.8) -- (545.0,615.2) -- (545.4,613.8) -- (545.9,612.2) -- (546.3,611.0) -- (546.5,610.1) -- (546.7,609.5) -- (546.8,609.2) -- (547.0,608.6) -- (547.3,607.9) -- (547.6,606.9) -- (547.9,605.8) -- (548.3,604.5) -- (548.8,603.2) -- (549.2,601.8) -- (549.7,600.2) -- (550.2,598.8) -- (550.6,597.2) -- (551.1,595.8) -- (551.6,594.2) -- (552.0,593.0) -- (552.3,591.9) -- (552.6,591.1) -- (552.7,590.5) -- (552.9,589.9) -- (553.2,589.1) -- (553.4,588.3) -- (553.7,587.4) -- (554.0,586.3) -- (554.4,585.1) -- (554.8,583.8) -- (555.2,582.2) -- (555.6,580.8) -- (556.1,579.2) -- (556.5,577.8) -- (556.9,576.2) -- (557.3,574.8) -- (557.7,573.3) -- (558.1,571.8) -- (558.4,570.4) -- (558.7,569.3) -- (558.9,568.5) -- (559.0,568.1) -- (559.0,568.1) -- (559.1,567.7) -- (559.3,566.9) -- (559.6,565.8) -- (560.0,564.2) -- (560.3,562.8) -- (560.7,561.2) -- (561.1,559.8) -- (561.4,558.2) -- (561.7,556.8) -- (562.1,555.2) -- (562.4,553.8) -- (562.7,552.2) -- (563.0,550.8) -- (563.3,549.2) -- (563.6,547.8) -- (563.9,546.2) -- (564.2,544.9) -- (564.4,543.7) -- (564.6,542.6) -- (564.8,541.7) -- (564.9,540.9) -- (565.0,540.2) -- (565.2,539.5) -- (565.3,538.9) -- (565.4,538.1) -- (565.6,537.0) -- (565.8,535.8) -- (566.1,534.2) -- (566.3,532.8) -- (566.6,531.2) -- (566.8,529.8) -- (567.1,528.2) -- (567.3,527.1) -- (567.4,526.4) -- (567.4,526.0);
\draw[cloudline] (554.8,1036.0) -- (555.1,1035.9) -- (555.6,1035.7) -- (556.4,1035.4) -- (557.4,1035.0) -- (558.6,1034.6) -- (559.9,1034.2) -- (561.2,1033.8) -- (562.8,1033.5) -- (564.2,1033.2) -- (565.8,1032.9) -- (567.2,1032.7) -- (568.8,1032.6) -- (570.2,1032.5) -- (571.8,1032.5) -- (573.2,1032.5) -- (574.8,1032.6) -- (576.2,1032.8) -- (577.8,1033.0) -- (579.2,1033.3) -- (580.8,1033.6) -- (582.2,1034.0) -- (583.6,1034.4) -- (585.0,1034.8) -- (586.4,1035.3) -- (587.4,1035.7) -- (588.2,1035.9) -- (588.6,1036.1) -- (588.8,1036.1) -- (589.3,1036.4) -- (590.1,1036.8) -- (591.2,1037.4) -- (592.8,1038.1) -- (594.2,1038.8) -- (595.5,1039.6) -- (596.7,1040.3) -- (597.8,1041.0) -- (598.7,1041.5) -- (599.5,1042.0) -- (600.0,1042.4) -- (600.4,1042.6) -- (601.1,1043.1) -- (602.0,1043.8) -- (603.2,1044.7) -- (604.8,1045.9) -- (605.9,1046.7) -- (606.7,1047.4) -- (607.2,1047.7) -- (607.3,1047.8) -- (607.7,1048.2) -- (608.5,1048.9) -- (609.5,1049.9) -- (610.9,1051.2) -- (612.0,1052.2) -- (612.8,1052.9) -- (613.3,1053.4) -- (613.6,1053.6) -- (614.1,1054.2) -- (614.8,1055.0) -- (615.8,1056.1) -- (617.1,1057.4) -- (618.1,1058.5) -- (618.8,1059.3) -- (619.2,1059.7) -- (619.3,1059.8) -- (619.6,1060.3) -- (620.3,1061.1) -- (621.2,1062.2) -- (622.4,1063.8) -- (623.3,1064.9) -- (624.0,1065.8) -- (624.5,1066.4) -- (624.7,1066.7) -- (625.0,1067.3) -- (625.5,1068.0) -- (626.2,1069.0) -- (627.0,1070.2) -- (627.8,1071.4) -- (628.5,1072.5) -- (629.3,1073.7) -- (630.0,1074.8) -- (630.5,1075.7) -- (631.0,1076.5) -- (631.3,1077.0) -- (631.5,1077.4) -- (631.9,1078.1) -- (632.4,1079.0) -- (633.1,1080.2) -- (633.9,1081.8) -- (634.6,1083.1) -- (635.2,1084.4) -- (635.8,1085.5) -- (636.3,1086.5) -- (636.7,1087.4) -- (637.0,1088.1) -- (637.3,1088.7) -- (637.6,1089.2) -- (637.9,1090.0) -- (638.3,1091.0) -- (638.9,1092.2) -- (639.5,1093.8) -- (640.1,1095.2) -- (640.7,1096.8) -- (641.3,1098.2) -- (641.9,1099.8) -- (642.3,1100.9) -- (642.7,1101.8) -- (642.8,1102.3) -- (642.9,1102.4) -- (643.1,1102.9) -- (643.3,1103.6) -- (643.7,1104.7) -- (644.1,1106.0) -- (644.6,1107.4) -- (645.0,1108.8) -- (645.5,1110.2) -- (645.9,1111.8) -- (646.4,1113.2) -- (646.8,1114.8) -- (647.3,1116.2) -- (647.7,1117.8) -- (648.1,1119.0) -- (648.4,1120.1) -- (648.6,1120.9) -- (648.8,1121.5) -- (648.9,1122.2) -- (649.1,1122.9) -- (649.3,1123.7) -- (649.6,1124.6) -- (649.7,1125.3) -- (649.7,1125.8) -- (649.6,1126.0) -- (649.3,1126.0) -- (648.8,1126.0) -- (647.9,1126.1) -- (646.8,1126.1) -- (645.2,1126.2) -- (643.8,1126.6) -- (642.3,1127.3) -- (640.9,1128.4) -- (639.6,1129.8) -- (638.5,1131.0) -- (637.7,1131.7) -- (637.3,1132.2) -- (637.2,1132.3) -- (636.7,1132.4) -- (635.9,1132.5) -- (634.8,1132.5) -- (633.2,1132.6) -- (631.8,1132.6) -- (630.2,1132.7) -- (628.8,1132.7) -- (627.2,1132.7) -- (625.8,1132.8) -- (624.2,1132.8) -- (622.8,1132.8) -- (621.2,1132.8) -- (619.8,1132.9) -- (618.2,1132.9) -- (616.8,1132.9) -- (615.2,1132.9) -- (613.8,1133.0) -- (612.2,1133.0) -- (610.8,1133.0) -- (609.2,1133.0) -- (607.8,1133.0) -- (606.2,1133.0) -- (604.8,1133.1) -- (603.2,1133.1) -- (601.8,1133.1) -- (600.2,1133.1) -- (598.8,1133.1) -- (597.2,1133.1) -- (595.8,1133.1) -- (594.2,1133.1) -- (592.8,1133.2) -- (591.2,1133.2) -- (590.1,1133.1) -- (589.2,1133.0) -- (588.6,1132.7) -- (588.3,1132.4) -- (587.7,1131.9) -- (587.0,1131.2) -- (586.0,1130.2) -- (584.8,1129.0) -- (583.5,1128.1) -- (582.2,1127.5) -- (580.8,1127.2) -- (579.2,1127.2) -- (577.8,1127.2) -- (576.2,1127.2) -- (574.8,1127.2) -- (573.2,1127.2) -- (571.8,1127.5) -- (570.5,1128.1) -- (569.2,1129.0) -- (568.0,1130.2) -- (567.0,1131.2) -- (566.3,1131.9) -- (565.8,1132.5) -- (565.5,1132.8) -- (564.9,1133.3) -- (564.2,1134.0) -- (563.2,1135.0) -- (562.0,1136.2) -- (561.0,1137.2) -- (560.3,1137.9) -- (559.8,1138.5) -- (559.5,1138.8) -- (558.9,1139.3) -- (558.2,1140.0) -- (557.2,1141.0) -- (556.0,1142.2) -- (555.1,1143.2) -- (554.3,1143.9) -- (553.8,1144.5) -- (553.5,1144.8) -- (552.9,1145.0) -- (552.0,1145.2) -- (550.8,1145.2) -- (549.2,1145.2) -- (547.8,1145.2) -- (546.2,1145.2) -- (544.8,1145.2) -- (543.2,1145.2) -- (541.8,1145.5) -- (540.5,1146.1) -- (539.2,1147.0) -- (538.0,1148.2) -- (537.0,1149.2) -- (536.3,1149.9) -- (535.7,1150.4) -- (535.4,1150.7) -- (534.8,1151.0) -- (533.9,1151.1) -- (532.8,1151.2) -- (531.2,1151.2) -- (529.8,1151.2) -- (528.2,1151.1) -- (526.8,1151.1) -- (525.2,1151.1) -- (523.8,1151.1) -- (522.2,1151.1) -- (520.8,1151.1) -- (519.2,1151.1) -- (517.8,1151.1) -- (516.2,1151.1) -- (514.8,1151.0) -- (513.2,1151.0) -- (511.8,1151.0) -- (510.2,1151.0) -- (508.8,1151.0) -- (507.2,1151.0) -- (505.8,1151.0) -- (504.2,1151.0) -- (502.8,1150.9) -- (501.2,1150.9) -- (499.8,1150.9) -- (498.2,1150.9) -- (496.8,1150.9) -- (495.2,1150.9) -- (493.8,1150.8) -- (492.2,1150.8) -- (490.8,1150.8) -- (489.2,1150.8) -- (487.8,1150.8) -- (486.2,1150.8) -- (484.8,1150.7) -- (483.2,1150.7) -- (481.8,1150.7) -- (480.2,1150.7) -- (478.8,1150.7) -- (477.2,1150.6) -- (476.1,1150.6) -- (475.3,1150.5) -- (474.8,1150.4) -- (474.6,1150.2) -- (474.2,1149.8) -- (473.4,1149.0) -- (472.4,1147.9) -- (471.0,1146.5) -- (469.6,1145.5) -- (468.2,1144.7) -- (466.8,1144.4) -- (465.2,1144.3) -- (464.1,1144.3) -- (463.3,1144.3) -- (462.9,1144.2) -- (462.8,1144.1) -- (462.4,1143.7) -- (461.6,1142.9) -- (460.6,1141.8) -- (459.1,1140.3) -- (458.1,1139.2) -- (457.3,1138.4) -- (457.0,1138.0) -- (457.0,1138.0) -- (457.0,1137.9) -- (457.0,1137.6) -- (457.0,1137.1) -- (457.0,1136.5) -- (457.2,1135.8) -- (457.6,1135.1) -- (458.2,1134.2) -- (459.0,1133.3) -- (459.8,1132.5) -- (460.6,1131.7) -- (461.3,1130.9) -- (462.0,1130.1) -- (462.7,1129.3) -- (463.3,1128.6) -- (464.0,1127.8) -- (464.7,1127.1) -- (465.4,1126.3) -- (466.1,1125.5) -- (466.9,1124.6) -- (467.7,1123.7) -- (468.5,1122.8) -- (469.2,1122.1) -- (469.7,1121.4) -- (470.2,1120.8) -- (470.9,1120.1) -- (471.6,1119.3) -- (472.5,1118.3) -- (473.5,1117.2) -- (474.3,1116.2) -- (475.0,1115.5) -- (475.5,1114.9) -- (475.8,1114.5) -- (476.3,1113.9) -- (477.1,1113.1) -- (478.1,1112.0) -- (479.2,1110.6) -- (480.2,1109.6) -- (480.8,1108.8) -- (481.2,1108.4) -- (481.3,1108.2) -- (481.8,1107.7) -- (482.5,1106.9) -- (483.5,1105.8) -- (484.8,1104.2) -- (485.8,1103.1) -- (486.5,1102.3) -- (486.8,1101.9) -- (486.9,1101.8) -- (487.3,1101.4) -- (487.9,1100.7) -- (488.9,1099.6) -- (490.1,1098.1) -- (491.1,1097.0) -- (491.8,1096.1) -- (492.3,1095.5) -- (492.6,1095.2) -- (493.1,1094.7) -- (493.7,1094.0) -- (494.5,1093.0) -- (495.6,1091.8) -- (496.5,1090.8) -- (497.2,1089.9) -- (497.8,1089.2) -- (498.3,1088.6) -- (498.9,1088.0) -- (499.5,1087.2) -- (500.2,1086.4) -- (501.1,1085.4) -- (501.8,1084.5) -- (502.6,1083.6) -- (503.3,1082.8) -- (504.0,1082.0) -- (504.6,1081.3) -- (505.3,1080.5) -- (505.9,1079.8) -- (506.5,1079.1) -- (507.2,1078.3) -- (508.0,1077.4) -- (508.8,1076.5) -- (509.7,1075.5) -- (510.4,1074.6) -- (511.1,1073.8) -- (511.6,1073.2) -- (512.1,1072.7) -- (512.7,1072.1) -- (513.4,1071.2) -- (514.3,1070.2) -- (515.4,1069.0) -- (516.3,1068.0) -- (516.9,1067.3) -- (517.4,1066.7) -- (517.7,1066.4) -- (518.2,1065.9) -- (518.9,1065.1) -- (519.9,1064.0) -- (521.2,1062.7) -- (522.1,1061.7) -- (522.8,1060.9) -- (523.3,1060.4) -- (523.4,1060.3) -- (523.9,1059.8) -- (524.6,1059.0) -- (525.7,1058.0) -- (527.0,1056.6) -- (528.0,1055.6) -- (528.8,1054.8) -- (529.2,1054.4) -- (529.4,1054.2) -- (529.9,1053.8) -- (530.6,1053.1) -- (531.7,1052.2) -- (533.0,1050.9) -- (534.1,1049.9) -- (534.9,1049.2) -- (535.5,1048.7) -- (535.8,1048.4) -- (536.3,1048.0) -- (537.1,1047.4) -- (538.1,1046.6) -- (539.2,1045.6) -- (540.3,1044.8) -- (541.3,1044.0) -- (542.2,1043.4) -- (543.0,1042.8) -- (543.7,1042.3) -- (544.5,1041.8) -- (545.2,1041.3) -- (545.9,1040.8) -- (546.8,1040.2) -- (547.9,1039.6) -- (549.2,1038.8) -- (550.8,1038.0) -- (552.0,1037.3) -- (553.0,1036.8) -- (553.7,1036.5) -- (554.1,1036.3) -- (554.5,1036.1) -- (554.7,1036.0) -- (554.8,1036.0);
\path[precipfill,even odd rule] (541.0,526.0) -- (541.0,526.0) -- (541.0,526.0) -- (541.0,526.0) -- (541.0,526.0) -- (541.4,525.9) -- (542.1,525.7) -- (543.2,525.5) -- (544.8,525.1) -- (546.0,525.0) -- (546.9,525.0) -- (547.5,525.1) -- (547.9,525.5) -- (548.4,526.0) -- (549.1,526.6) -- (550.1,527.4) -- (551.3,528.3) -- (552.3,528.7) -- (553.3,528.7) -- (554.2,528.3) -- (554.9,527.4) -- (555.7,526.5) -- (556.4,525.7) -- (557.2,524.9) -- (557.9,524.2) -- (558.8,523.6) -- (559.9,523.2) -- (561.2,522.9) -- (562.8,522.8) -- (564.2,522.7) -- (565.8,522.6) -- (567.2,522.5) -- (568.8,522.5) -- (570.2,522.4) -- (571.8,522.3) -- (573.2,522.3) -- (574.8,522.2) -- (576.2,522.2) -- (577.8,522.2) -- (579.2,522.1) -- (580.8,522.1) -- (582.2,522.1) -- (583.8,522.0) -- (585.2,522.0) -- (586.8,522.0) -- (588.2,522.0) -- (589.8,522.0) -- (591.2,521.9) -- (592.8,521.9) -- (594.2,521.9) -- (595.8,521.9) -- (597.2,521.9) -- (598.8,521.9) -- (600.2,521.9) -- (601.8,521.9) -- (603.2,521.9) -- (604.8,521.9) -- (606.2,521.8) -- (607.8,521.8) -- (609.2,521.8) -- (610.8,521.8) -- (612.2,521.8) -- (613.8,521.8) -- (615.2,521.8) -- (616.8,521.8) -- (618.2,521.8) -- (619.8,521.8) -- (621.2,521.8) -- (622.8,521.8) -- (624.2,521.8) -- (625.8,521.8) -- (627.2,521.8) -- (628.8,521.8) -- (630.2,521.8) -- (631.8,521.8) -- (633.2,521.8) -- (634.8,521.8) -- (636.2,521.8) -- (637.8,521.8) -- (639.2,521.8) -- (640.8,521.8) -- (642.2,521.8) -- (643.8,521.8) -- (645.2,521.8) -- (646.8,521.8) -- (648.2,521.8) -- (649.8,521.8) -- (651.2,521.8) -- (652.8,521.8) -- (654.2,521.8) -- (655.8,521.8) -- (657.2,521.8) -- (658.8,521.8) -- (660.2,521.8) -- (661.8,521.8) -- (663.2,521.8) -- (664.8,521.8) -- (666.2,521.8) -- (667.8,521.8) -- (669.2,521.8) -- (670.8,521.8) -- (672.2,521.8) -- (673.8,521.8) -- (675.2,521.8) -- (676.8,521.8) -- (678.2,521.8) -- (679.8,521.8) -- (681.2,521.8) -- (682.8,521.8) -- (684.2,521.8) -- (685.8,521.9) -- (687.2,521.9) -- (688.8,521.9) -- (690.1,522.1) -- (691.4,522.7) -- (692.5,523.4) -- (693.6,524.5) -- (694.5,525.3) -- (695.2,526.1) -- (695.8,526.7) -- (696.3,527.2) -- (696.9,527.8) -- (697.6,528.6) -- (698.5,529.4) -- (699.6,530.5) -- (700.4,531.4) -- (701.2,532.1) -- (701.8,532.7) -- (702.3,533.2) -- (703.0,533.6) -- (704.0,533.9) -- (705.2,534.0) -- (706.8,534.0) -- (708.0,533.9) -- (709.0,533.7) -- (709.8,533.3) -- (710.4,532.8) -- (711.0,532.1) -- (711.8,531.4) -- (712.6,530.6) -- (713.6,529.6) -- (714.7,528.9) -- (715.9,528.5) -- (717.2,528.2) -- (718.8,528.3) -- (720.2,528.3) -- (721.8,528.4) -- (723.2,528.4) -- (724.8,528.5) -- (726.2,528.5) -- (727.8,528.6) -- (729.2,528.6) -- (730.8,528.7) -- (732.2,528.8) -- (733.8,528.9) -- (735.2,528.9) -- (736.8,529.0) -- (738.2,529.1) -- (739.8,529.2) -- (741.2,529.3) -- (742.8,529.4) -- (744.1,529.3) -- (745.4,528.9) -- (746.6,528.2) -- (747.7,527.3) -- (748.7,526.6) -- (749.4,525.9) -- (750.0,525.4) -- (750.4,525.0) -- (751.1,524.7) -- (752.0,524.6) -- (753.2,524.7) -- (754.8,524.9) -- (756.2,525.1) -- (757.6,525.3) -- (759.0,525.5) -- (760.3,525.7) -- (761.3,526.0) -- (762.0,526.3) -- (762.5,526.8) -- (762.7,527.3) -- (763.0,527.9) -- (763.3,528.7) -- (763.6,529.5) -- (764.0,530.5) -- (764.5,531.6) -- (765.0,532.9) -- (765.6,534.2) -- (766.2,535.8) -- (766.8,537.2) -- (767.3,538.5) -- (767.8,539.8) -- (768.3,541.0) -- (768.7,542.0) -- (769.0,542.8) -- (769.2,543.3) -- (769.3,543.6) -- (769.6,544.2) -- (769.9,545.1) -- (770.4,546.2) -- (771.0,547.8) -- (771.7,549.2) -- (772.3,550.8) -- (772.9,552.2) -- (773.5,553.8) -- (773.9,555.0) -- (774.3,555.9) -- (774.6,556.6) -- (774.8,557.0) -- (775.0,557.6) -- (775.3,558.4) -- (775.7,559.3) -- (776.1,560.4) -- (776.6,561.6) -- (777.2,562.9) -- (777.7,564.2) -- (778.3,565.8) -- (778.9,567.1) -- (779.4,568.4) -- (779.9,569.6) -- (780.3,570.7) -- (780.7,571.7) -- (781.0,572.4) -- (781.3,573.0) -- (781.4,573.4) -- (781.7,574.1) -- (782.1,575.0) -- (782.7,576.2) -- (783.3,577.8) -- (784.0,579.2) -- (784.6,580.8) -- (785.2,582.2) -- (785.9,583.8) -- (786.4,584.9) -- (786.7,585.7) -- (786.9,586.2) -- (786.9,586.3) -- (787.1,586.7) -- (787.5,587.5) -- (788.0,588.5) -- (788.7,589.9) -- (789.4,591.3) -- (790.1,592.8) -- (790.8,594.2) -- (791.5,595.8) -- (792.1,596.9) -- (792.5,597.8) -- (792.7,598.3) -- (792.8,598.5) -- (793.1,599.0) -- (793.5,599.7) -- (794.1,600.7) -- (794.9,602.0) -- (795.7,603.3) -- (796.4,604.6) -- (797.2,605.8) -- (797.9,606.9) -- (798.5,607.9) -- (799.0,608.7) -- (799.4,609.2) -- (799.6,609.5) -- (800.1,610.1) -- (800.8,611.0) -- (801.8,612.2) -- (803.0,613.8) -- (804.0,614.9) -- (804.6,615.7) -- (804.9,616.1) -- (805.0,616.1) -- (805.4,616.5) -- (806.1,617.1) -- (807.2,618.1) -- (808.8,619.5) -- (809.9,620.5) -- (810.8,621.2) -- (811.3,621.7) -- (811.5,621.8) -- (812.0,622.1) -- (812.7,622.6) -- (813.8,623.2) -- (815.1,624.0) -- (816.4,624.7) -- (817.8,625.4) -- (819.2,626.1) -- (820.8,626.7) -- (821.9,627.2) -- (822.8,627.6) -- (823.4,627.8) -- (823.7,627.9) -- (824.3,628.0) -- (825.0,628.2) -- (826.0,628.5) -- (827.2,628.8) -- (828.5,629.1) -- (829.8,629.3) -- (831.2,629.5) -- (832.8,629.8) -- (834.2,629.9) -- (835.8,630.1) -- (837.2,630.2) -- (838.8,630.2) -- (840.2,630.3) -- (841.8,630.3) -- (843.2,630.2) -- (844.8,630.2) -- (846.2,630.1) -- (847.8,630.0) -- (849.2,629.8) -- (850.8,629.7) -- (852.2,629.5) -- (853.8,629.3) -- (855.2,629.0) -- (856.8,628.8) -- (858.0,628.5) -- (858.9,628.4) -- (859.6,628.2) -- (860.0,628.1) -- (860.5,628.0) -- (861.3,627.8) -- (862.2,627.6) -- (863.3,627.3) -- (864.5,627.0) -- (865.8,626.7) -- (867.2,626.3) -- (868.8,625.8) -- (870.2,625.4) -- (871.8,624.9) -- (873.2,624.4) -- (874.8,623.9) -- (876.1,623.4) -- (877.2,623.0) -- (878.1,622.7) -- (878.8,622.4) -- (879.5,622.1) -- (880.3,621.8) -- (881.1,621.5) -- (881.8,621.2) -- (882.8,620.8) -- (883.9,620.4) -- (885.2,619.8) -- (886.8,619.2) -- (888.2,618.5) -- (889.5,617.9) -- (890.7,617.4) -- (891.9,616.8) -- (892.9,616.4) -- (893.6,616.0) -- (894.2,615.8) -- (894.5,615.6) -- (895.1,615.3) -- (896.0,614.8) -- (897.2,614.2) -- (898.8,613.4) -- (900.1,612.7) -- (901.4,612.1) -- (902.6,611.4) -- (903.7,610.9) -- (904.6,610.4) -- (905.4,610.0) -- (905.9,609.7) -- (906.4,609.4) -- (907.1,609.0) -- (908.0,608.5) -- (909.2,607.8) -- (910.8,607.0) -- (911.9,606.6) -- (912.8,606.7) -- (913.3,607.3) -- (913.5,608.4) -- (913.6,609.6) -- (913.8,610.9) -- (914.0,612.2) -- (914.2,613.8) -- (914.7,615.1) -- (915.3,616.3) -- (916.2,617.3) -- (917.3,618.2) -- (918.5,618.9) -- (919.8,619.3) -- (921.2,619.4) -- (922.8,619.4) -- (924.2,619.3) -- (925.8,619.2) -- (927.2,619.2) -- (928.8,619.1) -- (930.2,619.0) -- (931.8,619.0) -- (933.2,618.9) -- (934.8,618.9) -- (936.1,619.0) -- (937.2,619.4) -- (938.2,620.0) -- (939.0,620.8) -- (939.8,621.5) -- (940.5,622.2) -- (941.3,622.9) -- (942.0,623.5) -- (942.9,623.9) -- (944.0,624.2) -- (945.2,624.4) -- (946.8,624.3) -- (948.1,624.5) -- (949.3,625.0) -- (950.4,625.7) -- (951.3,626.6) -- (952.2,627.4) -- (952.9,628.2) -- (953.6,628.8) -- (954.1,629.3) -- (954.9,629.7) -- (956.0,630.0) -- (957.2,630.1) -- (958.8,630.1) -- (960.1,630.3) -- (961.4,630.8) -- (962.5,631.5) -- (963.5,632.5) -- (964.3,633.4) -- (965.1,634.1) -- (965.7,634.7) -- (966.2,635.2) -- (966.9,635.8) -- (967.6,636.6) -- (968.5,637.5) -- (969.5,638.5) -- (970.4,639.4) -- (971.2,640.1) -- (971.8,640.7) -- (972.3,641.2) -- (972.9,641.8) -- (973.6,642.5) -- (974.5,643.4) -- (975.6,644.5) -- (976.5,645.3) -- (977.2,646.1) -- (977.8,646.7) -- (978.3,647.2) -- (978.9,647.8) -- (979.7,648.5) -- (980.6,649.4) -- (981.6,650.4) -- (982.5,651.3) -- (983.2,652.1) -- (983.8,652.7) -- (984.3,653.1) -- (984.9,653.8) -- (985.7,654.5) -- (986.6,655.4) -- (987.6,656.4) -- (988.4,657.6) -- (988.9,658.9) -- (989.2,660.2) -- (989.2,661.8) -- (989.3,663.0) -- (989.5,664.0) -- (989.9,664.7) -- (990.3,665.1) -- (990.9,665.7) -- (991.7,666.5) -- (992.6,667.4) -- (993.6,668.4) -- (994.4,669.6) -- (994.9,670.9) -- (995.2,672.2) -- (995.2,673.8) -- (995.2,675.2) -- (995.2,676.8) -- (995.2,678.2) -- (995.2,679.8) -- (995.2,681.2) -- (995.2,682.8) -- (995.2,684.2) -- (995.2,685.8) -- (995.2,687.2) -- (995.2,688.8) -- (995.2,690.2) -- (995.2,691.8) -- (995.2,693.2) -- (995.2,694.8) -- (995.2,696.2) -- (995.2,697.8) -- (995.2,699.2) -- (995.2,700.8) -- (995.2,702.2) -- (995.2,703.8) -- (995.2,705.2) -- (995.2,706.8) -- (995.2,708.2) -- (995.2,709.8) -- (995.2,711.2) -- (995.2,712.8) -- (995.2,714.2) -- (995.2,715.8) -- (995.2,717.2) -- (995.2,718.8) -- (995.2,720.2) -- (995.2,721.8) -- (995.2,723.2) -- (995.2,724.8) -- (995.2,726.2) -- (995.2,727.8) -- (995.2,729.2) -- (995.2,730.8) -- (995.2,732.2) -- (995.2,733.8) -- (995.2,735.2) -- (995.2,736.8) -- (995.2,738.2) -- (995.2,739.8) -- (995.2,741.2) -- (995.2,742.8) -- (995.2,744.2) -- (995.2,745.8) -- (995.2,747.2) -- (995.2,748.8) -- (995.2,750.2) -- (995.2,751.8) -- (994.9,753.1) -- (994.4,754.4) -- (993.6,755.6) -- (992.6,756.6) -- (991.7,757.5) -- (990.9,758.3) -- (990.3,758.9) -- (989.9,759.3) -- (989.5,760.0) -- (989.3,761.0) -- (989.2,762.2) -- (989.2,763.8) -- (988.9,765.1) -- (988.4,766.4) -- (987.6,767.6) -- (986.6,768.6) -- (985.7,769.5) -- (984.9,770.2) -- (984.3,770.8) -- (983.8,771.3) -- (983.2,771.9) -- (982.5,772.7) -- (981.6,773.5) -- (980.5,774.6) -- (979.7,775.5) -- (978.9,776.2) -- (978.3,776.8) -- (977.8,777.3) -- (977.2,777.9) -- (976.4,778.7) -- (975.6,779.5) -- (974.5,780.6) -- (973.6,781.4) -- (972.9,782.2) -- (972.3,782.8) -- (971.8,783.3) -- (971.1,783.9) -- (970.4,784.6) -- (969.5,785.5) -- (968.5,786.5) -- (967.6,787.4) -- (966.9,788.1) -- (966.2,788.8) -- (965.7,789.3) -- (965.1,789.9) -- (964.3,790.6) -- (963.4,791.5) -- (962.5,792.4) -- (961.4,793.2) -- (960.1,793.6) -- (958.8,793.9) -- (957.2,793.9) -- (956.0,794.0) -- (954.9,794.2) -- (954.1,794.6) -- (953.5,795.2) -- (952.8,795.8) -- (952.1,796.5) -- (951.2,797.3) -- (950.3,798.2) -- (949.3,798.9) -- (948.1,799.3) -- (946.8,799.5) -- (945.2,799.5) -- (943.9,799.6) -- (942.8,799.9) -- (941.8,800.4) -- (941.1,801.0) -- (940.3,801.7) -- (939.5,802.4) -- (938.8,803.1) -- (938.1,803.8) -- (937.2,804.3) -- (936.1,804.7) -- (934.8,804.8) -- (933.2,804.7) -- (931.8,804.6) -- (930.2,804.5) -- (928.8,804.4) -- (927.2,804.3) -- (925.8,804.4) -- (924.4,804.9) -- (923.0,805.6) -- (921.7,806.5) -- (920.6,807.5) -- (919.9,808.4) -- (919.4,809.3) -- (919.3,810.2) -- (918.8,810.7) -- (918.2,810.7) -- (917.3,810.2) -- (916.1,809.3) -- (915.2,808.6) -- (914.4,808.0) -- (913.9,807.6) -- (913.5,807.3) -- (912.9,806.8) -- (912.0,806.1) -- (910.8,805.1) -- (909.2,803.9) -- (908.1,803.0) -- (907.4,802.3) -- (907.0,802.0) -- (907.0,802.0) -- (906.6,801.7) -- (905.8,801.1) -- (904.7,800.2) -- (903.2,798.9) -- (902.0,797.9) -- (901.1,797.2) -- (900.5,796.7) -- (900.2,796.4) -- (899.7,795.9) -- (898.9,795.3) -- (898.0,794.5) -- (896.8,793.4) -- (895.7,792.6) -- (894.9,791.8) -- (894.2,791.2) -- (893.6,790.7) -- (893.0,790.1) -- (892.2,789.5) -- (891.4,788.7) -- (890.4,787.9) -- (889.5,787.1) -- (888.7,786.3) -- (887.9,785.6) -- (887.2,785.0) -- (886.5,784.3) -- (885.7,783.6) -- (884.9,782.9) -- (884.2,782.2) -- (883.4,781.5) -- (882.6,780.7) -- (881.7,779.9) -- (880.9,779.2) -- (880.1,778.4) -- (879.3,777.7) -- (878.6,777.1) -- (878.0,776.5) -- (877.2,775.8) -- (876.4,775.0) -- (875.5,774.2) -- (874.6,773.3) -- (873.7,772.5) -- (872.9,771.9) -- (872.3,771.3) -- (871.8,770.8) -- (871.1,770.1) -- (870.3,769.4) -- (869.3,768.5) -- (868.2,767.5) -- (867.3,766.7) -- (866.5,766.0) -- (866.0,765.5) -- (865.6,765.1) -- (865.0,764.5) -- (864.1,763.8) -- (863.0,762.9) -- (861.7,761.7) -- (860.7,760.8) -- (859.9,760.2) -- (859.5,759.7) -- (859.3,759.6) -- (858.8,759.1) -- (857.9,758.4) -- (856.8,757.5) -- (855.2,756.2) -- (854.1,755.3) -- (853.3,754.6) -- (852.8,754.2) -- (852.7,754.1) -- (852.2,753.8) -- (851.5,753.2) -- (850.4,752.4) -- (849.1,751.4) -- (847.9,750.5) -- (846.8,749.8) -- (846.0,749.1) -- (845.3,748.7) -- (844.6,748.2) -- (843.8,747.7) -- (843.1,747.1) -- (842.2,746.6) -- (841.2,745.9) -- (840.1,745.2) -- (838.8,744.4) -- (837.2,743.4) -- (836.1,742.8) -- (835.4,742.3) -- (835.0,742.0) -- (834.9,742.0) -- (834.5,741.8) -- (833.8,741.4) -- (832.7,740.8) -- (831.2,740.0) -- (829.7,739.3) -- (828.2,738.5) -- (826.8,737.8) -- (825.2,737.1) -- (824.1,736.6) -- (823.3,736.2) -- (822.9,736.0) -- (822.9,736.0) -- (822.5,735.9) -- (821.8,735.6) -- (820.6,735.1) -- (819.2,734.5) -- (817.7,734.0) -- (816.2,733.4) -- (814.8,732.9) -- (813.2,732.3) -- (811.8,731.9) -- (810.4,731.4) -- (809.0,731.0) -- (807.7,730.6) -- (806.6,730.3) -- (805.9,730.1) -- (805.4,729.9) -- (805.3,729.9) -- (804.8,729.7) -- (803.9,729.5) -- (802.8,729.3) -- (801.2,728.9) -- (799.8,728.5) -- (798.2,728.2) -- (796.8,727.9) -- (795.2,727.6) -- (793.8,727.4) -- (792.2,727.2) -- (790.8,726.9) -- (789.2,726.8) -- (787.8,726.6) -- (786.2,726.5) -- (784.8,726.4) -- (783.2,726.3) -- (781.8,726.3) -- (780.2,726.3) -- (778.8,726.4) -- (777.2,726.5) -- (775.8,726.7) -- (774.2,726.9) -- (772.8,727.2) -- (771.2,727.6) -- (769.9,727.9) -- (768.6,728.4) -- (767.5,728.8) -- (766.5,729.3) -- (765.7,729.7) -- (764.9,730.1) -- (764.3,730.5) -- (763.8,730.9) -- (763.1,731.4) -- (762.4,732.2) -- (761.5,733.1) -- (760.4,734.3) -- (759.6,735.4) -- (758.8,736.4) -- (758.2,737.4) -- (757.7,738.4) -- (757.3,739.2) -- (757.0,740.0) -- (756.7,740.6) -- (756.5,741.2) -- (756.3,742.0) -- (756.0,743.0) -- (755.7,744.2) -- (755.4,745.8) -- (755.1,747.2) -- (754.9,748.8) -- (754.7,750.2) -- (754.6,751.8) -- (754.5,753.2) -- (754.4,754.8) -- (754.3,756.2) -- (754.2,757.8) -- (754.2,759.2) -- (754.1,760.8) -- (754.1,762.2) -- (754.1,763.8) -- (754.1,765.2) -- (754.1,766.8) -- (754.0,768.2) -- (754.0,769.8) -- (754.0,771.2) -- (754.0,772.8) -- (754.0,774.2) -- (753.9,775.8) -- (753.9,777.2) -- (753.9,778.8) -- (753.8,780.2) -- (753.8,781.8) -- (753.7,783.2) -- (753.6,784.8) -- (753.5,786.2) -- (753.5,787.8) -- (753.4,789.2) -- (753.2,790.8) -- (753.1,792.2) -- (753.0,793.8) -- (752.8,795.2) -- (752.7,796.8) -- (752.5,798.2) -- (752.3,799.8) -- (752.1,801.2) -- (751.9,802.7) -- (751.6,804.2) -- (751.4,805.6) -- (751.2,806.8) -- (751.1,807.5) -- (751.0,807.9) -- (751.0,807.9) -- (750.9,808.3) -- (750.7,809.1) -- (750.5,810.2) -- (750.2,811.8) -- (749.9,813.2) -- (749.5,814.8) -- (749.1,816.2) -- (748.7,817.8) -- (748.3,819.2) -- (747.9,820.8) -- (747.4,822.2) -- (746.9,823.8) -- (746.4,825.0) -- (746.0,826.1) -- (745.7,827.0) -- (745.4,827.6) -- (745.1,828.3) -- (744.8,829.1) -- (744.5,829.9) -- (744.1,830.7) -- (743.6,831.7) -- (743.1,832.9) -- (742.4,834.2) -- (741.6,835.8) -- (740.9,837.0) -- (740.3,838.0) -- (739.9,838.8) -- (739.5,839.3) -- (739.1,839.9) -- (738.6,840.7) -- (738.0,841.5) -- (737.3,842.5) -- (736.6,843.5) -- (735.8,844.5) -- (735.0,845.4) -- (734.2,846.4) -- (733.5,847.2) -- (732.8,848.0) -- (732.3,848.6) -- (731.8,849.2) -- (731.1,849.8) -- (730.4,850.6) -- (729.5,851.4) -- (728.5,852.3) -- (727.6,853.2) -- (726.7,853.9) -- (726.0,854.6) -- (725.3,855.1) -- (724.6,855.7) -- (723.9,856.3) -- (723.1,857.0) -- (722.2,857.6) -- (721.3,858.3) -- (720.3,859.0) -- (719.2,859.8) -- (718.0,860.7) -- (717.0,861.4) -- (716.2,861.9) -- (715.7,862.3) -- (715.4,862.5) -- (714.8,862.9) -- (713.9,863.5) -- (712.8,864.3) -- (711.2,865.2) -- (710.0,866.1) -- (709.0,866.7) -- (708.2,867.2) -- (707.7,867.5) -- (707.1,867.9) -- (706.4,868.4) -- (705.5,869.0) -- (704.5,869.6) -- (703.4,870.3) -- (702.2,871.1) -- (701.0,871.9) -- (699.7,872.7) -- (698.6,873.4) -- (697.9,873.9) -- (697.4,874.2) -- (697.2,874.3) -- (696.7,874.6) -- (695.9,875.2) -- (694.8,876.0) -- (693.2,877.0) -- (692.0,877.8) -- (691.0,878.5) -- (690.2,879.1) -- (689.6,879.4) -- (688.9,879.9) -- (688.2,880.4) -- (687.3,881.0) -- (686.4,881.6) -- (685.4,882.4) -- (684.3,883.1) -- (683.1,883.9) -- (681.9,884.7) -- (680.9,885.4) -- (680.1,885.9) -- (679.6,886.3) -- (679.4,886.4) -- (678.8,886.8) -- (677.9,887.4) -- (676.8,888.2) -- (675.2,889.2) -- (674.0,890.0) -- (673.0,890.7) -- (672.2,891.2) -- (671.6,891.5) -- (671.0,891.9) -- (670.2,892.4) -- (669.4,892.9) -- (668.4,893.5) -- (667.3,894.1) -- (666.1,894.8) -- (664.8,895.6) -- (663.2,896.4) -- (662.1,897.0) -- (661.2,897.5) -- (660.7,897.8) -- (660.5,897.9) -- (660.0,898.1) -- (659.3,898.5) -- (658.3,899.0) -- (657.0,899.6) -- (655.6,900.2) -- (654.2,900.9) -- (652.8,901.5) -- (651.2,902.1) -- (649.9,902.6) -- (648.8,903.1) -- (647.9,903.4) -- (647.2,903.6) -- (646.5,903.9) -- (645.8,904.2) -- (645.0,904.4) -- (644.2,904.7) -- (643.2,905.1) -- (642.1,905.4) -- (640.8,905.8) -- (639.2,906.3) -- (637.8,906.7) -- (636.2,907.1) -- (634.8,907.5) -- (633.2,907.9) -- (631.8,908.3) -- (630.2,908.6) -- (628.8,908.9) -- (627.2,909.2) -- (626.0,909.4) -- (624.9,909.6) -- (624.1,909.8) -- (623.6,909.9) -- (622.9,910.0) -- (622.2,910.1) -- (621.3,910.2) -- (620.4,910.4) -- (619.3,910.6) -- (618.1,910.7) -- (616.8,910.9) -- (615.2,911.1) -- (613.8,911.3) -- (612.2,911.4) -- (610.8,911.6) -- (609.2,911.7) -- (607.8,911.8) -- (606.2,911.9) -- (604.8,911.9) -- (603.2,912.0) -- (601.8,912.0) -- (600.2,912.0) -- (598.8,912.0) -- (597.2,912.0) -- (595.8,911.9) -- (594.2,911.8) -- (592.8,911.7) -- (591.2,911.6) -- (589.8,911.4) -- (588.2,911.2) -- (586.8,911.0) -- (585.2,910.7) -- (584.0,910.5) -- (583.1,910.3) -- (582.5,910.2) -- (582.2,910.1) -- (581.7,910.0) -- (580.9,909.8) -- (580.0,909.5) -- (578.8,909.2) -- (577.5,908.7) -- (576.2,908.3) -- (574.8,907.7) -- (573.2,907.0) -- (571.9,906.4) -- (570.8,905.8) -- (569.8,905.3) -- (568.9,904.8) -- (568.1,904.3) -- (567.4,903.8) -- (566.7,903.3) -- (566.0,902.8) -- (565.2,902.2) -- (564.3,901.5) -- (563.3,900.6) -- (562.2,899.6) -- (561.3,898.7) -- (560.6,897.9) -- (560.0,897.3) -- (559.6,896.9) -- (559.1,896.3) -- (558.5,895.5) -- (557.8,894.6) -- (557.1,893.6) -- (556.3,892.5) -- (555.6,891.3) -- (554.8,890.0) -- (554.1,888.7) -- (553.5,887.7) -- (553.1,887.0) -- (552.9,886.5) -- (552.8,886.3) -- (552.5,885.8) -- (552.2,884.9) -- (551.7,883.8) -- (551.1,882.2) -- (550.6,880.8) -- (550.1,879.2) -- (549.7,877.8) -- (549.2,876.2) -- (548.9,874.8) -- (548.6,873.2) -- (548.3,871.8) -- (548.0,870.2) -- (547.8,868.8) -- (547.6,867.2) -- (547.5,865.8) -- (547.4,864.2) -- (547.3,862.8) -- (547.2,861.2) -- (547.2,859.8) -- (547.2,858.2) -- (547.2,856.8) -- (547.3,855.2) -- (547.4,853.8) -- (547.5,852.2) -- (547.6,850.8) -- (547.8,849.2) -- (548.0,847.8) -- (548.2,846.2) -- (548.4,844.8) -- (548.6,843.2) -- (548.9,841.8) -- (549.2,840.2) -- (549.5,838.8) -- (549.8,837.2) -- (550.2,835.8) -- (550.6,834.2) -- (550.9,832.8) -- (551.3,831.2) -- (551.7,829.8) -- (552.2,828.2) -- (552.5,827.1) -- (552.7,826.2) -- (552.9,825.7) -- (552.9,825.5) -- (553.1,825.1) -- (553.3,824.3) -- (553.6,823.3) -- (553.9,822.0) -- (554.3,820.6) -- (554.7,819.2) -- (555.1,817.8) -- (555.5,816.2) -- (555.9,814.8) -- (556.3,813.2) -- (556.6,811.8) -- (557.0,810.2) -- (557.3,808.8) -- (557.6,807.2) -- (557.9,805.8) -- (558.2,804.2) -- (558.4,802.8) -- (558.6,801.4) -- (558.7,800.0) -- (558.8,798.7) -- (558.9,797.6) -- (559.0,796.9) -- (559.0,796.4) -- (559.0,796.2) -- (559.0,796.0) -- (559.0,795.7) -- (559.0,795.3) -- (559.0,794.9) -- (559.0,794.3) -- (559.0,793.5) -- (559.0,792.6) -- (558.9,791.6) -- (558.8,790.4) -- (558.7,789.1) -- (558.5,787.8) -- (558.3,786.2) -- (558.0,784.8) -- (557.6,783.2) -- (557.3,781.8) -- (556.8,780.2) -- (556.3,778.8) -- (555.8,777.2) -- (555.2,775.8) -- (554.5,774.2) -- (554.0,773.1) -- (553.6,772.2) -- (553.4,771.6) -- (553.2,771.4) -- (552.9,770.9) -- (552.5,770.1) -- (551.9,769.1) -- (551.1,767.9) -- (550.4,766.7) -- (549.6,765.5) -- (548.9,764.4) -- (548.1,763.4) -- (547.5,762.5) -- (546.9,761.8) -- (546.5,761.2) -- (546.1,760.7) -- (545.6,760.0) -- (544.8,759.1) -- (543.9,758.0) -- (542.7,756.7) -- (541.8,755.7) -- (541.2,755.0) -- (540.7,754.5) -- (540.6,754.3) -- (540.1,753.8) -- (539.3,753.1) -- (538.3,752.0) -- (537.0,750.7) -- (536.0,749.7) -- (535.2,748.9) -- (534.7,748.4) -- (534.6,748.3) -- (534.1,747.8) -- (533.4,747.1) -- (532.3,746.1) -- (531.0,744.9) -- (529.9,743.9) -- (529.1,743.2) -- (528.6,742.7) -- (528.3,742.4) -- (527.8,741.9) -- (527.1,741.2) -- (526.1,740.3) -- (524.8,739.2) -- (523.8,738.2) -- (523.0,737.5) -- (522.4,736.9) -- (522.1,736.5) -- (521.5,736.0) -- (520.8,735.3) -- (519.8,734.3) -- (518.7,733.2) -- (517.8,732.3) -- (517.0,731.6) -- (516.4,731.0) -- (516.0,730.6) -- (515.5,730.0) -- (514.7,729.2) -- (513.8,728.2) -- (512.7,727.0) -- (511.8,726.0) -- (511.1,725.2) -- (510.6,724.7) -- (510.4,724.4) -- (509.9,723.8) -- (509.2,722.9) -- (508.2,721.8) -- (507.1,720.2) -- (506.2,719.1) -- (505.5,718.3) -- (505.2,717.8) -- (505.1,717.7) -- (504.8,717.3) -- (504.3,716.6) -- (503.6,715.5) -- (502.6,714.1) -- (501.8,712.8) -- (501.0,711.6) -- (500.4,710.5) -- (499.8,709.5) -- (499.3,708.6) -- (498.9,707.9) -- (498.6,707.3) -- (498.3,706.8) -- (497.9,706.0) -- (497.4,705.0) -- (496.9,703.8) -- (496.2,702.2) -- (495.5,700.8) -- (494.9,699.2) -- (494.4,697.8) -- (493.9,696.2) -- (493.5,695.1) -- (493.2,694.3) -- (493.1,693.9) -- (493.0,693.8) -- (492.9,693.4) -- (492.7,692.6) -- (492.4,691.5) -- (492.0,690.1) -- (491.7,688.7) -- (491.4,687.2) -- (491.1,685.8) -- (490.8,684.2) -- (490.6,682.8) -- (490.3,681.2) -- (490.2,679.8) -- (490.0,678.2) -- (489.9,676.8) -- (489.8,675.2) -- (489.8,673.8) -- (489.7,672.2) -- (489.7,670.8) -- (489.7,669.2) -- (489.8,667.8) -- (489.9,666.2) -- (490.0,664.8) -- (490.1,663.2) -- (490.2,661.8) -- (490.4,660.2) -- (490.6,658.8) -- (490.8,657.2) -- (491.1,655.8) -- (491.4,654.2) -- (491.7,652.8) -- (492.0,651.5) -- (492.3,650.3) -- (492.6,649.1) -- (492.8,648.2) -- (493.0,647.4) -- (493.1,646.9) -- (493.2,646.5) -- (493.4,645.9) -- (493.7,645.0) -- (494.0,643.8) -- (494.5,642.2) -- (495.0,640.8) -- (495.5,639.2) -- (496.0,637.8) -- (496.6,636.2) -- (497.1,634.9) -- (497.6,633.7) -- (498.0,632.7) -- (498.4,631.8) -- (498.8,631.0) -- (499.1,630.2) -- (499.4,629.6) -- (499.7,628.9) -- (500.0,628.1) -- (500.5,627.0) -- (501.1,625.8) -- (501.8,624.2) -- (502.5,622.9) -- (503.1,621.7) -- (503.7,620.6) -- (504.2,619.6) -- (504.7,618.8) -- (505.1,618.0) -- (505.4,617.4) -- (505.7,616.8) -- (506.1,616.0) -- (506.7,615.0) -- (507.4,613.8) -- (508.2,612.2) -- (508.9,611.0) -- (509.5,609.9) -- (510.0,609.0) -- (510.4,608.4) -- (510.8,607.7) -- (511.3,607.0) -- (511.7,606.1) -- (512.2,605.3) -- (512.8,604.3) -- (513.5,603.1) -- (514.2,601.8) -- (515.1,600.2) -- (515.8,599.1) -- (516.3,598.2) -- (516.7,597.7) -- (516.8,597.4) -- (517.1,596.9) -- (517.5,596.2) -- (518.0,595.2) -- (518.7,593.9) -- (519.5,592.6) -- (520.2,591.2) -- (521.0,589.8) -- (521.8,588.4) -- (522.4,587.3) -- (522.8,586.5) -- (523.0,586.1) -- (523.1,586.1) -- (523.3,585.7) -- (523.6,584.9) -- (524.2,583.8) -- (524.9,582.2) -- (525.6,580.8) -- (526.3,579.2) -- (527.0,577.8) -- (527.7,576.2) -- (528.2,575.1) -- (528.6,574.2) -- (528.8,573.7) -- (528.9,573.6) -- (529.1,573.1) -- (529.4,572.4) -- (529.7,571.3) -- (530.2,570.0) -- (530.7,568.6) -- (531.2,567.2) -- (531.8,565.8) -- (532.3,564.2) -- (532.8,562.8) -- (533.3,561.2) -- (533.8,559.8) -- (534.2,558.2) -- (534.6,557.1) -- (534.8,556.3) -- (535.0,555.9) -- (535.0,555.9) -- (535.1,555.5) -- (535.3,554.7) -- (535.5,553.6) -- (535.9,552.2) -- (536.2,550.7) -- (536.6,549.2) -- (536.9,547.8) -- (537.2,546.2) -- (537.6,544.8) -- (537.9,543.2) -- (538.2,541.8) -- (538.5,540.2) -- (538.8,538.8) -- (539.0,537.2) -- (539.3,535.8) -- (539.6,534.2) -- (539.9,532.8) -- (540.1,531.2) -- (540.4,529.8) -- (540.6,528.2) -- (540.8,527.1) -- (540.9,526.4) -- (541.0,526.0) -- cycle ;
\draw[isobar] (637.1,658.0) -- (638.6,657.7) -- (641.5,657.1) -- (644.5,656.6) -- (647.5,656.1) -- (650.5,655.7) -- (653.5,655.5) -- (656.5,655.3) -- (659.5,655.2) -- (662.5,655.2) -- (665.5,655.3) -- (668.5,655.5) -- (671.5,655.8) -- (674.5,656.3) -- (677.5,656.8) -- (679.8,657.3) -- (681.4,657.8) -- (682.9,658.2) -- (684.3,658.6) -- (686.5,659.4) -- (689.5,660.6) -- (692.4,661.9) -- (695.2,663.3) -- (696.7,664.1) -- (696.9,664.2) -- (698.5,665.2) -- (701.5,667.2) -- (703.6,668.6) -- (704.8,669.5) -- (706.3,670.8) -- (708.1,672.4) -- (709.7,673.9) -- (711.2,675.3) -- (712.7,676.9) -- (714.2,678.6) -- (715.5,680.1) -- (716.6,681.4) -- (718.1,683.4) -- (720.0,686.1) -- (721.1,687.6) -- (721.2,687.9) -- (722.2,689.5) -- (723.8,692.5) -- (725.2,695.3) -- (726.4,697.8) -- (727.1,699.3) -- (727.3,699.8) -- (727.9,701.5) -- (729.0,704.5) -- (729.8,707.5) -- (730.6,710.5) -- (731.2,713.5) -- (731.6,716.5) -- (731.9,719.5) -- (732.0,722.5) -- (732.0,725.5) -- (731.7,728.5) -- (731.4,731.5) -- (730.8,734.5) -- (730.1,737.5) -- (729.2,740.5) -- (728.3,743.2) -- (727.4,745.5) -- (726.9,747.0) -- (726.6,747.7) -- (725.6,749.5) -- (724.1,752.5) -- (722.7,755.0) -- (721.6,756.9) -- (720.6,758.4) -- (719.9,759.5) -- (718.5,761.5) -- (716.3,764.5) -- (715.1,766.0) -- (715.0,766.1) -- (713.6,767.6) -- (710.9,770.5) -- (709.4,772.1) -- (709.1,772.4) -- (707.5,773.8) -- (704.5,776.2) -- (702.7,777.5) -- (702.2,777.8) -- (700.7,778.6) -- (698.2,779.9) -- (695.5,780.9) -- (692.5,781.8) -- (689.5,782.5) -- (686.5,783.0) -- (683.5,783.2) -- (680.5,783.3) -- (677.5,783.1) -- (674.5,782.8) -- (671.5,782.3) -- (668.5,781.5) -- (665.5,780.6) -- (662.5,779.5) -- (660.5,778.8) -- (659.4,778.3) -- (657.9,777.5) -- (656.0,776.6) -- (653.5,775.2) -- (650.5,773.4) -- (648.8,772.4) -- (648.4,772.1) -- (646.9,771.2) -- (644.3,769.6) -- (641.9,768.1) -- (639.8,766.7) -- (638.3,765.8) -- (637.4,765.3) -- (635.5,764.4) -- (632.5,763.0) -- (629.8,761.7) -- (627.3,760.6) -- (625.8,759.9) -- (625.3,759.7) -- (623.5,758.9) -- (620.5,757.4) -- (617.8,756.0) -- (615.5,754.7) -- (614.0,753.8) -- (613.3,753.5) -- (611.5,752.3) -- (608.5,750.4) -- (606.5,749.1) -- (605.6,748.4) -- (604.1,747.2) -- (602.0,745.6) -- (600.2,744.1) -- (598.7,742.7) -- (597.2,741.3) -- (595.7,739.9) -- (594.3,738.4) -- (592.9,736.8) -- (591.4,735.0) -- (589.8,732.9) -- (588.7,731.4) -- (588.1,730.5) -- (587.0,728.5) -- (585.5,725.5) -- (584.3,722.7) -- (583.4,720.1) -- (582.9,718.6) -- (582.8,718.2) -- (582.6,716.5) -- (582.2,713.5) -- (582.1,710.5) -- (582.2,707.5) -- (582.4,705.1) -- (582.8,703.4) -- (583.1,701.9) -- (583.4,700.6) -- (584.2,698.5) -- (585.4,695.5) -- (586.7,692.6) -- (588.2,689.9) -- (589.1,688.4) -- (589.3,688.1) -- (590.5,686.5) -- (592.9,683.5) -- (594.3,681.7) -- (594.8,681.2) -- (596.3,679.7) -- (598.8,677.2) -- (600.3,675.8) -- (600.8,675.4) -- (602.5,674.1) -- (605.5,671.9) -- (607.3,670.6) -- (607.9,670.2) -- (609.4,669.3) -- (611.8,667.8) -- (614.5,666.4) -- (617.5,664.9) -- (619.1,664.1) -- (619.2,664.0) -- (620.7,663.4) -- (623.6,662.3) -- (626.5,661.2) -- (629.5,660.2) -- (632.5,659.3) -- (635.5,658.4) -- (637.0,658.0) -- (637.1,658.0) -- (637.1,658.0);
\draw[isobar] (644.3,610.0) -- (645.5,609.8) -- (647.8,609.4) -- (650.5,609.0) -- (653.5,608.5) -- (656.5,608.1) -- (659.5,607.8) -- (662.5,607.5) -- (665.5,607.3) -- (668.5,607.2) -- (671.5,607.2) -- (674.5,607.2) -- (677.5,607.3) -- (680.5,607.4) -- (683.5,607.7) -- (686.5,608.0) -- (689.5,608.5) -- (692.5,609.0) -- (695.5,609.6) -- (697.1,609.9) -- (697.3,610.0) -- (698.8,610.4) -- (701.6,611.1) -- (704.5,612.0) -- (707.5,613.0) -- (710.5,614.1) -- (713.5,615.3) -- (715.0,616.0) -- (715.1,616.0) -- (716.6,616.7) -- (719.5,618.2) -- (722.3,619.7) -- (725.0,621.2) -- (726.5,622.1) -- (726.8,622.3) -- (728.5,623.4) -- (731.5,625.5) -- (733.5,626.9) -- (734.5,627.6) -- (736.0,628.8) -- (738.0,630.5) -- (739.7,632.0) -- (741.1,633.3) -- (742.6,634.8) -- (744.2,636.4) -- (745.6,637.9) -- (746.9,639.3) -- (748.4,641.0) -- (750.1,643.0) -- (751.4,644.5) -- (752.2,645.5) -- (753.6,647.5) -- (755.8,650.5) -- (756.9,652.1) -- (757.0,652.2) -- (757.9,653.7) -- (759.7,656.6) -- (761.2,659.1) -- (762.4,661.2) -- (763.2,662.7) -- (763.7,663.6) -- (764.7,665.5) -- (766.2,668.5) -- (767.5,671.1) -- (768.5,673.3) -- (769.2,674.8) -- (769.5,675.6) -- (770.3,677.5) -- (771.6,680.5) -- (772.8,683.5) -- (774.0,686.5) -- (774.7,688.3) -- (774.9,688.9) -- (775.5,690.4) -- (776.4,692.8) -- (777.3,695.5) -- (778.4,698.5) -- (779.5,701.5) -- (780.5,704.5) -- (781.0,706.0) -- (781.0,706.0) -- (781.5,707.5) -- (782.4,710.5) -- (783.4,713.5) -- (784.2,716.5) -- (785.0,719.5) -- (785.8,722.5) -- (786.4,725.2) -- (786.8,727.6) -- (787.1,729.1) -- (787.2,729.7) -- (787.4,731.5) -- (787.7,734.5) -- (787.8,737.5) -- (787.7,740.5) -- (787.5,743.0) -- (787.2,745.0) -- (786.9,746.5) -- (786.7,747.5) -- (786.2,749.5) -- (785.2,752.5) -- (784.0,755.5) -- (782.5,758.5) -- (781.6,760.3) -- (781.2,761.0) -- (780.3,762.5) -- (778.8,764.8) -- (777.3,767.1) -- (775.8,769.2) -- (774.7,770.7) -- (774.1,771.6) -- (772.6,773.5) -- (770.3,776.5) -- (769.1,778.1) -- (769.0,778.2) -- (767.9,779.7) -- (765.7,782.6) -- (764.2,784.5) -- (763.4,785.5) -- (762.2,787.0) -- (760.7,789.0) -- (759.2,790.9) -- (757.7,792.8) -- (756.6,794.3) -- (755.7,795.4) -- (754.2,797.3) -- (752.1,799.9) -- (750.8,801.4) -- (750.5,801.8) -- (749.0,803.5) -- (746.3,806.4) -- (745.0,807.9) -- (745.0,808.0) -- (743.5,809.5) -- (740.5,812.5) -- (737.5,815.4) -- (734.5,818.1) -- (732.8,819.6) -- (732.5,819.9) -- (731.0,821.1) -- (728.3,823.2) -- (726.4,824.7) -- (725.3,825.6) -- (723.8,826.7) -- (721.9,828.0) -- (719.8,829.5) -- (717.3,831.2) -- (715.8,832.2) -- (715.3,832.5) -- (713.5,833.7) -- (710.5,835.6) -- (708.4,836.9) -- (707.1,837.6) -- (705.6,838.5) -- (703.9,839.5) -- (701.5,840.9) -- (698.5,842.6) -- (696.7,843.6) -- (696.2,843.9) -- (694.7,844.7) -- (692.2,846.1) -- (689.5,847.5) -- (686.5,849.2) -- (685.0,850.0) -- (684.9,850.0) -- (683.4,850.8) -- (680.5,852.3) -- (677.6,853.8) -- (674.8,855.3) -- (673.3,856.1) -- (673.1,856.2) -- (671.5,857.0) -- (668.5,858.5) -- (665.6,859.9) -- (662.8,861.3) -- (661.3,862.0) -- (661.1,862.1) -- (659.5,862.9) -- (656.5,864.3) -- (653.5,865.7) -- (650.5,867.0) -- (648.8,867.7) -- (648.3,867.9) -- (646.8,868.5) -- (644.3,869.5) -- (641.5,870.5) -- (638.5,871.4) -- (635.5,872.3) -- (632.5,873.1) -- (630.2,873.6) -- (628.5,873.9) -- (627.0,874.1) -- (625.7,874.3) -- (623.5,874.5) -- (620.5,874.6) -- (617.5,874.5) -- (614.5,874.3) -- (612.8,874.1) -- (612.5,874.0) -- (611.0,873.7) -- (608.3,873.0) -- (605.5,871.9) -- (602.5,870.5) -- (600.3,869.4) -- (598.9,868.5) -- (597.4,867.4) -- (595.8,866.1) -- (594.1,864.6) -- (592.3,862.9) -- (590.8,861.3) -- (589.6,859.9) -- (588.4,858.4) -- (587.2,856.8) -- (585.7,854.7) -- (583.9,852.0) -- (582.9,850.5) -- (582.7,850.2) -- (581.8,848.5) -- (580.1,845.5) -- (578.7,842.8) -- (577.6,840.5) -- (576.8,839.0) -- (576.5,838.3) -- (575.7,836.5) -- (574.4,833.5) -- (573.1,830.5) -- (571.9,827.5) -- (571.2,825.8) -- (571.1,825.5) -- (570.5,824.0) -- (569.4,821.3) -- (568.3,818.5) -- (567.1,815.5) -- (566.2,813.0) -- (565.4,811.0) -- (564.8,809.5) -- (564.4,808.5) -- (563.6,806.5) -- (562.4,803.5) -- (561.2,800.5) -- (560.0,797.5) -- (559.3,795.8) -- (559.1,795.4) -- (558.4,793.9) -- (557.2,791.3) -- (555.9,788.5) -- (554.5,785.5) -- (553.5,783.7) -- (553.2,783.0) -- (552.4,781.5) -- (551.1,779.2) -- (549.6,776.5) -- (547.9,773.6) -- (547.0,772.1) -- (546.9,772.0) -- (546.0,770.5) -- (544.2,767.5) -- (542.8,764.9) -- (541.6,762.8) -- (540.8,761.3) -- (540.3,760.4) -- (539.4,758.5) -- (538.0,755.5) -- (536.8,752.5) -- (535.7,749.5) -- (535.1,747.8) -- (535.0,747.5) -- (534.7,746.0) -- (534.0,743.3) -- (533.4,740.5) -- (532.9,737.5) -- (532.5,734.5) -- (532.3,731.5) -- (532.2,728.5) -- (532.2,725.5) -- (532.3,722.5) -- (532.6,719.5) -- (533.0,716.5) -- (533.4,713.5) -- (534.0,710.6) -- (534.7,707.7) -- (535.0,706.2) -- (535.0,706.1) -- (535.5,704.5) -- (536.3,701.5) -- (537.3,698.5) -- (538.3,695.5) -- (539.4,692.7) -- (540.5,690.1) -- (541.1,688.6) -- (541.3,688.2) -- (542.0,686.5) -- (543.4,683.5) -- (544.9,680.7) -- (546.3,678.0) -- (547.1,676.5) -- (547.2,676.2) -- (548.2,674.5) -- (549.8,671.5) -- (551.2,669.0) -- (552.4,667.0) -- (553.3,665.5) -- (553.8,664.5) -- (555.0,662.5) -- (556.8,659.5) -- (558.0,657.5) -- (558.7,656.4) -- (559.6,654.9) -- (560.9,653.0) -- (562.4,650.8) -- (564.1,648.4) -- (565.2,646.9) -- (565.7,646.3) -- (567.2,644.5) -- (569.7,641.6) -- (571.0,640.1) -- (571.1,640.0) -- (572.6,638.6) -- (575.5,635.9) -- (577.2,634.4) -- (577.5,634.1) -- (579.0,633.0) -- (581.7,631.0) -- (583.8,629.5) -- (585.3,628.5) -- (586.8,627.6) -- (588.3,626.7) -- (590.5,625.5) -- (593.5,624.0) -- (595.7,622.9) -- (597.2,622.3) -- (598.7,621.7) -- (600.2,621.1) -- (602.5,620.2) -- (605.5,619.2) -- (608.5,618.2) -- (611.5,617.3) -- (613.9,616.7) -- (615.6,616.2) -- (617.1,615.8) -- (618.4,615.5) -- (620.5,615.0) -- (623.5,614.3) -- (626.5,613.6) -- (629.5,613.0) -- (632.5,612.3) -- (635.5,611.7) -- (638.5,611.1) -- (641.5,610.5) -- (643.3,610.2) -- (644.0,610.1) -- (644.3,610.0);
\draw[isobar] (415.0,520.0) -- (414.0,521.5) -- (412.2,524.5) -- (410.7,527.3) -- (409.6,530.0) -- (408.9,531.5) -- (408.8,531.8) -- (408.4,533.5) -- (407.7,536.5) -- (407.2,539.5) -- (406.9,542.5) -- (406.9,545.5) -- (407.0,548.5) -- (407.3,551.5) -- (407.8,554.5) -- (408.3,556.9) -- (408.8,558.6) -- (409.2,560.1) -- (409.5,561.4) -- (410.2,563.5) -- (411.4,566.5) -- (412.6,569.5) -- (414.1,572.5) -- (414.9,574.1) -- (415.0,574.2) -- (415.8,575.7) -- (417.5,578.6) -- (419.0,580.9) -- (420.3,582.8) -- (421.4,584.3) -- (422.1,585.4) -- (423.6,587.4) -- (425.9,590.1) -- (427.1,591.6) -- (427.3,591.9) -- (428.8,593.5) -- (431.6,596.5) -- (433.0,598.0) -- (433.0,598.0) -- (434.5,599.5) -- (437.5,602.3) -- (439.1,603.8) -- (439.2,603.9) -- (440.7,605.2) -- (443.6,607.7) -- (445.3,609.2) -- (446.0,609.7) -- (447.5,610.9) -- (449.8,612.8) -- (451.7,614.3) -- (453.1,615.4) -- (454.6,616.6) -- (456.2,617.9) -- (458.0,619.4) -- (460.0,621.1) -- (461.5,622.4) -- (462.5,623.3) -- (464.1,624.8) -- (466.3,626.9) -- (467.8,628.5) -- (468.6,629.4) -- (469.8,630.9) -- (471.3,633.0) -- (472.8,635.3) -- (474.3,637.8) -- (475.1,639.3) -- (475.3,639.8) -- (475.9,641.5) -- (476.8,644.5) -- (477.4,647.5) -- (477.7,650.5) -- (477.7,653.5) -- (477.4,656.5) -- (476.9,659.5) -- (476.2,662.5) -- (475.6,664.5) -- (475.2,665.6) -- (474.6,667.1) -- (473.8,669.0) -- (472.5,671.5) -- (471.0,674.5) -- (469.9,676.5) -- (469.3,677.5) -- (468.3,679.0) -- (467.0,681.0) -- (465.5,683.2) -- (463.8,685.6) -- (462.8,687.1) -- (462.3,687.7) -- (461.0,689.5) -- (458.7,692.5) -- (457.4,694.2) -- (457.1,694.6) -- (455.9,696.1) -- (453.8,698.7) -- (452.3,700.6) -- (451.4,701.7) -- (450.2,703.2) -- (448.7,705.1) -- (447.2,707.0) -- (445.7,708.9) -- (444.6,710.4) -- (443.7,711.5) -- (442.2,713.4) -- (440.1,716.3) -- (438.9,717.8) -- (438.8,717.9) -- (437.7,719.5) -- (435.5,722.5) -- (434.1,724.5) -- (433.4,725.6) -- (432.3,727.1) -- (431.0,729.0) -- (429.5,731.4) -- (427.8,734.1) -- (426.9,735.6) -- (426.8,735.9) -- (425.8,737.5) -- (424.1,740.5) -- (422.7,743.2) -- (421.6,745.7) -- (420.9,747.2) -- (420.6,747.7) -- (419.8,749.5) -- (418.6,752.5) -- (417.6,755.5) -- (416.7,758.5) -- (415.9,761.4) -- (415.3,764.2) -- (415.0,765.7) -- (414.9,765.9) -- (414.7,767.5) -- (414.3,770.5) -- (414.1,773.5) -- (414.0,776.5) -- (414.0,779.5) -- (414.2,782.5) -- (414.5,785.1) -- (414.8,787.3) -- (415.1,788.8) -- (415.2,789.6) -- (415.6,791.5) -- (416.4,794.5) -- (417.3,797.5) -- (418.4,800.5) -- (419.5,803.0) -- (420.5,805.1) -- (421.2,806.6) -- (421.7,807.5) -- (422.8,809.5) -- (424.6,812.5) -- (425.9,814.5) -- (426.6,815.5) -- (427.7,817.0) -- (429.2,819.0) -- (430.7,820.9) -- (432.2,822.6) -- (433.6,824.1) -- (434.7,825.4) -- (436.2,826.9) -- (438.1,828.7) -- (439.6,830.2) -- (440.9,831.4) -- (442.4,832.7) -- (444.1,834.2) -- (445.9,835.7) -- (447.8,837.2) -- (449.3,838.5) -- (450.4,839.4) -- (452.4,840.9) -- (455.2,843.0) -- (456.7,844.1) -- (456.9,844.2) -- (458.5,845.4) -- (461.5,847.6) -- (463.4,849.0) -- (464.3,849.7) -- (465.8,850.8) -- (467.9,852.4) -- (469.9,853.9) -- (471.8,855.3) -- (473.3,856.4) -- (474.4,857.3) -- (476.3,858.8) -- (479.0,860.9) -- (480.5,862.1) -- (480.8,862.4) -- (482.5,863.9) -- (485.5,866.5) -- (487.1,867.8) -- (487.2,867.9) -- (488.7,869.3) -- (491.6,872.0) -- (493.1,873.5) -- (493.4,873.8) -- (494.9,875.4) -- (497.6,878.1) -- (499.1,879.6) -- (499.3,879.9) -- (500.8,881.5) -- (503.5,884.5) -- (504.9,886.0) -- (505.0,886.1) -- (506.2,887.6) -- (508.5,890.5) -- (510.0,892.4) -- (510.7,893.3) -- (511.8,894.8) -- (513.3,896.9) -- (514.8,899.1) -- (516.3,901.4) -- (517.2,902.9) -- (517.7,903.6) -- (518.9,905.5) -- (520.7,908.5) -- (521.9,910.6) -- (522.6,911.9) -- (523.5,913.4) -- (524.6,915.1) -- (526.0,917.5) -- (527.7,920.5) -- (528.7,922.2) -- (528.9,922.5) -- (529.9,924.0) -- (531.6,926.7) -- (533.1,928.9) -- (534.4,930.7) -- (535.5,932.2) -- (536.4,933.4) -- (537.9,935.2) -- (540.0,937.5) -- (541.3,939.0) -- (542.0,939.7) -- (543.5,941.0) -- (545.8,943.1) -- (547.6,944.6) -- (548.9,945.5) -- (550.4,946.6) -- (552.1,947.7) -- (554.5,949.0) -- (557.5,950.6) -- (559.4,951.5) -- (560.1,951.8) -- (561.6,952.4) -- (563.9,953.3) -- (566.5,954.2) -- (569.5,955.1) -- (572.5,955.9) -- (575.5,956.6) -- (578.5,957.2) -- (581.5,957.7) -- (583.0,958.0) -- (583.1,958.0) -- (584.6,958.2) -- (587.5,958.5) -- (590.5,958.8) -- (593.5,959.1) -- (596.5,959.2) -- (599.5,959.3) -- (602.5,959.4) -- (605.5,959.4) -- (608.5,959.4) -- (611.5,959.3) -- (614.5,959.1) -- (617.5,958.9) -- (620.5,958.6) -- (623.5,958.2) -- (625.0,958.0) -- (625.1,958.0) -- (626.6,957.7) -- (629.5,957.2) -- (632.5,956.5) -- (635.5,955.7) -- (638.5,954.8) -- (641.5,953.6) -- (643.5,952.8) -- (644.6,952.3) -- (646.1,951.5) -- (648.0,950.5) -- (650.5,949.0) -- (653.5,947.0) -- (655.0,946.0) -- (655.0,946.0) -- (656.5,944.8) -- (659.5,942.5) -- (661.4,941.0) -- (662.1,940.3) -- (663.6,939.0) -- (665.9,937.0) -- (667.5,935.5) -- (668.6,934.5) -- (670.1,933.1) -- (672.0,931.2) -- (673.6,929.7) -- (674.7,928.6) -- (676.2,927.1) -- (678.1,925.2) -- (679.6,923.7) -- (680.7,922.6) -- (682.2,921.1) -- (684.1,919.2) -- (685.5,917.7) -- (686.6,916.6) -- (688.1,915.0) -- (690.0,913.1) -- (691.5,911.6) -- (692.6,910.5) -- (694.1,909.1) -- (696.0,907.3) -- (697.6,905.8) -- (698.8,904.6) -- (700.3,903.2) -- (702.1,901.6) -- (703.8,900.1) -- (705.4,898.7) -- (706.9,897.4) -- (708.3,896.3) -- (710.2,894.8) -- (712.5,892.9) -- (714.0,891.8) -- (714.7,891.3) -- (716.5,890.0) -- (719.5,887.8) -- (721.3,886.6) -- (721.8,886.2) -- (723.3,885.2) -- (725.8,883.6) -- (728.1,882.1) -- (730.3,880.7) -- (731.8,879.8) -- (732.6,879.3) -- (734.5,878.1) -- (737.5,876.3) -- (739.6,875.1) -- (740.9,874.4) -- (742.4,873.5) -- (744.1,872.5) -- (746.5,871.1) -- (749.5,869.4) -- (751.2,868.4) -- (751.7,868.1) -- (753.2,867.3) -- (755.7,865.9) -- (758.5,864.4) -- (761.4,862.8) -- (762.9,862.0) -- (763.0,861.9) -- (764.5,861.1) -- (767.5,859.5) -- (770.3,858.0) -- (773.0,856.7) -- (774.5,855.9) -- (774.8,855.7) -- (776.5,854.9) -- (779.5,853.4) -- (782.5,852.0) -- (785.5,850.7) -- (787.0,850.0) -- (787.1,850.0) -- (788.6,849.4) -- (791.5,848.2) -- (794.5,847.1) -- (797.5,846.1) -- (800.3,845.2) -- (802.9,844.4) -- (804.4,843.9) -- (804.8,843.8) -- (806.5,843.3) -- (809.5,842.5) -- (812.5,841.7) -- (815.5,840.8) -- (818.5,840.0) -- (821.5,839.0) -- (823.4,838.4) -- (824.2,838.1) -- (825.7,837.6) -- (827.9,836.7) -- (830.5,835.5) -- (833.5,834.0) -- (835.5,832.9) -- (836.4,832.3) -- (837.9,831.2) -- (840.0,829.7) -- (841.8,828.2) -- (843.4,826.7) -- (844.9,825.2) -- (846.3,823.7) -- (847.6,822.2) -- (848.7,820.7) -- (850.2,818.5) -- (852.0,815.5) -- (852.9,813.9) -- (853.0,813.8) -- (853.6,812.3) -- (854.9,809.4) -- (856.0,806.5) -- (857.0,803.5) -- (857.8,800.5) -- (858.5,797.5) -- (858.9,795.7) -- (859.0,795.1) -- (859.2,793.6) -- (859.5,791.2) -- (859.8,788.5) -- (860.0,785.5) -- (860.1,782.5) -- (860.1,779.5) -- (860.0,776.5) -- (859.8,773.5) -- (859.5,770.5) -- (859.2,767.5) -- (859.0,766.0) -- (859.0,765.9) -- (858.7,764.4) -- (858.2,761.5) -- (857.7,758.5) -- (857.0,755.5) -- (856.2,752.5) -- (855.4,749.5) -- (854.5,746.6) -- (853.5,743.8) -- (853.0,742.3) -- (852.9,742.1) -- (852.3,740.5) -- (851.1,737.5) -- (849.7,734.5) -- (848.3,731.5) -- (847.5,729.7) -- (847.2,729.1) -- (846.4,727.6) -- (845.1,725.2) -- (843.6,722.6) -- (841.9,719.7) -- (841.0,718.2) -- (840.9,718.1) -- (839.8,716.5) -- (837.9,713.5) -- (836.4,711.3) -- (835.5,709.9) -- (834.5,708.4) -- (833.4,706.8) -- (831.9,704.6) -- (830.0,701.9) -- (828.9,700.4) -- (828.7,700.1) -- (827.7,698.5) -- (825.7,695.5) -- (824.3,693.3) -- (823.4,691.9) -- (822.6,690.4) -- (821.7,688.8) -- (820.6,686.5) -- (819.2,683.5) -- (818.1,680.9) -- (817.4,678.7) -- (816.9,677.2) -- (816.6,676.4) -- (816.1,674.5) -- (815.3,671.5) -- (814.5,668.5) -- (813.7,665.5) -- (813.0,662.5) -- (812.2,659.5) -- (811.6,657.2) -- (811.2,655.7) -- (810.8,654.2) -- (810.4,652.7) -- (809.7,650.5) -- (808.8,647.5) -- (807.8,644.5) -- (806.6,641.5) -- (805.8,639.4) -- (805.3,638.2) -- (804.6,636.7) -- (803.7,634.9) -- (802.5,632.5) -- (800.9,629.5) -- (799.8,627.5) -- (799.3,626.6) -- (798.4,625.1) -- (797.1,623.0) -- (795.6,620.7) -- (793.9,618.2) -- (792.8,616.7) -- (792.5,616.2) -- (791.2,614.5) -- (788.9,611.5) -- (787.6,609.8) -- (787.2,609.3) -- (786.0,607.8) -- (783.9,605.3) -- (782.4,603.5) -- (781.5,602.4) -- (780.1,600.9) -- (778.3,599.0) -- (776.8,597.4) -- (775.6,596.1) -- (774.1,594.6) -- (772.4,592.9) -- (770.9,591.4) -- (769.6,590.2) -- (768.0,588.7) -- (766.1,586.9) -- (764.6,585.5) -- (763.5,584.6) -- (761.7,583.1) -- (759.1,581.0) -- (757.6,579.8) -- (757.2,579.5) -- (755.5,578.2) -- (752.5,576.0) -- (750.7,574.7) -- (750.0,574.2) -- (748.5,573.2) -- (746.2,571.6) -- (743.8,570.1) -- (741.5,568.7) -- (740.0,567.8) -- (739.3,567.4) -- (737.5,566.4) -- (734.5,564.7) -- (732.0,563.4) -- (730.1,562.5) -- (728.6,561.7) -- (727.5,561.2) -- (725.5,560.3) -- (722.5,559.0) -- (719.5,557.7) -- (716.6,556.6) -- (715.1,556.0) -- (715.0,556.0) -- (713.5,555.4) -- (710.5,554.4) -- (707.5,553.4) -- (704.5,552.5) -- (701.5,551.6) -- (698.5,550.9) -- (696.5,550.4) -- (695.4,550.1) -- (693.9,549.8) -- (692.0,549.3) -- (689.5,548.8) -- (686.5,548.2) -- (683.5,547.7) -- (680.5,547.2) -- (677.5,546.8) -- (674.5,546.4) -- (671.5,546.1) -- (668.5,545.8) -- (665.5,545.5) -- (662.5,545.3) -- (659.5,545.1) -- (656.5,544.9) -- (653.5,544.8) -- (650.5,544.6) -- (647.5,544.5) -- (644.5,544.4) -- (641.5,544.2) -- (638.5,544.1) -- (636.8,544.0) -- (636.5,544.0) -- (635.0,543.9) -- (632.3,543.8) -- (629.5,543.6) -- (626.5,543.4) -- (623.5,543.2) -- (620.5,542.9) -- (617.5,542.6) -- (614.5,542.3) -- (611.5,541.9) -- (608.5,541.5) -- (605.5,541.0) -- (602.5,540.5) -- (599.5,540.0) -- (596.5,539.4) -- (593.7,538.9) -- (591.0,538.3) -- (589.5,538.0) -- (589.2,537.9) -- (587.5,537.5) -- (584.5,536.7) -- (581.5,535.9) -- (578.5,535.1) -- (575.5,534.2) -- (572.5,533.3) -- (570.4,532.6) -- (569.2,532.2) -- (567.7,531.7) -- (565.9,531.1) -- (563.5,530.2) -- (560.5,529.1) -- (557.5,528.0) -- (554.5,526.8) -- (552.8,526.2) -- (552.5,526.1) -- (551.0,525.4) -- (548.3,524.3) -- (545.5,523.2) -- (542.5,521.9) -- (540.3,520.9) -- (538.8,520.3) -- (538.1,520.0);
\draw[isobar] (307.0,1012.5) -- (308.5,1012.9) -- (311.5,1013.7) -- (314.5,1014.5) -- (317.5,1015.2) -- (320.5,1015.8) -- (323.5,1016.4) -- (326.5,1016.9) -- (329.5,1017.4) -- (331.8,1017.7) -- (333.5,1017.9) -- (335.0,1018.1) -- (336.3,1018.2) -- (338.5,1018.5) -- (341.5,1018.7) -- (344.5,1018.9) -- (347.5,1019.1) -- (350.5,1019.2) -- (353.5,1019.2) -- (356.5,1019.2) -- (359.5,1019.1) -- (362.5,1019.0) -- (365.5,1018.8) -- (368.5,1018.5) -- (371.5,1018.2) -- (373.2,1018.1) -- (373.6,1018.0) -- (375.1,1017.8) -- (377.7,1017.5) -- (380.5,1017.1) -- (383.5,1016.7) -- (386.5,1016.2) -- (389.5,1015.7) -- (392.5,1015.1) -- (395.5,1014.5) -- (398.5,1013.8) -- (401.5,1013.2) -- (403.9,1012.6) -- (405.6,1012.2) -- (407.1,1011.9) -- (408.4,1011.6) -- (410.5,1011.1) -- (413.5,1010.5) -- (416.5,1009.9) -- (419.5,1009.2) -- (422.5,1008.7) -- (425.5,1008.1) -- (428.5,1007.6) -- (431.5,1007.2) -- (434.5,1006.9) -- (437.5,1006.7) -- (440.5,1006.7) -- (443.5,1007.0) -- (446.5,1007.5) -- (449.5,1008.4) -- (452.1,1009.7) -- (454.2,1011.2) -- (455.7,1012.5) -- (456.6,1013.6) -- (457.5,1015.1) -- (458.6,1017.0) -- (459.6,1019.5) -- (460.5,1022.5) -- (461.1,1025.5) -- (461.4,1028.5) -- (461.4,1031.5) -- (461.1,1034.5) -- (460.6,1037.5) -- (460.0,1040.5) -- (459.1,1043.5) -- (458.1,1046.5) -- (457.4,1048.3) -- (457.1,1049.0) -- (456.4,1050.5) -- (455.3,1052.8) -- (453.8,1055.5) -- (452.1,1058.5) -- (451.2,1060.1) -- (451.1,1060.3) -- (450.0,1061.8) -- (448.1,1064.6) -- (446.6,1066.7) -- (445.5,1068.2) -- (444.4,1069.7) -- (443.2,1071.2) -- (441.7,1073.1) -- (439.9,1075.3) -- (438.7,1076.8) -- (438.0,1077.6) -- (436.5,1079.4) -- (434.2,1082.1) -- (432.9,1083.6) -- (432.6,1083.9) -- (431.2,1085.5) -- (428.5,1088.5) -- (427.1,1090.0) -- (427.0,1090.1) -- (425.7,1091.6) -- (423.0,1094.5) -- (421.5,1096.2) -- (421.2,1096.5) -- (419.8,1098.0) -- (417.3,1100.7) -- (415.8,1102.3) -- (415.3,1102.9) -- (413.9,1104.4) -- (411.7,1106.8) -- (410.2,1108.5) -- (409.4,1109.4) -- (409.0,1109.8);
\draw[isobar] (784.3,520.0) -- (785.0,520.6) -- (786.3,521.9) -- (787.9,523.4) -- (789.6,525.1) -- (791.1,526.6) -- (792.4,527.9) -- (793.8,529.4) -- (795.5,531.1) -- (797.0,532.7) -- (798.3,534.1) -- (799.7,535.6) -- (801.1,537.2) -- (802.6,538.9) -- (804.2,540.6) -- (805.5,542.1) -- (806.6,543.4) -- (808.1,545.2) -- (810.0,547.5) -- (811.3,549.0) -- (811.8,549.7) -- (813.3,551.5) -- (815.7,554.5) -- (816.9,556.0) -- (817.0,556.1) -- (818.1,557.6) -- (820.4,560.5) -- (821.9,562.5) -- (822.6,563.5) -- (823.7,565.0) -- (825.2,567.0) -- (826.7,569.0) -- (828.2,571.1) -- (829.3,572.6) -- (830.0,573.5) -- (831.5,575.5) -- (833.6,578.5) -- (834.8,580.1) -- (834.9,580.3) -- (836.0,581.8) -- (838.0,584.6) -- (839.5,586.7) -- (840.5,588.1) -- (841.5,589.6) -- (842.6,591.2) -- (844.1,593.5) -- (846.0,596.4) -- (847.0,597.9) -- (847.1,598.0) -- (848.0,599.5) -- (849.8,602.5) -- (851.3,605.0) -- (852.4,607.0) -- (853.3,608.5) -- (853.9,609.5) -- (854.9,611.5) -- (856.5,614.5) -- (857.7,616.9) -- (858.6,618.7) -- (859.3,620.2) -- (859.9,621.4) -- (860.9,623.5) -- (862.2,626.5) -- (863.4,629.2) -- (864.5,631.7) -- (865.1,633.2) -- (865.4,633.7) -- (866.1,635.5) -- (867.4,638.5) -- (868.6,641.5) -- (869.8,644.5) -- (870.6,646.4) -- (870.9,647.1) -- (871.5,648.6) -- (872.5,650.9) -- (873.6,653.5) -- (874.9,656.5) -- (875.9,658.8) -- (876.6,660.5) -- (877.3,662.0) -- (877.9,663.3) -- (879.0,665.5) -- (880.4,668.5) -- (881.6,671.0) -- (882.5,672.9) -- (883.3,674.4) -- (883.8,675.5) -- (884.9,677.5) -- (886.4,680.5) -- (887.6,682.9) -- (888.5,684.6) -- (889.3,686.1) -- (890.0,687.4) -- (891.1,689.5) -- (892.7,692.5) -- (893.9,694.7) -- (894.6,696.1) -- (895.4,697.6) -- (896.2,699.2) -- (897.4,701.5) -- (898.9,704.5) -- (900.0,706.7) -- (900.7,708.0) -- (901.4,709.5) -- (902.3,711.2) -- (903.4,713.5) -- (904.9,716.5) -- (906.0,718.7) -- (906.7,720.2) -- (907.4,721.7) -- (908.1,723.2) -- (909.2,725.5) -- (910.5,728.5) -- (911.6,731.1) -- (912.5,733.3) -- (913.2,734.8) -- (913.5,735.6) -- (914.2,737.5) -- (915.4,740.5) -- (916.5,743.5) -- (917.5,746.5) -- (918.3,748.8) -- (918.8,750.5) -- (919.2,752.0) -- (919.6,753.3) -- (920.2,755.5) -- (921.0,758.5) -- (921.7,761.5) -- (922.3,764.5) -- (922.9,767.5) -- (923.4,770.5) -- (923.9,773.5) -- (924.2,776.5) -- (924.5,779.5) -- (924.8,782.5) -- (924.9,784.9) -- (925.0,786.6) -- (925.0,788.1) -- (925.1,789.4) -- (925.1,791.2) -- (925.0,793.6) -- (925.0,795.1) -- (925.0,795.7) -- (924.9,797.5) -- (924.7,800.5) -- (924.4,803.5) -- (924.0,806.5) -- (923.6,809.5) -- (923.1,812.5) -- (922.5,815.5) -- (921.8,818.5) -- (921.0,821.5) -- (920.1,824.5) -- (919.5,826.5) -- (919.2,827.4) -- (918.6,828.9) -- (917.9,831.0) -- (916.9,833.5) -- (915.6,836.5) -- (914.5,839.0) -- (913.5,841.0) -- (912.7,842.5) -- (912.2,843.5) -- (911.1,845.5) -- (909.3,848.5) -- (908.1,850.5) -- (907.4,851.6) -- (906.3,853.1) -- (905.0,855.0) -- (903.5,857.1) -- (901.8,859.2) -- (900.6,860.7) -- (899.9,861.6) -- (898.4,863.3) -- (896.1,865.8) -- (894.8,867.3) -- (894.3,867.8) -- (892.8,869.2) -- (890.3,871.6) -- (888.7,873.1) -- (888.0,873.7) -- (886.5,875.0) -- (884.2,876.9) -- (882.3,878.4) -- (880.8,879.5) -- (879.3,880.6) -- (877.8,881.7) -- (875.5,883.2) -- (872.6,885.1) -- (871.1,886.0) -- (871.0,886.1) -- (869.5,887.0) -- (866.5,888.7) -- (863.8,890.2) -- (861.5,891.4) -- (860.0,892.2) -- (859.3,892.5) -- (857.5,893.4) -- (854.5,894.8) -- (851.6,896.2) -- (848.8,897.4) -- (847.3,898.0) -- (847.1,898.1) -- (845.5,898.8) -- (842.5,900.1) -- (839.5,901.3) -- (836.5,902.5) -- (834.5,903.4) -- (833.4,903.8) -- (831.9,904.4) -- (830.0,905.2) -- (827.5,906.2) -- (824.5,907.4) -- (821.8,908.5) -- (819.4,909.5) -- (817.9,910.1) -- (817.3,910.4) -- (815.5,911.2) -- (812.5,912.4) -- (809.5,913.7) -- (806.5,915.0) -- (804.8,915.7) -- (804.3,915.9) -- (802.8,916.6) -- (800.3,917.7) -- (797.5,919.0) -- (794.5,920.3) -- (792.5,921.3) -- (791.4,921.8) -- (789.9,922.5) -- (788.0,923.4) -- (785.5,924.7) -- (782.5,926.2) -- (780.5,927.2) -- (779.5,927.7) -- (778.0,928.6) -- (776.0,929.7) -- (773.5,931.2) -- (770.5,933.0) -- (769.0,933.9) -- (768.9,934.0) -- (767.4,935.0) -- (764.5,937.1) -- (762.4,938.6) -- (761.2,939.5) -- (759.7,940.7) -- (757.9,942.2) -- (756.2,943.7) -- (754.6,945.2) -- (753.1,946.7) -- (751.7,948.1) -- (750.3,949.6) -- (748.8,951.2) -- (747.3,952.9) -- (745.8,954.7) -- (744.5,956.2) -- (743.6,957.4) -- (742.1,959.4) -- (740.0,962.1) -- (738.9,963.6) -- (738.7,963.9) -- (737.6,965.5) -- (735.5,968.5) -- (734.1,970.5) -- (733.4,971.6) -- (732.3,973.1) -- (731.0,975.0) -- (729.5,977.2) -- (727.8,979.7) -- (726.8,981.2) -- (726.5,981.7) -- (725.3,983.5) -- (723.2,986.5) -- (721.9,988.4) -- (721.3,989.2) -- (720.2,990.7) -- (718.6,992.9) -- (717.1,994.9) -- (715.7,996.8) -- (714.5,998.3) -- (713.6,999.4) -- (712.1,1001.2) -- (710.0,1003.7) -- (708.8,1005.2) -- (708.3,1005.7) -- (706.8,1007.3) -- (704.3,1010.0) -- (702.8,1011.5) -- (702.5,1011.8) -- (701.0,1013.2) -- (698.3,1015.7) -- (696.7,1017.2) -- (696.1,1017.7) -- (694.6,1019.0) -- (692.2,1021.0) -- (690.3,1022.5) -- (688.9,1023.5) -- (687.4,1024.6) -- (685.8,1025.8) -- (683.6,1027.3) -- (680.9,1029.1) -- (679.4,1030.1) -- (679.1,1030.3) -- (677.5,1031.3) -- (674.5,1033.1) -- (672.1,1034.5) -- (670.3,1035.5) -- (668.8,1036.4) -- (667.6,1037.1) -- (665.5,1038.2) -- (662.5,1039.8) -- (660.3,1041.0) -- (658.9,1041.7) -- (657.4,1042.4) -- (655.8,1043.3) -- (653.5,1044.5) -- (650.5,1046.0) -- (648.4,1047.1) -- (647.2,1047.7) -- (645.7,1048.5) -- (643.9,1049.5) -- (641.5,1050.8) -- (638.5,1052.6) -- (636.8,1053.6) -- (636.3,1053.9) -- (634.8,1054.9) -- (632.3,1056.7) -- (630.2,1058.2) -- (628.7,1059.4) -- (627.2,1060.7) -- (625.7,1062.0) -- (624.2,1063.5) -- (622.7,1065.2) -- (621.2,1066.9) -- (619.7,1068.7) -- (618.6,1070.2) -- (617.7,1071.4) -- (616.2,1073.5) -- (614.3,1076.5) -- (613.2,1078.1) -- (613.1,1078.4) -- (612.2,1079.9) -- (610.7,1082.6) -- (609.2,1085.5) -- (607.7,1088.4) -- (607.0,1089.9) -- (606.9,1090.0) -- (606.2,1091.5) -- (604.9,1094.5) -- (603.5,1097.5) -- (602.2,1100.5) -- (601.4,1102.4) -- (601.1,1103.1) -- (600.5,1104.6) -- (599.6,1106.9) -- (598.6,1109.5) -- (597.4,1112.5) -- (596.3,1115.2) -- (595.4,1117.5) -- (594.9,1119.0) -- (594.6,1119.7) -- (593.9,1121.5) -- (592.8,1124.5) -- (591.7,1127.5) -- (590.6,1130.5) -- (589.7,1132.7) -- (589.2,1134.0) -- (588.7,1135.5) -- (588.1,1137.2) -- (587.2,1139.5) -- (586.1,1142.5) -- (585.6,1144.0);
\draw[isobar] (882.8,520.0) -- (882.9,520.1) -- (883.0,520.2) -- (883.9,521.7) -- (885.8,524.6) -- (887.3,526.9) -- (888.4,528.7) -- (889.4,530.2) -- (890.2,531.4) -- (891.5,533.5) -- (893.4,536.5) -- (894.5,538.3) -- (894.8,538.8) -- (895.8,540.3) -- (897.3,542.8) -- (898.8,545.2) -- (900.3,547.7) -- (901.2,549.2) -- (901.5,549.7) -- (902.6,551.5) -- (904.3,554.5) -- (905.7,556.8) -- (906.6,558.3) -- (907.4,559.8) -- (908.3,561.3) -- (909.6,563.5) -- (911.3,566.5) -- (912.4,568.4) -- (912.8,569.1) -- (913.6,570.6) -- (914.9,572.9) -- (916.4,575.5) -- (918.0,578.5) -- (918.9,580.1) -- (919.0,580.2) -- (919.8,581.7) -- (921.3,584.6) -- (922.8,587.4) -- (924.3,590.1) -- (925.1,591.6) -- (925.2,591.9) -- (926.1,593.5) -- (927.6,596.5) -- (929.1,599.3) -- (930.4,601.8) -- (931.0,603.0);
\draw[isobar] (931.5,604.0) -- (932.3,605.5) -- (933.8,608.5) -- (935.2,611.2) -- (936.4,613.7) -- (937.1,615.2) -- (937.4,615.7) -- (938.3,617.5) -- (939.8,620.5) -- (941.1,623.3) -- (942.4,625.8) -- (943.1,627.3) -- (943.3,627.8) -- (944.2,629.5) -- (945.6,632.5) -- (947.0,635.4) -- (948.3,638.1) -- (949.1,639.6) -- (949.2,639.9) -- (950.0,641.5) -- (951.4,644.5) -- (952.8,647.5) -- (954.2,650.5) -- (955.0,652.0) -- (955.0,652.1) -- (955.7,653.6) -- (957.1,656.5) -- (958.5,659.5) -- (959.9,662.5) -- (960.7,664.2) -- (960.9,664.7) -- (961.6,666.2) -- (962.7,668.7) -- (964.0,671.5) -- (965.3,674.5) -- (966.2,676.6) -- (966.7,677.7) -- (967.4,679.2) -- (968.2,681.1) -- (969.3,683.5) -- (970.6,686.5) -- (971.6,689.1) -- (972.5,691.2) -- (973.2,692.7) -- (973.5,693.6) -- (974.3,695.5) -- (975.5,698.5) -- (976.7,701.5) -- (977.9,704.5) -- (978.6,706.4) -- (978.9,707.1) -- (979.4,708.6) -- (980.3,710.9) -- (981.2,713.5) -- (982.2,716.5) -- (983.3,719.5) -- (984.3,722.5) -- (984.8,724.2) -- (984.9,724.6) -- (985.4,726.1) -- (986.2,728.7) -- (987.0,731.5) -- (987.9,734.5) -- (988.7,737.5) -- (989.5,740.5) -- (990.2,743.1) -- (990.7,745.4) -- (991.1,746.9) -- (991.3,747.6) -- (991.7,749.5) -- (992.3,752.5) -- (992.9,755.5) -- (993.5,758.5) -- (994.0,761.5) -- (994.5,764.5) -- (994.9,767.5) -- (995.3,770.5) -- (995.7,773.5) -- (996.0,776.5) -- (996.1,778.0);
\draw[isobar] (758.2,1132.0) -- (759.2,1130.5) -- (761.0,1127.5) -- (762.2,1125.6) -- (762.7,1124.8) -- (763.7,1123.3) -- (765.0,1121.1) -- (766.5,1118.7) -- (768.2,1116.1) -- (769.1,1114.6) -- (769.4,1114.2) -- (770.5,1112.5) -- (772.4,1109.5) -- (773.8,1107.4) -- (774.6,1106.1) -- (775.6,1104.6) -- (776.7,1102.9) -- (778.2,1100.6) -- (780.1,1097.7) -- (781.0,1096.2) -- (781.1,1096.1) -- (782.1,1094.5) -- (784.0,1091.5) -- (785.5,1089.2) -- (786.5,1087.7) -- (787.5,1086.2) -- (788.4,1084.7) -- (789.8,1082.5) -- (791.8,1079.5) -- (792.8,1077.9) -- (792.9,1077.7) -- (793.9,1076.2) -- (795.7,1073.4) -- (797.2,1071.1) -- (798.4,1069.2) -- (799.4,1067.7) -- (800.1,1066.6) -- (801.5,1064.5) -- (803.4,1061.5) -- (804.6,1059.8) -- (804.9,1059.3) -- (805.8,1057.8) -- (807.5,1055.3) -- (809.0,1053.0) -- (810.3,1051.1) -- (811.4,1049.6) -- (812.1,1048.5) -- (813.5,1046.5) -- (815.6,1043.5) -- (816.8,1041.9) -- (816.9,1041.7) -- (818.0,1040.2) -- (820.1,1037.4) -- (821.6,1035.4) -- (822.5,1034.1) -- (823.7,1032.6) -- (825.1,1030.9) -- (826.6,1029.0) -- (828.2,1027.1) -- (829.4,1025.6) -- (830.3,1024.5) -- (831.8,1022.8) -- (833.9,1020.5) -- (835.3,1019.0) -- (836.0,1018.3) -- (837.5,1016.8) -- (839.8,1014.5) -- (841.3,1013.0) -- (842.0,1012.3) -- (843.5,1011.0) -- (845.8,1008.9) -- (847.5,1007.4) -- (848.5,1006.5) -- (850.0,1005.2) -- (852.0,1003.5) -- (853.8,1002.0) -- (855.5,1000.7) -- (857.0,999.5) -- (858.3,998.4) -- (860.3,996.9) -- (862.8,995.0) -- (864.3,993.8) -- (864.8,993.5) -- (866.5,992.2) -- (869.5,990.0) -- (871.3,988.7) -- (871.9,988.2) -- (873.4,987.1) -- (875.8,985.4) -- (877.9,983.9) -- (879.6,982.6) -- (881.1,981.6) -- (882.4,980.7) -- (884.5,979.2) -- (887.4,977.1) -- (888.9,976.0) -- (889.0,975.9) -- (890.5,974.8) -- (893.5,972.7) -- (895.5,971.2) -- (896.6,970.4) -- (898.1,969.3) -- (900.0,967.9) -- (902.0,966.4) -- (904.1,964.8) -- (905.6,963.7) -- (906.5,963.0) -- (908.5,961.5) -- (911.4,959.2) -- (912.9,958.0) -- (913.0,957.9) -- (914.5,956.7) -- (917.5,954.2) -- (919.0,953.0);
\node[isolabel,rotate=-10] at (673.0,782.6) {1004};
\node[isolabel,rotate=25] at (661.0,862.2) {1008};
\node[isolabel,rotate=45] at (685.0,918.2) {1012};
\node[isolabel,rotate=35] at (877.0,882.2) {1016};
\node[isolabel,rotate=51] at (825.8,1030.0) {1020};
\draw[fblue,line width=1.10pt] (655.0,702.0) -- (620.0,740.0) -- (600.0,770.0) -- (572.0,815.0) -- (553.0,865.0) -- (550.0,915.0) -- (540.0,970.0) -- (535.0,1020.0) -- (525.0,1072.0);
\fill[fblue] (651.6,705.7) -- (654.4,721.8) -- (638.1,720.4) -- cycle;
\fill[fblue] (631.3,727.7) -- (634.1,743.9) -- (617.7,742.5) -- cycle;
\fill[fblue] (612.6,751.1) -- (617.9,766.6) -- (601.5,767.7) -- cycle;
\fill[fblue] (596.2,776.2) -- (601.9,791.5) -- (585.6,793.2) -- cycle;
\fill[fblue] (580.3,801.7) -- (586.1,817.0) -- (569.7,818.6) -- cycle;
\fill[fblue] (566.9,828.4) -- (575.5,842.3) -- (559.8,847.0) -- cycle;
\fill[fblue] (553.6,855.8) -- (565.9,866.6) -- (552.4,875.8) -- cycle;
\fill[fblue] (551.8,885.8) -- (564.1,896.5) -- (550.6,905.7) -- cycle;
\fill[fblue] (549.9,915.7) -- (560.9,927.9) -- (546.3,935.4) -- cycle;
\fill[fblue] (544.5,945.2) -- (555.5,957.4) -- (540.9,964.9) -- cycle;
\fill[fblue] (539.5,974.8) -- (551.5,986.0) -- (537.5,994.7) -- cycle;
\fill[fblue] (536.5,1004.6) -- (548.5,1015.9) -- (534.5,1024.5) -- cycle;
\fill[fblue] (532.3,1034.3) -- (543.1,1046.6) -- (528.5,1053.9) -- cycle;
\draw[fred,line width=1.10pt] (655.0,702.0) -- (690.0,704.0) -- (722.0,710.0) -- (760.0,735.0) -- (818.0,768.0) -- (870.0,772.0) -- (940.0,768.0);
\fill[fred] (661.5,702.4) -- (661.8,700.5) -- (662.5,698.7) -- (663.6,697.2) -- (665.1,695.9) -- (666.7,695.0) -- (668.6,694.5) -- (670.5,694.4) -- (672.3,694.7) -- (674.1,695.4) -- (675.6,696.5) -- (676.9,697.9) -- (677.8,699.6) -- (678.4,701.4) -- (678.5,703.3) -- cycle;
\fill[fred] (691.4,704.3) -- (692.0,702.4) -- (692.9,700.8) -- (694.2,699.4) -- (695.8,698.3) -- (697.6,697.6) -- (699.4,697.3) -- (701.3,697.5) -- (703.2,698.0) -- (704.8,699.0) -- (706.2,700.3) -- (707.3,701.8) -- (708.0,703.6) -- (708.3,705.5) -- (708.1,707.4) -- cycle;
\fill[fred] (721.1,709.4) -- (722.3,707.9) -- (723.8,706.8) -- (725.5,706.0) -- (727.4,705.6) -- (729.3,705.6) -- (731.1,706.1) -- (732.8,707.0) -- (734.3,708.2) -- (735.5,709.7) -- (736.2,711.4) -- (736.6,713.3) -- (736.6,715.2) -- (736.1,717.0) -- (735.3,718.7) -- cycle;
\fill[fred] (746.1,725.9) -- (747.3,724.4) -- (748.9,723.3) -- (750.6,722.5) -- (752.5,722.1) -- (754.4,722.1) -- (756.2,722.6) -- (757.9,723.4) -- (759.4,724.7) -- (760.5,726.2) -- (761.3,727.9) -- (761.7,729.8) -- (761.7,731.7) -- (761.2,733.5) -- (760.3,735.2) -- cycle;
\fill[fred] (771.6,741.6) -- (772.8,740.1) -- (774.2,738.8) -- (775.9,737.9) -- (777.7,737.4) -- (779.6,737.3) -- (781.5,737.7) -- (783.2,738.4) -- (784.8,739.6) -- (786.0,741.0) -- (786.9,742.7) -- (787.4,744.5) -- (787.5,746.4) -- (787.2,748.3) -- (786.4,750.0) -- cycle;
\fill[fred] (797.7,756.5) -- (798.8,754.9) -- (800.3,753.7) -- (802.0,752.8) -- (803.8,752.3) -- (805.7,752.2) -- (807.6,752.5) -- (809.3,753.3) -- (810.9,754.4) -- (812.1,755.8) -- (813.0,757.5) -- (813.5,759.3) -- (813.6,761.2) -- (813.2,763.1) -- (812.5,764.9) -- cycle;
\fill[fred] (824.6,768.5) -- (825.0,766.6) -- (825.8,764.9) -- (826.9,763.4) -- (828.3,762.1) -- (830.0,761.2) -- (831.9,760.8) -- (833.8,760.7) -- (835.6,761.0) -- (837.4,761.8) -- (838.9,762.9) -- (840.2,764.4) -- (841.0,766.1) -- (841.5,767.9) -- (841.6,769.8) -- cycle;
\fill[fred] (854.6,770.8) -- (854.9,768.9) -- (855.7,767.2) -- (856.8,765.7) -- (858.3,764.4) -- (859.9,763.5) -- (861.8,763.1) -- (863.7,763.0) -- (865.6,763.3) -- (867.3,764.1) -- (868.8,765.2) -- (870.1,766.7) -- (871.0,768.4) -- (871.4,770.2) -- (871.5,772.1) -- cycle;
\fill[fred] (884.5,771.2) -- (884.6,769.3) -- (885.1,767.4) -- (886.0,765.8) -- (887.3,764.4) -- (888.9,763.3) -- (890.6,762.5) -- (892.5,762.2) -- (894.4,762.3) -- (896.2,762.8) -- (897.9,763.7) -- (899.3,765.0) -- (900.4,766.6) -- (901.1,768.3) -- (901.5,770.2) -- cycle;
\fill[fred] (914.4,769.5) -- (914.5,767.6) -- (915.1,765.7) -- (916.0,764.1) -- (917.3,762.6) -- (918.8,761.5) -- (920.6,760.8) -- (922.4,760.5) -- (924.3,760.6) -- (926.2,761.1) -- (927.8,762.0) -- (929.3,763.3) -- (930.4,764.9) -- (931.1,766.6) -- (931.4,768.5) -- cycle;
\node[lowmark] at (655.0,706.0) {L};
\draw[black,line width=0.51pt] (476.0,583.0) -- (456.5,551.8);
\draw[black,line width=0.51pt] (456.5,551.8) -- (464.2,539.2);
\draw[black,line width=0.51pt,fill=white] (476.0,583.0) circle (7.8);
\fill[black] (476.0,583.0) -- (476.0,575.2) -- (477.0,575.2) -- (478.0,575.4) -- (479.0,575.8) -- (479.9,576.2) -- (480.8,576.8) -- (481.5,577.5) -- (482.2,578.2) -- (482.8,579.1) -- (483.2,580.0) -- (483.6,581.0) -- (483.8,582.0) -- (483.8,583.0) -- cycle;
\node[anchor=east,inner sep=0pt,font=\fontsize{5.4}{5.4}\selectfont,text=black] at (465.9,572.9) {34};
\node[anchor=east,inner sep=0pt,font=\fontsize{5.4}{5.4}\selectfont,text=black] at (465.9,601.4) {22};
\node[anchor=west,inner sep=0pt,font=\fontsize{5.4}{5.4}\selectfont,text=black] at (486.1,572.9) {110};
\draw[black,line width=0.51pt] (626.0,593.0) -- (662.7,589.8);
\draw[black,line width=0.51pt] (662.7,589.8) -- (670.4,602.4);
\draw[black,line width=0.51pt,fill=white] (626.0,593.0) circle (7.8);
\fill[black] (626.0,593.0) circle (7.8);
\node[anchor=east,inner sep=0pt,font=\fontsize{5.4}{5.4}\selectfont,text=black] at (615.9,582.9) {45};
\node[anchor=east,inner sep=0pt,font=\fontsize{5.4}{5.4}\selectfont,text=black] at (615.9,611.4) {43};
\node[anchor=west,inner sep=0pt,font=\fontsize{5.4}{5.4}\selectfont,text=black] at (636.1,582.9) {096};
\fill[black] (605.2,597.6) circle (2.0);
\fill[black] (611.8,597.6) circle (2.0);
\draw[black,line width=0.51pt] (824.0,581.0) -- (859.9,572.7);
\draw[black,line width=0.51pt] (859.9,572.7) -- (869.3,584.1);
\draw[black,line width=0.51pt] (854.5,574.0) -- (859.2,579.7);
\draw[black,line width=0.51pt,fill=white] (824.0,581.0) circle (7.8);
\fill[black] (824.0,581.0) -- (824.0,573.2) -- (825.0,573.2) -- (826.0,573.4) -- (827.0,573.8) -- (827.9,574.2) -- (828.8,574.8) -- (829.5,575.5) -- (830.2,576.2) -- (830.8,577.1) -- (831.2,578.0) -- (831.6,579.0) -- (831.8,580.0) -- (831.8,581.0) -- (831.8,582.0) -- (831.6,583.0) -- (831.2,584.0) -- (830.8,584.9) -- (830.2,585.8) -- (829.5,586.5) -- (828.8,587.2) -- (827.9,587.8) -- (827.0,588.2) -- (826.0,588.6) -- (825.0,588.8) -- (824.0,588.8) -- (823.0,588.8) -- (822.0,588.6) -- (821.0,588.2) -- (820.1,587.8) -- (819.2,587.2) -- (818.5,586.5) -- (817.8,585.8) -- (817.2,584.9) -- (816.8,584.0) -- (816.4,583.0) -- (816.2,582.0) -- (816.2,581.0) -- cycle;
\node[anchor=east,inner sep=0pt,font=\fontsize{5.4}{5.4}\selectfont,text=black] at (813.9,570.9) {39};
\node[anchor=east,inner sep=0pt,font=\fontsize{5.4}{5.4}\selectfont,text=black] at (813.9,599.4) {39};
\node[anchor=west,inner sep=0pt,font=\fontsize{5.4}{5.4}\selectfont,text=black] at (834.1,570.9) {154};
\draw[black,line width=0.51pt] (465.0,654.0) -- (441.3,625.8);
\draw[black,line width=0.51pt] (441.3,625.8) -- (447.2,612.2);
\draw[black,line width=0.51pt,fill=white] (465.0,654.0) circle (7.8);
\fill[black] (465.0,654.0) -- (465.0,646.2) -- (466.0,646.2) -- (467.0,646.4) -- (468.0,646.8) -- (468.9,647.2) -- (469.8,647.8) -- (470.5,648.5) -- (471.2,649.2) -- (471.8,650.1) -- (472.2,651.0) -- (472.6,652.0) -- (472.8,653.0) -- (472.8,654.0) -- (472.8,655.0) -- (472.6,656.0) -- (472.2,657.0) -- (471.8,657.9) -- (471.2,658.8) -- (470.5,659.5) -- (469.8,660.2) -- (468.9,660.8) -- (468.0,661.2) -- (467.0,661.6) -- (466.0,661.8) -- (465.0,661.8) -- cycle;
\node[anchor=east,inner sep=0pt,font=\fontsize{5.4}{5.4}\selectfont,text=black] at (454.9,643.9) {36};
\node[anchor=east,inner sep=0pt,font=\fontsize{5.4}{5.4}\selectfont,text=black] at (454.9,672.4) {28};
\node[anchor=west,inner sep=0pt,font=\fontsize{5.4}{5.4}\selectfont,text=black] at (475.1,643.9) {124};
\draw[black,line width=0.51pt] (544.0,662.0) -- (519.4,634.7);
\draw[black,line width=0.51pt] (519.4,634.7) -- (524.8,620.9);
\draw[black,line width=0.51pt] (523.1,638.8) -- (528.5,625.0);
\draw[black,line width=0.51pt,fill=white] (544.0,662.0) circle (7.8);
\fill[black] (544.0,662.0) circle (7.8);
\node[anchor=east,inner sep=0pt,font=\fontsize{5.4}{5.4}\selectfont,text=black] at (533.9,651.9) {40};
\node[anchor=east,inner sep=0pt,font=\fontsize{5.4}{5.4}\selectfont,text=black] at (533.9,680.4) {38};
\node[anchor=west,inner sep=0pt,font=\fontsize{5.4}{5.4}\selectfont,text=black] at (554.1,651.9) {088};
\fill[black] (526.5,666.6) circle (2.0);
\draw[black,line width=0.51pt] (738.0,624.0) -- (767.0,646.7);
\draw[black,line width=0.51pt] (767.0,646.7) -- (764.1,661.2);
\draw[black,line width=0.51pt] (762.6,643.3) -- (759.7,657.8);
\draw[black,line width=0.51pt,fill=white] (738.0,624.0) circle (7.8);
\fill[black] (738.0,624.0) circle (7.8);
\node[anchor=east,inner sep=0pt,font=\fontsize{5.4}{5.4}\selectfont,text=black] at (727.9,613.9) {39};
\node[anchor=east,inner sep=0pt,font=\fontsize{5.4}{5.4}\selectfont,text=black] at (727.9,642.4) {39};
\node[anchor=west,inner sep=0pt,font=\fontsize{5.4}{5.4}\selectfont,text=black] at (748.1,613.9) {084};
\fill[black] (717.2,628.6) circle (2.0);
\fill[black] (723.8,628.6) circle (2.0);
\draw[black,line width=0.51pt] (822.0,635.0) -- (857.0,646.4);
\draw[black,line width=0.51pt] (857.0,646.4) -- (859.2,661.0);
\draw[black,line width=0.51pt] (851.7,644.7) -- (854.0,659.3);
\draw[black,line width=0.51pt,fill=white] (822.0,635.0) circle (7.8);
\fill[black] (822.0,635.0) -- (822.0,627.2) -- (823.0,627.2) -- (824.0,627.4) -- (825.0,627.8) -- (825.9,628.2) -- (826.8,628.8) -- (827.5,629.5) -- (828.2,630.2) -- (828.8,631.1) -- (829.2,632.0) -- (829.6,633.0) -- (829.8,634.0) -- (829.8,635.0) -- (829.8,636.0) -- (829.6,637.0) -- (829.2,638.0) -- (828.8,638.9) -- (828.2,639.8) -- (827.5,640.5) -- (826.8,641.2) -- (825.9,641.8) -- (825.0,642.2) -- (824.0,642.6) -- (823.0,642.8) -- (822.0,642.8) -- cycle;
\draw[black,line width=0.51pt] (814.2,635.0) -- (829.8,635.0);
\draw[black,line width=0.51pt] (822.0,627.2) -- (822.0,642.8);
\node[anchor=east,inner sep=0pt,font=\fontsize{5.4}{5.4}\selectfont,text=black] at (811.9,624.9) {43};
\node[anchor=east,inner sep=0pt,font=\fontsize{5.4}{5.4}\selectfont,text=black] at (811.9,653.4) {41};
\node[anchor=west,inner sep=0pt,font=\fontsize{5.4}{5.4}\selectfont,text=black] at (832.1,624.9) {132};
\draw[black,line width=0.51pt] (672.0,681.0) -- (701.0,658.3);
\draw[black,line width=0.51pt] (701.0,658.3) -- (714.4,664.7);
\draw[black,line width=0.51pt,fill=white] (672.0,681.0) circle (7.8);
\fill[black] (672.0,681.0) circle (7.8);
\node[anchor=east,inner sep=0pt,font=\fontsize{5.4}{5.4}\selectfont,text=black] at (661.9,670.9) {48};
\node[anchor=east,inner sep=0pt,font=\fontsize{5.4}{5.4}\selectfont,text=black] at (661.9,699.4) {45};
\node[anchor=west,inner sep=0pt,font=\fontsize{5.4}{5.4}\selectfont,text=black] at (682.1,670.9) {024};
\fill[black] (651.2,685.6) circle (2.0);
\fill[black] (657.8,685.6) circle (2.0);
\draw[black,line width=0.51pt] (600.0,710.0) -- (565.9,696.2);
\draw[black,line width=0.51pt] (565.9,696.2) -- (564.7,681.4);
\draw[black,line width=0.51pt,fill=white] (600.0,710.0) circle (7.8);
\fill[black] (600.0,710.0) circle (7.8);
\node[anchor=east,inner sep=0pt,font=\fontsize{5.4}{5.4}\selectfont,text=black] at (589.9,699.9) {48};
\node[anchor=west,inner sep=0pt,font=\fontsize{5.4}{5.4}\selectfont,text=black] at (610.1,699.9) {026};
\fill[black] (579.2,714.6) circle (2.0);
\fill[black] (585.8,714.6) circle (2.0);
\draw[black,line width=0.51pt] (824.0,683.0) -- (860.4,688.1);
\draw[black,line width=0.51pt] (855.0,687.4) -- (857.3,694.4);
\draw[black,line width=0.51pt,fill=white] (824.0,683.0) circle (7.8);
\fill[black] (824.0,683.0) circle (7.8);
\node[anchor=east,inner sep=0pt,font=\fontsize{5.4}{5.4}\selectfont,text=black] at (813.9,672.9) {46};
\node[anchor=east,inner sep=0pt,font=\fontsize{5.4}{5.4}\selectfont,text=black] at (813.9,701.4) {43};
\node[anchor=west,inner sep=0pt,font=\fontsize{5.4}{5.4}\selectfont,text=black] at (834.1,672.9) {124};
\fill[black] (803.2,687.6) circle (2.0);
\fill[black] (809.8,687.6) circle (2.0);
\draw[black,line width=0.51pt] (908.0,712.0) -- (943.5,721.5);
\draw[black,line width=0.51pt,fill=white] (908.0,712.0) circle (7.8);
\fill[black] (908.0,712.0) circle (7.8);
\node[anchor=east,inner sep=0pt,font=\fontsize{5.4}{5.4}\selectfont,text=black] at (897.9,701.9) {45};
\node[anchor=east,inner sep=0pt,font=\fontsize{5.4}{5.4}\selectfont,text=black] at (897.9,730.4) {43};
\node[anchor=west,inner sep=0pt,font=\fontsize{5.4}{5.4}\selectfont,text=black] at (918.1,701.9) {163};
\fill[black] (890.5,716.6) circle (2.0);
\draw[black,line width=0.51pt] (419.0,722.0) -- (394.4,694.7);
\draw[black,line width=0.51pt] (394.4,694.7) -- (399.8,680.9);
\draw[black,line width=0.51pt,fill=white] (419.0,722.0) circle (7.8);
\node[anchor=east,inner sep=0pt,font=\fontsize{5.4}{5.4}\selectfont,text=black] at (408.9,711.9) {39};
\node[anchor=east,inner sep=0pt,font=\fontsize{5.4}{5.4}\selectfont,text=black] at (408.9,740.4) {31};
\node[anchor=west,inner sep=0pt,font=\fontsize{5.4}{5.4}\selectfont,text=black] at (429.1,711.9) {124};
\draw[black,line width=0.51pt] (714.0,722.0) -- (698.4,755.4);
\draw[black,line width=0.51pt] (698.4,755.4) -- (683.6,755.8);
\draw[black,line width=0.51pt,fill=white] (714.0,722.0) circle (7.8);
\fill[black] (714.0,722.0) -- (714.0,714.2) -- (715.0,714.2) -- (716.0,714.4) -- (717.0,714.8) -- (717.9,715.2) -- (718.8,715.8) -- (719.5,716.5) -- (720.2,717.2) -- (720.8,718.1) -- (721.2,719.0) -- (721.6,720.0) -- (721.8,721.0) -- (721.8,722.0) -- (721.8,723.0) -- (721.6,724.0) -- (721.2,725.0) -- (720.8,725.9) -- (720.2,726.8) -- (719.5,727.5) -- (718.8,728.2) -- (717.9,728.8) -- (717.0,729.2) -- (716.0,729.6) -- (715.0,729.8) -- (714.0,729.8) -- (713.0,729.8) -- (712.0,729.6) -- (711.0,729.2) -- (710.1,728.8) -- (709.2,728.2) -- (708.5,727.5) -- (707.8,726.8) -- (707.2,725.9) -- (706.8,725.0) -- (706.4,724.0) -- (706.2,723.0) -- (706.2,722.0) -- cycle;
\node[anchor=east,inner sep=0pt,font=\fontsize{5.4}{5.4}\selectfont,text=black] at (703.9,711.9) {57};
\node[anchor=east,inner sep=0pt,font=\fontsize{5.4}{5.4}\selectfont,text=black] at (703.9,740.4) {54};
\node[anchor=west,inner sep=0pt,font=\fontsize{5.4}{5.4}\selectfont,text=black] at (724.1,711.9) {029};
\draw[black,line width=0.51pt] (818.0,742.0) -- (854.2,748.4);
\draw[black,line width=0.51pt] (854.2,748.4) -- (858.5,762.6);
\draw[black,line width=0.51pt,fill=white] (818.0,742.0) circle (7.8);
\fill[black] (818.0,742.0) circle (7.8);
\node[anchor=east,inner sep=0pt,font=\fontsize{5.4}{5.4}\selectfont,text=black] at (807.9,731.9) {54};
\node[anchor=east,inner sep=0pt,font=\fontsize{5.4}{5.4}\selectfont,text=black] at (807.9,760.4) {45};
\node[anchor=west,inner sep=0pt,font=\fontsize{5.4}{5.4}\selectfont,text=black] at (828.1,731.9) {097};
\draw[black,line width=0.51pt] (469.0,770.0) -- (438.5,749.4);
\draw[black,line width=0.51pt] (438.5,749.4) -- (440.4,734.7);
\draw[black,line width=0.51pt,fill=white] (469.0,770.0) circle (7.8);
\draw[black,line width=0.51pt] (469.0,762.2) -- (469.0,777.8);
\node[anchor=east,inner sep=0pt,font=\fontsize{5.4}{5.4}\selectfont,text=black] at (458.9,759.9) {41};
\node[anchor=east,inner sep=0pt,font=\fontsize{5.4}{5.4}\selectfont,text=black] at (458.9,788.4) {30};
\node[anchor=west,inner sep=0pt,font=\fontsize{5.4}{5.4}\selectfont,text=black] at (479.1,759.9) {103};
\draw[black,line width=0.51pt] (570.0,768.0) -- (544.9,741.1);
\draw[black,line width=0.51pt] (544.9,741.1) -- (550.1,727.2);
\draw[black,line width=0.51pt] (548.7,745.1) -- (553.8,731.2);
\draw[black,line width=0.51pt] (552.4,749.2) -- (555.0,742.2);
\draw[black,line width=0.51pt,fill=white] (570.0,768.0) circle (7.8);
\fill[black] (570.0,768.0) -- (570.0,760.2) -- (571.0,760.2) -- (572.0,760.4) -- (573.0,760.8) -- (573.9,761.2) -- (574.8,761.8) -- (575.5,762.5) -- (576.2,763.2) -- (576.8,764.1) -- (577.2,765.0) -- (577.6,766.0) -- (577.8,767.0) -- (577.8,768.0) -- (577.8,769.0) -- (577.6,770.0) -- (577.2,771.0) -- (576.8,771.9) -- (576.2,772.8) -- (575.5,773.5) -- (574.8,774.2) -- (573.9,774.8) -- (573.0,775.2) -- (572.0,775.6) -- (571.0,775.8) -- (570.0,775.8) -- cycle;
\node[anchor=east,inner sep=0pt,font=\fontsize{5.4}{5.4}\selectfont,text=black] at (559.9,757.9) {45};
\node[anchor=east,inner sep=0pt,font=\fontsize{5.4}{5.4}\selectfont,text=black] at (559.9,786.4) {40};
\node[anchor=west,inner sep=0pt,font=\fontsize{5.4}{5.4}\selectfont,text=black] at (580.1,757.9) {069};
\draw[black,line width=0.51pt] (637.0,765.0) -- (618.0,796.5);
\draw[black,line width=0.51pt] (618.0,796.5) -- (603.3,795.4);
\draw[black,line width=0.51pt] (620.9,791.8) -- (606.1,790.7);
\draw[black,line width=0.51pt] (623.7,787.1) -- (616.3,786.5);
\draw[black,line width=0.51pt,fill=white] (637.0,765.0) circle (7.8);
\fill[black] (637.0,765.0) circle (7.8);
\node[anchor=east,inner sep=0pt,font=\fontsize{5.4}{5.4}\selectfont,text=black] at (626.9,754.9) {60};
\node[anchor=east,inner sep=0pt,font=\fontsize{5.4}{5.4}\selectfont,text=black] at (626.9,783.4) {55};
\node[anchor=west,inner sep=0pt,font=\fontsize{5.4}{5.4}\selectfont,text=black] at (647.1,754.9) {040};
\draw[black,line width=0.43pt] (616.4,763.3) -- (616.4,775.9);
\draw[black,line width=0.43pt] (616.4,763.3) -- (621.1,763.3) .. controls (625.4,763.3) and (625.4,770.0) .. (621.1,770.0) -- (616.4,770.0);
\draw[black,line width=0.43pt] (620.3,770.0) -- (625.8,775.9);
\fill[black] (626.9,777.0) -- (623.1,776.1) -- (625.6,773.4) -- cycle;
\fill[black] (618.7,759.4) circle (1.9);
\draw[black,line width=0.51pt] (704.0,785.0) -- (685.0,816.5);
\draw[black,line width=0.51pt] (685.0,816.5) -- (670.3,815.4);
\draw[black,line width=0.51pt,fill=white] (704.0,785.0) circle (7.8);
\fill[black] (704.0,785.0) -- (704.0,777.2) -- (705.0,777.2) -- (706.0,777.4) -- (707.0,777.8) -- (707.9,778.2) -- (708.8,778.8) -- (709.5,779.5) -- (710.2,780.2) -- (710.8,781.1) -- (711.2,782.0) -- (711.6,783.0) -- (711.8,784.0) -- (711.8,785.0) -- (711.8,786.0) -- (711.6,787.0) -- (711.2,788.0) -- (710.8,788.9) -- (710.2,789.8) -- (709.5,790.5) -- (708.8,791.2) -- (707.9,791.8) -- (707.0,792.2) -- (706.0,792.6) -- (705.0,792.8) -- (704.0,792.8) -- cycle;
\node[anchor=east,inner sep=0pt,font=\fontsize{5.4}{5.4}\selectfont,text=black] at (693.9,774.9) {61};
\node[anchor=east,inner sep=0pt,font=\fontsize{5.4}{5.4}\selectfont,text=black] at (693.9,803.4) {57};
\node[anchor=west,inner sep=0pt,font=\fontsize{5.4}{5.4}\selectfont,text=black] at (714.1,774.9) {044};
\fill[black] (683.2,789.6) circle (2.0);
\fill[black] (689.8,789.6) circle (2.0);
\draw[black,line width=0.51pt] (787.0,798.0) -- (787.0,834.8);
\draw[black,line width=0.51pt] (787.0,834.8) -- (773.8,841.4);
\draw[black,line width=0.51pt,fill=white] (787.0,798.0) circle (7.8);
\fill[black] (787.0,798.0) -- (787.0,790.2) -- (788.0,790.2) -- (789.0,790.4) -- (790.0,790.8) -- (790.9,791.2) -- (791.8,791.8) -- (792.5,792.5) -- (793.2,793.2) -- (793.8,794.1) -- (794.2,795.0) -- (794.6,796.0) -- (794.8,797.0) -- (794.8,798.0) -- (794.8,799.0) -- (794.6,800.0) -- (794.2,801.0) -- (793.8,801.9) -- (793.2,802.8) -- (792.5,803.5) -- (791.8,804.2) -- (790.9,804.8) -- (790.0,805.2) -- (789.0,805.6) -- (788.0,805.8) -- (787.0,805.8) -- (786.0,805.8) -- (785.0,805.6) -- (784.0,805.2) -- (783.1,804.8) -- (782.2,804.2) -- (781.5,803.5) -- (780.8,802.8) -- (780.2,801.9) -- (779.8,801.0) -- (779.4,800.0) -- (779.2,799.0) -- (779.2,798.0) -- cycle;
\node[anchor=east,inner sep=0pt,font=\fontsize{5.4}{5.4}\selectfont,text=black] at (776.9,787.9) {66};
\node[anchor=east,inner sep=0pt,font=\fontsize{5.4}{5.4}\selectfont,text=black] at (776.9,816.4) {53};
\node[anchor=west,inner sep=0pt,font=\fontsize{5.4}{5.4}\selectfont,text=black] at (797.1,787.9) {099};
\draw[black,line width=0.51pt] (828.0,818.0) -- (828.0,854.8);
\draw[black,line width=0.51pt] (828.0,854.8) -- (814.8,861.4);
\draw[black,line width=0.51pt,fill=white] (828.0,818.0) circle (7.8);
\fill[black] (828.0,818.0) -- (828.0,810.2) -- (829.0,810.2) -- (830.0,810.4) -- (831.0,810.8) -- (831.9,811.2) -- (832.8,811.8) -- (833.5,812.5) -- (834.2,813.2) -- (834.8,814.1) -- (835.2,815.0) -- (835.6,816.0) -- (835.8,817.0) -- (835.8,818.0) -- cycle;
\node[anchor=east,inner sep=0pt,font=\fontsize{5.4}{5.4}\selectfont,text=black] at (817.9,807.9) {66};
\node[anchor=east,inner sep=0pt,font=\fontsize{5.4}{5.4}\selectfont,text=black] at (817.9,836.4) {55};
\node[anchor=west,inner sep=0pt,font=\fontsize{5.4}{5.4}\selectfont,text=black] at (838.1,807.9) {111};
\draw[black,line width=0.51pt] (405.0,838.0) -- (374.9,816.9);
\draw[black,line width=0.51pt] (374.9,816.9) -- (377.0,802.2);
\draw[black,line width=0.51pt,fill=white] (405.0,838.0) circle (7.8);
\node[anchor=east,inner sep=0pt,font=\fontsize{5.4}{5.4}\selectfont,text=black] at (394.9,827.9) {41};
\node[anchor=east,inner sep=0pt,font=\fontsize{5.4}{5.4}\selectfont,text=black] at (394.9,856.4) {33};
\node[anchor=west,inner sep=0pt,font=\fontsize{5.4}{5.4}\selectfont,text=black] at (415.1,827.9) {126};
\draw[black,line width=0.51pt] (496.0,831.0) -- (467.8,807.3);
\draw[black,line width=0.51pt] (467.8,807.3) -- (471.3,792.9);
\draw[black,line width=0.51pt,fill=white] (496.0,831.0) circle (7.8);
\node[anchor=east,inner sep=0pt,font=\fontsize{5.4}{5.4}\selectfont,text=black] at (485.9,820.9) {42};
\node[anchor=east,inner sep=0pt,font=\fontsize{5.4}{5.4}\selectfont,text=black] at (485.9,849.4) {36};
\node[anchor=west,inner sep=0pt,font=\fontsize{5.4}{5.4}\selectfont,text=black] at (506.1,820.9) {109};
\draw[black,line width=0.51pt] (608.0,861.0) -- (588.5,892.2);
\draw[black,line width=0.51pt] (588.5,892.2) -- (573.8,890.8);
\draw[black,line width=0.51pt,fill=white] (608.0,861.0) circle (7.8);
\fill[black] (608.0,861.0) circle (7.8);
\node[anchor=east,inner sep=0pt,font=\fontsize{5.4}{5.4}\selectfont,text=black] at (597.9,850.9) {61};
\node[anchor=east,inner sep=0pt,font=\fontsize{5.4}{5.4}\selectfont,text=black] at (597.9,879.4) {57};
\node[anchor=west,inner sep=0pt,font=\fontsize{5.4}{5.4}\selectfont,text=black] at (618.1,850.9) {075};
\fill[black] (587.2,865.6) circle (2.0);
\fill[black] (593.8,865.6) circle (2.0);
\draw[black,line width=0.51pt] (828.0,885.0) -- (828.0,921.8);
\draw[black,line width=0.51pt] (828.0,921.8) -- (814.8,928.4);
\draw[black,line width=0.51pt,fill=white] (828.0,885.0) circle (7.8);
\fill[black] (828.0,885.0) -- (828.0,877.2) -- (829.0,877.2) -- (830.0,877.4) -- (831.0,877.8) -- (831.9,878.2) -- (832.8,878.8) -- (833.5,879.5) -- (834.2,880.2) -- (834.8,881.1) -- (835.2,882.0) -- (835.6,883.0) -- (835.8,884.0) -- (835.8,885.0) -- (835.8,886.0) -- (835.6,887.0) -- (835.2,888.0) -- (834.8,888.9) -- (834.2,889.8) -- (833.5,890.5) -- (832.8,891.2) -- (831.9,891.8) -- (831.0,892.2) -- (830.0,892.6) -- (829.0,892.8) -- (828.0,892.8) -- cycle;
\node[anchor=east,inner sep=0pt,font=\fontsize{5.4}{5.4}\selectfont,text=black] at (817.9,874.9) {71};
\node[anchor=east,inner sep=0pt,font=\fontsize{5.4}{5.4}\selectfont,text=black] at (817.9,903.4) {60};
\node[anchor=west,inner sep=0pt,font=\fontsize{5.4}{5.4}\selectfont,text=black] at (838.1,874.9) {149};
\draw[black,line width=0.51pt] (498.0,909.0) -- (465.5,891.7);
\draw[black,line width=0.51pt] (465.5,891.7) -- (465.9,876.9);
\draw[black,line width=0.51pt,fill=white] (498.0,909.0) circle (7.8);
\node[anchor=east,inner sep=0pt,font=\fontsize{5.4}{5.4}\selectfont,text=black] at (487.9,898.9) {48};
\node[anchor=east,inner sep=0pt,font=\fontsize{5.4}{5.4}\selectfont,text=black] at (487.9,927.4) {38};
\node[anchor=west,inner sep=0pt,font=\fontsize{5.4}{5.4}\selectfont,text=black] at (508.1,898.9) {126};
\draw[black,line width=0.51pt] (729.0,914.0) -- (729.0,950.8);
\draw[black,line width=0.51pt] (729.0,950.8) -- (715.8,957.4);
\draw[black,line width=0.51pt,fill=white] (729.0,914.0) circle (7.8);
\fill[black] (729.0,914.0) -- (729.0,906.2) -- (730.0,906.2) -- (731.0,906.4) -- (732.0,906.8) -- (732.9,907.2) -- (733.8,907.8) -- (734.5,908.5) -- (735.2,909.2) -- (735.8,910.1) -- (736.2,911.0) -- (736.6,912.0) -- (736.8,913.0) -- (736.8,914.0) -- (736.8,915.0) -- (736.6,916.0) -- (736.2,917.0) -- (735.8,917.9) -- (735.2,918.8) -- (734.5,919.5) -- (733.8,920.2) -- (732.9,920.8) -- (732.0,921.2) -- (731.0,921.6) -- (730.0,921.8) -- (729.0,921.8) -- cycle;
\node[anchor=east,inner sep=0pt,font=\fontsize{5.4}{5.4}\selectfont,text=black] at (718.9,903.9) {70};
\node[anchor=east,inner sep=0pt,font=\fontsize{5.4}{5.4}\selectfont,text=black] at (718.9,932.4) {59};
\node[anchor=west,inner sep=0pt,font=\fontsize{5.4}{5.4}\selectfont,text=black] at (739.1,903.9) {140};
\draw[black,line width=0.51pt] (390.0,946.0) -- (356.6,930.4);
\draw[black,line width=0.51pt] (356.6,930.4) -- (356.2,915.6);
\draw[black,line width=0.51pt,fill=white] (390.0,946.0) circle (7.8);
\node[anchor=east,inner sep=0pt,font=\fontsize{5.4}{5.4}\selectfont,text=black] at (379.9,935.9) {49};
\node[anchor=east,inner sep=0pt,font=\fontsize{5.4}{5.4}\selectfont,text=black] at (379.9,964.4) {38};
\node[anchor=west,inner sep=0pt,font=\fontsize{5.4}{5.4}\selectfont,text=black] at (400.1,935.9) {128};
\draw[black,line width=0.51pt] (603.0,945.0) -- (573.2,966.6);
\draw[black,line width=0.51pt] (573.2,966.6) -- (560.1,959.8);
\draw[black,line width=0.51pt,fill=white] (603.0,945.0) circle (7.8);
\fill[black] (603.0,945.0) -- (603.0,937.2) -- (604.0,937.2) -- (605.0,937.4) -- (606.0,937.8) -- (606.9,938.2) -- (607.8,938.8) -- (608.5,939.5) -- (609.2,940.2) -- (609.8,941.1) -- (610.2,942.0) -- (610.6,943.0) -- (610.8,944.0) -- (610.8,945.0) -- (610.8,946.0) -- (610.6,947.0) -- (610.2,948.0) -- (609.8,948.9) -- (609.2,949.8) -- (608.5,950.5) -- (607.8,951.2) -- (606.9,951.8) -- (606.0,952.2) -- (605.0,952.6) -- (604.0,952.8) -- (603.0,952.8) -- (602.0,952.8) -- (601.0,952.6) -- (600.0,952.2) -- (599.1,951.8) -- (598.2,951.2) -- (597.5,950.5) -- (596.8,949.8) -- (596.2,948.9) -- (595.8,948.0) -- (595.4,947.0) -- (595.2,946.0) -- (595.2,945.0) -- cycle;
\node[anchor=east,inner sep=0pt,font=\fontsize{5.4}{5.4}\selectfont,text=black] at (592.9,934.9) {66};
\node[anchor=east,inner sep=0pt,font=\fontsize{5.4}{5.4}\selectfont,text=black] at (592.9,963.4) {59};
\node[anchor=west,inner sep=0pt,font=\fontsize{5.4}{5.4}\selectfont,text=black] at (613.1,934.9) {114};
\draw[black,line width=0.51pt] (648.0,960.0) -- (631.3,992.8);
\draw[black,line width=0.51pt] (631.3,992.8) -- (616.5,992.7);
\draw[black,line width=0.51pt,fill=white] (648.0,960.0) circle (7.8);
\node[anchor=east,inner sep=0pt,font=\fontsize{5.4}{5.4}\selectfont,text=black] at (637.9,949.9) {66};
\node[anchor=east,inner sep=0pt,font=\fontsize{5.4}{5.4}\selectfont,text=black] at (637.9,978.4) {60};
\node[anchor=west,inner sep=0pt,font=\fontsize{5.4}{5.4}\selectfont,text=black] at (658.1,949.9) {124};
\draw[black,line width=0.51pt] (452.0,1009.0) -- (423.8,985.3);
\draw[black,line width=0.51pt] (423.8,985.3) -- (427.3,970.9);
\draw[black,line width=0.51pt,fill=white] (452.0,1009.0) circle (7.8);
\node[anchor=east,inner sep=0pt,font=\fontsize{5.4}{5.4}\selectfont,text=black] at (441.9,998.9) {50};
\node[anchor=east,inner sep=0pt,font=\fontsize{5.4}{5.4}\selectfont,text=black] at (441.9,1027.4) {38};
\node[anchor=west,inner sep=0pt,font=\fontsize{5.4}{5.4}\selectfont,text=black] at (462.1,998.9) {120};
\draw[black,line width=0.51pt] (507.0,1061.0) -- (471.5,1070.5);
\draw[black,line width=0.51pt] (471.5,1070.5) -- (461.6,1059.4);
\draw[black,line width=0.51pt,fill=white] (507.0,1061.0) circle (7.8);
\fill[black] (507.0,1061.0) circle (7.8);
\node[anchor=east,inner sep=0pt,font=\fontsize{5.4}{5.4}\selectfont,text=black] at (496.9,1050.9) {52};
\node[anchor=east,inner sep=0pt,font=\fontsize{5.4}{5.4}\selectfont,text=black] at (496.9,1079.4) {50};
\node[anchor=west,inner sep=0pt,font=\fontsize{5.4}{5.4}\selectfont,text=black] at (517.1,1050.9) {128};
\draw[black,line width=0.51pt] (613.0,1041.0) -- (606.0,1077.1);
\draw[black,line width=0.51pt] (606.0,1077.1) -- (591.7,1081.1);
\draw[black,line width=0.51pt,fill=white] (613.0,1041.0) circle (7.8);
\fill[black] (613.0,1041.0) circle (7.8);
\node[anchor=east,inner sep=0pt,font=\fontsize{5.4}{5.4}\selectfont,text=black] at (602.9,1030.9) {63};
\node[anchor=east,inner sep=0pt,font=\fontsize{5.4}{5.4}\selectfont,text=black] at (602.9,1059.4) {61};
\node[anchor=west,inner sep=0pt,font=\fontsize{5.4}{5.4}\selectfont,text=black] at (623.1,1030.9) {153};
\draw[black,line width=0.51pt] (725.0,1031.0) -- (710.6,1064.9);
\draw[black,line width=0.51pt] (710.6,1064.9) -- (695.8,1065.8);
\draw[black,line width=0.51pt,fill=white] (725.0,1031.0) circle (7.8);
\fill[black] (725.0,1031.0) -- (725.0,1023.2) -- (726.0,1023.2) -- (727.0,1023.4) -- (728.0,1023.8) -- (728.9,1024.2) -- (729.8,1024.8) -- (730.5,1025.5) -- (731.2,1026.2) -- (731.8,1027.1) -- (732.2,1028.0) -- (732.6,1029.0) -- (732.8,1030.0) -- (732.8,1031.0) -- cycle;
\node[anchor=east,inner sep=0pt,font=\fontsize{5.4}{5.4}\selectfont,text=black] at (714.9,1020.9) {69};
\node[anchor=east,inner sep=0pt,font=\fontsize{5.4}{5.4}\selectfont,text=black] at (714.9,1049.4) {61};
\node[anchor=west,inner sep=0pt,font=\fontsize{5.4}{5.4}\selectfont,text=black] at (735.1,1020.9) {171};
\draw[black,line width=0.51pt] (822.0,1006.0) -- (806.4,1039.4);
\draw[black,line width=0.51pt] (806.4,1039.4) -- (791.6,1039.8);
\draw[black,line width=0.51pt,fill=white] (822.0,1006.0) circle (7.8);
\node[anchor=east,inner sep=0pt,font=\fontsize{5.4}{5.4}\selectfont,text=black] at (811.9,995.9) {72};
\node[anchor=east,inner sep=0pt,font=\fontsize{5.4}{5.4}\selectfont,text=black] at (811.9,1024.4) {64};
\node[anchor=west,inner sep=0pt,font=\fontsize{5.4}{5.4}\selectfont,text=black] at (832.1,995.9) {194};
}
\caption{Question 14 complete: isobars every 4\,mb (orange), the low (L), the cold front (blue triangles) and the warm front (red semicircles), the area with more than 75\,\% cloud (dashed outline) and the precipitation (blue shading), over the 34 stations. The state outlines are real boundaries. The fronts are an interpretation, drawn from the temperature, dew-point and wind contrasts.}\label{p3-fig:final}
\end{figure}
\end{landscape}

\section{Reading the finished picture}

The map is a text-book mid-latitude depression.
\begin{itemize}
\item The low is over Iowa and northern Illinois, at about 1002\,hPa.
\item The \textbf{warm sector} is the tongue of warm, humid air in the south and east, between the two fronts: temperatures in the 60s and 70s \degF\ and dew points of 55 to 64\,\degF.
\item The cold air in the north-west is dry (dew points of 22 to 38\,\degF) and mostly clear.
\item The cloud and rain are on the fronts and near the low, where the air is being lifted.
\item The winds circle the low anticlockwise (Chapter~\ref{p3-ch:stream}) and cross the isobars towards low pressure (Chapter~\ref{p3-ch:fric}).
\end{itemize}

\begin{workedbox}[Worked example: is the wind consistent with the isobars?]
Between stations 9 (1002.6\,hPa) and 5 (1008.8\,hPa), about 216\,km apart, the pressure changes by 6.2\,hPa. The geostrophic wind for that gradient at 41.5\deg N ($f=9.7\times10^{-5}$\,s$^{-1}$) is
\[
V_g=\frac{1}{\rho f}\frac{\Delta p}{\Delta n}=\frac{620}{1.2\times9.7\times10^{-5}\times2.16\times10^{5}}=25\ \ms\approx48\ \text{knots}.
\]
The observed surface winds are 10 to 25 knots. That is about what friction does: it slows the surface wind to a half or two thirds of the geostrophic value and turns it across the isobars (Chapter~\ref{p3-ch:fric}).
\end{workedbox}

\section{Questions that come up}

\begin{askbox}
\textbf{Where exactly should I put the front?}\\
There is no single right answer: a front analysis is an interpretation. Put the line along the zone where the temperature, dew point and wind change fastest, bend it to follow the stations, and begin it at the low. Check that the stations on the cold side are colder and drier than those on the warm side.
\end{askbox}

\begin{askbox}
\textbf{Why do the isobars bend sharply where a front crosses them?}\\
Because the wind changes direction at a front, and the wind follows the isobars, the isobars must change direction there too. A front lies along a trough of low pressure, with the isobars forming a V pointing away from the low.
\end{askbox}

\begin{askbox}
\textbf{Which stations count as ``cloud above 75\,\%''?}\\
Those whose circle is filled to at least three quarters: 7--8 tenths, 9 tenths (filled with a gap) or overcast. Stations with half or less are outside the outline.
\end{askbox}

\section{Try it yourself: Practice Set~\thechapter}

\begin{tryit}[Practice Set~\thechapter: fronts and isobars]
\begin{enumerate}[label=\textbf{\thechapter.\arabic*},leftmargin=*,itemsep=3pt]
\item Stations X (1011.1\,hPa) and Y (1014.9\,hPa) are neighbours. Where does the 1012\,mb isobar cross the line between them?
\item A station reports a temperature of 61\,\degF, a dew point of 57\,\degF\ and a south-south-west wind of 10 knots. A neighbour to its west reports 42\,\degF, 36\,\degF\ and a north-west wind. What kind of front lies between them?
\item Station P is 30\,\degF\ with dew point 28\,\degF, east wind, overcast and snow. Station Q to the south is 41\,\degF\ with dew point 38\,\degF\ and south-east wind. Is a warm or cold front likely to lie between them? Why?
\item Why is there no 1000\,mb isobar on Figure~\ref{p3-fig:fig10}?
\end{enumerate}
\end{tryit}

\begin{answerbox}[Answers to Practice Set~\thechapter]
\begin{enumerate}[label=\textbf{\thechapter.\arabic*},leftmargin=*,itemsep=3pt]
\item $(1012-1011.1)/(1014.9-1011.1)=0.9/3.8=0.24$: 24\,\% of the way from X to Y.
\item A \textbf{cold front}: the temperature falls 19\,\degF, the dew point 21\,\degF, and the wind turns clockwise from south-south-west to north-west as you cross it from east to west.
\item A \textbf{warm front}: the warmer, moister air (41\,\degF, 38\,\degF) lies to the south with a south-east wind, riding over the colder air with east wind to the north, where it is snowing. The wind turns clockwise from east to south-east.
\item No station reports a pressure as low as 1000\,hPa (the lowest is 1002.4), and an isobar should only be drawn where the data support it.
\end{enumerate}
\end{answerbox}
